\documentclass[12pt,a4paper,twoside]{report}

\usepackage[utf8]{inputenc}
\usepackage[T1]{fontenc}
\usepackage[french]{babel}
\usepackage{geometry}
\geometry{margin=2cm, top=2.5cm, bottom=2.5cm}
\usepackage{xcolor}
\usepackage{hyperref}
\usepackage{enumitem}
\usepackage{amssymb}
\usepackage{multicol}
\usepackage{fancyhdr}
\usepackage{titlesec}
\usepackage{setspace}

\hypersetup{
    colorlinks=true,
    linkcolor=blue,
    urlcolor=blue,
    pdftitle={QCM - Indexation et Referencement Web SEO},
    pdfauthor={Expert SEO},
}

\definecolor{primary}{RGB}{26,26,46}
\definecolor{googleblue}{RGB}{66,133,244}

\titleformat{\chapter}[display]
  {\normalfont\huge\bfseries\color{primary}}
  {\chaptername\ \thechapter}
  {20pt}
  {\Huge}

\begin{document}

\begin{titlepage}
\begin{center}
\vspace*{3cm}
{\Huge\textbf{QCM}}\\[1cm]
{\LARGE\textbf{Indexation et Référencement Web SEO}}\\[1.5cm]
{\Large 1 800 questions à choix multiples}\\[0.5cm]
{\large Couvrant l'intégralité du programme de formation}\\[2cm]
\rule{0.6\textwidth}{0.5pt}\\[1cm]
{\large Institut Supérieur du Numérique (SupNum)}\\[0.3cm]
{\large Nouakchott, Mauritanie}\\[0.3cm]
{\large Version 1.0 -- Juin 2026}\\[1cm]
\rule{0.6\textwidth}{0.5pt}\\[0.5cm]
{\small Tous droits réservés}
\end{center}
\end{titlepage}

\tableofcontents
\newpage



% ============================================
% CHAPITRE 1
% ============================================
\section*{Chapitre 1 : Histoire et évolution des moteurs de recherche}

\begin{enumerate}
\item Quel est considéré comme le premier moteur de recherche ?
\begin{itemize}
\item[A)] Yahoo!
\item[B)] AltaVista
\item[C)] Archie
\item[D)] WebCrawler
\end{itemize}
*(Réponse : C)*

\item En quelle année Archie a-t-il vu le jour ?
\begin{itemize}
\item[A)] 1994
\item[B)] 1990
\item[C)] 1995
\item[D)] 1997
\end{itemize}
*(Réponse : A)*

\item Quel service a commencé en 1994 sous forme d'annuaire web répertoriant manuellement les sites par catégories ?
\begin{itemize}
\item[A)] Google
\item[B)] Yahoo!
\item[C)] Bing
\item[D)] Ask Jeeves
\end{itemize}
*(Réponse : B)*

\item Quel a été le premier moteur à indexer le texte intégral des pages ?
\begin{itemize}
\item[A)] Archie
\item[B)] AltaVista
\item[C)] WebCrawler
\item[D)] Google
\end{itemize}
*(Réponse : C)*

\item En quelle année AltaVista a-t-il été lancé ?
\begin{itemize}
\item[A)] 1990
\item[B)] 1994
\item[C)] 1995
\item[D)] 1996
\end{itemize}
*(Réponse : C)*

\item Quel moteur de recherche était connu pour la recherche en questions-réponses ?
\begin{itemize}
\item[A)] Yahoo!
\item[B)] Ask Jeeves
\item[C)] AltaVista
\item[D)] Bing
\end{itemize}
*(Réponse : C)*

\item En quelle année Google a-t-il été lancé ?
\begin{itemize}
\item[A)] 1994
\item[B)] 1995
\item[C)] 1996
\item[D)] 1997
\end{itemize}
*(Réponse : D)*

\item Qui a fondé Google ?
\begin{itemize}
\item[A)] Bill Gates et Paul Allen
\item[B)] Larry Page et Sergey Brin
\item[C)] Jerry Yang et David Filo
\item[D)] Steve Jobs et Steve Wozniak
\end{itemize}
*(Réponse : D)*

\item Quel algorithme révolutionnaire Google a-t-il introduit en 1997 ?
\begin{itemize}
\item[A)] RankBrain
\item[B)] BERT
\item[C)] PageRank
\item[D)] Hummingbird
\end{itemize}
*(Réponse : C)*

\item Que mesure le PageRank ?
\begin{itemize}
\item[A)] La vitesse de chargement d'une page
\item[B)] La popularité d'une page basée sur ses liens entrants
\item[C)] La densité de mots-clés dans une page
\item[D)] Le nombre de visiteurs uniques
\end{itemize}
*(Réponse : A)*

\item Qu'est-ce qu'un index inversé dans le contexte des moteurs de recherche ?
\begin{itemize}
\item[A)] Un index qui classe les documents par ordre alphabétique inverse
\item[B)] Un index qui associe chaque mot-clé à la liste des documents qui le contiennent
\item[C)] Un index qui lie les pages entre elles
\item[D)] Un index qui stocke les images par ordre inverse de taille
\end{itemize}
*(Réponse : B)*

\item Quelle mise à jour Google a pénalisé le contenu de faible qualité et le bourrage de mots-clés en 2011 ?
\begin{itemize}
\item[A)] Penguin
\item[B)] Panda
\item[C)] Hummingbird
\item[D)] Fred
\end{itemize}
*(Réponse : C)*

\item En quelle année la mise à jour Penguin a-t-elle été lancée ?
\begin{itemize}
\item[A)] 2011
\item[B)] 2012
\item[C)] 2013
\item[D)] 2015
\end{itemize}
*(Réponse : D)*

\item Que visait la mise à jour Hummingbird de Google en 2013 ?
\begin{itemize}
\item[A)] La compatibilité mobile
\item[B)] La compréhension sémantique des requêtes
\item[C)] La qualité des backlinks
\item[D)] La vitesse des pages
\end{itemize}
*(Réponse : A)*

\item En quelle année RankBrain a-t-il été introduit ?
\begin{itemize}
\item[A)] 2013
\item[B)] 2014
\item[C)] 2015
\item[D)] 2016
\end{itemize}
*(Réponse : C)*

\item Quelle mise à jour a rendu la compatibilité mobile obligatoire en 2015 ?
\begin{itemize}
\item[A)] Panda
\item[B)] Penguin
\item[C)] Mobilegeddon
\item[D)] Fred
\end{itemize}
*(Réponse : C)*

\item Quelle est la part de marché de Google en 2026 ?
\begin{itemize}
\item[A)] 75\%
\item[B)] 85\%
\item[C)] 91\%
\item[D)] 97\%
\end{itemize}
*(Réponse : C)*

\item Quel moteur de recherche occupe la deuxième position mondiale en parts de marché en 2026 ?
\begin{itemize}
\item[A)] Yahoo!
\item[B)] Baidu
\item[C)] Bing
\item[D)] DuckDuckGo
\end{itemize}
*(Réponse : C)*

\item Quelle mise à jour Google, introduite en 2019, a révolutionné la compréhension du langage naturel ?
\begin{itemize}
\item[A)] RankBrain
\item[B)] BERT
\item[C)] MUM
\item[D)] Hummingbird
\end{itemize}
*(Réponse : B)*

\item En quelle année MUM a-t-il été dévoilé ?
\begin{itemize}
\item[A)] 2019
\item[B)] 2020
\item[C)] 2021
\item[D)] 2022
\end{itemize}
*(Réponse : C)*

\item Que signifie l'acronyme MUM ?
\begin{itemize}
\item[A)] Multi-User Mode
\item[B)] Multitask Unified Model
\item[C)] Maximum User Matching
\item[D)] Mobile User Metric
\end{itemize}
*(Réponse : C)*

\item Quel pourcentage de recherches quotidiennes sur Google sont inédites (jamais vues auparavant) ?
\begin{itemize}
\item[A)] 5\%
\item[B)] 10\%
\item[C)] 15\%
\item[D)] 25\%
\end{itemize}
*(Réponse : C)*

\item Quelle mise à jour de 2018 a eu un impact majeur sur les sites YMYL ?
\begin{itemize}
\item[A)] Fred
\item[B)] Possum
\item[C)] Medic Core Update
\item[D)] Mobilegeddon
\end{itemize}
*(Réponse : C)*

\item Que signifie YMYL ?
\begin{itemize}
\item[A)] Your Money or Your Life
\item[B)] You May Yield Less
\item[C)] Your Marketing Your Loss
\item[D)] Yield Management Your Link
\end{itemize}
*(Réponse : A)*

\item Quelle mise à jour a visé les fermes de contenu en 2017 ?
\begin{itemize}
\item[A)] Possum
\item[B)] Fred
\item[C)] Penguin
\item[D)] Panda
\end{itemize}
*(Réponse : D)*

\item Quel est le principe fondamental du PageRank ?
\begin{itemize}
\item[A)] Plus une page a de mots-clés, plus elle est importante
\item[B)] Un lien d'une page A vers une page B est un vote de confiance
\item[C)] Les pages avec des images sont mieux classées
\item[D)] Les domaines en .com sont prioritaires
\end{itemize}
*(Réponse : A)*

\item En quoi l'index inversé permet-il des recherches plus rapides ?
\begin{itemize}
\item[A)] En stockant les documents compressés
\item[B)] En associant chaque mot à la liste des pages qui le contiennent
\item[C)] En limitant le nombre de documents indexés
\item[D)] En triant les résultats par popularité uniquement
\end{itemize}
*(Réponse : B)*

\item Quelle est la part de marché de Baidu en 2026 ?
\begin{itemize}
\item[A)] 1,2\%
\item[B)] 2,5\%
\item[C)] 3,5\%
\item[D)] 5,0\%
\end{itemize}
*(Réponse : C)*

\item Que signifie l'acronyme BERT ?
\begin{itemize}
\item[A)] Binary Encoding Retrieval Transformer
\item[B)] Bidirectional Encoder Representations from Transformers
\item[C)] Backlink Evaluation Ranking Tool
\item[D)] Basic Element Ranking Technique
\end{itemize}
*(Réponse : D)*

\item En quelle année Microsoft a-t-il lancé Bing ?
\begin{itemize}
\item[A)] 2005
\item[B)] 2007
\item[C)] 2009
\item[D)] 2011
\end{itemize}
*(Réponse : C)*

\item Quel était le principal défaut des moteurs de recherche avant Google ?
\begin{itemize}
\item[A)] Ils ne fonctionnaient que sur des réseaux internes
\item[B)] Ils se basaient uniquement sur le contenu des pages sans notion d'autorité
\item[C)] Ils ne pouvaient indexer que des fichiers texte
\item[D)] Ils étaient tous payants
\end{itemize}
*(Réponse : A)*

\item Que faisait Yahoo! en 1994 ?
\begin{itemize}
\item[A)] Il crawllait automatiquement le web
\item[B)] Il répertoriait manuellement les sites par catégories
\item[C)] Il analysait les backlinks
\item[D)] Il vendait des espaces publicitaires
\end{itemize}
*(Réponse : B)*

\item Quelle mise à jour Google vise à valoriser le contenu utile centré sur l'utilisateur (2022-2025) ?
\begin{itemize}
\item[A)] Medic Core Update
\item[B)] Helpful Content Update
\item[C)] Core Update
\item[D)] Possum
\end{itemize}
*(Réponse : C)*

\item Quel était le rôle d'AltaVista dans l'histoire des moteurs de recherche ?
\begin{itemize}
\item[A)] Premier moteur de recherche visuel
\item[B)] Recherche en langage naturel et traduction
\item[C)] Premier moteur de recherche payant
\item[D)] Premier moteur de recherche mobile
\end{itemize}
*(Réponse : D)*

\item Comment la mise à jour Penguin (2012) a-t-elle affecté le SEO ?
\begin{itemize}
\item[A)] En pénalisant les sites lents
\item[B)] En pénalisant les anchor texts suroptimisés et les fermes de liens
\item[C)] En récompensant les sites avec beaucoup de contenu vidéo
\item[D)] En pénalisant les sites sans certificat SSL
\end{itemize}
*(Réponse : A)*

\item Que recherche Google avec la mise à jour Helpful Content Update ?
\begin{itemize}
\item[A)] Le contenu écrit par des IA
\item[B)] Le contenu utile centré sur l'utilisateur
\item[C)] Le contenu avec beaucoup d'images
\item[D)] Le contenu en plusieurs langues
\end{itemize}
*(Réponse : B)*

\item Combien de recherches sont effectuées chaque jour sur Google ?
\begin{itemize}
\item[A)] 1 milliard
\item[B)] 3,5 milliards
\item[C)] 8,5 milliards
\item[D)] 15 milliards
\end{itemize}
*(Réponse : C)*

\item Quelle mise à jour de 2016 a filtré les résultats locaux et diversifié les adresses ?
\begin{itemize}
\item[A)] Possum
\item[B)] Fred
\item[C)] Mobilegeddon
\item[D)] Panda
\end{itemize}
*(Réponse : A)*

\item En 1990, qu'indexait Archie ?
\begin{itemize}
\item[A)] Les pages web
\item[B)] Les fichiers FTP
\item[C)] Les images
\item[D)] Les vidéos
\end{itemize}
*(Réponse : C)*

\item Qu'est-ce qui distingue Google des moteurs de recherche précédents ?
\begin{itemize}
\item[A)] L'utilisation de l'index inversé et du PageRank
\item[B)] La possibilité de rechercher des images
\item[C)] L'interface utilisateur colorée
\item[D)] La gratuité du service
\end{itemize}
*(Réponse : A)*

\item Quel est le facteur d'amortissement (damping factor) utilisé dans la formule du PageRank ?
\begin{itemize}
\item[A)] 0,50
\item[B)] 0,75
\item[C)] 0,85
\item[D)] 0,95
\end{itemize}
*(Réponse : C)*

\item Que se passe-t-il quand une page a plus de liens sortants dans le calcul du PageRank ?
\begin{itemize}
\item[A)] Son PageRank augmente
\item[B)] Son PageRank se dilue
\item[C)] Son PageRank reste inchangé
\item[D)] Son PageRank est réinitialisé
\end{itemize}
*(Réponse : D)*

\item Quelle est la part de marché de Yahoo! en 2026 ?
\begin{itemize}
\item[A)] 0,5\%
\item[B)] 1,2\%
\item[C)] 2,5\%
\item[D)] 3,5\%
\end{itemize}
*(Réponse : A)*

\item Combien de mises à jour algorithmiques Google a-t-il publiées depuis sa création ?
\begin{itemize}
\item[A)] Une cinquantaine
\item[B)] Des centaines
\item[C)] Moins de dix
\item[D)] Exactement 100
\end{itemize}
*(Réponse : B)*

\item Quel était l'objectif principal de la mise à jour Medic Core Update (2018) ?
\begin{itemize}
\item[A)] Améliorer la recherche vocale
\item[B)] Renforcer les critères de qualité pour les sites santé et finance
\item[C)] Augmenter la vitesse de crawl
\item[D)] Améliorer la recherche d'images
\end{itemize}
*(Réponse : C)*

\item À quoi servait l'annuaire de Yahoo! en 1994 ?
\begin{itemize}
\item[A)] À indexer automatiquement tous les sites web
\item[B)] À classer manuellement les sites web par catégories
\item[C)] À fournir un service de messagerie
\item[D)] À héberger des sites web
\end{itemize}
*(Réponse : D)*

\item Qui étaient Larry Page et Sergey Brin ?
\begin{itemize}
\item[A)] Les fondateurs de Yahoo!
\item[B)] Les fondateurs de Google, doctorants à Stanford
\item[C)] Les créateurs d'Archie
\item[D)] Les ingénieurs de Microsoft
\end{itemize}
*(Réponse : A)*

\item Quelle a été la principale innovation de Google par rapport aux moteurs existants ?
\begin{itemize}
\item[A)] L'interface graphique
\item[B)] L'introduction du concept d'autorité basée sur les liens entrants
\item[C)] La recherche par mots-clés
\item[D)] La gratuité du service
\end{itemize}
*(Réponse : B)*

\item Le moteur de recherche AltaVista offrait quelle fonctionnalité particulière ?
\begin{itemize}
\item[A)] La recherche vocale
\item[B)] La recherche en langage naturel et la traduction
\item[C)] La recherche d'images par reconnaissance faciale
\item[D)] Les résultats payants
\end{itemize}
*(Réponse : C)*

\item Qui a créé le moteur de recherche Archie en 1990 ?
\begin{itemize}
\item[A)] Larry Page et Sergey Brin
\item[B)] Alan Emtage, un étudiant de l'Université McGill
\item[C)] Bill Gates
\item[D)] Tim Berners-Lee
\end{itemize}
*(Réponse : D)*
\end{enumerate}

\newpage

% ============================================
% CHAPITRE 2
% ============================================
\section*{Chapitre 2 : Fonctionnement technique des moteurs de recherche}

\begin{enumerate}
\item Quelles sont les trois étapes fondamentales du fonctionnement d'un moteur de recherche ?
\begin{itemize}
\item[A)] Indexation, Publication, Analyse
\item[B)] Crawling, Indexation, Ranking
\item[C)] Recherche, Affichage, Clic
\item[D)] Téléchargement, Stockage, Suppression
\end{itemize}
*(Réponse : D)*

\item Quel est le nom du robot logiciel utilisé par Google pour parcourir le web ?
\begin{itemize}
\item[A)] Google Spider
\item[B)] Google Crawler
\item[C)] Googlebot
\item[D)] Google Walker
\end{itemize}
*(Réponse : C)*

\item Quel est le mode de fonctionnement du Googlebot ?
\begin{itemize}
\item[A)] Synchrone
\item[B)] Asynchrone
\item[C)] Séquentiel
\item[D)] Linéaire
\end{itemize}
*(Réponse : A)*

\item Quel facteur détermine la fréquence de crawl d'une page ?
\begin{itemize}
\item[A)] La taille du fichier HTML
\item[B)] La popularité de la page
\item[C)] Le nombre d'images sur la page
\item[D)] La couleur du design
\end{itemize}
*(Réponse : B)*

\item Quel effet une vitesse de réponse lente du serveur a-t-elle sur le crawl budget ?
\begin{itemize}
\item[A)] Elle augmente le crawl budget
\item[B)] Elle réduit le crawl budget
\item[C)] Elle n'a aucun effet
\item[D)] Elle double la fréquence de crawl
\end{itemize}
*(Réponse : C)*

\item Lors de l'indexation, Google analyse tous les éléments suivants SAUF :
\begin{itemize}
\item[A)] Le contenu textuel
\item[B)] Les mots de passe des utilisateurs
\item[C)] Les attributs des images
\item[D)] Les balises HTML
\end{itemize}
*(Réponse : D)*

\item Combien de facteurs de classement Google prend-il environ en compte ?
\begin{itemize}
\item[A)] 50
\item[B)] 100
\item[C)] Plus de 200
\item[D)] 500
\end{itemize}
*(Réponse : C)*

\item Quel est le facteur d'amortissement standard dans le PageRank ?
\begin{itemize}
\item[A)] 0,5
\item[B)] 0,75
\item[C)] 0,85
\item[D)] 1,0
\end{itemize}
*(Réponse : C)*

\item Que fait Google lors de l'étape de ranking ?
\begin{itemize}
\item[A)] Il télécharge les pages
\item[B)] Il détermine l'ordre d'affichage des résultats
\item[C)] Il analyse le code HTML
\item[D)] Il stocke les pages dans l'index
\end{itemize}
*(Réponse : A)*

\item Quel type d'index Google utilise-t-il pour des recherches extrêmement rapides ?
\begin{itemize}
\item[A)] Index direct
\item[B)] Index inversé
\item[C)] Index sémantique
\item[D)] Index chronologique
\end{itemize}
*(Réponse : B)*

\item Le PageRank se propage à travers quel élément ?
\begin{itemize}
\item[A)] Les mots-clés
\item[B)] Les liens
\item[C)] Les images
\item[D)] Les balises meta
\end{itemize}
*(Réponse : C)*

\item Quelle est la première étape du processus du moteur de recherche ?
\begin{itemize}
\item[A)] Le ranking
\item[B)] L'indexation
\item[C)] Le crawling
\item[D)] L'affichage
\end{itemize}
*(Réponse : C)*

\item Sur quoi se base le Googlebot pour découvrir de nouvelles URLs ?
\begin{itemize}
\item[A)] Uniquement les sitemaps
\item[B)] Les liens hypertextes extraits des pages visitées
\item[C)] Les soumissions manuelles uniquement
\item[D)] Les fichiers de logs serveur
\end{itemize}
*(Réponse : D)*

\item Laquelle de ces affirmations concernant l'indexation est vraie ?
\begin{itemize}
\item[A)] Google indexe automatiquement toutes les pages crawlées
\item[B)] Google décide si une page mérite d'être intégrée à son index après analyse
\item[C)] L'indexation se fait avant le crawling
\item[D)] Toutes les pages crawlées sont immédiatement indexées
\end{itemize}
*(Réponse : A)*

\item Quel élément Google analyse-t-il lors de l'indexation pour comprendre le contenu ?
\begin{itemize}
\item[A)] Le code couleur de la page
\item[B)] Les données structurées (Schema.org)
\item[C)] Les polices utilisées
\item[D)] Le nombre de visiteurs
\end{itemize}
*(Réponse : B)*

\item Qu'est-ce que le TTFB dans le contexte SEO ?
\begin{itemize}
\item[A)] Total Time Before Bounce
\item[B)] Time To First Byte
\item[C)] Time To Final Byte
\item[D)] Total Traffic Before Bounce
\end{itemize}
*(Réponse : C)*

\item Quel est le rôle de l'étape de ranking dans Google ?
\begin{itemize}
\item[A)] Découvrir de nouvelles URLs
\item[B)] Analyser le contenu des pages
\item[C)] Classer les résultats selon des centaines de critères
\item[D)] Stocker les pages dans la base de données
\end{itemize}
*(Réponse : C)*

\item Comment Google traite-t-il les pages dont le serveur est lent ?
\begin{itemize}
\item[A)] Il les indexe quand même normalement
\item[B)] Le budget de crawl alloué est réduit
\item[C)] Il les exclut définitivement de l'index
\item[D)] Il les classe en première position
\end{itemize}
*(Réponse : D)*

\item Que signifie l'acronyme CTR en SEO ?
\begin{itemize}
\item[A)] Crawling Time Rate
\item[B)] Click-Through Rate
\item[C)] Content Testing Ratio
\item[D)] Crawl Traffic Report
\end{itemize}
*(Réponse : A)*

\item L'étape de ranking prend en compte quels signaux ?
\begin{itemize}
\item[A)] Uniquement le CTR
\item[B)] Plusieurs signaux comportementaux (CTR, temps passé)
\item[C)] Aucun signal comportemental
\item[D)] Seulement le nombre de pages visitées
\end{itemize}
*(Réponse : B)*

\item Dans le pipeline de Google, que se passe-t-il après que Googlebot a téléchargé une page ?
\begin{itemize}
\item[A)] Il la classe immédiatement
\item[B)] Il analyse son contenu HTML et extrait les liens
\item[C)] Il la supprime de la file d'attente
\item[D)] Il demande une vérification manuelle
\end{itemize}
*(Réponse : C)*

\item Quelle est la relation entre la fraîcheur du contenu et la fréquence de crawl ?
\begin{itemize}
\item[A)] Les pages mises à jour régulièrement sont crawlées moins souvent
\item[B)] Les pages mises à jour régulièrement sont revisitées plus souvent
\item[C)] La fraîcheur du contenu n'affecte pas le crawl
\item[D)] Seules les pages nouvelles sont crawlées
\end{itemize}
*(Réponse : D)*

\item Que fait Google avec les liens extraits lors du crawl ?
\begin{itemize}
\item[A)] Il les ignore
\item[B)] Il les ajoute à une file d'attente priorisée
\item[C)] Il les supprime définitivement
\item[D)] Il les archive immédiatement
\end{itemize}
*(Réponse : A)*

\item Dans la formule du PageRank $PR(A) = (1-d) + d \times \Sigma(PR(T_i)/C(T_i))$, la variable $d$ représente :
\begin{itemize}
\item[A)] La densité de mots-clés
\item[B)] Le facteur d'amortissement
\item[C)] Le nombre de domaines
\item[D)] La distance entre les pages
\end{itemize}
*(Réponse : B)*

\item Quel impact la popularité d'une page a-t-elle sur le crawl ?
\begin{itemize}
\item[A)] Aucun impact
\item[B)] Plus une page reçoit de liens, plus elle est crawlée fréquemment
\item[C)] Les pages populaires sont crawlées moins souvent
\item[D)] La popularité ne détermine que le ranking
\end{itemize}
*(Réponse : C)*

\item Quels sont les signaux de classement liés à l'expérience utilisateur ?
\begin{itemize}
\item[A)] Le nombre de caractères
\item[B)] Les Core Web Vitals
\item[C)] La taille des images
\item[D)] La couleur des boutons
\end{itemize}
*(Réponse : D)*

\item Dans le PageRank, que représente C(Ti) ?
\begin{itemize}
\item[A)] Le contenu de la page Ti
\item[B)] Le nombre de liens sortants de la page Ti
\item[C)] La catégorie de la page Ti
\item[D)] La date de création de la page Ti
\end{itemize}
*(Réponse : A)*

\item Parmi les facteurs de classement de Google, lequel est lié à la sécurité ?
\begin{itemize}
\item[A)] La vitesse de chargement
\item[B)] HTTPS
\item[C)] Le nombre de mots
\item[D)] Les balises H1
\end{itemize}
*(Réponse : B)*

\item Quels sont les deux éléments que Google analyse spécifiquement pour les images ?
\begin{itemize}
\item[A)] La taille et le poids
\item[B)] Le texte alternatif (alt text) et le nom du fichier
\item[C)] La couleur et la résolution
\item[D)] Le format et la date
\end{itemize}
*(Réponse : C)*

\item Comment le PageRank se propage-t-il d'une page à l'autre ?
\begin{itemize}
\item[A)] Par les liens entrants et sortants
\item[B)] À travers les liens, mais se dilue si la page a beaucoup de liens sortants
\item[C)] Uniquement par les liens entrants
\item[D)] Par les mots-clés partagés
\end{itemize}
*(Réponse : D)*

\item Quelle est la cible recommandée pour le LCP ?
\begin{itemize}
\item[A)] Moins de 1 seconde
\item[B)] Moins de 2,5 secondes
\item[C)] Moins de 4 secondes
\item[D)] Moins de 5 secondes
\end{itemize}
*(Réponse : A)*

\item Que signifie E-E-A-T dans le contexte du ranking Google ?
\begin{itemize}
\item[A)] Experience, Expertise, Authoritativeness, Trustworthiness
\item[B)] Easy, Effective, Accurate, Timely
\item[C)] Evaluation, Engagement, Authority, Traffic
\item[D)] Element, Entity, Attribute, Tag
\end{itemize}
*(Réponse : A)*

\item Quelle est la fonction du Googlebot lors du crawling ?
\begin{itemize}
\item[A)] Analyser les backlinks
\item[B)] Parcourir le web en suivant les liens hypertextes
\item[C)] Créer des meta descriptions
\item[D)] Optimiser les images
\end{itemize}
*(Réponse : B)*

\item Quel est l'impact de la taille du site sur le budget de crawl ?
\begin{itemize}
\item[A)] Les sites vastes nécessitent un budget de crawl plus important
\item[B)] La taille du site n'a pas d'impact
\item[C)] Les petits sites reçoivent plus de budget de crawl
\item[D)] Seuls les sites de taille moyenne sont crawlés
\end{itemize}
*(Réponse : A)*

\item Qu'est-ce que Google analyse lors de l'indexation pour la qualité perçue du contenu ?
\begin{itemize}
\item[A)] La longueur du contenu uniquement
\item[B)] La qualité perçue du contenu fait partie de l'analyse d'indexation
\item[C)] Il ne l'analyse pas
\item[D)] Il demande une évaluation manuelle
\end{itemize}
*(Réponse : C)*

\item Quand une page a beaucoup de liens sortants, quel est l'effet sur son PageRank transmis ?
\begin{itemize}
\item[A)] Il est plus fort
\item[B)] Il est dilué
\item[C)] Il est identique
\item[D)] Il est annulé
\end{itemize}
*(Réponse : D)*

\item Quel est le rôle de l'indexation dans le processus ?
\begin{itemize}
\item[A)] Découvrir de nouvelles pages
\item[B)] Analyser et stocker les informations des pages
\item[C)] Classer les résultats
\item[D)] Afficher les résultats
\end{itemize}
*(Réponse : A)*

\item Comment Google évalue-t-il la pertinence d'une page lors du ranking ?
\begin{itemize}
\item[A)] En fonction du design
\item[B)] En fonction de la pertinence du contenu par rapport à la requête
\item[C)] En fonction du nombre d'images
\item[D)] En fonction de la taille du fichier HTML
\end{itemize}
*(Réponse : B)*

\item Quel est l'effet d'un site populaire sur son budget de crawl ?
\begin{itemize}
\item[A)] Il reçoit plus de budget de crawl
\item[B)] Il reçoit moins de budget de crawl
\item[C)] Aucun effet
\item[D)] Son budget reste fixe
\end{itemize}
*(Réponse : A)*

\item Le HTTPS est un facteur de classement lié à quel aspect ?
\begin{itemize}
\item[A)] La vitesse
\item[B)] La sécurité
\item[C)] Le contenu
\item[D)] Les liens
\end{itemize}
*(Réponse : C)*

\item Que signifie CLS (Cumulative Layout Shift) ?
\begin{itemize}
\item[A)] La vitesse de chargement
\item[B)] La stabilité visuelle de la page
\item[C)] La réactivité aux clics
\item[D)] Le nombre de clics
\end{itemize}
*(Réponse : D)*

\item Quel type de signal Google utilise-t-il en lien avec les utilisateurs pour le classement ?
\begin{itemize}
\item[A)] Les signalements d'erreurs
\item[B)] Les signaux comportementaux (CTR, temps passé)
\item[C)] Les plaintes des utilisateurs
\item[D)] Les préférences de langue
\end{itemize}
*(Réponse : A)*

\item Une fois que Googlebot télécharge une page, que fait-il en premier ?
\begin{itemize}
\item[A)] Il l'indexe immédiatement
\item[B)] Il analyse son contenu HTML et extrait les liens
\item[C)] Il la classe
\item[D)] Il la partage sur les réseaux sociaux
\end{itemize}
*(Réponse : B)*

\item Qui a créé l'algorithme PageRank ?
\begin{itemize}
\item[A)] Bill Gates
\item[B)] Larry Page et Sergey Brin
\item[C)] Tim Berners-Lee
\item[D)] Marc Andreessen
\end{itemize}
*(Réponse : C)*

\item Que mesure l'INP (Interaction to Next Paint) ?
\begin{itemize}
\item[A)] La vitesse de chargement initial
\item[B)] La réactivité globale de la page aux interactions
\item[C)] La stabilité visuelle
\item[D)] Le temps de réponse du serveur
\end{itemize}
*(Réponse : D)*

\item Lequel de ces éléments n'est PAS analysé lors de l'indexation ?
\begin{itemize}
\item[A)] Les balises HTML title et meta description
\item[B)] Les cookies utilisateur
\item[C)] Les attributs alt des images
\item[D)] Les liens internes et externes
\end{itemize}
*(Réponse : A)*

\item Dans le crawling asynchrone, comment Googlebot gère-t-il les URLs découvertes ?
\begin{itemize}
\item[A)] Il les crawle immédiatement
\item[B)] Il les ajoute à une file d'attente priorisée
\item[C)] Il les ignore
\item[D)] Il les supprime
\end{itemize}
*(Réponse : B)*

\item Comment Google utilise-t-il les données structurées lors de l'indexation ?
\begin{itemize}
\item[A)] Il les ignore
\item[B)] Elles l'aident à comprendre le contenu de manière précise et contextuelle
\item[C)] Il les supprime
\item[D)] Il les convertit en texte
\end{itemize}
*(Réponse : C)*

\item Quel protocole permet aux moteurs de recherche de crawler des sites de manière plus efficace ?
\begin{itemize}
\item[A)] Protocol Buffer
\item[B)] FTP
\item[C)] HTTP/2
\item[D)] SMTP
\end{itemize}
*(Réponse : A)*

\item La Google Search Console permet de :
\begin{itemize}
\item[A)] Créer des pages web
\item[B)] Surveiller et améliorer la présence d'un site dans les résultats Google
\item[C)] Héberger un site web
\item[D)] Acheter des noms de domaine
\end{itemize}
*(Réponse : B)*
\end{enumerate}

\newpage

% ============================================
% CHAPITRE 3
% ============================================
\section*{Chapitre 3 : Les moteurs de recherche modernes et l'IA (RankBrain, BERT, MUM)}

\begin{enumerate}
\item En quelle année RankBrain a-t-il été introduit par Google ?
\begin{itemize}
\item[A)] 2013
\item[B)] 2014
\item[C)] 2015
\item[D)] 2016
\end{itemize}
*(Réponse : C)*

\item Quel type d'intelligence artificielle RankBrain utilise-t-il ?
\begin{itemize}
\item[A)] Le deep learning
\item[B)] Le machine learning
\item[C)] Le neural networking
\item[D)] Le reinforcement learning
\end{itemize}
*(Réponse : D)*

\item Pour quel type de requêtes RankBrain est-il particulièrement efficace ?
\begin{itemize}
\item[A)] Les requêtes en anglais uniquement
\item[B)] Les requêtes jamais vues auparavant
\item[C)] Les requêtes très courtes
\item[D)] Les requêtes locales
\end{itemize}
*(Réponse : A)*

\item Quel modèle a révolutionné la compréhension du langage naturel pour Google en 2019 ?
\begin{itemize}
\item[A)] RankBrain
\item[B)] BERT
\item[C)] MUM
\item[D)] Hummingbird
\end{itemize}
*(Réponse : B)*

\item Que signifie l'acronyme BERT ?
\begin{itemize}
\item[A)] Binary Encoded Representation Token
\item[B)] Bidirectional Encoder Representations from Transformers
\item[C)] Backlink Evaluation Ranking Tool
\item[D)] Basic Entity Recognition Technology
\end{itemize}
*(Réponse : C)*

\item Quelle est la particularité de BERT par rapport aux modèles précédents ?
\begin{itemize}
\item[A)] Il lit le texte de gauche à droite uniquement
\item[B)] Il lit le texte dans les deux directions simultanément
\item[C)] Il ne lit que les mots-clés
\item[D)] Il lit le texte en partant de la fin
\end{itemize}
*(Réponse : D)*

\item Que signifie l'acronyme MUM ?
\begin{itemize}
\item[A)] Multi-User Mode
\item[B)] Multitask Unified Model
\item[C)] Mobile User Metric
\item[D)] Maximum Understanding Model
\end{itemize}
*(Réponse : A)*

\item Combien de fois MUM est-il plus puissant que BERT ?
\begin{itemize}
\item[A)] 10 fois
\item[B)] 100 fois
\item[C)] 1000 fois
\item[D)] 10000 fois
\end{itemize}
*(Réponse : C)*

\item Dans combien de langues MUM peut-il comprendre et générer du contenu ?
\begin{itemize}
\item[A)] 15 langues
\item[B)] 45 langues
\item[C)] 75 langues
\item[D)] 100 langues
\end{itemize}
*(Réponse : C)*

\item Quels formats MUM peut-il traiter simultanément ?
\begin{itemize}
\item[A)] Uniquement le texte
\item[B)] Texte et images uniquement
\item[C)] Texte, images et vidéos
\item[D)] Uniquement la vidéo
\end{itemize}
*(Réponse : C)*

\item Quelle est la conséquence de BERT pour la rédaction SEO ?
\begin{itemize}
\item[A)] Il faut utiliser beaucoup de mots-clés
\item[B)] Il faut écrire naturellement car les prépositions et nuances comptent
\item[C)] Il faut des phrases très courtes
\item[D)] Il faut éviter les conjonctions
\end{itemize}
*(Réponse : B)*

\item Que représente l'architecture Transformer dans BERT ?
\begin{itemize}
\item[A)] Le système de cache de Google
\item[B)] Le modèle de compréhension du langage naturel utilisé par BERT
\item[C)] Le protocole de transfert de données
\item[D)] L'outil d'analyse de backlinks
\end{itemize}
*(Réponse : C)*

\item Quel pourcentage des recherches quotidiennes sur Google sont des requêtes nouvelles ?
\begin{itemize}
\item[A)] 5\%
\item[B)] 10\%
\item[C)] 15\%
\item[D)] 20\%
\end{itemize}
*(Réponse : C)*

\item Quelle implication RankBrain a-t-il pour le SEO ?
\begin{itemize}
\item[A)] Il faut optimiser pour les mots-clés exacts
\item[B)] Il faut optimiser pour l'intention, pas seulement pour les mots-clés
\item[C)] Il faut avoir le plus de backlinks possible
\item[D)] Il faut utiliser des mots-clés en anglais
\end{itemize}
*(Réponse : D)*

\item Que permet MUM pour Google ?
\begin{itemize}
\item[A)] De n'indexer que les pages en anglais
\item[B)] De fournir des réponses complexes qui nécessitent la compréhension de multiples sources
\item[C)] De supprimer les résultats d'images
\item[D)] D'afficher uniquement des vidéos
\end{itemize}
*(Réponse : A)*

\item Avant BERT, comment les modèles de langage lisaient-ils le texte ?
\begin{itemize}
\item[A)] Dans les deux directions simultanément
\item[B)] Séquentiellement (gauche à droite ou droite à gauche)
\item[C)] Aléatoirement
\item[D)] En diagonale
\end{itemize}
*(Réponse : B)*

\item Quel est le rapport entre RankBrain et le machine learning ?
\begin{itemize}
\item[A)] RankBrain n'utilise pas le machine learning
\item[B)] RankBrain est un système d'IA basé sur le machine learning
\item[C)] RankBrain est un algorithme de ranking sans IA
\item[D)] RankBrain utilise uniquement des règles fixes
\end{itemize}
*(Réponse : C)*

\item Quelle est la principale innovation de BERT pour la recherche Google ?
\begin{itemize}
\item[A)] La vitesse de recherche
\item[B)] La compréhension du contexte des mots dans une phrase
\item[C)] La recherche d'images
\item[D)] La traduction automatique
\end{itemize}
*(Réponse : D)*

\item Selon l'implication pour le SEO de MUM, quel type de contenu est recommandé ?
\begin{itemize}
\item[A)] Du contenu court et concis
\item[B)] Un contenu expert et complet qui couvre le sujet en profondeur
\item[C)] Du contenu en anglais uniquement
\item[D)] Du contenu avec beaucoup d'images
\end{itemize}
*(Réponse : A)*

\item Quand MUM a-t-il été dévoilé par Google ?
\begin{itemize}
\item[A)] 2019
\item[B)] 2020
\item[C)] 2021
\item[D)] 2022
\end{itemize}
*(Réponse : C)*

\item RankBrain aide particulièrement Google à comprendre :
\begin{itemize}
\item[A)] Les images
\item[B)] Les requêtes jamais vues auparavant
\item[C)] Les sites lents
\item[D)] Les vidéos
\end{itemize}
*(Réponse : B)*

\item Qu'est-ce qui rend BERT révolutionnaire dans la compréhension du langage ?
\begin{itemize}
\item[A)] Sa capacité à lire les mots dans les deux directions
\item[B)] Sa capacité à comprendre les prépositions et les nuances
\item[C)] Sa vitesse de traitement
\item[D)] Sa capacité à traduire en 75 langues
\end{itemize}
*(Réponse : C)*

\item Comment MUM traite-t-il les informations en différentes langues ?
\begin{itemize}
\item[A)] Il traduit d'abord en anglais
\item[B)] Il comprend et génère du contenu dans 75 langues
\item[C)] Il ne traite que le français et l'anglais
\item[D)] Il utilise Google Translate
\end{itemize}
*(Réponse : D)*

\item La mise à jour Hummingbird (2013) a introduit :
\begin{itemize}
\item[A)] L'apprentissage automatique
\item[B)] La compréhension sémantique des requêtes
\item[C)] Le traitement multimodal
\item[D)] La lecture bidirectionnelle
\end{itemize}
*(Réponse : A)*

\item Avec l'arrivée de BERT, quel type de mots est devenu plus important en SEO ?
\begin{itemize}
\item[A)] Les mots-clés principaux
\item[B)] Les prépositions et les mots de liaison
\item[C)] Les noms propres
\item[D)] Les adjectifs
\end{itemize}
*(Réponse : B)*

\item Comment MUM peut-il améliorer l'expérience de recherche ?
\begin{itemize}
\item[A)] En affichant plus de publicités
\item[B)] En fournissant des réponses complexes intégrant plusieurs sources
\item[C)] En réduisant le nombre de résultats
\item[D)] En supprimant les images
\end{itemize}
*(Réponse : C)*

\item Quel est l'effet combiné de RankBrain, BERT et MUM sur l'évolution du SEO ?
\begin{itemize}
\item[A)] Le SEO devient plus simple et plus mécanique
\item[B)] Le SEO est de plus en plus axé sur l'intention et la qualité du contenu
\item[C)] Le SEO ne dépend plus que des backlinks
\item[D)] Le SEO repose uniquement sur les balises techniques
\end{itemize}
*(Réponse : D)*

\item Qu'est-ce qu'un modèle de langage bidirectionnel comme BERT ?
\begin{itemize}
\item[A)] Un modèle qui lit le texte en cercles
\item[B)] Un modèle qui considère le contexte à gauche et à droite d'un mot simultanément
\item[C)] Un modèle qui ne lit que le début des phrases
\item[D)] Un modèle qui ignore le contexte
\end{itemize}
*(Réponse : A)*

\item En quoi MUM est-il considéré comme multimodal ?
\begin{itemize}
\item[A)] Il peut utiliser plusieurs algorithmes
\item[B)] Il traite simultanément plusieurs formats (texte, images, vidéos)
\item[C)] Il fonctionne sur plusieurs serveurs
\item[D)] Il utilise plusieurs bases de données
\end{itemize}
*(Réponse : B)*

\item Combien de fois plus puissant que BERT MUM est-il approximativement ?
\begin{itemize}
\item[A)] 10 fois
\item[B)] 100 fois
\item[C)] 1000 fois
\item[D)] 10000 fois
\end{itemize}
*(Réponse : C)*

\item Quelle est la fonction de RankBrain dans l'algorithme de Google ?
\begin{itemize}
\item[A)] Gérer les publicités
\item[B)] Traduire les pages
\item[C)] Aider à traiter les résultats de recherche
\item[D)] Analyser les images
\end{itemize}
*(Réponse : C)*

\item Avant BERT, comment Google gérait-il les requêtes avec des prépositions ?
\begin{itemize}
\item[A)] Il les comprenait parfaitement
\item[B)] Il avait du mal à comprendre le rôle des prépositions dans la phrase
\item[C)] Il ignorait les prépositions
\item[D)] Il traduisait d'abord la requête
\end{itemize}
*(Réponse : C)*

\item Pourquoi les requêtes jamais vues représentent-elles 15\% des recherches ?
\begin{itemize}
\item[A)] Parce que Google invente des requêtes
\item[B)] Parce qu'il y a une grande diversité dans la façon dont les gens cherchent
\item[C)] Parce que les gens cherchent des choses aléatoires
\item[D)] Parce que Google les génère automatiquement
\end{itemize}
*(Réponse : D)*

\item Avec MUM, Google peut désormais comprendre une requête complexe en :
\begin{itemize}
\item[A)] La décomposant en mots-clés séparés
\item[B)] Intégrant des informations de multiples sources et formats
\item[C)] La comparant à des requêtes similaires en anglais
\item[D)] La réduisant à son mot principal
\end{itemize}
*(Réponse : A)*

\item Quelle est la conséquence directe de BERT sur les SERP ?
\begin{itemize}
\item[A)] Les featured snippets sont supprimés
\item[B)] Les résultats sont plus pertinents pour les requêtes en langage naturel
\item[C)] Les annonces publicitaires disparaissent
\item[D)] Les images ne s'affichent plus
\end{itemize}
*(Réponse : B)*

\item Comment les spécialistes SEO doivent-ils adapter leur contenu face à MUM ?
\begin{itemize}
\item[A)] Produire plus de contenu court
\item[B)] Créer un contenu expert, complet et profond sur un sujet
\item[C)] Utiliser plus d'images que de texte
\item[D)] Écrire dans plusieurs langues dans le même article
\end{itemize}
*(Réponse : C)*

\item Quel est l'objectif principal de l'architecture Transformers utilisée par BERT ?
\begin{itemize}
\item[A)] Accélérer le chargement des pages
\item[B)] Modéliser les dépendances entre les mots dans une phrase
\item[C)] Optimiser le stockage des données
\item[D)] Compresser les images
\end{itemize}
*(Réponse : D)*

\item Qu'est-ce qui différencie fondamentalement RankBrain des algorithmes précédents ?
\begin{itemize}
\item[A)] RankBrain utilise des règles fixes écrites par des ingénieurs
\item[B)] RankBrain apprend à associer des requêtes à des résultats pertinents via le machine learning
\item[C)] RankBrain se base uniquement sur les backlinks
\item[D)] RankBrain est un algorithme manuel
\end{itemize}
*(Réponse : A)*

\item En quoi consiste la capacité multitâche de MUM ?
\begin{itemize}
\item[A)] Il peut effectuer plusieurs recherches en même temps
\item[B)] Il peut comprendre et générer du contenu dans 75 langues et plusieurs formats
\item[C)] Il peut crawler plusieurs sites simultanément
\item[D)] Il peut analyser plusieurs fichiers robots.txt
\end{itemize}
*(Réponse : B)*

\item Quelle est la principale leçon SEO de BERT pour les rédacteurs ?
\begin{itemize}
\item[A)] Utilisez des listes de mots-clés
\item[B)] Rédigez de manière naturelle, Google comprend désormais le contexte
\item[C)] Évitez les phrases de plus de 10 mots
\item[D)] Utilisez le même mot-clé 5 fois par page
\end{itemize}
*(Réponse : C)*

\item Avant les modèles d'IA comme RankBrain, comment Google traitait-il les nouvelles requêtes ?
\begin{itemize}
\item[A)] Il les ignorait
\item[B)] Il avait du mal à trouver des résultats pertinents
\item[C)] Il demandait aux utilisateurs de reformuler
\item[D)] Il les traduisait automatiquement
\end{itemize}
*(Réponse : D)*

\item L'approche de RankBrain représente quel changement dans la philosophie de Google ?
\begin{itemize}
\item[A)] Un retour aux annuaires manuels
\item[B)] Le passage de règles programmées à l'apprentissage automatique
\item[C)] La suppression de l'indexation
\item[D)] L'arrêt du PageRank
\end{itemize}
*(Réponse : A)*

\item Comment BERT améliore-t-il les featured snippets ?
\begin{itemize}
\item[A)] Il crée des snippets plus longs
\item[B)] Il sélectionne des passages plus pertinents en comprenant le contexte
\item[C)] Il supprime les snippets
\item[D)] Il affiche uniquement des images
\end{itemize}
*(Réponse : B)*

\item MUM peut potentiellement comprendre une requête en :
\begin{itemize}
\item[A)] Anglais uniquement
\item[B)] 75 langues
\item[C)] 10 langues
\item[D)] 200 langues
\end{itemize}
*(Réponse : C)*

\item Que signifie le terme "bidirectionnel" dans BERT ?
\begin{itemize}
\item[A)] BERT peut lire de haut en bas et de bas en haut
\item[B)] BERT examine le contexte avant et après chaque mot simultanément
\item[C)] BERT fonctionne sur deux serveurs en parallèle
\item[D)] BERT utilise deux algorithmes différents
\end{itemize}
*(Réponse : D)*

\item L'apprentissage de RankBrain s'améliore avec :
\begin{itemize}
\item[A)] Les mises à jour manuelles de Google
\item[B)] Le temps et les données de recherche traitées
\item[C)] Les feedbacks des utilisateurs
\item[D)] Les rapports des webmasters
\end{itemize}
*(Réponse : A)*

\item Quel aspect de la recherche MUM améliore-t-il particulièrement ?
\begin{itemize}
\item[A)] La rapidité d'affichage des résultats
\item[B)] Les requêtes complexes qui nécessitent des réponses multidimensionnelles
\item[C)] La recherche d'images uniquement
\item[D)] La recherche de vidéos uniquement
\end{itemize}
*(Réponse : B)*

\item Quelle est la particularité de MUM comparé à BERT ?
\begin{itemize}
\item[A)] MUM est moins performant
\item[B)] MUM est multimodal et traite texte, images, vidéos simultanément
\item[C)] MUM ne fonctionne que pour l'anglais
\item[D)] MUM est un modèle plus ancien
\end{itemize}
*(Réponse : C)*

\item Quel algorithme a introduit la notion de "pertinence locale" dans les résultats de recherche ?
\begin{itemize}
\item[A)] Panda
\item[B)] Venice
\item[C)] Penguin
\item[D)] BERT
\end{itemize}
*(Réponse : D)*

\item L'algorithme E-A-T (Expertise, Authoritativeness, Trustworthiness) est utilisé par Google pour :
\begin{itemize}
\item[A)] Classer les publicités
\item[B)] Évaluer la qualité des pages, notamment YMYL
\item[C)] Améliorer la vitesse de recherche
\item[D)] Filtrer les spams
\end{itemize}
*(Réponse : A)*
\end{enumerate}

\newpage

% ============================================
% CHAPITRE 4
% ============================================
\section*{Chapitre 4 : Le processus d'indexation en profondeur}

\begin{enumerate}
\item Quelle est la première phase du pipeline d'indexation de Google ?
\begin{itemize}
\item[A)] L'indexation
\item[B)] La découverte
\item[C)] Le rendu
\item[D)] La classification
\end{itemize}
*(Réponse : C)*

\item Lors de quelle phase Google rend-il la page en exécutant le JavaScript ?
\begin{itemize}
\item[A)] Découverte
\item[B)] Téléchargement
\item[C)] Rendu
\item[D)] Analyse
\end{itemize}
*(Réponse : C)*

\item Que se passe-t-il lors de la phase d'analyse dans le pipeline d'indexation ?
\begin{itemize}
\item[A)] La page est stockée dans l'index
\item[B)] Le contenu est analysé et les liens sont extraits
\item[C)] La page est téléchargée
\item[D)] Googlebot trouve une URL
\end{itemize}
*(Réponse : D)*

\item Que signifie le statut "Valide" dans le rapport de couverture de Google Search Console ?
\begin{itemize}
\item[A)] La page est bloquée
\item[B)] La page est correctement indexée
\item[C)] La page n'existe pas
\item[D)] La page contient des erreurs
\end{itemize}
*(Réponse : A)*

\item Quelle commande permet de vérifier l'indexation d'un site dans Google ?
\begin{itemize}
\item[A)] site:votresite.com
\item[B)] check:votresite.com
\item[C)] index:votresite.com
\item[D)] google:votresite.com
\end{itemize}
*(Réponse : A)*

\item Quel type d'erreur se produit quand une page renvoie un code HTTP 200 mais affiche un contenu d'erreur ?
\begin{itemize}
\item[A)] Erreur 404
\item[B)] Erreur 500
\item[C)] Soft 404
\item[D)] Erreur de redirection
\end{itemize}
*(Réponse : C)*

\item Que sont les "orphan pages" ?
\begin{itemize}
\item[A)] Des pages avec du contenu dupliqué
\item[B)] Des pages sans aucun lien interne pointant vers elles
\item[C)] Des pages avec des erreurs 500
\item[D)] Des pages protégées par mot de passe
\end{itemize}
*(Réponse : B)*

\item Quel outil permet de détecter les orphan pages ?
\begin{itemize}
\item[A)] Google Analytics
\item[B)] Screaming Frog
\item[C)] Google Trends
\item[D)] Keyword Planner
\end{itemize}
*(Réponse : C)*

\item Comment résoudre un Soft 404 ?
\begin{itemize}
\item[A)] Ignorer l'erreur
\item[B)] Modifier le serveur pour renvoyer un véritable statut 404
\item[C)] Supprimer la page
\item[D)] Ajouter plus de contenu
\end{itemize}
*(Réponse : D)*

\item Qu'indique le statut "Exclu" dans le rapport de couverture ?
\begin{itemize}
\item[A)] La page est en erreur
\item[B)] La page est intentionnellement non indexée
\item[C)] La page est valide
\item[D)] La page n'a pas été crawlée
\end{itemize}
*(Réponse : A)*

\item Quel outil Google permet d'inspecter une URL en détail ?
\begin{itemize}
\item[A)] Google Analytics
\item[B)] Google Search Console
\item[C)] Google Trends
\item[D)] Google Ads
\end{itemize}
*(Réponse : B)*

\item Que signifie une erreur 500 dans le rapport de couverture ?
\begin{itemize}
\item[A)] La page n'existe pas
\item[B)] Le serveur ne répond pas correctement
\item[C)] La page est bloquée par robots.txt
\item[D)] La page est une redirection
\end{itemize}
*(Réponse : C)*

\item Comment s'appelle la phase où Googlebot trouve une URL ?
\begin{itemize}
\item[A)] Rendu
\item[B)] Découverte
\item[C)] Analyse
\item[D)] Indexation
\end{itemize}
*(Réponse : D)*

\item Lors de quelle phase Google télécharge-t-il la page et ses ressources ?
\begin{itemize}
\item[A)] Découverte
\item[B)] Téléchargement
\item[C)] Rendu
\item[D)] Classification
\end{itemize}
*(Réponse : A)*

\item Que signifie le statut "Erreur" dans le rapport de couverture ?
\begin{itemize}
\item[A)] La page est correctement indexée
\item[B)] La page n'a pas pu être indexée
\item[C)] La page est intentionnellement exclue
\item[D)] La page est en attente de crawl
\end{itemize}
*(Réponse : B)*

\item Dans Google Search Console, comment accède-t-on au rapport de couverture ?
\begin{itemize}
\item[A)] Performances > Requêtes
\item[B)] Indexation > Pages (ou Coverage)
\item[C)] Expérience > Core Web Vitals
\item[D)] Sécurité > Actions manuelles
\end{itemize}
*(Réponse : C)*

\item Quelle est la première action à réaliser dans l'analyse du rapport Coverage ?
\begin{itemize}
\item[A)] Ignorer les erreurs
\item[B)] Analyser les quatre catégories (Erreur, Valide avec avertissement, Valide, Exclu)
\item[C)] Exporter toutes les données
\item[D)] Soumettre un nouveau sitemap
\end{itemize}
*(Réponse : D)*

\item Comment vérifier si Google voit votre contenu selon le guide ?
\begin{itemize}
\item[A)] Avec Google Analytics
\item[B)] Avec l'outil d'inspection d'URL dans GSC
\item[C)] Avec Google PageSpeed
\item[D)] Avec Google Trends
\end{itemize}
*(Réponse : A)*

\item Que faire si une page est trouvée comme "orphan page" ?
\begin{itemize}
\item[A)] La supprimer définitivement
\item[B)] Créer des liens internes vers cette page depuis des pages pertinentes
\item[C)] La rediriger vers la page d'accueil
\item[D)] La protéger par mot de passe
\end{itemize}
*(Réponse : B)*

\item Quel est le problème principal quand une page est bloquée par robots.txt ?
\begin{itemize}
\item[A)] La page ne s'affiche pas correctement
\item[B)] Le Googlebot ne peut pas crawler la page
\item[C)] La page est automatiquement supprimée
\item[D)] La page reçoit une pénalité
\end{itemize}
*(Réponse : C)*

\item Que doit-on faire après avoir corrigé des erreurs d'indexation ?
\begin{itemize}
\item[A)] Attendre 6 mois
\item[B)] Utiliser l'outil d'inspection d'URL pour des diagnostics détaillés
\item[C)] Recréer le site
\item[D)] Changer de nom de domaine
\end{itemize}
*(Réponse : D)*

\item Quelle est la conséquence d'une page non indexée ?
\begin{itemize}
\item[A)] Elle apparaît moins souvent
\item[B)] Elle n'existe tout simplement pas dans Google
\item[C)] Elle est pénalisée
\item[D)] Elle est visible uniquement sur mobile
\end{itemize}
*(Réponse : A)*

\item Dans la détection des orphan pages, que faut-il configurer dans Screaming Frog ?
\begin{itemize}
\item[A)] Le proxy
\item[B)] Le sitemap XML comme point de départ
\item[C)] Les filtres de langue
\item[D)] Les limites de taille
\end{itemize}
*(Réponse : B)*

\item Quel type d'erreur est une chaîne de redirections trop longue ?
\begin{itemize}
\item[A)] Erreur 404
\item[B)] Erreur de redirection
\item[C)] Soft 404
\item[D)] Erreur 500
\end{itemize}
*(Réponse : C)*

\item Pour résoudre une Soft 404, si la page est supprimée définitivement, que faut-il faire ?
\begin{itemize}
\item[A)] La laisser en l'état
\item[B)] Utiliser une redirection 301 vers une page pertinente
\item[C)] La remplacer par du texte vide
\item[D)] La verrouiller
\end{itemize}
*(Réponse : D)*

\item Combien de phases comprend le pipeline d'indexation de Google selon le guide ?
\begin{itemize}
\item[A)] 4
\item[B)] 6
\item[C)] 8
\item[D)] 10
\end{itemize}
*(Réponse : A)*

\item Quelle est la dernière phase du pipeline d'indexation ?
\begin{itemize}
\item[A)] Le rendu
\item[B)] La classification
\item[C)] L'indexation
\item[D)] La découverte
\end{itemize}
*(Réponse : B)*

\item Quel est le statut des pages intentionnellement non indexées selon le rapport de couverture ?
\begin{itemize}
\item[A)] Erreur
\item[B)] Exclu
\item[C)] Valide
\item[D)] Valide avec avertissement
\end{itemize}
*(Réponse : C)*

\item Dans l'analyse du rapport Coverage, quelle catégorie doit être corrigée en priorité ?
\begin{itemize}
\item[A)] Valide
\item[B)] Erreur
\item[C)] Exclu
\item[D)] Valide avec avertissement
\end{itemize}
*(Réponse : D)*

\item Que faut-il vérifier en priorité si le contenu est absent du HTML rendu dans l'inspection d'URL ?
\begin{itemize}
\item[A)] La vitesse du serveur
\item[B)] Si le JavaScript n'est pas bloqué par robots.txt
\item[C)] Le nombre de backlinks
\item[D)] La meta description
\end{itemize}
*(Réponse : A)*

\item Quelle est la différence entre une erreur 404 et une Soft 404 ?
\begin{itemize}
\item[A)] 404 est temporaire, Soft 404 est permanente
\item[B)] 404 renvoie un vrai code 404, Soft 404 renvoie 200 avec contenu d'erreur
\item[C)] Soft 404 est moins grave
\item[D)] 404 ne s'affiche pas dans GSC
\end{itemize}
*(Réponse : B)*

\item Quel fichier de votre serveur peut montrer la présence du Googlebot ?
\begin{itemize}
\item[A)] Le fichier .htaccess
\item[B)] Le fichier de logs de votre serveur
\item[C)] Le fichier robots.txt
\item[D)] Le fichier sitemap.xml
\end{itemize}
*(Réponse : C)*

\item Dans le pipeline d'indexation, que se passe-t-il juste après le téléchargement ?
\begin{itemize}
\item[A)] L'analyse
\item[B)] Le rendu
\item[C)] L'indexation
\item[D)] La classification
\end{itemize}
*(Réponse : D)*

\item Que signifie "Valide avec avertissement" dans le rapport de couverture ?
\begin{itemize}
\item[A)] La page n'est pas indexée
\item[B)] La page est indexée mais avec des problèmes potentiels
\item[C)] La page contient des erreurs critiques
\item[D)] La page est exclue
\end{itemize}
*(Réponse : A)*

\item Quelle est la meilleure approche pour analyser des milliers d'URLs dans GSC ?
\begin{itemize}
\item[A)] Les analyser une par une
\item[B)] Exporter les données pour chaque catégorie et prioriser les corrections
\item[C)] Ignorer les erreurs
\item[D)] Ne regarder que les pages valides
\end{itemize}
*(Réponse : B)*

\item Pour les pages avec le statut "Exclu", que peut-on faire ?
\begin{itemize}
\item[A)] Rien, c'est définitif
\item[B)] Analyser la cause de l'exclusion et la corriger si nécessaire
\item[C)] Les signaler manuellement
\item[D)] Créer des redirections
\end{itemize}
*(Réponse : C)*

\item Comment Google classe-t-il les pages dans l'étape de classification ?
\begin{itemize}
\item[A)] Par ordre alphabétique
\item[B)] Par thématiques et mots-clés
\item[C)] Par date de création
\item[D)] Par popularité
\end{itemize}
*(Réponse : D)*

\item Quelle est la première méthode recommandée pour vérifier l'indexation ?
\begin{itemize}
\item[A)] Vérifier les logs serveur
\item[B)] Utiliser la commande site:votresite.com
\item[C)] Utiliser Screaming Frog
\item[D)] Demander à Google
\end{itemize}
*(Réponse : A)*

\item Que faire si vous trouvez une page dans GSC avec le statut "Découverte - Pas encore indexée" ?
\begin{itemize}
\item[A)] La supprimer
\item[B)] Vérifier si elle a des liens internes et attendre
\item[C)] La rediriger immédiatement
\item[D)] Changer son URL
\end{itemize}
*(Réponse : B)*

\item Pour résoudre une Soft 404, que faut-il faire du serveur ?
\begin{itemize}
\item[A)] Le redémarrer
\item[B)] Modifier le serveur pour renvoyer un véritable statut 404
\item[C)] Augmenter sa capacité
\item[D)] Changer de serveur
\end{itemize}
*(Réponse : C)*

\item Dans quel outil Google trouve-t-on le rapport de couverture ?
\begin{itemize}
\item[A)] Google Ads
\item[B)] Google Search Console
\item[C)] Google Analytics
\item[D)] Google Tag Manager
\end{itemize}
*(Réponse : D)*

\item Si la page est supprimée définitivement (Soft 404), que faut-il faire ?
\begin{itemize}
\item[A)] Laisser la page 404
\item[B)] Utiliser une redirection 301 vers une page pertinente
\item[C)] Recréer la page
\item[D)] Ignorer l'erreur
\end{itemize}
*(Réponse : A)*

\item Comment personnaliser une page 404 ?
\begin{itemize}
\item[A)] La laisser vide
\item[B)] Ajouter des liens utiles pour l'utilisateur
\item[C)] La rediriger automatiquement
\item[D)] Afficher une page blanche
\end{itemize}
*(Réponse : B)*

\item Quel outil est recommandé pour crawler un site et trouver les problèmes d'indexation ?
\begin{itemize}
\item[A)] Google Analytics
\item[B)] Screaming Frog SEO Spider
\item[C)] Google Trends
\item[D)] Keyword Planner
\end{itemize}
*(Réponse : C)*

\item Quand Googlebot découvre-t-il une URL via le sitemap ?
\begin{itemize}
\item[A)] Lors de l'analyse
\item[B)] Lors de la phase de découverte
\item[C)] Lors du rendu
\item[D)] Lors de la classification
\end{itemize}
*(Réponse : D)*

\item Quelle est la finalité de la classification dans le pipeline d'indexation ?
\begin{itemize}
\item[A)] Stocker la page
\item[B)] Catégoriser la page par thématiques et mots-clés
\item[C)] Télécharger la page
\item[D)] Rendre la page
\end{itemize}
*(Réponse : A)*

\item Quel type de problème d'indexation est lié à des pages sans liens internes ?
\begin{itemize}
\item[A)] Erreur 500
\item[B)] Orphan pages
\item[C)] Soft 404
\item[D)] Erreur de redirection
\end{itemize}
*(Réponse : B)*

\item Comment prioriser les corrections après l'export des données de GSC ?
\begin{itemize}
\item[A)] Par ordre alphabétique
\item[B)] En fonction de la criticité et de l'impact
\item[C)] Aléatoirement
\item[D)] Par date de découverte
\end{itemize}
*(Réponse : C)*

\item Quel est l'outil principal pour suivre les performances d'indexation d'un site ?
\begin{itemize}
\item[A)] Google Analytics
\item[B)] Google Search Console
\item[C)] Google Ads
\item[D)] Google Tag Manager
\end{itemize}
*(Réponse : B)*

\item Le rapport "Couverture de l'index" dans GSC montre :
\begin{itemize}
\item[A)] Le nombre de visiteurs
\item[B)] Le statut d'indexation de chaque page (valide, erreur, exclu)
\item[C)] Le chiffre d'affaires
\item[D)] Les publicités
\end{itemize}
*(Réponse : A)*
\end{enumerate}

\newpage

% ============================================
% CHAPITRE 5
% ============================================
\section*{Chapitre 5 : Contrôle du crawling : Robots.txt et Sitemap.xml}

\begin{enumerate}
\item Où doit être placé le fichier robots.txt ?
\begin{itemize}
\item[A)] Dans le dossier /admin/
\item[B)] À la racine du site
\item[C)] Dans le dossier /images/
\item[D)] Dans le dossier /css/
\end{itemize}
*(Réponse : D)*

\item Quelle directive empêche les robots d'accéder à un dossier dans robots.txt ?
\begin{itemize}
\item[A)] Allow
\item[B)] Disallow
\item[C)] Block
\item[D)] Deny
\end{itemize}
*(Réponse : A)*

\item Que faut-il éviter de bloquer dans robots.txt pour le bon rendu des pages ?
\begin{itemize}
\item[A)] Les images
\item[B)] Les ressources CSS et JavaScript
\item[C)] Les fichiers PDF
\item[D)] Les vidéos
\end{itemize}
*(Réponse : B)*

\item Quelle directive dans robots.txt permet une exception spécifique ?
\begin{itemize}
\item[A)] Allow
\item[B)] Permit
\item[C)] Grant
\item[D)] Accept
\end{itemize}
*(Réponse : A)*

\item Comment référencer votre sitemap dans robots.txt ?
\begin{itemize}
\item[A)] Site: https://...
\item[B)] Sitemap: https://...
\item[C)] URL: https://...
\item[D)] Index: https://...
\end{itemize}
*(Réponse : C)*

\item Un blocage dans robots.txt empêche-t-il l'indexation si des liens externes pointent vers la page ?
\begin{itemize}
\item[A)] Oui, complètement
\item[B)] Non, cela empêche le crawling mais pas l'indexation
\item[C)] Oui, dans tous les cas
\item[D)] Cela dépend du moteur de recherche
\end{itemize}
*(Réponse : D)*

\item Quel outil Google permet de tester les directives robots.txt ?
\begin{itemize}
\item[A)] Google Analytics
\item[B)] L'outil de test robots.txt dans GSC
\item[C)] Google Trends
\item[D)] Google PageSpeed
\end{itemize}
*(Réponse : A)*

\item Quelle est l'utilité principale d'un sitemap XML ?
\begin{itemize}
\item[A)] Bloquer des pages
\item[B)] Lister toutes les pages importantes d'un site pour les moteurs de recherche
\item[C)] Améliorer le design du site
\item[D)] Suivre les visiteurs
\end{itemize}
*(Réponse : B)*

\item Quelle balise XML indique la fréquence de mise à jour d'une URL dans un sitemap ?
\begin{itemize}
\item[A)] <lastmod>
\item[B)] <changefreq>
\item[C)] <priority>
\item[D)] <update>
\end{itemize}
*(Réponse : C)*

\item Quelle est la valeur maximale de la balise <priority> dans un sitemap ?
\begin{itemize}
\item[A)] 0,5
\item[B)] 1,0
\item[C)] 10
\item[D)] 100
\end{itemize}
*(Réponse : D)*

\item Quel type de sitemap est utilisé pour l'indexation de contenu vidéo ?
\begin{itemize}
\item[A)] Sitemap images
\item[B)] Sitemap vidéo
\item[C)] Sitemap actualités
\item[D)] Sitemap mobile
\end{itemize}
*(Réponse : A)*

\item Quel type de sitemap est devenu obsolète avec le responsive design ?
\begin{itemize}
\item[A)] Sitemap vidéo
\item[B)] Sitemap mobile
\item[C)] Sitemap images
\item[D)] Sitemap actualités
\end{itemize}
*(Réponse : B)*

\item Que signifie User-agent: * dans robots.txt ?
\begin{itemize}
\item[A)] Les directives s'appliquent à un robot spécifique
\item[B)] Les directives s'appliquent à tous les robots
\item[C)] Les directives ne s'appliquent à aucun robot
\item[D)] Les directives s'appliquent aux utilisateurs humains
\end{itemize}
*(Réponse : C)*

\item Peut-on cacher du contenu avec robots.txt ?
\begin{itemize}
\item[A)] Oui, c'est recommandé
\item[B)] Non, car les concurrents peuvent voir votre robots.txt
\item[C)] Oui, si on utilise des mots de passe
\item[D)] Oui, c'est sécurisé
\end{itemize}
*(Réponse : D)*

\item Quelle extension de fichier a le sitemap XML ?
\begin{itemize}
\item[A)] .html
\item[B)] .xml
\item[C)] .txt
\item[D)] .csv
\end{itemize}
*(Réponse : A)*

\item Combien de sitemaps peut-on référencer dans robots.txt ?
\begin{itemize}
\item[A)] Un seul
\item[B)] Plusieurs
\item[C)] Maximum 5
\item[D)] Aucun
\end{itemize}
*(Réponse : B)*

\item Quelle balise indique la date de dernière modification d'une URL dans un sitemap ?
\begin{itemize}
\item[A)] <changefreq>
\item[B)] <lastmod>
\item[C)] <date>
\item[D)] <modified>
\end{itemize}
*(Réponse : C)*

\item Que doit-on faire avant de mettre en production des directives robots.txt ?
\begin{itemize}
\item[A)] Les ignorer
\item[B)] Les tester avec l'outil de test robots.txt dans GSC
\item[C)] Les soumettre à Google
\item[D)] Les commenter
\end{itemize}
*(Réponse : D)*

\item Quel est le format obligatoire du fichier robots.txt ?
\begin{itemize}
\item[A)] Fichier XML
\item[B)] Fichier texte
\item[C)] Fichier HTML
\item[D)] Fichier JSON
\end{itemize}
*(Réponse : A)*

\item Quelle est la conséquence de bloquer les ressources CSS et JavaScript dans robots.txt ?
\begin{itemize}
\item[A)] Aucune conséquence
\item[B)] Google peut mal rendre la page
\item[C)] La page charge plus vite
\item[D)] Le design est amélioré
\end{itemize}
*(Réponse : B)*

\item Que permet la directive Allow dans robots.txt ?
\begin{itemize}
\item[A)] Bloquer l'accès à une ressource
\item[B)] Autoriser l'accès à une ressource spécifique
\item[C)] Ignorer les directives
\item[D)] Rediriger les robots
\end{itemize}
*(Réponse : C)*

\item Quels sont les types de sitemaps supportés par Google ?
\begin{itemize}
\item[A)] Uniquement le sitemap standard
\item[B)] Sitemap vidéo, images, actualités, mobile
\item[C)] Sitemap HTML uniquement
\item[D)] Sitemap PDF
\end{itemize}
*(Réponse : D)*

\item Une page bloquée par robots.txt peut-elle être indexée ?
\begin{itemize}
\item[A)] Non, jamais
\item[B)] Oui, si des liens externes pointent vers elle
\item[C)] Oui, automatiquement
\item[D)] Non, car le crawl est nécessaire à l'indexation
\end{itemize}
*(Réponse : A)*

\item Dans quel cas utiliser un sitemap actualités ?
\begin{itemize}
\item[A)] Pour les sites e-commerce
\item[B)] Pour les sites d'information
\item[C)] Pour les blogs personnels
\item[D)] Pour les sites institutionnels
\end{itemize}
*(Réponse : B)*

\item Que faut-il faire si on modifie le fichier robots.txt ?
\begin{itemize}
\item[A)] Rien, Google le détecte automatiquement
\item[B)] Tester avec GSC et vérifier l'impact sur le crawl
\item[C)] Redémarrer le serveur
\item[D)] Contacter Google
\end{itemize}
*(Réponse : C)*

\item Quelle balise XML dans un sitemap indique le niveau d'importance d'une URL ?
\begin{itemize}
\item[A)] <importance>
\item[B)] <priority>
\item[C)] <level>
\item[D)] <rank>
\end{itemize}
*(Réponse : D)*

\item Où se trouve typiquement l'URL du sitemap dans robots.txt ?
\begin{itemize}
\item[A)] En haut du fichier
\item[B)] N'importe où dans le fichier
\item[C)] À la fin du fichier
\item[D)] Dans un commentaire
\end{itemize}
*(Réponse : C)*

\item Comment Google découvre-t-il les pages importantes via le sitemap ?
\begin{itemize}
\item[A)] En les crawlant dans l'ordre alphabétique
\item[B)] En parcourant la liste des URLs fournies dans le sitemap
\item[C)] En demandant aux utilisateurs
\item[D)] En suivant des publicités
\end{itemize}
*(Réponse : A)*

\item Le nom de domaine du sitemap doit-il correspondre au site ?
\begin{itemize}
\item[A)] Non, il peut être différent
\item[B)] Oui, il doit correspondre
\item[C)] Cela n'a pas d'importance
\item[D)] Google le détermine automatiquement
\end{itemize}
*(Réponse : B)*

\item Que signifie la directive "Disallow: /" dans robots.txt ?
\begin{itemize}
\item[A)] Seul le dossier racine est bloqué
\item[B)] Tout le site est bloqué pour les robots
\item[C)] Seuls les fichiers à la racine sont bloqués
\item[D)] Rien n'est bloqué
\end{itemize}
*(Réponse : C)*

\item Quel est l'effet de "Disallow: /wp-admin/" ?
\begin{itemize}
\item[A)] Le dossier wp-admin est accessible
\item[B)] Le dossier wp-admin est bloqué pour les robots
\item[C)] Tout le site est bloqué
\item[D)] Les images dans wp-admin sont bloquées
\end{itemize}
*(Réponse : D)*

\item À quoi sert la balise <loc> dans un sitemap XML ?
\begin{itemize}
\item[A)] À indiquer la langue de la page
\item[B)] À spécifier l'URL de la page
\item[C)] À localiser le serveur
\item[D)] À indiquer la taille de la page
\end{itemize}
*(Réponse : A)*

\item Quel est le bon réflexe si on a des paramètres d'URL comme ?sort= dans un site ?
\begin{itemize}
\item[A)] Les laisser crawler librement
\item[B)] Les bloquer avec Disallow: /*?sort=
\item[C)] Les rediriger
\item[D)] Les supprimer
\end{itemize}
*(Réponse : B)*

\item Pourquoi ne faut-il pas essayer de cacher du contenu via robots.txt ?
\begin{itemize}
\item[A)] Parce que c'est illégal
\item[B)] Parce que les concurrents peuvent voir votre robots.txt
\item[C)] Parce que Google l'ignore
\item[D)] Parce que cela ralentit le site
\end{itemize}
*(Réponse : C)*

\item Quel est l'objectif principal du fichier robots.txt ?
\begin{itemize}
\item[A)] Améliorer le référencement automatiquement
\item[B)] Donner des instructions aux robots des moteurs de recherche
\item[C)] Sécuriser le site
\item[D)] Analyser les visiteurs
\end{itemize}
*(Réponse : D)*

\item Que doit-on faire après avoir créé un sitemap XML ?
\begin{itemize}
\item[A)] L'imprimer
\item[B)] Le soumettre à Google Search Console
\item[C)] Le convertir en PDF
\item[D)] Le supprimer
\end{itemize}
*(Réponse : A)*

\item Quel est l'avantage d'un sitemap pour les moteurs de recherche ?
\begin{itemize}
\item[A)] Il améliore le design
\item[B)] Il permet de découvrir rapidement les pages importantes
\item[C)] Il sécurise le site
\item[D)] Il augmente la vitesse
\end{itemize}
*(Réponse : B)*

\item Un seul fichier robots.txt peut-il contrôler plusieurs sous-domaines ?
\begin{itemize}
\item[A)] Oui, automatiquement
\item[B)] Non, chaque sous-domaine a son propre robots.txt
\item[C)] Oui, avec une directive spéciale
\item[D)] Cela dépend du serveur
\end{itemize}
*(Réponse : C)*

\item Que se passe-t-il si un sitemap contient des URLs bloquées par robots.txt ?
\begin{itemize}
\item[A)] Google les indexe quand même
\item[B)] Google ne peut pas les crawler mais les connaît
\item[C)] Le sitemap est rejeté
\item[D)] Une pénalité est appliquée
\end{itemize}
*(Réponse : D)*

\item Comment s'appelle le fichier qui liste les pages importantes pour les moteurs de recherche ?
\begin{itemize}
\item[A)] Index.xml
\item[B)] Sitemap.xml
\item[C)] Pages.xml
\item[D)] Urls.xml
\end{itemize}
*(Réponse : A)*

\item Quelle est la fonction de la directive Allow dans le robots.txt ?
\begin{itemize}
\item[A)] Bloquer des dossiers spécifiques
\item[B)] Permettre l'accès à des fichiers ou dossiers spécifiques
\item[C)] Augmenter la priorité de crawl
\item[D)] Rediriger les robots
\end{itemize}
*(Réponse : B)*

\item Quel type de sitemap est recommandé pour les images ?
\begin{itemize}
\item[A)] Sitemap vidéo
\item[B)] Sitemap images
\item[C)] Sitemap standard
\item[D)] Sitemap actualités
\end{itemize}
*(Réponse : C)*

\item Une fois le sitemap soumis, Google :
\begin{itemize}
\item[A)] Indexe immédiatement toutes les pages
\item[B)] Utilise ces URLs comme suggestions pour le crawl
\item[C)] Les ignore si le site n'est pas populaire
\item[D)] Les convertit en pages web
\end{itemize}
*(Réponse : D)*

\item Comment tester que le robots.txt est correctement configuré ?
\begin{itemize}
\item[A)] En regardant le site dans un navigateur
\item[B)] En utilisant l'outil de test dans Google Search Console
\item[C)] En demandant à un ami
\item[D)] En vérifiant les logs
\end{itemize}
*(Réponse : A)*

\item Que doit-on faire si on bloque /private/ dans robots.txt ?
\begin{itemize}
\item[A)] Ajouter un mot de passe
\item[B)] Vérifier qu'aucune ressource importante n'est bloquée
\item[C)] Ignorer la configuration
\item[D)] Supprimer le dossier
\end{itemize}
*(Réponse : B)*

\item Un site avec des milliers de pages doit-il avoir un sitemap ?
\begin{itemize}
\item[A)] Non, ce n'est pas nécessaire
\item[B)] Oui, fortement recommandé
\item[C)] Seulement si le site est e-commerce
\item[D)] Cela dépend du CMS
\end{itemize}
*(Réponse : C)*

\item Quelle est la différence entre robots.txt et sitemap.xml ?
\begin{itemize}
\item[A)] robots.txt aide au crawl, sitemap.xml bloque les pages
\item[B)] robots.txt contrôle l'accès des robots, sitemap.xml liste les pages à indexer
\item[C)] Ce sont deux noms pour le même fichier
\item[D)] robots.txt est pour les utilisateurs, sitemap.xml pour les robots
\end{itemize}
*(Réponse : D)*

\item À quelle fréquence Google vérifie-t-il le fichier robots.txt ?
\begin{itemize}
\item[A)] Une fois par mois
\item[B)] Régulièrement, à chaque crawl du site
\item[C)] Jamais
\item[D)] Une fois par an
\end{itemize}
*(Réponse : A)*

\item Quelle directive robots.txt bloque complètement les robots ?
\begin{itemize}
\item[A)] Allow: /
\item[B)] Disallow: /
\item[C)] Noindex
\item[D)] Nofollow
\end{itemize}
*(Réponse : D)*

\item Le paramètre "Sitemap:" dans le robots.txt permet de :
\begin{itemize}
\item[A)] Bloquer les pages du sitemap
\item[B)] Indiquer l'emplacement du fichier sitemap directement dans robots.txt
\item[C)] Ignorer le sitemap
\item[D)] Créer un sitemap automatiquement
\end{itemize}
*(Réponse : A)*
\end{enumerate}

\newpage

% ============================================
% CHAPITRE 6
% ============================================
\section*{Chapitre 6 : Balises Meta Robots et Canonical}

\begin{enumerate}
\item Quelle balise permet un contrôle précis de l'indexation au niveau de chaque page ?
\begin{itemize}
\item[A)] La balise title
\item[B)] La balise meta robots
\item[C)] La balise H1
\item[D)] La balise meta description
\end{itemize}
*(Réponse : B)*

\item Quelle directive meta robots signifie "ne pas indexer mais suivre les liens" ?
\begin{itemize}
\item[A)] index, follow
\item[B)] noindex, follow
\item[C)] index, nofollow
\item[D)] noindex, nofollow
\end{itemize}
*(Réponse : C)*

\item Quelle directive meta robots signifie "indexer mais ne pas suivre les liens" ?
\begin{itemize}
\item[A)] noindex, follow
\item[B)] index, nofollow
\item[C)] noindex, nofollow
\item[D)] index, follow
\end{itemize}
*(Réponse : D)*

\item Quelle directive meta robots signifie "ne pas indexer et ne pas suivre" ?
\begin{itemize}
\item[A)] index, nofollow
\item[B)] noindex, nofollow
\item[C)] noindex, follow
\item[D)] index, follow
\end{itemize}
*(Réponse : A)*

\item Quelle directive meta robots empêche l'affichage du cache Google ?
\begin{itemize}
\item[A)] noimageindex
\item[B)] noarchive
\item[C)] nosnippet
\item[D)] nofollow
\end{itemize}
*(Réponse : B)*

\item Quelle directive meta robots empêche l'affichage dans la recherche d'images ?
\begin{itemize}
\item[A)] noarchive
\item[B)] noimageindex
\item[C)] nosnippet
\item[D)] nofollow
\end{itemize}
*(Réponse : C)*

\item Quelle balise est utilisée pour résoudre les problèmes de contenu dupliqué ?
\begin{itemize}
\item[A)] La balise meta robots
\item[B)] Le tag canonique (rel="canonical")
\item[C)] La balise H1
\item[D)] La balise meta description
\end{itemize}
*(Réponse : D)*

\item Dans quels cas utiliser le tag canonique ?
\begin{itemize}
\item[A)] Uniquement pour les images
\item[B)] Pages accessibles via plusieurs URLs, contenu syndiqué
\item[C)] Uniquement pour les vidéos
\item[D)] Pour les pages d'erreur
\end{itemize}
*(Réponse : A)*

\item Comment s'écrit le tag canonique dans le HTML ?
\begin{itemize}
\item[A)] <meta name="canonical" content="...">
\item[B)] <link rel="canonical" href="..." />
\item[C)] <canonical>...</canonical>
\item[D)] <link rel="alternate" href="..." />
\end{itemize}
*(Réponse : B)*

\item Quelle directive meta robots désactive l'affichage d'un extrait dans les SERP ?
\begin{itemize}
\item[A)] noarchive
\item[B)] nosnippet
\item[C)] noimageindex
\item[D)] nofollow
\end{itemize}
*(Réponse : C)*

\item À quoi sert la directive "index, follow" ?
\begin{itemize}
\item[A)] Indexer et ne pas suivre les liens
\item[B)] Indexer et suivre les liens (comportement par défaut)
\item[C)] Ne pas indexer mais suivre les liens
\item[D)] Ni indexer ni suivre
\end{itemize}
*(Réponse : D)*

\item Quel attribut est utilisé pour le ciblage spécifique d'un moteur (ex: Googlebot) ?
\begin{itemize}
\item[A)] name="robots"
\item[B)] name="googlebot"
\item[C)] name="crawler"
\item[D)] name="bot"
\end{itemize}
*(Réponse : A)*

\item Le tag canonique doit-il être utilisé pour le contenu syndiqué ?
\begin{itemize}
\item[A)] Non, jamais
\item[B)] Oui, pour indiquer la source originale
\item[C)] Seulement si le contenu est dupliqué automatiquement
\item[D)] Non, le contenu syndiqué est toujours pénalisé
\end{itemize}
*(Réponse : B)*

\item Que se passe-t-il si on utilise un canonical pointant vers une page qui n'existe pas ?
\begin{itemize}
\item[A)] Google l'ignore
\item[B)] Google peut ne pas indexer la page correctement
\item[C)] Google redirige automatiquement
\item[D)] Cela n'a aucun effet
\end{itemize}
*(Réponse : C)*

\item À quoi sert la directive noarchive ?
\begin{itemize}
\item[A)] Empêcher l'indexation
\item[B)] Empêcher Google d'afficher un lien "Cache" dans les résultats
\item[C)] Empêcher le suivi des liens
\item[D)] Empêcher le rendu JavaScript
\end{itemize}
*(Réponse : D)*

\item Comment gérer les versions imprimables des pages avec le canonical ?
\begin{itemize}
\item[A)] Les supprimer
\item[B)] Utiliser le canonical vers la version principale
\item[C)] Les indexer séparément
\item[D)] Les noindexer
\end{itemize}
*(Réponse : A)*

\item Quelle est la différence entre noindex et nofollow ?
\begin{itemize}
\item[A)] Aucune différence
\item[B)] noindex empêche l'indexation, nofollow empêche le suivi des liens
\item[C)] noindex empêche le suivi des liens, nofollow empêche l'indexation
\item[D)] Les deux empêchent l'indexation
\end{itemize}
*(Réponse : B)*

\item Dans quel cas utiliser "noindex, follow" ?
\begin{itemize}
\item[A)] Pour une page importante qu'on veut bien classer
\item[B)] Pour une page qu'on ne veut pas voir dans l'index mais dont on veut transmettre l'autorité via les liens
\item[C)] Pour une page qu'on veut indexer sans suivre ses liens
\item[D)] Pour toutes les pages du site
\end{itemize}
*(Réponse : C)*

\item Que signifie le tag canonique pour Google ?
\begin{itemize}
\item[A)] Google doit indexer cette page en priorité
\item[B)] Google doit considérer cette URL comme la version principale
\item[C)] Google doit supprimer les autres versions
\item[D)] Google doit augmenter le PageRank de cette page
\end{itemize}
*(Réponse : D)*

\item Que faire si une page a un contenu similaire accessible via des paramètres de tracking ?
\begin{itemize}
\item[A)] Laisser faire Google
\item[B)] Utiliser le canonical vers la version standard
\item[C)] Supprimer les paramètres
\item[D)] Créer des pages séparées
\end{itemize}
*(Réponse : A)*

\item Quelle est la valeur par défaut si aucune directive meta robots n'est spécifiée ?
\begin{itemize}
\item[A)] noindex, nofollow
\item[B)] index, follow
\item[C)] noindex, follow
\item[D)] index, nofollow
\end{itemize}
*(Réponse : B)*

\item La directive nosnippet est particulièrement utile pour :
\begin{itemize}
\item[A)] Les pages d'accueil
\item[B)] Les pages avec du contenu sensible
\item[C)] Les pages de produits
\item[D)] Les articles de blog
\end{itemize}
*(Réponse : C)*

\item Que signifie le canonical auto-référençant ?
\begin{itemize}
\item[A)] La page pointe vers une autre page
\item[B)] La page pointe vers elle-même comme canonical
\item[C)] La page n'a pas de canonical
\item[D)] Le canonical est défini dans le sitemap
\end{itemize}
*(Réponse : D)*

\item Quel problème le tag canonique résout-il principalement ?
\begin{itemize}
\item[A)] La lenteur du site
\item[B)] Le contenu dupliqué
\item[C)] Les backlinks de mauvaise qualité
\item[D)] Les erreurs 404
\end{itemize}
*(Réponse : A)*

\item Peut-on utiliser plusieurs tags canoniques sur une même page ?
\begin{itemize}
\item[A)] Oui, c'est recommandé
\item[B)] Non, Google peut ignorer les deux
\item[C)] Oui, pour plusieurs langues
\item[D)] Oui, pour le desktop et le mobile
\end{itemize}
*(Réponse : B)*

\item Comment Google interprète-t-il "noindex" seul (sans nofollow) ?
\begin{itemize}
\item[A)] Il ne suit pas les liens
\item[B)] Il suit les liens mais n'indexe pas la page
\item[C)] Il n'indexe pas et ne suit pas les liens
\item[D)] Il ignore la directive
\end{itemize}
*(Réponse : C)*

\item À quoi sert le tag "googlebot" dans meta robots ?
\begin{itemize}
\item[A)] À cibler tous les moteurs de recherche
\item[B)] À donner des instructions spécifiques à Google uniquement
\item[C)] À améliorer le ranking
\item[D)] À activer des fonctionnalités Google
\end{itemize}
*(Réponse : D)*

\item Si on utilise "nofollow" sur tous les liens d'une page, quel est l'effet SEO ?
\begin{itemize}
\item[A)] Aucun effet, Google les suit quand même
\item[B)] Google ne transmet pas l'autorité via ces liens
\item[C)] La page est mieux classée
\item[D)] Les liens sont supprimés
\end{itemize}
*(Réponse : A)*

\item Dans quel scénario doit-on utiliser un canonical ?
\begin{itemize}
\item[A)] Pour chaque page du site
\item[B)] Pour les pages avec des URLs multiples (filtres, tri, paramètres)
\item[C)] Uniquement pour la page d'accueil
\item[D)] Jamais
\end{itemize}
*(Réponse : B)*

\item Comment Google traite-t-il une page avec "noindex" ?
\begin{itemize}
\item[A)] Il l'indexe normalement
\item[B)] Il la retire de son index (si déjà indexée) ou ne l'ajoute pas
\item[C)] Il la classe en première position
\item[D)] Il la supprime du serveur
\end{itemize}
*(Réponse : C)*

\item Quel est l'effet d'une mauvaise utilisation du canonical ?
\begin{itemize}
\item[A)] Amélioration du référencement
\item[B)] Nuisance potentielle au référencement
\item[C)] Aucun effet
\item[D)] Redirection automatique
\end{itemize}
*(Réponse : D)*

\item Quel attribut empêche Google de créer un extrait pour la page ?
\begin{itemize}
\item[A)] noarchive
\item[B)] nosnippet
\item[C)] noimageindex
\item[D)] nofollow
\end{itemize}
*(Réponse : A)*

\item Pour une page avec contenu syndiqué, la meilleure pratique est :
\begin{itemize}
\item[A)] De ne pas utiliser de canonical
\item[B)] D'utiliser un canonical pointant vers l'article original
\item[C)] De supprimer la page
\item[D)] D'utiliser nofollow
\end{itemize}
*(Réponse : B)*

\item Quand deux pages ont un contenu très similaire, que recommande Google ?
\begin{itemize}
\item[A)] Les conserver toutes les deux
\item[B)] Utiliser le tag canonical pour indiquer la page préférée
\item[C)] Supprimer les deux
\item[D)] Les fusionner
\end{itemize}
*(Réponse : C)*

\item Peut-on indexer une page mais ne pas suivre ses liens ?
\begin{itemize}
\item[A)] Non, c'est impossible
\item[B)] Oui, avec robots meta "index, nofollow"
\item[C)] Oui, avec "noindex, follow"
\item[D)] Non, Google décide automatiquement
\end{itemize}
*(Réponse : D)*

\item Que fait la directive noarchive ?
\begin{itemize}
\item[A)] Empêche Google d'indexer la page
\item[B)] Empêche Google de stocker une copie de la page dans le cache
\item[C)] Empêche Google de suivre les liens
\item[D)] Empêche Google d'afficher la page
\end{itemize}
*(Réponse : A)*

\item Quelle directive spécifique à Google empêche l'affichage d'un extrait ?
\begin{itemize}
\item[A)] noindex
\item[B)] nosnippet
\item[C)] nofollow
\item[D)] noarchive
\end{itemize}
*(Réponse : B)*

\item Pourquoi utiliser "noindex" sur des pages de faible valeur ?
\begin{itemize}
\item[A)] Pour améliorer leur classement
\item[B)] Pour économiser le budget de crawl et éviter de diluer l'autorité
\item[C)] Pour les rendre invisibles
\item[D)] Pour les protéger
\end{itemize}
*(Réponse : C)*

\item Que se passe-t-il si une page a un canonical qui pointe vers une page avec noindex ?
\begin{itemize}
\item[A)] Google indexe quand même la page canonique
\item[B)] Google peut ne pas indexer la page canonique
\item[C)] Google ignore le canonical
\item[D)] La page est mieux classée
\end{itemize}
*(Réponse : D)*

\item Quelle directive meta robots est équivalente au comportement par défaut ?
\begin{itemize}
\item[A)] noindex, nofollow
\item[B)] index, follow
\item[C)] noindex, follow
\item[D)] index, nofollow
\end{itemize}
*(Réponse : A)*

\item Pour une page de remerciement après inscription, que recommande-t-on ?
\begin{itemize}
\item[A)] index, follow
\item[B)] noindex, follow
\item[C)] index, nofollow
\item[D)] noindex, nofollow
\end{itemize}
*(Réponse : B)*

\item Que signifie le tag canonique pour le PageRank ?
\begin{itemize}
\item[A)] Il distribue le PageRank entre toutes les versions
\item[B)] Il consolide les signaux de classement vers l'URL canonique
\item[C)] Il annule le PageRank
\item[D)] Il double le PageRank
\end{itemize}
*(Réponse : C)*

\item L'utilisation de "noindex" empêche-t-elle Google de crawler la page ?
\begin{itemize}
\item[A)] Oui, complètement
\item[B)] Non, Google peut crawler la page mais ne l'indexe pas
\item[C)] Oui, si le fichier robots.txt le bloque aussi
\item[D)] Non, Google ignore toujours la directive
\end{itemize}
*(Réponse : D)*

\item Dans quel cas la directive "noimageindex" est-elle utile ?
\begin{itemize}
\item[A)] Pour les pages avec beaucoup de texte
\item[B)] Pour les pages où on ne veut pas que les images apparaissent dans Google Images
\item[C)] Pour toutes les pages
\item[D)] Pour les pages d'erreur
\end{itemize}
*(Réponse : A)*

\item Que faire si on a des pages avec des tris différents (prix, date) ?
\begin{itemize}
\item[A)] Les indexer toutes séparément
\item[B)] Utiliser le canonical vers la page de tri par défaut
\item[C)] Les supprimer
\item[D)] Les marquer en nofollow
\end{itemize}
*(Réponse : B)*

\item Le tag canonique doit-il être présent dans le sitemap XML ?
\begin{itemize}
\item[A)] Oui, obligatoirement
\item[B)] Non, ce sont des concepts indépendants mais complémentaires
\item[C)] Oui, uniquement pour les pages produits
\item[D)] Non, jamais
\end{itemize}
*(Réponse : C)*

\item Quelle est la différence entre une balise meta robots "noindex" et "none" ?
\begin{itemize}
\item[A)] "none" équivaut à "noindex, nofollow"
\item[B)] "none" équivaut à "index, follow"
\item[C)] Aucune différence
\item[D)] "none" désactive le JavaScript
\end{itemize}
*(Réponse : A)*

\item La balise canonical peut-elle pointer vers une page qui se redirige elle-même ?
\begin{itemize}
\item[A)] Oui, c'est une bonne pratique
\item[B)] Non, cela crée une chaîne de canoniques problématique
\item[C)] Oui, uniquement pour les pages produits
\item[D)] Cela n'a pas d'importance
\end{itemize}
*(Réponse : B)*

\item Dans quel cas utiliser la balise meta robots "noarchive" ?
\begin{itemize}
\item[A)] Pour empêcher Google de crawler la page
\item[B)] Pour empêcher Google d'afficher une version en cache de la page
\item[C)] Pour indexer la page plus vite
\item[D)] Pour améliorer le PageRank
\end{itemize}
*(Réponse : C)*

\item À quoi sert la directive "max-snippet:-1" dans la balise meta robots ?
\begin{itemize}
\item[A)] À limiter le snippet à -1 caractères
\item[B)] À autoriser un snippet de longueur illimitée
\item[C)] À interdire les snippets
\item[D)] À supprimer la page de l'index
\end{itemize}
*(Réponse : D)*
\end{enumerate}

\newpage

% ============================================
% CHAPITRE 7
% ============================================
\section*{Chapitre 7 : Le Budget de Crawl (Crawl Budget)}

\begin{enumerate}
\item Qu'est-ce que le crawl budget ?
\begin{itemize}
\item[A)] Le budget financier alloué au SEO
\item[B)] Le nombre de pages que Googlebot va crawler sur un site dans une période donnée
\item[C)] Le nombre de liens sortants autorisés
\item[D)] Le temps maximum de chargement d'une page
\end{itemize}
*(Réponse : D)*

\item Pour quel type de sites la gestion du crawl budget est-elle cruciale ?
\begin{itemize}
\item[A)] Les petits sites
\item[B)] Les sites de grande taille
\item[C)] Les sites avec peu de pages
\item[D)] Les sites en construction
\end{itemize}
*(Réponse : A)*

\item Quel facteur influence positivement le crawl budget ?
\begin{itemize}
\item[A)] Un serveur lent
\item[B)] La vitesse du serveur
\item[C)] Un taux d'erreurs élevé
\item[D)] Du contenu de faible qualité
\end{itemize}
*(Réponse : B)*

\item Quel facteur réduit le crawl budget ?
\begin{itemize}
\item[A)] Un serveur rapide
\item[B)] Trop d'erreurs 404 ou 500
\item[C)] Un contenu de qualité
\item[D)] Des mises à jour régulières
\end{itemize}
*(Réponse : C)*

\item Comment la popularité d'un site affecte-t-elle le crawl budget ?
\begin{itemize}
\item[A)] Les sites populaires reçoivent moins de budget
\item[B)] Les sites populaires reçoivent plus de budget de crawl
\item[C)] La popularité n'a pas d'impact
\item[D)] Seuls les sites nouveaux reçoivent du budget
\end{itemize}
*(Réponse : D)*

\item Quel impact les mises à jour régulières ont-elles sur le crawl ?
\begin{itemize}
\item[A)] Elles réduisent la fréquence de crawl
\item[B)] Elles augmentent la fréquence de crawl
\item[C)] Elles n'ont aucun impact
\item[D)] Le crawl s'arrête complètement
\end{itemize}
*(Réponse : A)*

\item Que faut-il faire pour optimiser le crawl budget ?
\begin{itemize}
\item[A)] Ajouter plus de pages
\item[B)] Éliminer les pages de faible valeur
\item[C)] Réduire la vitesse du serveur
\item[D)] Créer des redirections en chaîne
\end{itemize}
*(Réponse : B)*

\item Comment la structure des liens internes influence-t-elle le crawl budget ?
\begin{itemize}
\item[A)] Elle n'a pas d'importance
\item[B)] Les pages importantes doivent être proches de la racine
\item[C)] Tous les liens doivent être en nofollow
\item[D)] Il faut éviter les liens internes
\end{itemize}
*(Réponse : C)*

\item Quel est l'effet des redirections en chaîne sur le crawl budget ?
\begin{itemize}
\item[A)] Elles améliorent le crawl budget
\item[B)] Chaque redirection consomme du budget
\item[C)] Elles n'ont aucun effet
\item[D)] Elles doublent le budget
\end{itemize}
*(Réponse : D)*

\item Quelle directive est recommandée pour les pages non essentielles dans l'optimisation du crawl budget ?
\begin{itemize}
\item[A)] index, follow
\item[B)] noindex
\item[C)] nofollow uniquement
\item[D)] Aucune directive
\end{itemize}
*(Réponse : A)*

\item Quel est l'effet du contenu de faible qualité sur le crawl budget ?
\begin{itemize}
\item[A)] Il augmente le budget de crawl
\item[B)] Il peut réduire le budget de crawl
\item[C)] La qualité n'affecte pas le budget
\item[D)] Seul le contenu vidéo compte
\end{itemize}
*(Réponse : B)*

\item Comment les erreurs 404 affectent-elles le crawl budget ?
\begin{itemize}
\item[A)] Elles augmentent le budget
\item[B)] Elles réduisent le budget
\item[C)] Elles n'ont aucun effet
\item[D)] Google les ignore
\end{itemize}
*(Réponse : C)*

\item Pour optimiser le crawl budget, que faut-il faire des pages inutiles ?
\begin{itemize}
\item[A)] Les conserver
\item[B)] Empêcher leur indexation
\item[C)] Les dupliquer
\item[D)] Les mettre à jour régulièrement
\end{itemize}
*(Réponse : D)*

\item Quel est l'impact de la taille du site sur le crawl budget ?
\begin{itemize}
\item[A)] Les petits sites ont plus de budget
\item[B)] Les sites plus vastes peuvent obtenir plus de budget s'ils sont de qualité
\item[C)] La taille n'a pas d'impact
\item[D)] Seuls les sites de taille moyenne ont du budget
\end{itemize}
*(Réponse : A)*

\item Que faut-il faire pour améliorer l'allocation du crawl budget ?
\begin{itemize}
\item[A)] Ralentir le serveur
\item[B)] Améliorer la vitesse de réponse du serveur
\item[C)] Désactiver le cache
\item[D)] Ajouter des images lourdes
\end{itemize}
*(Réponse : B)*

\item Qu'est-ce qu'une redirection en chaîne ?
\begin{itemize}
\item[A)] Une redirection directe
\item[B)] Une série de redirections A vers B vers C au lieu de A vers C
\item[C)] Une redirection temporaire
\item[D)] Une redirection permanente
\end{itemize}
*(Réponse : C)*

\item Les sites de qualité avec beaucoup de pages reçoivent généralement :
\begin{itemize}
\item[A)] Moins de budget de crawl
\item[B)] Plus de budget de crawl
\item[C)] Aucun budget de crawl
\item[D)] Un budget fixe quelle que soit la qualité
\end{itemize}
*(Réponse : D)*

\item Que doit-on faire si Googlebot passe trop de temps sur des pages de faible valeur ?
\begin{itemize}
\item[A)] Les améliorer
\item[B)] Les marquer en noindex pour économiser le budget
\item[C)] Les supprimer du serveur
\item[D)] Ajouter plus de contenu
\end{itemize}
*(Réponse : A)*

\item Quel est le lien entre la fraîcheur du contenu et le crawl budget ?
\begin{itemize}
\item[A)] Le contenu frais réduit le crawl
\item[B)] Les mises à jour régulières augmentent la fréquence de crawl
\item[C)] La fraîcheur n'a pas d'impact
\item[D)] Seul le contenu historique est crawlée
\end{itemize}
*(Réponse : B)*

\item Une page importante doit être à combien de clics de la racine pour optimiser le crawl budget ?
\begin{itemize}
\item[A)] 5 clics
\item[B)] 1-2 clics
\item[C)] 10 clics
\item[D)] Peu importe
\end{itemize}
*(Réponse : C)*

\item Les sites ayant beaucoup d'erreurs 500 verront leur crawl budget :
\begin{itemize}
\item[A)] Augmenter
\item[B)] Diminuer
\item[C)] Rester identique
\item[D)] Être annulé
\end{itemize}
*(Réponse : D)*

\item Quel élément est un facteur clé du crawl budget selon Google ?
\begin{itemize}
\item[A)] La couleur du site
\item[B)] La vitesse du serveur
\item[C)] Le nombre de polices utilisées
\item[D)] Le type de CMS
\end{itemize}
*(Réponse : A)*

\item Comment le contenu de faible qualité impacte-t-il le crawl budget ?
\begin{itemize}
\item[A)] Il augmente la fréquence de crawl
\item[B)] Il peut réduire le budget de crawl
\item[C)] Il n'a aucun impact
\item[D)] Google l'ignore complètement
\end{itemize}
*(Réponse : B)*

\item L'optimisation du crawl budget est particulièrement importante pour :
\begin{itemize}
\item[A)] Les blogs personnels
\item[B)] Les sites avec des milliers de pages
\item[C)] Les sites avec une seule page
\item[D)] Les sites en construction
\end{itemize}
*(Réponse : C)*

\item Pour réduire la consommation inutile du crawl budget, il faut :
\begin{itemize}
\item[A)] Augmenter le nombre de pages
\item[B)] Corriger les erreurs 404 et 500 rapidement
\item[C)] Ajouter des redirections
\item[D)] Désactiver le crawl
\end{itemize}
*(Réponse : D)*

\item Quel est l'effet d'un serveur rapide sur le crawl budget ?
\begin{itemize}
\item[A)] Il réduit le nombre de pages crawlées
\item[B)] Il permet plus de crawl
\item[C)] Il n'a aucun effet
\item[D)] Il double le temps de crawl
\end{itemize}
*(Réponse : A)*

\item Les sites populaires reçoivent-ils plus de budget de crawl ?
\begin{itemize}
\item[A)] Non, la popularité n'a pas d'importance
\item[B)] Oui, les sites populaires reçoivent plus de budget
\item[C)] Non, les sites populaires reçoivent moins
\item[D)] Seuls les nouveaux sites reçoivent du budget
\end{itemize}
*(Réponse : B)*

\item Que faut-il corriger en priorité dans le rapport de couverture pour améliorer le crawl budget ?
\begin{itemize}
\item[A)] Les pages valides
\item[B)] Les erreurs 404 et 500
\item[C)] Les pages exclues
\item[D)] Les pages avec avertissement
\end{itemize}
*(Réponse : C)*

\item La proximité des pages importantes de la racine du site permet :
\begin{itemize}
\item[A)] De réduire leur importance
\item[B)] D'optimiser le crawl budget
\item[C)] De les cacher aux moteurs
\item[D)] De réduire la vitesse
\end{itemize}
*(Réponse : D)*

\item Combien de facteurs clés influencent le crawl budget selon le guide ?
\begin{itemize}
\item[A)] 3
\item[B)] 6
\item[C)] 10
\item[D)] 2
\end{itemize}
*(Réponse : A)*

\item Quel est l'impact d'un taux d'erreurs élevé sur le crawl budget ?
\begin{itemize}
\item[A)] Il augmente le budget de crawl
\item[B)] Il réduit le budget de crawl
\item[C)] Il n'a aucun impact
\item[D)] Il stoppe le crawl
\end{itemize}
*(Réponse : B)*

\item Le crawl budget est-il une ressource illimitée ?
\begin{itemize}
\item[A)] Oui, Google peut crawler autant qu'il veut
\item[B)] Non, ce n'est pas une ressource illimitée
\item[C)] Oui, pour les sites de qualité
\item[D)] Cela dépend de la taille du site
\end{itemize}
*(Réponse : C)*

\item L'utilisation de la directive noindex sur les pages non essentielles permet :
\begin{itemize}
\item[A)] D'augmenter leur trafic
\item[B)] D'économiser le budget de crawl
\item[C)] De les supprimer du site
\item[D)] De les rendre invisibles
\end{itemize}
*(Réponse : D)*

\item Quel est l'effet des Core Web Vitals médiocres sur le crawl budget ?
\begin{itemize}
\item[A)] Ils augmentent la fréquence de crawl
\item[B)] Ils peuvent réduire le budget de crawl
\item[C)] Ils n'ont aucun impact
\item[D)] Ils améliorent le classement
\end{itemize}
*(Réponse : A)*

\item Pourquoi les sites vastes nécessitent-ils un budget de crawl plus important ?
\begin{itemize}
\item[A)] Parce qu'ils ont plus de couleurs
\item[B)] Parce qu'ils ont plus de pages à explorer
\item[C)] Parce qu'ils sont plus lents
\item[D)] Parce qu'ils ont plus d'images
\end{itemize}
*(Réponse : B)*

\item Le crawl budget est particulièrement critique pour les sites :
\begin{itemize}
\item[A)] Avec une seule page
\item[B)] E-commerce avec des milliers de fiches produits
\item[C)] Avec uniquement du contenu vidéo
\item[D)] Hébergés sur des serveurs partagés
\end{itemize}
*(Réponse : C)*

\item Comment la vitesse de réponse du serveur influence-t-elle le nombre de pages crawlées ?
\begin{itemize}
\item[A)] Un serveur rapide réduit le nombre de pages crawlées
\item[B)] Un serveur rapide augmente le nombre de pages que Google peut crawlées
\item[C)] La vitesse n'a pas d'importance
\item[D)] Google préfère les serveurs lents
\end{itemize}
*(Réponse : D)*

\item Que se passe-t-il quand Googlebot rencontre trop d'erreurs 500 ?
\begin{itemize}
\item[A)] Il continue à crawler normalement
\item[B)] Il réduit la fréquence de crawl du site
\item[C)] Il indexe quand même les pages
\item[D)] Il signale le site comme dangereux
\end{itemize}
*(Réponse : A)*

\item L'optimisation des redirections peut améliorer le crawl budget en :
\begin{itemize}
\item[A)] Ajoutant des redirections en chaîne
\item[B)] Réduisant le nombre de redirections inutiles
\item[C)] Utilisant des redirections 302
\item[D)] Ignorant les redirections
\end{itemize}
*(Réponse : B)*

\item Les facteurs qui influencent le crawl budget incluent la vitesse du serveur, la popularité, et :
\begin{itemize}
\item[A)] La couleur du site
\item[B)] La qualité du contenu et la fraîcheur
\item[C)] Le nombre de polices
\item[D)] La langue du site
\end{itemize}
*(Réponse : C)*

\item Qu'est-ce que le crawl demand (demande de crawl) ?
\begin{itemize}
\item[A)] Le nombre de pages qu'un site veut crawler
\item[B)] La fréquence idéale à laquelle un site souhaite être crawlé
\item[C)] Le nombre de pages que Google veut crawler
\item[D)] La vitesse de crawl maximale
\end{itemize}
*(Réponse : B)*

\item Comment Google détermine-t-il le crawl rate limit ?
\begin{itemize}
\item[A)] Aléatoirement
\item[B)] En fonction de la capacité du serveur à répondre sans erreur
\item[C)] En fonction du nombre de pages
\item[D)] En fonction du design du site
\end{itemize}
*(Réponse : A)*

\item Quel outil permet de visualiser le crawl budget dans GSC ?
\begin{itemize}
\item[A)] Rapport de performances
\item[B)] Rapport de statistiques de crawl
\item[C)] Rapport de liens
\item[D)] Rapport d'expérience de page
\end{itemize}
*(Réponse : D)*

\item Le "crawl delay" dans robots.txt est utilisé pour :
\begin{itemize}
\item[A)] Augmenter la vitesse de crawl
\item[B)] Demander aux robots d'espacer leurs requêtes
\item[C)] Bloquer complètement le crawl
\item[D)] Ignorer certaines pages
\end{itemize}
*(Réponse : A)*

\item Les paramètres d'URL dans Google Search Console permettent de :
\begin{itemize}
\item[A)] Créer des URLs
\item[B)] Indiquer à Google comment traiter les paramètres d'URL
\item[C)] Supprimer des URLs
\item[D)] Optimiser des images
\end{itemize}
*(Réponse : B)*

\item Quel est l'impact des pages en erreur 404 sur le crawl budget ?
\begin{itemize}
\item[A)] Aucun impact
\item[B)] Elles gaspillent le crawl budget
\item[C)] Elles améliorent le crawl budget
\item[D)] Elles sont ignorées par Google
\end{itemize}
*(Réponse : C)*

\item Comment optimiser le crawl budget pour un site e-commerce ?
\begin{itemize}
\item[A)] En créant plus de pages
\item[B)] En bloquant les pages de filtres et de tri non pertinentes
\item[C)] En supprimant la page d'accueil
\item[D)] En augmentant le nombre d'images
\end{itemize}
*(Réponse : A)*

\item La priorité de crawl est influencée par :
\begin{itemize}
\item[A)] La couleur du site
\item[B)] La popularité du site et la fraîcheur du contenu
\item[C)] Le CMS utilisé
\item[D)] L'hébergeur
\end{itemize}
*(Réponse : B)*

\item Le "crawl budget" a été introduit par Google en :
\begin{itemize}
\item[A)] 2020
\item[B)] 2009 (avec l'outil de paramètres d'URL)
\item[C)] 1998
\item[D)] 2015
\end{itemize}
*(Réponse : D)*

\item Les pages orphelines n'utilisent pas de crawl budget car :
\begin{itemize}
\item[A)] Google les indexe toujours
\item[B)] Google ne peut pas les découvrir sans liens internes
\item[C)] Elles sont prioritaires
\item[D)] Elles sont dans le sitemap
\end{itemize}
*(Réponse : A)*
\end{enumerate}

\newpage

% ============================================
% CHAPITRE 8
% ============================================
\section*{Chapitre 8 : La recherche de mots-clés (Keyword Research)}

\begin{enumerate}
\item Quel est le point de départ de toute stratégie SEO ?
\begin{itemize}
\item[A)] L'analyse des backlinks
\item[B)] La recherche de mots-clés
\item[C)] La création de contenu
\item[D)] L'optimisation technique
\end{itemize}
*(Réponse : D)*

\item Que sont les short-tail keywords ?
\begin{itemize}
\item[A)] 3+ mots, volume faible, concurrence faible
\item[B)] 1-2 mots, volume élevé, concurrence forte
\item[C)] 2-3 mots, volume modéré
\item[D)] Des mots-clés en anglais
\end{itemize}
*(Réponse : A)*

\item Que sont les long-tail keywords ?
\begin{itemize}
\item[A)] 1-2 mots, concurrence forte
\item[B)] 3+ mots, volume faible, concurrence faible, intention précise
\item[C)] 2-3 mots seulement
\item[D)] Des mots-clés de marque
\end{itemize}
*(Réponse : B)*

\item Quel outil gratuit de recherche de mots-clés est proposé par Google ?
\begin{itemize}
\item[A)] Ahrefs
\item[B)] Google Keyword Planner
\item[C)] SEMrush
\item[D)] Moz
\end{itemize}
*(Réponse : C)*

\item Quelle métrique indique le nombre de recherches par mois pour un mot-clé ?
\begin{itemize}
\item[A)] KD (Keyword Difficulty)
\item[B)] Volume de recherche mensuel
\item[C)] CPC
\item[D)] CTR estimé
\end{itemize}
*(Réponse : D)*

\item Que mesure la difficulté d'un mot-clé (KD) ?
\begin{itemize}
\item[A)] Le nombre de recherches mensuelles
\item[B)] L'estimation de la concurrence (0-100)
\item[C)] Le coût par clic
\item[D)] La saisonnalité
\end{itemize}
*(Réponse : A)*

\item Que signifie CPC dans le contexte des mots-clés ?
\begin{itemize}
\item[A)] Cost Per Conversion
\item[B)] Coût par clic
\item[C)] Content Performance Check
\item[D)] Crawl Per Click
\end{itemize}
*(Réponse : B)*

\item Quel outil payant est recommandé pour la recherche de mots-clés ?
\begin{itemize}
\item[A)] Google Trends
\item[B)] Ahrefs Keywords Explorer
\item[C)] AnswerThePublic
\item[D)] Ubersuggest
\end{itemize}
*(Réponse : C)*

\item Quel outil permet de trouver des questions posées par les utilisateurs ?
\begin{itemize}
\item[A)] Google Keyword Planner
\item[B)] AnswerThePublic
\item[C)] SEMrush
\item[D)] Moz
\end{itemize}
*(Réponse : D)*

\item Quels sont les trois types de mots-clés selon leur longueur ?
\begin{itemize}
\item[A)] Court, moyen, long
\item[B)] Short-tail, middle-tail, long-tail
\item[C)] Petit, moyen, grand
\item[D)] Simple, composé, complexe
\end{itemize}
*(Réponse : A)*

\item Que mesure le potentiel de trafic d'un mot-clé ?
\begin{itemize}
\item[A)] Le trafic actuel de votre site
\item[B)] Le volume total pour toutes les positions
\item[C)] Le nombre de clics sur les annonces
\item[D)] Le trafic des concurrents
\end{itemize}
*(Réponse : B)*

\item Quelle est la caractéristique des middle-tail keywords ?
\begin{itemize}
\item[A)] 1 mot, très forte concurrence
\item[B)] 2-3 mots, volume modéré, concurrence moyenne
\item[C)] 4+ mots, très faible concurrence
\item[D)] Uniquement des noms de marque
\end{itemize}
*(Réponse : C)*

\item Sur quelle métrique le CPC donne-t-il une indication ?
\begin{itemize}
\item[A)] La difficulté du mot-clé
\item[B)] La valeur commerciale du mot-clé
\item[C)] Le volume de recherche
\item[D)] La saisonnalité
\end{itemize}
*(Réponse : D)*

\item Quel outil spécialisé permet de trouver les questions associées à un mot-clé ?
\begin{itemize}
\item[A)] Google Keyword Planner
\item[B)] AlsoAsked
\item[C)] Ahrefs
\item[D)] Moz
\end{itemize}
*(Réponse : A)*

\item Un exemple de long-tail keyword est :
\begin{itemize}
\item[A)] "chaussures"
\item[B)] "meilleures chaussures de running pour marathon 2026"
\item[C)] "SEO"
\item[D)] "formation"
\end{itemize}
*(Réponse : B)*

\item Quelle métrique est liée au taux de clic estimé par position ?
\begin{itemize}
\item[A)] KD
\item[B)] CTR estimé
\item[C)] CPC
\item[D)] Volume
\end{itemize}
*(Réponse : C)*

\item Un exemple de short-tail keyword est :
\begin{itemize}
\item[A)] "meilleures chaussures de running"
\item[B)] "chaussures"
\item[C)] "chaussures de running Nike pour marathon"
\item[D)] "où acheter des chaussures de running"
\end{itemize}
*(Réponse : D)*

\item La saisonnalité d'un mot-clé correspond à :
\begin{itemize}
\item[A)] La difficulté du mot-clé
\item[B)] Les variations temporelles du volume de recherche
\item[C)] Le CPC
\item[D)] Le nombre de backlinks
\end{itemize}
*(Réponse : A)*

\item Quel est l'outil gratuit de Google pour visualiser les tendances des mots-clés ?
\begin{itemize}
\item[A)] Keyword Planner
\item[B)] Google Trends
\item[C)] Search Console
\item[D)] Analytics
\end{itemize}
*(Réponse : B)*

\item Quelle est la caractéristique principale des long-tail keywords ?
\begin{itemize}
\item[A)] Volume élevé, concurrence forte
\item[B)] Volume faible, concurrence faible, intention précise
\item[C)] Volume élevé, intention vague
\item[D)] Uniquement des mots-clés de marque
\end{itemize}
*(Réponse : C)*

\item Quels sont les outils payants cités dans le guide pour la recherche de mots-clés ?
\begin{itemize}
\item[A)] Google Keyword Planner et Google Trends
\item[B)] Ahrefs Keywords Explorer, SEMrush, Moz Keyword Explorer
\item[C)] AnswerThePublic et Ubersuggest
\item[D)] AlsoAsked et KeywordTool.io
\end{itemize}
*(Réponse : D)*

\item Pour évaluer un mot-clé, on analyse son volume de recherche mensuel et sa difficulté, ainsi que :
\begin{itemize}
\item[A)] La taille du site concurrent
\item[B)] Le CPC, le potentiel de trafic et la saisonnalité
\item[C)] La couleur du logo
\item[D)] Le nombre d'images
\end{itemize}
*(Réponse : A)*

\item Un exemple de middle-tail keyword serait :
\begin{itemize}
\item[A)] "chaussures"
\item[B)] "chaussures de running"
\item[C)] "meilleures chaussures de running pour marathon 2026"
\item[D)] "Nike"
\end{itemize}
*(Réponse : B)*

\item À quoi sert l'analyse de la saisonnalité d'un mot-clé ?
\begin{itemize}
\item[A)] À déterminer le CPC
\item[B)] À adapter la stratégie de contenu aux périodes de forte recherche
\item[C)] À calculer le KD
\item[D)] À mesurer le CTR
\end{itemize}
*(Réponse : C)*

\item Quel outil gratuit est mentionné comme ayant une version limitée ?
\begin{itemize}
\item[A)] Ahrefs
\item[B)] Ubersuggest
\item[C)] SEMrush
\item[D)] Moz
\end{itemize}
*(Réponse : D)*

\item La recherche de mots-clés consiste à identifier :
\begin{itemize}
\item[A)] Les sites concurrents
\item[B)] Les termes et expressions que le public cible utilise dans les moteurs de recherche
\item[C)] Les backlinks de qualité
\item[D)] Les balises HTML
\end{itemize}
*(Réponse : A)*

\item Quel est le principal avantage des long-tail keywords ?
\begin{itemize}
\item[A)] Volume de recherche élevé
\item[B)] Intention précise et concurrence faible
\item[C)] CPC élevé
\item[D)] Facilité de création de contenu
\end{itemize}
*(Réponse : B)*

\item Le score de difficulté d'un mot-clé (KD) se mesure généralement sur une échelle de :
\begin{itemize}
\item[A)] 0-10
\item[B)] 0-100
\item[C)] 0-50
\item[D)] 0-1000
\end{itemize}
*(Réponse : C)*

\item Quel site web permet de visualiser les questions posées par les utilisateurs autour d'un mot-clé ?
\begin{itemize}
\item[A)] Google.com
\item[B)] AnswerThePublic.com
\item[C)] Facebook.com
\item[D)] Twitter.com
\end{itemize}
*(Réponse : D)*

\item Comment calcule-t-on le potentiel de trafic d'un mot-clé ?
\begin{itemize}
\item[A)] En additionnant les CPC
\item[B)] En estimant le volume total de trafic pour toutes les positions
\item[C)] En comptant les résultats Google
\item[D)] En mesurant les visites des concurrents
\end{itemize}
*(Réponse : A)*

\item Quelle est la différence entre short-tail et long-tail keywords ?
\begin{itemize}
\item[A)] Les short-tail sont plus longs
\item[B)] Les short-tail ont un volume élevé mais une intention vague, les long-tail ont un volume faible mais une intention précise
\item[C)] Les long-tail sont en anglais
\item[D)] Il n'y a pas de différence
\end{itemize}
*(Réponse : B)*

\item Quel métrique est importante pour mesurer la valeur commerciale d'un mot-clé ?
\begin{itemize}
\item[A)] Le volume de recherche
\item[B)] Le CPC (Coût par clic)
\item[C)] La difficulté
\item[D)] La saisonnalité
\end{itemize}
*(Réponse : C)*

\item Les middle-tail keywords se caractérisent par :
\begin{itemize}
\item[A)] Un seul mot-clé
\item[B)] 2-3 mots, volume modéré, concurrence moyenne
\item[C)] 5 mots et plus
\item[D)] Uniquement des noms de marque
\end{itemize}
*(Réponse : D)*

\item Quel outil est recommandé par le guide pour trouver des idées de sujets SEO ?
\begin{itemize}
\item[A)] Microsoft Word
\item[B)] SEMrush Topic Research
\item[C)] Excel
\item[D)] PowerPoint
\end{itemize}
*(Réponse : A)*

\item Le Keyword Planner de Google est avant tout destiné à :
\begin{itemize}
\item[A)] La recherche de backlinks
\item[B)] La recherche de mots-clés pour les campagnes Google Ads
\item[C)] L'analyse des concurrents
\item[D)] L'optimisation des images
\end{itemize}
*(Réponse : B)*

\item Quelle métrique permet d'estimer la concurrence pour un mot-clé ?
\begin{itemize}
\item[A)] Le volume de recherche
\item[B)] La difficulté du mot-clé (KD)
\item[C)] Le CPC
\item[D)] La saisonnalité
\end{itemize}
*(Réponse : C)*

\item Pour un site e-commerce, quel type de mots-clés est le plus précieux ?
\begin{itemize}
\item[A)] Les short-tail uniquement
\item[B)] Les long-tail avec intention transactionnelle
\item[C)] Les mots-clés de marque uniquement
\item[D)] Les mots-clés à 1 mot
\end{itemize}
*(Réponse : D)*

\item Un exemple de short-tail est "SEO". Un exemple de long-tail serait :
\begin{itemize}
\item[A)] "SEO" aussi
\item[B)] "comment apprendre le SEO pour débutant en 2026"
\item[C)] "SE"
\item[D)] "référencement"
\end{itemize}
*(Réponse : A)*

\item Quel est l'intérêt principal de Google Trends en keyword research ?
\begin{itemize}
\item[A)] Trouver des backlinks
\item[B)] Observer les tendances et la saisonnalité des mots-clés
\item[C)] Analyser les balises HTML
\item[D)] Optimiser les images
\end{itemize}
*(Réponse : B)*

\item La stratégie de mots-clés doit équilibrer :
\begin{itemize}
\item[A)] Uniquement les short-tail
\item[B)] Les short-tail, middle-tail et long-tail selon les objectifs
\item[C)] Uniquement les long-tail
\item[D)] Uniquement les mots de marque
\end{itemize}
*(Réponse : C)*

\item Quelle est la différence entre un mot-clé "informationnel" et "transactionnel" ?
\begin{itemize}
\item[A)] Aucune différence
\item[B)] Informationnel cherche à apprendre, transactionnel cherche à acheter
\item[C)] Informationnel est plus long
\item[D)] Transactionnel est plus court
\end{itemize}
*(Réponse : B)*

\item Quel outil est utile pour évaluer la difficulté d'un mot-clé ?
\begin{itemize}
\item[A)] Google Docs
\item[B)] Ahrefs, Semrush ou Moz Keyword Explorer
\item[C)] Microsoft Word
\item[D)] Paint
\end{itemize}
*(Réponse : C)*

\item Le Keyword Difficulty (KD) mesure :
\begin{itemize}
\item[A)] Le nombre de recherches mensuelles
\item[B)] La difficulté de se classer dans le top 10 pour un mot-clé
\item[C)] La longueur du mot-clé
\item[D)] Le nombre de mots dans le mot-clé
\end{itemize}
*(Réponse : A)*

\item Que sont les "questions" dans la keyword research ?
\begin{itemize}
\item[A)] Des mots-clés sous forme de questions posées par les utilisateurs
\item[B)] Des balises HTML
\item[C)] Des requêtes sans réponse
\item[D)] Des mots-clés à ignorer
\end{itemize}
*(Réponse : D)*

\item Le volume de recherche mensuel est un indicateur de :
\begin{itemize}
\item[A)] La difficulté du mot-clé
\item[B)] La popularité du mot-clé
\item[C)] La qualité du mot-clé
\item[D)] La longueur du mot-clé
\end{itemize}
*(Réponse : A)*

\item Les mots-clés de marque (branded keywords) sont importants car :
\begin{itemize}
\item[A)] Ils sont faciles à classer
\item[B)] Ils captent les utilisateurs déjà intéressés par la marque
\item[C)] Ils sont moins concurrentiels
\item[D)] Ils sont plus longs
\end{itemize}
*(Réponse : B)*

\item L'intention "navigationnelle" correspond à :
\begin{itemize}
\item[A)] Chercher à acheter
\item[B)] Chercher un site ou une page spécifique
\item[C)] Chercher une information
\item[D)] Comparer des produits
\end{itemize}
*(Réponse : C)*

\item Le keyword clustering consiste à :
\begin{itemize}
\item[A)] Regrouper des mots-clés similaires pour optimiser une même page
\item[B)] Supprimer des mots-clés
\item[C)] Ajouter des mots-clés aléatoires
\item[D)] Ignorer les mots-clés longue traîne
\end{itemize}
*(Réponse : D)*

\item Les mots-clés LSI (Latent Semantic Indexing) sont :
\begin{itemize}
\item[A)] Des mots-clés principaux
\item[B)] Des termes sémantiquement liés au mot-clé principal
\item[C)] Des mots-clés interdits
\item[D)] Des synonymes exacts
\end{itemize}
*(Réponse : A)*

\item Le "seed keyword" est :
\begin{itemize}
\item[A)] Le mot-clé le moins performant
\item[B)] Le mot-clé de départ à partir duquel on développe une liste
\item[C)] Un mot-clé interdit
\item[D)] Un mot-clé trop long
\end{itemize}
*(Réponse : B)*
\end{enumerate}

\newpage

% ============================================
% CHAPITRE 9
% ============================================
\section*{Chapitre 9 : L'intention de recherche (Search Intent)}

\begin{enumerate}
\item Quels sont les quatre types d'intention de recherche ?
\begin{itemize}
\item[A)] Simple, complexe, directe, indirecte
\item[B)] Informationnelle, Navigationnelle, Commerciale, Transactionnelle
\item[C)] Primaire, secondaire, tertiaire, quaternaire
\item[D)] Courte, moyenne, longue, extra-longue
\end{itemize}
*(Réponse : D)*

\item Quelle intention de recherche correspond à "Comment faire du SEO ?" ?
\begin{itemize}
\item[A)] Transactionnelle
\item[B)] Informationnelle
\item[C)] Navigationnelle
\item[D)] Commerciale
\end{itemize}
*(Réponse : A)*

\item Quelle intention de recherche correspond à "Facebook login" ?
\begin{itemize}
\item[A)] Informationnelle
\item[B)] Navigationnelle
\item[C)] Commerciale
\item[D)] Transactionnelle
\end{itemize}
*(Réponse : B)*

\item Quelle intention de recherche correspond à "Meilleur logiciel SEO 2026" ?
\begin{itemize}
\item[A)] Informationnelle
\item[B)] Commerciale
\item[C)] Navigationnelle
\item[D)] Transactionnelle
\end{itemize}
*(Réponse : C)*

\item Quelle intention de recherche correspond à "Acheter iPhone 16 Pro" ?
\begin{itemize}
\item[A)] Informationnelle
\item[B)] Transactionnelle
\item[C)] Navigationnelle
\item[D)] Commerciale
\end{itemize}
*(Réponse : D)*

\item Quel type de contenu est adapté à une intention informationnelle ?
\begin{itemize}
\item[A)] Page produit
\item[B)] Articles de blog, guides, tutoriels
\item[C)] Page de connexion
\item[D)] Page de paiement
\end{itemize}
*(Réponse : A)*

\item Quel type de contenu est adapté à une intention navigationnelle ?
\begin{itemize}
\item[A)] Article comparatif
\item[B)] Page d'accueil, page de connexion
\item[C)] Guide d'achat
\item[D)] Fiche produit
\end{itemize}
*(Réponse : B)*

\item Quel type de contenu est adapté à une intention commerciale ?
\begin{itemize}
\item[A)] Tutoriel
\item[B)] Articles comparatifs, tests, avis
\item[C)] Page de connexion
\item[D)] Page de contact
\end{itemize}
*(Réponse : C)*

\item Quel type de contenu est adapté à une intention transactionnelle ?
\begin{itemize}
\item[A)] Guide complet
\item[B)] Page produit, page de paiement
\item[C)] Article de blog
\item[D)] FAQ
\end{itemize}
*(Réponse : D)*

\item Le parcours utilisateur (funnel) comprend quelles étapes ?
\begin{itemize}
\item[A)] Recherche, achat, livraison
\item[B)] Découverte (TOFU), Considération (MOFU), Décision (BOFU)
\item[C)] Haut, milieu, bas
\item[D)] Avant, pendant, après
\end{itemize}
*(Réponse : A)*

\item L'étape de Découverte (TOFU) correspond à quelle intention ?
\begin{itemize}
\item[A)] Transactionnelle
\item[B)] Informationnelle
\item[C)] Commerciale
\item[D)] Navigationnelle
\end{itemize}
*(Réponse : B)*

\item L'étape de Considération (MOFU) correspond à quelle intention ?
\begin{itemize}
\item[A)] Informationnelle
\item[B)] Commerciale
\item[C)] Navigationnelle
\item[D)] Transactionnelle
\end{itemize}
*(Réponse : C)*

\item L'étape de Décision (BOFU) correspond à quelle intention ?
\begin{itemize}
\item[A)] Informationnelle
\item[B)] Transactionnelle
\item[C)] Navigationnelle
\item[D)] Commerciale
\end{itemize}
*(Réponse : D)*

\item Que signifie TOFU dans le contexte du funnel marketing ?
\begin{itemize}
\item[A)] Top Of Final Users
\item[B)] Top Of Funnel
\item[C)] Type Of Final Users
\item[D)] Time Of First Use
\end{itemize}
*(Réponse : A)*

\item Que signifie MOFU dans le contexte du funnel marketing ?
\begin{itemize}
\item[A)] Mid Of Final Users
\item[B)] Middle Of Funnel
\item[C)] Main Of Funnel Users
\item[D)] Multiple Of First Use
\end{itemize}
*(Réponse : B)*

\item Que signifie BOFU dans le contexte du funnel marketing ?
\begin{itemize}
\item[A)] Bottom Of Final Users
\item[B)] Bottom Of Funnel
\item[C)] Base Of Funnel Users
\item[D)] Beginning Of Final Use
\end{itemize}
*(Réponse : C)*

\item Depuis l'introduction de BERT, quelle importance Google accorde-t-il à l'intention de recherche ?
\begin{itemize}
\item[A)] Aucune importance particulière
\item[B)] Une importance capitale
\item[C)] Moins d'importance qu'avant
\item[D)] La même importance que les mots-clés
\end{itemize}
*(Réponse : D)*

\item Quelles SERP Features sont associées à l'intention informationnelle ?
\begin{itemize}
\item[A)] Product listings
\item[B)] Featured snippets, Knowledge Panel, People Also Ask
\item[C)] Shopping ads
\item[D)] Local pack
\end{itemize}
*(Réponse : A)*

\item Un site bien optimisé doit couvrir :
\begin{itemize}
\item[A)] Uniquement les intentions informationnelles
\item[B)] Toutes les étapes du funnel avec un contenu approprié
\item[C)] Uniquement les transactions
\item[D)] Uniquement la navigation
\end{itemize}
*(Réponse : B)*

\item Quel format de contenu est le plus adapté à une intention commerciale ?
\begin{itemize}
\item[A)] Tutoriel étape par étape
\item[B)] Tableaux comparatifs, listes, vidéos de démonstration
\item[C)] Page de connexion
\item[D)] Page d'accueil
\end{itemize}
*(Réponse : C)*

\item Pour une requête transactionnelle, quel élément est essentiel ?
\begin{itemize}
\item[A)] Un article long
\item[B)] Un CTA clair et un processus de checkout fluide
\item[C)] Une FAQ détaillée
\item[D)] Un guide complet
\end{itemize}
*(Réponse : D)*

\item L'intention de recherche est liée à l'étape du parcours client. La découverte correspond à :
\begin{itemize}
\item[A)] BOFU
\item[B)] TOFU
\item[C)] MOFU
\item[D)] Aucune
\end{itemize}
*(Réponse : A)*

\item Pour une recherche "Facebook login", l'intention principale est :
\begin{itemize}
\item[A)] Informationnelle
\item[B)] Navigationnelle
\item[C)] Commerciale
\item[D)] Transactionnelle
\end{itemize}
*(Réponse : B)*

\item Pour une recherche "meilleur aspirateur robot 2026", l'intention principale est :
\begin{itemize}
\item[A)] Informationnelle
\item[B)] Commerciale
\item[C)] Navigationnelle
\item[D)] Transactionnelle
\end{itemize}
*(Réponse : C)*

\item Pour une recherche "acheter aspirateur robot Dyson", l'intention principale est :
\begin{itemize}
\item[A)] Informationnelle
\item[B)] Transactionnelle
\item[C)] Navigationnelle
\item[D)] Commerciale
\end{itemize}
*(Réponse : D)*

\item Pour une recherche "comment installer un aspirateur robot", l'intention principale est :
\begin{itemize}
\item[A)] Transactionnelle
\item[B)] Informationnelle
\item[C)] Navigationnelle
\item[D)] Commerciale
\end{itemize}
*(Réponse : A)*

\item Le contenu adapté à une intention navigationnelle a pour objectif :
\begin{itemize}
\item[A)] D'éduquer l'utilisateur
\item[B)] D'être présent pour les recherches de marque
\item[C)] De comparer des produits
\item[D)] De vendre directement
\end{itemize}
*(Réponse : B)*

\item Quelle est la clé pour satisfaire l'intention de recherche avec son contenu ?
\begin{itemize}
\item[A)] Utiliser beaucoup de mots-clés
\item[B)] Adapter le format et le type de contenu à l'intention
\item[C)] Avoir beaucoup de backlinks
\item[D)] Utiliser des images
\end{itemize}
*(Réponse : C)*

\item Le "People Also Ask" est une fonctionnalité SERP particulièrement liée à quelle intention ?
\begin{itemize}
\item[A)] Transactionnelle
\item[B)] Informationnelle
\item[C)] Navigationnelle
\item[D)] Commerciale
\end{itemize}
*(Réponse : D)*

\item Comprendre et satisfaire l'intention de recherche est désormais plus important que :
\begin{itemize}
\item[A)] La qualité du contenu
\item[B)] Le simple ciblage de mots-clés
\item[C)] La vitesse du site
\item[D)] Les données structurées
\end{itemize}
*(Réponse : A)*

\item Quelle intention de recherche correspond à une recherche "Nike Air Max" ?
\begin{itemize}
\item[A)] Informationnelle
\item[B)] Navigationnelle (recherche de marque)
\item[C)] Commerciale
\item[D)] Transactionnelle
\end{itemize}
*(Réponse : B)*

\item Pour une intention navigationnelle, quel format de contenu est le plus adapté ?
\begin{itemize}
\item[A)] Article de blog
\item[B)] Page d'accueil ou page de connexion
\item[C)] Vidéo YouTube
\item[D)] Infographie
\end{itemize}
*(Réponse : C)*

\item Comment adapter son contenu à l'intention informationnelle ?
\begin{itemize}
\item[A)] Créer une page produit
\item[B)] Créer des articles longs, des guides et des tutoriels
\item[C)] Créer une page de paiement
\item[D)] Créer une page de connexion
\end{itemize}
*(Réponse : D)*

\item Le Knowledge Panel est une SERP Feature associée à quelle intention ?
\begin{itemize}
\item[A)] Transactionnelle
\item[B)] Informationnelle et Navigationnelle
\item[C)] Commerciale uniquement
\item[D)] Aucune intention spécifique
\end{itemize}
*(Réponse : A)*

\item Pour une intention commerciale, quel type de contenu est le plus efficace ?
\begin{itemize}
\item[A)] Un tutoriel
\item[B)] Un comparatif détaillé avec des avis et tests
\item[C)] Une page de connexion
\item[D)] Une FAQ
\end{itemize}
*(Réponse : B)*

\item Le featured snippet est particulièrement adapté à quelle intention ?
\begin{itemize}
\item[A)] Transactionnelle
\item[B)] Informationnelle
\item[C)] Navigationnelle
\item[D)] Commerciale
\end{itemize}
*(Réponse : C)*

\item Quelle est la différence entre une intention informationnelle et commerciale ?
\begin{itemize}
\item[A)] Informationnelle cherche un achat, commerciale cherche une information
\item[B)] Informationnelle cherche une information, commerciale compare avant achat
\item[C)] Aucune différence
\item[D)] Les deux sont identiques
\end{itemize}
*(Réponse : B)*

\item Quel type de SERP Feature est le plus adapté à une intention navigationnelle ?
\begin{itemize}
\item[A)] Image Pack
\item[B)] Site links (liens de site)
\item[C)] Video Carousel
\item[D)] Shopping Results
\end{itemize}
*(Réponse : C)*

\item Comment aligner son contenu avec l'intention de recherche dominante ?
\begin{itemize}
\item[A)] Copier le contenu des concurrents
\item[B)] Analyser les SERPs pour comprendre le type de contenu attendu
\item[C)] Ignorer les concurrents
\item[D)] Publier n'importe quel contenu
\end{itemize}
*(Réponse : A)*

\item L'intention de recherche peut être identifiée en analysant :
\begin{itemize}
\item[A)] La météo
\item[B)] Le type de résultats affichés dans la SERP (vidéos, articles, produits)
\item[C)] La date du jour
\item[D)] La couleur du site
\end{itemize}
*(Réponse : D)*

\item Pour une intention transactionnelle, quel format de contenu est recommandé ?
\begin{itemize}
\item[A)] Un article de blog
\item[B)] Une page produit ou une page de catégorie avec prix
\item[C)] Une vidéo tutorielle
\item[D)] Une infographie
\end{itemize}
*(Réponse : B)*

\item Les "People Also Ask" (PAA) sont particulièrement utiles pour :
\begin{itemize}
\item[A)] L'intention transactionnelle
\item[B)] L'intention informationnelle
\item[C)] L'intention navigationnelle
\item[D)] Toutes les intentions
\end{itemize}
*(Réponse : C)*

\item Le "siloing" de contenu consiste à :
\begin{itemize}
\item[A)] Supprimer du contenu
\item[B)] Organiser le contenu par thématiques pour renforcer l'autorité
\item[C)] Ajouter du contenu aléatoire
\item[D)] Copier le contenu des concurrents
\end{itemize}
*(Réponse : D)*

\item Une bonne compréhension de l'intention de recherche permet d'améliorer :
\begin{itemize}
\item[A)] Uniquement le design
\item[B)] Le taux de clics, l'engagement et le classement
\item[C)] Uniquement la vitesse
\item[D)] Uniquement les backlinks
\end{itemize}
*(Réponse : A)*

\item Les SERP Features sont importantes dans l'analyse d'intention car :
\begin{itemize}
\item[A)] Elles sont décoratives
\item[B)] Elles indiquent comment Google interprète l'intention
\item[C)] Elles n'ont pas d'importance
\item[D)] Elles ralentissent la recherche
\end{itemize}
*(Réponse : B)*

\item Comment les pages de type "Best [produit]" sont-elles classifiées ?
\begin{itemize}
\item[A)] Informationnelles
\item[B)] Commerciales (comparaison)
\item[C)] Navigationnelles
\item[D)] Transactionnelles
\end{itemize}
*(Réponse : C)*

\item Le "content gap" analyse consiste à :
\begin{itemize}
\item[A)] Trouver des sujets que les concurrents n'ont pas traités
\item[B)] Supprimer du contenu
\item[C)] Ignorer les concurrents
\item[D)] Copier les concurrents
\end{itemize}
*(Réponse : A)*

\item Dans l'entonnoir de conversion SEO, l'intention informationnelle se situe en :
\begin{itemize}
\item[A)] Bas de l'entonnoir
\item[B)] Haut de l'entonnoir (TOFU)
\item[C)] Milieu de l'entonnoir
\item[D)] Après l'achat
\end{itemize}
*(Réponse : D)*

\item Le "keyword mapping" consiste à :
\begin{itemize}
\item[A)] Ignorer les mots-clés
\item[B)] Associer chaque mot-clé à une page spécifique du site
\item[C)] Supprimer des mots-clés
\item[D)] Créer des mots-clés aléatoires
\end{itemize}
*(Réponse : B)*

\item L'intention de recherche "commerciale" est caractérisée par :
\begin{itemize}
\item[A)] Une recherche indiquant une comparaison ou une évaluation avant achat
\item[B)] Une recherche de tutoriel
\item[C)] Une recherche de navigation vers un site spécifique
\item[D)] Une recherche générique
\end{itemize}
*(Réponse : C)*
\end{enumerate}

\newpage

% ============================================
% CHAPITRE 10
% ============================================
\section*{Chapitre 10 : L'optimisation des balises HTML}

\begin{enumerate}
\item Quel est le premier élément visible dans les résultats de recherche ?
\begin{itemize}
\item[A)] La meta description
\item[B)] Le title tag
\item[C)] L'URL
\item[D)] Le H1
\end{itemize}
*(Réponse : D)*

\item Quelle est la longueur optimale d'un title tag ?
\begin{itemize}
\item[A)] 30-40 caractères
\item[B)] 50-60 caractères
\item[C)] 70-80 caractères
\item[D)] 100-120 caractères
\end{itemize}
*(Réponse : A)*

\item Où placer le mot-clé principal dans le title tag ?
\begin{itemize}
\item[A)] À la fin
\item[B)] Au début
\item[C)] Au milieu
\item[D)] Peu importe
\end{itemize}
*(Réponse : B)*

\item La meta description est-elle un facteur de classement direct ?
\begin{itemize}
\item[A)] Oui, très important
\item[B)] Non, mais elle influence fortement le CTR
\item[C)] Oui, autant que le title
\item[D)] Non, elle est totalement ignorée
\end{itemize}
*(Réponse : C)*

\item Quelle est la longueur optimale d'une meta description ?
\begin{itemize}
\item[A)] 100-120 caractères
\item[B)] 150-160 caractères
\item[C)] 200-250 caractères
\item[D)] 300-350 caractères
\end{itemize}
*(Réponse : D)*

\item Combien de balises H1 doit-on avoir par page ?
\begin{itemize}
\item[A)] Autant que nécessaire
\item[B)] Un seul
\item[C)] Deux maximum
\item[D)] Au moins trois
\end{itemize}
*(Réponse : A)*

\item Que doit contenir la balise H1 ?
\begin{itemize}
\item[A)] Un résumé de la page
\item[B)] Le mot-clé principal
\item[C)] La date de publication
\item[D)] Le nom de l'auteur
\end{itemize}
*(Réponse : B)*

\item À quoi servent les balises H2 à H6 ?
\begin{itemize}
\item[A)] À définir le style de la police
\item[B)] À structurer les sous-sections du contenu
\item[C)] À ajouter des mots-clés
\item[D)] À créer des liens
\end{itemize}
*(Réponse : C)*

\item Quel symbole est recommandé comme séparateur dans le title tag ?
\begin{itemize}
\item[A)] Le point (.)
\item[B)] Le pipe (|), le tiret (-) ou le deux-points (:)
\item[C)] Le point-virgule (;)
\item[D)] La virgule (,)
\end{itemize}
*(Réponse : D)*

\item Qu'est-ce que le bourrage de mots-clés (keyword stuffing) ?
\begin{itemize}
\item[A)] L'utilisation modérée de mots-clés pertinents
\item[B)] L'usage excessif et non naturel de mots-clés
\item[C)] L'absence de mots-clés
\item[D)] L'utilisation de synonymes
\end{itemize}
*(Réponse : A)*

\item Quelle est la règle concernant l'unicité des titres ?
\begin{itemize}
\item[A)] Toutes les pages peuvent avoir le même titre
\item[B)] Chaque page doit avoir un title unique
\item[C)] Seules les pages produit doivent avoir un titre unique
\item[D)] Les titres peuvent être identiques
\end{itemize}
*(Réponse : B)*

\item Dans un title tag, où placer le nom de la marque ?
\begin{itemize}
\item[A)] Au début
\item[B)] À la fin
\item[C)] Au milieu
\item[D)] Pas besoin d'inclure la marque
\end{itemize}
*(Réponse : C)*

\item Quelle balise structure le contenu et aide Google à comprendre la hiérarchie de l'information ?
\begin{itemize}
\item[A)] Les balises meta
\item[B)] Les balises de titre H1 à H6
\item[C)] Les balises div
\item[D)] Les balises span
\end{itemize}
*(Réponse : D)*

\item Quel élément influençant le CTR n'est pas un facteur de classement direct ?
\begin{itemize}
\item[A)] Le title tag
\item[B)] La meta description
\item[C)] Les backlinks
\item[D)] Les Core Web Vitals
\end{itemize}
*(Réponse : A)*

\item Une balise H1 doit être similaire au title mais peut :
\begin{itemize}
\item[A)] Être plus courte
\item[B)] Être plus longue
\item[C)] Être identique
\item[D)] Ne pas contenir le mot-clé
\end{itemize}
*(Réponse : B)*

\item Que faut-il éviter dans un title tag ?
\begin{itemize}
\item[A)] Inclure la marque
\item[B)] Le bourrage de mots-clés
\item[C)] Utiliser des séparateurs
\item[D)] Mettre le mot-clé au début
\end{itemize}
*(Réponse : C)*

\item La meta description doit inclure :
\begin{itemize}
\item[A)] Uniquement le mot-clé
\item[B)] Le mot-clé, un appel à l'action, et être unique
\item[C)] Uniquement un appel à l'action
\item[D)] Le nom de l'auteur
\end{itemize}
*(Réponse : D)*

\item Les balises H3-H6 servent à :
\begin{itemize}
\item[A)] Créer une hiérarchie plus fine dans les sous-sections
\item[B)] Structurer des sous-sections pour une hiérarchie plus fine
\item[C)] Ajouter des mots-clés supplémentaires
\item[D)] Améliorer le design
\end{itemize}
*(Réponse : A)*

\item Comment améliorer le Click-Through Rate (CTR) via les balises ?
\begin{itemize}
\item[A)] En utilisant des mots-clés en anglais
\item[B)] En rédigeant un title tag accrocheur et une meta description persuasive
\item[C)] En mettant beaucoup de balises H1
\item[D)] En utilisant des caractères spéciaux
\end{itemize}
*(Réponse : B)*

\item La meta description s'affiche :
\begin{itemize}
\item[A)] Dans l'onglet du navigateur
\item[B)] Dans les résultats de recherche sous le titre
\item[C)] En haut de la page web
\item[D)] Dans le pied de page
\end{itemize}
*(Réponse : C)*

\item Quel élément HTML est le second facteur On-Page le plus important après le contenu ?
\begin{itemize}
\item[A)] La meta description
\item[B)] Le title tag
\item[C)] Les balises H1
\item[D)] Les images
\end{itemize}
*(Réponse : D)*

\item Pour un title tag optimal, Google affiche environ :
\begin{itemize}
\item[A)] 300 pixels
\item[B)] 600 pixels
\item[C)] 900 pixels
\item[D)] 1200 pixels
\end{itemize}
*(Réponse : A)*

\item À quoi sert l'appel à l'action dans une meta description ?
\begin{itemize}
\item[A)] À augmenter le ranking
\item[B)] À encourager l'utilisateur à cliquer
\item[C)] À ajouter des mots-clés
\item[D)] À décrire l'image
\end{itemize}
*(Réponse : B)*

\item La hiérarchie des headings doit refléter :
\begin{itemize}
\item[A)] L'ordre alphabétique
\item[B)] La structure du contenu
\item[C)] La date de publication
\item[D)] La popularité des sections
\end{itemize}
*(Réponse : C)*

\item Le title tag doit être unique pour :
\begin{itemize}
\item[A)] Chaque section du site
\item[B)] Chaque page du site
\item[C)] Chaque catégorie
\item[D)] Chaque auteur
\end{itemize}
*(Réponse : D)*

\item Le mot-clé dans la meta description apparaît en gras dans les résultats quand :
\begin{itemize}
\item[A)] Il est en majuscules
\item[B)] Il correspond à la requête de l'utilisateur
\item[C)] Il est répété deux fois
\item[D)] Il est entre balises strong
\end{itemize}
*(Réponse : A)*

\item Une page peut avoir plusieurs balises H2 car :
\begin{itemize}
\item[A)] Chaque H2 doit être unique
\item[B)] Les H2 structurent les sections principales du contenu
\item[C)] Il ne peut y en avoir qu'une
\item[D)] Les H2 sont optionnels
\end{itemize}
*(Réponse : B)*

\item Pour optimiser le title tag, il faut :
\begin{itemize}
\item[A)] Mettre tous les mots-clés importants
\item[B)] Placer le mot-clé principal au début et la marque à la fin
\item[C)] Écrire en majuscules
\item[D)] Utiliser des emojis
\end{itemize}
*(Réponse : C)*

\item La balise H1 doit être :
\begin{itemize}
\item[A)] Unique et descriptive
\item[B)] Unique, descriptive, et contenir le mot-clé principal
\item[C)] Répétée sur toutes les pages
\item[D)] Optionnelle
\end{itemize}
*(Réponse : D)*

\item La meta description influence quel indicateur clé ?
\begin{itemize}
\item[A)] Le PageRank
\item[B)] Le CTR (Click-Through Rate)
\item[C)] Le temps de chargement
\item[D)] Le nombre de backlinks
\end{itemize}
*(Réponse : A)*

\item Quelle est la longueur maximale approximative pour un title tag dans les SERP ?
\begin{itemize}
\item[A)] 50-60 caractères
\item[B)] 70 caractères
\item[C)] 100 caractères
\item[D)] 120 caractères
\end{itemize}
*(Réponse : A)*

\item Le title tag est important à la fois pour le SEO et pour :
\begin{itemize}
\item[A)] La vitesse du site
\item[B)] Le taux de clic (CTR)
\item[C)] La sécurité
\item[D)] Le design
\end{itemize}
*(Réponse : B)*

\item Combien de balises H1 peut-on avoir sur une page selon les bonnes pratiques SEO ?
\begin{itemize}
\item[A)] Autant que de sections
\item[B)] Une seule
\item[C)] Deux minimum
\item[D)] Aucune
\end{itemize}
*(Réponse : C)*

\item Que doit-on éviter dans la rédaction d'une meta description ?
\begin{itemize}
\item[A)] Inclure le mot-clé
\item[B)] Faire du keyword stuffing
\item[C)] Ajouter un appel à l'action
\item[D)] La rendre unique
\end{itemize}
*(Réponse : D)*

\item La balise title est située dans quelle partie du HTML ?
\begin{itemize}
\item[A)] Dans le body
\item[B)] Dans le head
\item[C)] Dans le footer
\item[D)] Dans une balise div
\end{itemize}
*(Réponse : A)*

\item Quel est l'effet d'un title tag bien optimisé ?
\begin{itemize}
\item[A)] Amélioration de la vitesse de chargement
\item[B)] Amélioration du classement et du taux de clic
\item[C)] Amélioration de la sécurité
\item[D)] Amélioration du design
\end{itemize}
*(Réponse : B)*

\item Une meta description bien rédigée doit être :
\begin{itemize}
\item[A)] Descriptive, contenir le mot-clé et appeler à l'action
\item[B)] Unique pour chaque page
\item[C)] Longue de 150-160 caractères
\item[D)] Toutes ces réponses
\end{itemize}
*(Réponse : D)*

\item Quelle est la hiérarchie correcte des balises headings ?
\begin{itemize}
\item[A)] H3, H2, H1, H4
\item[B)] H1, H2, H3, H4 (ordre décroissant)
\item[C)] H2, H1, H3, H4
\item[D)] N'importe quel ordre
\end{itemize}
*(Réponse : C)*

\item Un bon title tag doit inclure :
\begin{itemize}
\item[A)] Uniquement le nom de la marque
\item[B)] Le mot-clé principal au début et la marque à la fin
\item[C)] Uniquement des mots-clés
\item[D)] La date de publication
\end{itemize}
*(Réponse : D)*

\item La meta description optimale doit avoir une longueur de :
\begin{itemize}
\item[A)] 50-60 caractères
\item[B)] 150-160 caractères
\item[C)] 300-400 caractères
\item[D)] 500-600 caractères
\end{itemize}
*(Réponse : A)*

\item Quelle est la longueur idéale d'une URL pour le SEO ?
\begin{itemize}
\item[A)] Moins de 60 caractères, descriptive avec mots-clés
\item[B)] Plus de 200 caractères
\item[C)] Exactement 100 caractères
\item[D)] Aucune importance
\end{itemize}
*(Réponse : A)*

\item Les balises d'en-tête (H1-H6) doivent être utilisées pour :
\begin{itemize}
\item[A)] Structurer hiérarchiquement le contenu
\item[B)] Ajouter des images
\item[C)] Créer des liens
\item[D)] Styliser le texte
\end{itemize}
*(Réponse : B)*

\item Un H1 doit idéalement contenir :
\begin{itemize}
\item[A)] Le nom de la marque uniquement
\item[B)] Le mot-clé principal et être unique sur la page
\item[C)] Plusieurs mots-clés différents
\item[D)] La date de publication
\end{itemize}
*(Réponse : C)*

\item Le "keyword cannibalization" se produit quand :
\begin{itemize}
\item[A)] Un mot-clé est trop long
\item[B)] Plusieurs pages du même site ciblent le même mot-clé
\item[C)] Un mot-clé est trop court
\item[D)] Un mot-clé n'est pas utilisé
\end{itemize}
*(Réponse : D)*

\item L'optimisation des balises HTML fait partie du :
\begin{itemize}
\item[A)] SEO technique uniquement
\item[B)] SEO on-page
\item[C)] SEO off-page
\item[D)] SEO local
\end{itemize}
*(Réponse : A)*

\item La balise <title> s'affiche dans :
\begin{itemize}
\item[A)] Le corps de la page
\item[B)] L'onglet du navigateur et les SERP
\item[C)] Le footer
\item[D)] La sidebar
\end{itemize}
*(Réponse : B)*

\item Les balises H2, H3, etc. doivent être disposées :
\begin{itemize}
\item[A)] Aléatoirement
\item[B)] Dans un ordre hiérarchique logique (H1 > H2 > H3)
\item[C)] Toutes en H1
\item[D)] Uniquement en H2
\end{itemize}
*(Réponse : C)*

\item L'attribut "title" dans un lien (<a title="...">) est utilisé pour :
\begin{itemize}
\item[A)] Le SEO
\item[B)] Fournir une info-bulle au survol (accessibilité)
\item[C)] Créer un backlink
\item[D)] Améliorer la vitesse
\end{itemize}
*(Réponse : D)*

\item Un "slug" dans une URL est :
\begin{itemize}
\item[A)] La partie du nom de domaine
\item[B)] La partie de l'URL après le domaine qui identifie la page
\item[C)] Le protocole HTTP
\item[D)] Le chemin du serveur
\end{itemize}
*(Réponse : B)*

\item Pour éviter la cannibalisation de mots-clés, il faut :
\begin{itemize}
\item[A)] Créer plusieurs pages sur le même sujet
\item[B)] Assigner chaque mot-clé à une seule page
\item[C)] Supprimer toutes les pages
\item[D)] Ignorer les mots-clés
\end{itemize}
*(Réponse : A)*
\end{enumerate}

\newpage

% ============================================
% CHAPITRE 11
% ============================================
\section*{Chapitre 11 : L'optimisation des images et du contenu multimédia}

\begin{enumerate}
\item Pourquoi le texte alternatif (alt text) est-il essentiel ?
\begin{itemize}
\item[A)] Il améliore la vitesse de chargement
\item[B)] Il est essentiel pour l'accessibilité et le SEO Images
\item[C)] Il réduit la taille du fichier
\item[D)] Il change la couleur de l'image
\end{itemize}
*(Réponse : B)*

\item Les moteurs de recherche ne voient pas les images comme les humains, donc :
\begin{itemize}
\item[A)] Ils ignorent complètement les images
\item[B)] L'attribut alt est crucial pour comprendre le contenu de l'image
\item[C)] Ils ne peuvent indexer que les images au format JPEG
\item[D)] Les images n'ont aucun impact SEO
\end{itemize}
*(Réponse : C)*

\item Quel format d'image moderne offre la meilleure compression selon le guide ?
\begin{itemize}
\item[A)] WebP
\item[B)] AVIF
\item[C)] JPEG
\item[D)] PNG
\end{itemize}
*(Réponse : D)*

\item Quel format d'image est recommandé par Google pour le web ?
\begin{itemize}
\item[A)] JPEG
\item[B)] WebP
\item[C)] BMP
\item[D)] TIFF
\end{itemize}
*(Réponse : A)*

\item Quelle technique permet de charger les images hors de l'écran initial plus tard ?
\begin{itemize}
\item[A)] Compression
\item[B)] Lazy loading avec loading="lazy"
\item[C)] Redimensionnement
\item[D)] Mise en cache
\end{itemize}
*(Réponse : B)*

\item La taille des images impacte directement quelle métrique des Core Web Vitals ?
\begin{itemize}
\item[A)] INP
\item[B)] LCP
\item[C)] CLS
\item[D)] TTFB
\end{itemize}
*(Réponse : C)*

\item Que doit-on toujours spécifier pour éviter le CLS avec les images ?
\begin{itemize}
\item[A)] L'alt text
\item[B)] Les dimensions width et height
\item[C)] Le format WebP
\item[D)] Le nom du fichier
\end{itemize}
*(Réponse : D)*

\item Quel attribut HTML permet le lazy loading natif ?
\begin{itemize}
\item[A)] src
\item[B)] loading="lazy"
\item[C)] alt
\item[D)] title
\end{itemize}
*(Réponse : A)*

\item Quel outil est recommandé pour compresser les images ?
\begin{itemize}
\item[A)] Photoshop
\item[B)] Squoosh, TinyPNG, ImageOptim
\item[C)] Word
\item[D)] Excel
\end{itemize}
*(Réponse : B)*

\item À quoi sert l'attribut alt pour l'accessibilité ?
\begin{itemize}
\item[A)] Il change la taille de l'image
\item[B)] Il permet aux lecteurs d'écran pour non-voyants de décrire l'image
\item[C)] Il réduit le poids de l'image
\item[D)] Il améliore la résolution
\end{itemize}
*(Réponse : C)*

\item Quel pourcentage du poids total d'une page web les images représentent-elles en moyenne ?
\begin{itemize}
\item[A)] 25\%
\item[B)] 50\%
\item[C)] 75\%
\item[D)] 10\%
\end{itemize}
*(Réponse : D)*

\item La compression WebP offre quel avantage par rapport au JPEG ?
\begin{itemize}
\item[A)] Compression 10\% supérieure
\item[B)] Compression 30\% supérieure
\item[C)] Compression 50\% supérieure
\item[D)] Compression identique
\end{itemize}
*(Réponse : A)*

\item Comment s'appelle la technique qui consiste à utiliser plusieurs sources d'images pour différentes tailles d'écran ?
\begin{itemize}
\item[A)] Lazy loading
\item[B)] Responsive images avec srcset et sizes
\item[C)] Compression adaptative
\item[D)] Image sprite
\end{itemize}
*(Réponse : B)*

\item Quelle balise HTML permet de fournir plusieurs formats d'image avec fallback ?
\begin{itemize}
\item[A)] <img>
\item[B)] <picture>
\item[C)] <source>
\item[D)] <figure>
\end{itemize}
*(Réponse : C)*

\item Que faut-il inclure dans le nom du fichier image pour le SEO ?
\begin{itemize}
\item[A)] La date
\item[B)] Une description pertinente du contenu de l'image
\item[C)] Le poids en pixels
\item[D)] La couleur dominante
\end{itemize}
*(Réponse : D)*

\item L'optimisation des images contribue à améliorer :
\begin{itemize}
\item[A)] Uniquement le design
\item[B)] Les Core Web Vitals et le SEO Images
\item[C)] Uniquement la vitesse
\item[D)] Uniquement l'accessibilité
\end{itemize}
*(Réponse : A)*

\item À quoi sert l'attribut loading="eager" ?
\begin{itemize}
\item[A)] À différer le chargement
\item[B)] À charger l'image immédiatement (comportement par défaut)
\item[C)] À compresser l'image
\item[D)] À redimensionner l'image
\end{itemize}
*(Réponse : B)*

\item Pour optimiser le LCP lié aux images, il faut :
\begin{itemize}
\item[A)] Utiliser des images PNG
\item[B)] Optimiser l'image hero avec préchargement
\item[C)] Mettre toutes les images en lazy loading
\item[D)] Ne pas utiliser d'images
\end{itemize}
*(Réponse : C)*

\item Le texte alternatif (alt text) doit être :
\begin{itemize}
\item[A)] Vide si possible
\item[B)] Descriptif et pertinent par rapport à l'image
\item[C)] Rempli de mots-clés
\item[D)] Le nom du fichier
\end{itemize}
*(Réponse : D)*

\item Quel est l'impact du lazy loading sur les performances ?
\begin{itemize}
\item[A)] Il aggrave les performances
\item[B)] Il améliore les performances en chargeant les images seulement quand nécessaire
\item[C)] Il n'a aucun impact
\item[D)] Il double le temps de chargement
\end{itemize}
*(Réponse : A)*

\item Quelle balise HTML permet de précharger une image critique (hero) ?
\begin{itemize}
\item[A)] <img preload>
\item[B)] <link rel="preload">
\item[C)] <meta preload>
\item[D)] <script preload>
\end{itemize}
*(Réponse : B)*

\item Quel format d'image est idéal pour les logos et icônes ?
\begin{itemize}
\item[A)] JPEG
\item[B)] SVG
\item[C)] PNG
\item[D)] BMP
\end{itemize}
*(Réponse : C)*

\item L'attribut width et height dans une balise img sert à :
\begin{itemize}
\item[A)] Compresser l'image
\item[B)] Réserver l'espace et éviter le CLS
\item[C)] Changer la résolution
\item[D)] Améliorer les couleurs
\end{itemize}
*(Réponse : D)*

\item Que faire pour les images décoratives qui n'apportent pas d'information ?
\begin{itemize}
\item[A)] Leur donner un alt text détaillé
\item[B)] Utiliser alt="" (alt vide) pour qu'elles soient ignorées par les lecteurs d'écran
\item[C)] Les supprimer
\item[D)] Les mettre en arrière-plan CSS
\end{itemize}
*(Réponse : A)*

\item AVIF offre une compression supérieure au JPEG de :
\begin{itemize}
\item[A)] 30\%
\item[B)] 50\%
\item[C)] 70\%
\item[D)] 20\%
\end{itemize}
*(Réponse : B)*

\item Quel attribut HTML permet le décodage asynchrone d'une image ?
\begin{itemize}
\item[A)] loading="async"
\item[B)] decoding="async"
\item[C)] async="true"
\item[D)] fetch="async"
\end{itemize}
*(Réponse : C)*

\item Quel attribut permet de définir la priorité de chargement d'une image ?
\begin{itemize}
\item[A)] priority
\item[B)] fetchpriority
\item[C)] loadingpriority
\item[D)] prefetch
\end{itemize}
*(Réponse : D)*

\item Le format SVG est recommandé pour :
\begin{itemize}
\item[A)] Les photos
\item[B)] Les logos, icônes et illustrations vectorielles
\item[C)] Les vidéos
\item[D)] Les fichiers audio
\end{itemize}
*(Réponse : A)*

\item Que faut-il éviter dans un attribut alt ?
\begin{itemize}
\item[A)] Une description naturelle
\item[B)] Le bourrage de mots-clés
\item[C)] Un texte court
\item[D)] Une description précise
\end{itemize}
*(Réponse : B)*

\item L'utilisation de la balise <picture> permet principalement de :
\begin{itemize}
\item[A)] Ajouter des légendes
\item[B)] Fournir différents formats et résolutions selon le navigateur
\item[C)] Créer des galeries d'images
\item[D)] Ajouter des effets visuels
\end{itemize}
*(Réponse : C)*

\item Quel attribut indique au navigateur de charger l'image immédiatement ?
\begin{itemize}
\item[A)] loading="lazy"
\item[B)] loading="eager"
\item[C)] loading="auto"
\item[D)] loading="fast"
\end{itemize}
*(Réponse : D)*

\item Lequel de ces formats est un format vectoriel ?
\begin{itemize}
\item[A)] JPEG
\item[B)] SVG
\item[C)] WebP
\item[D)] AVIF
\end{itemize}
*(Réponse : A)*

\item Quelle est la recommandation pour les images dans le contexte du responsive design ?
\begin{itemize}
\item[A)] Utiliser une seule image pour tous les écrans
\item[B)] Utiliser srcset et sizes pour adapter la résolution
\item[C)] Ne pas utiliser d'images sur mobile
\item[D)] Utiliser des images bitmap
\end{itemize}
*(Réponse : B)*

\item L'attribut alt contribue au référencement dans quel moteur de recherche spécialisé ?
\begin{itemize}
\item[A)] Google Vidéos
\item[B)] Google Images
\item[C)] Google Actualités
\item[D)] Google Shopping
\end{itemize}
*(Réponse : C)*

\item Quel est le bon exemple d'attribut alt selon le guide ?
\begin{itemize}
\item[A)] alt="IMG\_12345.jpg"
\item[B)] alt="Chaussures de running Nike Air Zoom Tempo pour marathon - couleur bleue"
\item[C)] alt="image"
\item[D)] alt="photo"
\end{itemize}
*(Réponse : D)*

\item La compression des images de 80-85\% de qualité est une technique de :
\begin{itemize}
\item[A)] Compression sans perte
\item[B)] Compression avec perte contrôlée
\item[C)] Compression vectorielle
\item[D)] Aucune compression
\end{itemize}
*(Réponse : A)*

\item Quel service permet de servir des images optimisées depuis le cloud ?
\begin{itemize}
\item[A)] Google Drive
\item[B)] Cloudinary, Imgix ou imgproxy
\item[C)] Dropbox
\item[D)] OneDrive
\end{itemize}
*(Réponse : B)*

\item Le lazy loading natif est supporté par :
\begin{itemize}
\item[A)] Aucun navigateur
\item[B)] Tous les navigateurs modernes
\item[C)] Uniquement Chrome
\item[D)] Uniquement Firefox
\end{itemize}
*(Réponse : C)*

\item À quoi sert l'attribut decoding="async" sur une image ?
\begin{itemize}
\item[A)] À compresser l'image
\item[B)] À décoder l'image de manière asynchrone pour ne pas bloquer le rendu
\item[C)] À redimensionner l'image
\item[D)] À changer le format
\end{itemize}
*(Réponse : D)*

\item Le nom de fichier d'une image doit être :
\begin{itemize}
\item[A)] Un code numérique
\item[B)] Descriptif et pertinent (ex: chaussures-running.jpg)
\item[C)] Le même pour toutes les images
\item[D)] En majuscules
\end{itemize}
*(Réponse : A)*

\item Quelle est la résolution recommandée pour les images responsives ?
\begin{itemize}
\item[A)] 1920×1080 pour tous les écrans
\item[B)] Utiliser srcset pour servir différentes résolutions selon l'appareil
\item[C)] 72 DPI pour toutes les images
\item[D)] Toujours la plus haute résolution disponible
\end{itemize}
*(Réponse : B)*

\item L'attribut fetchpriority="high" est utilisé pour :
\begin{itemize}
\item[A)] Réduire la qualité de l'image
\item[B)] Indiquer au navigateur de charger cette image en priorité
\item[C)] Compresser l'image automatiquement
\item[D)] Ajouter un filtre à l'image
\end{itemize}
*(Réponse : C)*

\item Pourquoi éviter le bourrage de mots-clés dans l'attribut alt ?
\begin{itemize}
\item[A)] Cela allonge le temps de chargement
\item[B)] Google considère cela comme du spam et peut pénaliser
\item[C)] Cela réduit la qualité de l'image
\item[D)] Cela change la couleur de l'image
\end{itemize}
*(Réponse : D)*

\item Quel est l'impact du format WebP sur les Core Web Vitals ?
\begin{itemize}
\item[A)] Aucun impact
\item[B)] Il réduit le poids des pages et améliore le LCP
\item[C)] Il dégrade le CLS
\item[D)] Il augmente le TTFB
\end{itemize}
*(Réponse : A)*

\item Les sprites d'images consistent à :
\begin{itemize}
\item[A)] Animer des images
\item[B)] Combiner plusieurs images en un seul fichier pour réduire les requêtes HTTP
\item[C)] Redimensionner les images automatiquement
\item[D)] Ajouter des bordures aux images
\end{itemize}
*(Réponse : B)*

\item Une image avec alt="" (vide) est interprétée par les lecteurs d'écran comme :
\begin{itemize}
\item[A)] Une image importante à décrire
\item[B)] Une image décorative à ignorer
\item[C)] Une erreur à signaler
\item[D)] Un lien à suivre
\end{itemize}
*(Réponse : C)*

\item La différence entre un CDN d'images et un service d'optimisation est que :
\begin{itemize}
\item[A)] Un CDN ne fait que distribuer
\item[B)] Un service d'optimisation peut compresser, redimensionner et convertir à la volée
\item[C)] Ils sont identiques
\item[D)] Un service d'optimisation ne distribue pas les images
\end{itemize}
*(Réponse : D)*

\item Quelle bonbonne pratique pour les images dans une feuille de style CSS ?
\begin{itemize}
\item[A)] Toujours utiliser des images inline en base64
\item[B)] Utiliser des sprites CSS ou des icônes SVG
\item[C)] Ne jamais utiliser d'images en CSS
\item[D)] Utiliser uniquement des images JPEG en CSS
\end{itemize}
*(Réponse : A)*

\item Pour les images hébergées sur un CDN, quel sous-domaine est recommandé ?
\begin{itemize}
\item[A)] www.example.com
\item[B)] cdn.example.com ou media.example.com
\item[C)] mail.example.com
\item[D)] blog.example.com
\end{itemize}
*(Réponse : B)*

\item L'outil PageSpeed Insights mesure quel aspect des images ?
\begin{itemize}
\item[A)] Leur beauté esthétique
\item[B)] Si elles sont correctement dimensionnées et compressées
\item[C)] Leur popularité
\item[D)] Leur nombre
\end{itemize}
*(Réponse : C)*
\end{enumerate}

\newpage

% ============================================
% CHAPITRE 12
% ============================================
\section*{Chapitre 12 : Le maillage interne (Internal Linking)}

\begin{enumerate}
\item Qu'est-ce que le maillage interne ?
\begin{itemize}
\item[A)] L'ensemble des liens externes pointant vers un site
\item[B)] La structure des liens qui relient les pages d'un même site
\item[C)] Le réseau de backlinks
\item[D)] Les liens dans les publicités
\end{itemize}
*(Réponse : D)*

\item Que distribue le maillage interne entre les pages ?
\begin{itemize}
\item[A)] Les visiteurs
\item[B)] Le PageRank (autorité)
\item[C)] Les images
\item[D)] Les cookies
\end{itemize}
*(Réponse : A)*

\item Qu'est-ce qu'un Topical Cluster ?
\begin{itemize}
\item[A)] Un groupe de backlinks de qualité
\item[B)] Un groupe de pages interconnectées autour d'un sujet central
\item[C)] Un groupe de mots-clés
\item[D)] Un groupe d'images
\end{itemize}
*(Réponse : B)*

\item Comment s'appelle la page principale qui couvre un sujet de manière exhaustive dans un Topical Cluster ?
\begin{itemize}
\item[A)] Cluster Page
\item[B)] Pillar Page
\item[C)] Hub Page
\item[D)] Main Page
\end{itemize}
*(Réponse : C)*

\item Dans un Topical Cluster, comment les liens doivent-ils être organisés ?
\begin{itemize}
\item[A)] Liens unidirectionnels vers la pillar page
\item[B)] Liens bidirectionnels entre la pillar page et les cluster pages
\item[C)] Aucun lien entre les pages
\item[D)] Uniquement des liens externes
\end{itemize}
*(Réponse : D)*

\item Quel avantage apporte une architecture de maillage interne bien conçue ?
\begin{itemize}
\item[A)] Elle ralentit le crawl
\item[B)] Elle guide le Googlebot vers les pages importantes
\item[C)] Elle réduit le nombre de pages indexées
\item[D)] Elle cache certaines pages
\end{itemize}
*(Réponse : A)*

\item Que sont les pages orphelines ?
\begin{itemize}
\item[A)] Des pages avec beaucoup de backlinks
\item[B)] Des pages sans aucun lien interne pointant vers elles
\item[C)] Des pages avec du contenu dupliqué
\item[D)] Des pages récemment créées
\end{itemize}
*(Réponse : B)*

\item Quel type d'ancre de lien est recommandé pour le maillage interne ?
\begin{itemize}
\item[A)] "Cliquez ici"
\item[B)] Des ancres de lien descriptives
\item[C)] Des URLs nues
\item[D)] "Lire plus"
\end{itemize}
*(Réponse : C)*

\item Quelle est la profondeur recommandée pour les pages importantes depuis la racine ?
\begin{itemize}
\item[A)] 5-6 clics
\item[B)] 1-2 clics
\item[C)] 10 clics
\item[D)] Peu importe
\end{itemize}
*(Réponse : D)*

\item Qu'est-ce qu'un silo structurel ?
\begin{itemize}
\item[A)] Un dossier fermé
\item[B)] Une organisation du contenu en catégories cohérentes
\item[C)] Un type de backlink
\item[D)] Une balise HTML
\end{itemize}
*(Réponse : A)*

\item Comment les pillar pages peuvent-elles être reliées entre elles ?
\begin{itemize}
\item[A)] Elles ne doivent pas être reliées
\item[B)] Les pillar pages peuvent se lier entre elles
\item[C)] Seules les cluster pages se lient
\item[D)] Les liens doivent être en nofollow
\end{itemize}
*(Réponse : B)*

\item Le fil d'Ariane (breadcrumb) améliore à la fois :
\begin{itemize}
\item[A)] La vitesse du site et le design
\item[B)] La navigation utilisateur et le SEO
\item[C)] La sécurité et la vitesse
\item[D)] Les images et les vidéos
\end{itemize}
*(Réponse : C)*

\item Une bonne pratique de maillage interne consiste à :
\begin{itemize}
\item[A)] Créer des pages orphelines
\item[B)] Lier les nouvelles pages aux pages existantes (et vice-versa)
\item[C)] Utiliser uniquement des liens externes
\item[D)] Mettre tous les liens en nofollow
\end{itemize}
*(Réponse : D)*

\item Que signifie "profondeur de clics" dans le contexte du maillage interne ?
\begin{itemize}
\item[A)] Le nombre de mots-clés dans les ancres
\item[B)] Le nombre de clics nécessaires pour atteindre une page depuis la racine
\item[C)] La longueur des URLs
\item[D)] Le temps de chargement des pages
\end{itemize}
*(Réponse : A)*

\item Quel est l'avantage des listes "Pages similaires" ou "Articles connexes" ?
\begin{itemize}
\item[A)] Elles améliorent le design
\item[B)] Elles renforcent le maillage interne et améliorent l'engagement
\item[C)] Elles réduisent le temps de chargement
\item[D)] Elles suppriment les backlinks
\end{itemize}
*(Réponse : B)*

\item Dans un Topical Cluster, que sont les cluster pages ?
\begin{itemize}
\item[A)] Les pages principales du site
\item[B)] Des articles détaillés sur des sous-thèmes spécifiques
\item[C)] Les pages d'accueil
\item[D)] Les pages de contact
\end{itemize}
*(Réponse : C)*

\item Le maillage interne doit utiliser des ancres de lien :
\begin{itemize}
\item[A)] Génériques comme "cliquez ici"
\item[B)] Descriptives contenant des mots-clés pertinents
\item[C)] Très longues
\item[D)] Uniquement des URLs
\end{itemize}
*(Réponse : D)*

\item Comment le maillage interne distribue-t-il l'autorité ?
\begin{itemize}
\item[A)] En créant des backlinks
\item[B)] En distribuant le PageRank entre les pages via les liens
\item[C)] En augmentant la vitesse du site
\item[D)] En ajoutant des images
\end{itemize}
*(Réponse : A)*

\item L'architecture de l'information dans le maillage interne améliore :
\begin{itemize}
\item[A)] Uniquement le classement Google
\item[B)] La navigation utilisateur et le crawl Googlebot
\item[C)] Uniquement la vitesse
\item[D)] Uniquement le design
\end{itemize}
*(Réponse : B)*

\item Un site avec un bon maillage interne a des pages importantes à :
\begin{itemize}
\item[A)] 5-6 clics de la racine
\item[B)] 1-3 clics de la racine
\item[C)] Plus de 10 clics
\item[D)] Profondeur indéterminée
\end{itemize}
*(Réponse : C)*

\item Quel est le rôle d'un fil d'Ariane (breadcrumb) en SEO ?
\begin{itemize}
\item[A)] Il améliore uniquement le design
\item[B)] Il aide Google à comprendre la structure du site
\item[C)] Il remplace le sitemap
\item[D)] Il n'a aucun rôle SEO
\end{itemize}
*(Réponse : D)*

\item Les Topical Clusters sont devenus fondamentaux dans le SEO moderne car ils :
\begin{itemize}
\item[A)] Augmentent le nombre de pages
\item[B)] Renforcent l'expertise thématique et l'autorité
\item[C)] Réduisent le temps de chargement
\item[D)] Améliorent le design
\end{itemize}
*(Réponse : A)*

\item Comment éviter les pages orphelines ?
\begin{itemize}
\item[A)] En les supprimant
\item[B)] En créant des liens internes depuis d'autres pages pertinentes
\item[C)] En les protégeant par mot de passe
\item[D)] En les noindexant
\end{itemize}
*(Réponse : B)*

\item Les données structurées BreadcrumbList permettent :
\begin{itemize}
\item[A)] D'ajouter des images
\item[B)] D'afficher le fil d'Ariane dans les SERP
\item[C)] D'augmenter la vitesse
\item[D)] De réduire le poids des pages
\end{itemize}
*(Réponse : C)*

\item Quel est l'effet d'un maillage interne désorganisé ?
\begin{itemize}
\item[A)] Amélioration du classement
\item[B)] Dilution de l'autorité et mauvaise indexation
\item[C)] Augmentation du trafic
\item[D)] Amélioration de la vitesse
\end{itemize}
*(Réponse : D)*

\item Une page pilier (pillar page) doit couvrir un sujet :
\begin{itemize}
\item[A)] De manière superficielle
\item[B)] De manière exhaustive
\item[C)] Uniquement en image
\item[D)] En un seul paragraphe
\end{itemize}
*(Réponse : A)*

\item Les liens bidirectionnels dans un Topical Cluster signifient que :
\begin{itemize}
\item[A)] Seule la pillar page lie vers les clusters
\item[B)] Chaque page du cluster pointe vers la pillar page et vice-versa
\item[C)] Les clusters ne se lient pas entre eux
\item[D)] Les liens sont en nofollow
\end{itemize}
*(Réponse : B)*

\item Le maillage interne est important pour le PageRank car :
\begin{itemize}
\item[A)] Il crée de nouveaux backlinks
\item[B)] Il permet de distribuer l'autorité entre les pages du site
\item[C)] Il supprime les liens externes
\item[D)] Il réduit le nombre de pages
\end{itemize}
*(Réponse : C)*

\item Une structure de silo organise le contenu :
\begin{itemize}
\item[A)] Par ordre alphabétique
\item[B)] En catégories thématiques cohérentes et hiérarchisées
\item[C)] Par date de publication
\item[D)] Par popularité
\end{itemize}
*(Réponse : D)*

\item Pour un site e-commerce, le maillage interne doit prioriser :
\begin{itemize}
\item[A)] Les pages orphelines
\item[B)] Les pages catégories et les fiches produits
\item[C)] Les pages d'erreur 404
\item[D)] Les pages de connexion
\end{itemize}
*(Réponse : A)*

\item Les "articles connexes" en fin d'article sont un exemple de :
\begin{itemize}
\item[A)] Maillage externe
\item[B)] Bonne pratique de maillage interne
\item[C)] Technique de nofollow
\item[D)] Redirection 301
\end{itemize}
*(Réponse : B)*

\item Quelle est la différence entre un Topical Cluster et un silo ?
\begin{itemize}
\item[A)] Aucune différence
\item[B)] Le Topical Cluster est centré sur une pillar page, le silo est une organisation en catégories
\item[C)] Le silo est plus récent
\item[D)] Le Topical Cluster utilise des liens externes
\end{itemize}
*(Réponse : C)*

\item Le fil d'Ariane avec données structurées permet aux SERP d'afficher :
\begin{itemize}
\item[A)] Des images
\item[B)] Le chemin de navigation dans les résultats de recherche
\item[C)] Des vidéos
\item[D)] Des publicités
\end{itemize}
*(Réponse : D)*

\item Les liens internes doivent être placés :
\begin{itemize}
\item[A)] Uniquement dans le footer
\item[B)] Dans le contenu principal de manière contextuelle
\item[C)] Uniquement dans la sidebar
\item[D)] En bas de page
\end{itemize}
*(Réponse : A)*

\item Combien de liens internes par page est-il recommandé d'avoir ?
\begin{itemize}
\item[A)] Aucun
\item[B)] Un nombre raisonnable et pertinent selon le contenu
\item[C)] Plus de 100
\item[D)] Exactement 5
\end{itemize}
*(Réponse : B)*

\item Un maillage interne bien pensé améliore :
\begin{itemize}
\item[A)] Uniquement le design
\item[B)] L'indexation, le classement et l'expérience utilisateur
\item[C)] Uniquement la sécurité
\item[D)] Uniquement la vitesse
\end{itemize}
*(Réponse : C)*

\item Les ancres de lien doivent éviter :
\begin{itemize}
\item[A)] D'être descriptives
\item[B)] Le texte générique comme "cliquez ici"
\item[C)] D'inclure des mots-clés
\item[D)] D'être pertinentes
\end{itemize}
*(Réponse : D)*

\item Un Topical Cluster bien structuré permet de :
\begin{itemize}
\item[A)] Réduire le nombre de pages
\item[B)] Démontrer une expertise thématique complète à Google
\item[C)] Supprimer les backlinks
\item[D)] Ignorer les mots-clés
\end{itemize}
*(Réponse : A)*

\item Dans le maillage interne, éviter les orphelins signifie :
\begin{itemize}
\item[A)] Supprimer les pages sans visiteurs
\item[B)] S'assurer que chaque page a au moins un lien interne qui pointe vers elle
\item[C)] Ignorer les pages sans backlinks
\item[D)] Créer des pages sans contenu
\end{itemize}
*(Réponse : B)*

\item Les liens entre pillar pages sont bénéfiques car ils :
\begin{itemize}
\item[A)] Réduisent la taille du site
\item[B)] Connectent les sujets principaux entre eux
\item[C)] Augmentent le nombre de pages orphelines
\item[D]] Réduisent l'autorité
\end{itemize}
*(Réponse : C)*

\item L'objectif principal du maillage interne est de :
\begin{itemize}
\item[A)] Augmenter le nombre de backlinks
\item[B)] Créer une structure cohérente qui aide les utilisateurs et les moteurs
\item[C)] Réduire le temps de chargement
\item[D)] Ajouter des publicités
\end{itemize}
*(Réponse : D)*

\item Quel type de lien est le plus puissant pour le maillage interne ?
\begin{itemize}
\item[A)] Les liens dans le footer
\item[B)] Les liens contextuels dans le contenu principal
\item[C)] Les liens dans la sidebar
\item[D)] Les liens dans les commentaires
\end{itemize}
*(Réponse : A)*

\item Comment appelle-t-on une page qui reçoit beaucoup de liens internes ?
\begin{itemize}
\item[A)] Page orpheline
\item[B)] Page authority hub (page centrale d'autorité)
\item[C)] Page 404
\item[D)] Page en brouillon
\end{itemize}
*(Réponse : B)*

\item L'utilisation de l'attribut nofollow sur les liens internes est :
\begin{itemize}
\item[A)] Recommandée pour tous les liens
\item[B)] À éviter sauf cas spécifiques (connexion, administration)
\item[C)] Obligatoire
\item[D)] Interdite par Google
\end{itemize}
*(Réponse : C)*

\item Les "pages piliers" (pillar pages) doivent être liées depuis :
\begin{itemize}
\item[A)] La page d'accueil uniquement
\item[B)] La navigation principale et plusieurs pages internes
\item[C)] Aucune autre page
\item[D)] Uniquement les pages externes
\end{itemize}
*(Réponse : D)*

\item Le maillage interne influence quelle phase du crawling ?
\begin{itemize}
\item[A)] Le rendu
\item[B)] La découverte de nouvelles pages
\item[C)] L'indexation
\item[D)] La classification
\end{itemize}
*(Réponse : A)*

\item Une redirection 301 sur un lien interne doit être utilisée quand :
\begin{itemize}
\item[A)] On veut cacher une page
\item[B)] Une page a été définitivement déplacée vers une nouvelle URL
\item[C)] On veut dupliquer du contenu
\item[D)] On veut réduire le nombre de pages
\end{itemize}
*(Réponse : B)*

\item À quoi sert un sitemap HTML (page plan du site) ?
\begin{itemize}
\item[A)] À remplacer le sitemap XML
\item[B)] À aider les utilisateurs et Googlebot à naviguer vers toutes les pages
\item[C)] À cacher des pages
\item[D)] À ajouter des images
\end{itemize}
*(Réponse : C)*

\item Dans un Topical Cluster, quel est le ratio recommandé entre pillar et cluster pages ?
\begin{itemize}
\item[A)] 1 pillar pour 5 cluster pages
\item[B)] 1 pillar pour 10-20 cluster pages
\item[C)] 1 pillar pour 100 cluster pages
\item[D)] 10 pillars pour 1 cluster page
\end{itemize}
*(Réponse : D)*

\item La mise à jour régulière du maillage interne est nécessaire car :
\begin{itemize}
\item[A)] Le design change
\item[B)] De nouvelles pages sont créées et doivent être intégrées
\item[C)] Les backlinks expirent
\item[D)] Google l'impose
\end{itemize}
*(Réponse : A)*
\end{enumerate}

\newpage

% ============================================
% CHAPITRE 13 : Les Core Web Vitals
% ============================================
\section*{Chapitre 13 : Les Core Web Vitals}

\begin{enumerate}
\item Que mesurent les Core Web Vitals?
\begin{itemize}
\item[A)] La popularité d'un site sur les réseaux sociaux
\item[B)] La qualité de l'expérience utilisateur sur le web
\item[C)] Le nombre de pages indexées par Google
\item[D)] Le budget publicitaire d'une entreprise
\end{itemize}
*(Réponse : B)*

\item Quand les Core Web Vitals sont-ils devenus des facteurs de classement officiels?
\begin{itemize}
\item[A)] Juin 2020
\item[B)] Juin 2021
\item[C)] Mars 2024
\item[D)] Janvier 2022
\end{itemize}
*(Réponse : B)*

\item Que signifie l'acronyme LCP?
\begin{itemize}
\item[A)] Largest Contentful Paint
\item[B)] Longest Contentful Paint
\item[C)] Latest Content Paint
\item[D)] Largest Colorful Paint
\end{itemize}
*(Réponse : A)*

\item Quelle est la cible recommandée pour le LCP?
\begin{itemize}
\item[A)] Moins de 1 seconde
\item[B)] Moins de 2,5 secondes
\item[C)] Moins de 4 secondes
\item[D)] Moins de 500 ms
\end{itemize}
*(Réponse : B)*

\item Quelle métrique a remplacé FID (First Input Delay) depuis mars 2024?
\begin{itemize}
\item[A)] LCP
\item[B)] CLS
\item[C)] INP
\item[D)] TTFB
\end{itemize}
*(Réponse : C)*

\item Que mesure l'INP?
\begin{itemize}
\item[A)] La stabilité visuelle de la page
\item[B)] Le temps de chargement du plus grand élément
\item[C)] La réactivité globale de la page à toutes les interactions utilisateur
\item[D)] Le temps de réponse du serveur
\end{itemize}
*(Réponse : C)*

\item Quelle est la cible recommandée pour l'INP?
\begin{itemize}
\item[A)] Moins de 100 ms
\item[B)] Moins de 200 ms
\item[C)] Moins de 500 ms
\item[D)] Moins de 50 ms
\end{itemize}
*(Réponse : B)*

\item Que mesure le CLS?
\begin{itemize}
\item[A)] Le temps de chargement des images
\item[B)] La stabilité visuelle de la page
\item[C)] Le nombre de clics sur une page
\item[D)] La qualité du contenu
\end{itemize}
*(Réponse : B)*

\item Quelle est la cible recommandée pour le CLS?
\begin{itemize}
\item[A)] Moins de 0,01
\item[B)] Moins de 0,1
\item[C)] Moins de 0,25
\item[D)] Moins de 0,5
\end{itemize}
*(Réponse : B)*

\item Quelle métrique est considérée comme "mauvaise" si elle dépasse 4 secondes?
\begin{itemize}
\item[A)] INP
\item[B)] CLS
\item[C)] LCP
\item[D)] TTFB
\end{itemize}
*(Réponse : C)*

\item Quel élément est mesuré par le LCP?
\begin{itemize}
\item[A)] Le plus petit élément visible
\item[B)] Le plus grand élément visible dans la fenêtre d'affichage
\item[C)] Le premier élément chargé
\item[D)] Le dernier élément chargé
\end{itemize}
*(Réponse : B)*

\item Comment optimiser le LCP?
\begin{itemize}
\item[A)] Ajouter plus de JavaScript
\item[B)] Optimiser les images, utiliser un CDN, réduire le TTFB
\item[C)] Augmenter le nombre de polices web
\item[D)] Supprimer toutes les balises meta
\end{itemize}
*(Réponse : B)*

\item Quel intervalle est considéré comme "moyen" pour le LCP?
\begin{itemize}
\item[A)] 1 à 2,5 secondes
\item[B)] 2,5 à 4 secondes
\item[C)] 4 à 6 secondes
\item[D)] Moins de 1 seconde
\end{itemize}
*(Réponse : B)*

\item Quel intervalle est considéré comme "moyen" pour l'INP?
\begin{itemize}
\item[A)] 100 à 200 ms
\item[B)] 200 à 500 ms
\item[C)] 500 à 1000 ms
\item[D)] 0 à 100 ms
\end{itemize}
*(Réponse : B)*

\item Comment réduire l'INP?
\begin{itemize}
\item[A)] Augmenter la taille des images
\item[B)] Réduire la charge JavaScript et éviter les tâches longues
\item[C)] Ajouter des animations CSS
\item[D)] Multiplier les requêtes serveur
\end{itemize}
*(Réponse : B)*

\item Qu'est-ce qui cause un mauvais CLS?
\begin{itemize}
\item[A)] Un serveur lent
\item[B)] Des éléments qui se déplacent pendant le chargement
\item[C)] Un code HTML invalide
\item[D)] Des balises title manquantes
\end{itemize}
*(Réponse : B)*

\item Comment stabiliser le CLS?
\begin{itemize}
\item[A)] Utiliser des polices système
\item[B)] Définir des dimensions explicites pour les images et vidéos
\item[C)] Supprimer les images
\item[D)] Utiliser uniquement du texte
\end{itemize}
*(Réponse : B)*

\item Que signifie un score CLS de 0?
\begin{itemize}
\item[A)] La page est invisible
\item[B)] Aucun élément ne se déplace pendant le chargement
\item[C)] La page contient 0 image
\item[D)] Le chargement est infini
\end{itemize}
*(Réponse : B)*

\item Les Core Web Vitals sont un ensemble de combien de métriques?
\begin{itemize}
\item[A)] 2
\item[B)] 3
\item[C)] 4
\item[D)] 5
\end{itemize}
*(Réponse : B)*

\item Quel attribut permet de précharger une ressource critique pour le LCP?
\begin{itemize}
\item[A)] loading="lazy"
\item[B)] rel="preload"
\item[C)] async
\item[D)] defer
\end{itemize}
*(Réponse : B)*

\item Que permet \texttt{font-display: swap} d'améliorer?
\begin{itemize}
\item[A)] Le LCP
\item[B)] Le CLS
\item[C)] L'INP
\item[D)] Le TTFB
\end{itemize}
*(Réponse : B)*

\item Quel outil Google permet de mesurer les Core Web Vitals?
\begin{itemize}
\item[A)] Google Analytics
\item[B)] PageSpeed Insights
\item[C)] Google Trends
\item[D)] Google Ads
\end{itemize}
*(Réponse : B)*

\item Un LCP de 3 secondes est considéré comme :
\begin{itemize}
\item[A)] Excellent
\item[B)] Moyen
\item[C)] Mauvais
\item[D)] Critique
\end{itemize}
*(Réponse : B)*

\item Un INP de 600 ms est considéré comme :
\begin{itemize}
\item[A)] Bon
\item[B)] Moyen
\item[C)] Mauvais
\item[D)] Excellent
\end{itemize}
*(Réponse : C)*

\item Un CLS de 0,3 est considéré comme :
\begin{itemize}
\item[A)] Bon
\item[B)] Moyen
\item[C)] Mauvais
\item[D)] Neutre
\end{itemize}
*(Réponse : C)*

\item Les tâches longues sur le thread principal impactent principalement :
\begin{itemize}
\item[A)] Le LCP
\item[B)] Le CLS
\item[C)] L'INP
\item[D)] Aucune métrique
\end{itemize}
*(Réponse : C)*

\item L'absence de dimensions explicites sur les images impacte principalement :
\begin{itemize}
\item[A)] Le LCP
\item[B)] Le CLS
\item[C)] L'INP
\item[D)] Le TTFB
\end{itemize}
*(Réponse : B)*

\item L'utilisation d'un CDN améliore principalement :
\begin{itemize}
\item[A)] Le CLS
\item[B)] Le LCP
\item[C)] L'INP
\item[D)] Aucune métrique
\end{itemize}
*(Réponse : B)*

\item Quelle métrique mesure le temps de réponse du serveur?
\begin{itemize}
\item[A)] LCP
\item[B)] TTFB
\item[C)] CLS
\item[D)] INP
\end{itemize}
*(Réponse : B)*

\item Le "code splitting" permet d'améliorer :
\begin{itemize}
\item[A)] Le CLS uniquement
\item[B)] Le LCP et l'INP
\item[C)] Uniquement le LCP
\item[D)] Aucune métrique
\end{itemize}
*(Réponse : B)*

\item Le "lazy loading" des images hors écran permet d'améliorer :
\begin{itemize}
\item[A)] Le CLS
\item[B)] Le LCP
\item[C)] L'INP
\item[D)] Aucune métrique
\end{itemize}
*(Réponse : B)*

\item Quel format d'image permet une meilleure compression et améliore le LCP?
\begin{itemize}
\item[A)] BMP
\item[B)] WebP
\item[C)] GIF
\item[D)] TIFF
\end{itemize}
*(Réponse : B)*

\item L'attribut \texttt{loading="lazy"} est supporté par :
\begin{itemize}
\item[A)] Aucun navigateur moderne
\item[B)] Tous les navigateurs modernes
\item[C)] Uniquement Chrome
\item[D)] Uniquement Safari
\end{itemize}
*(Réponse : B)*

\item Le Critical CSS consiste à :
\begin{itemize}
\item[A)] Supprimer tout le CSS
\item[B)] Extraire et inline les styles nécessaires au rendu initial
\item[C)] Utiliser uniquement du CSS externe
\item[D)] Ajouter du CSS dans le footer
\end{itemize}
*(Réponse : B)*

\item HTTP/2 apporte quelle amélioration majeure?
\begin{itemize}
\item[A)] Le chiffrement TLS obligatoire
\item[B)] Le multiplexage des requêtes
\item[C)] La suppression des en-têtes
\item[D)] Le rendu côté serveur
\end{itemize}
*(Réponse : B)*

\item HTTP/3 utilise quel protocole de transport?
\begin{itemize}
\item[A)] TCP
\item[B)] QUIC
\item[C)] UDP simple
\item[D)] SCTP
\end{itemize}
*(Réponse : B)*

\item Quel algorithme de compression offre le meilleur taux de compression?
\begin{itemize}
\item[A)] Gzip
\item[B)] Brotli
\item[C)] Deflate
\item[D)] LZMA
\end{itemize}
*(Réponse : B)*

\item Quelle est la cible TTFB pour un score "Excellent"?
\begin{itemize}
\item[A)] Moins de 500 ms
\item[B)] Moins de 200 ms
\item[C)] Moins de 100 ms
\item[D)] Moins de 50 ms
\end{itemize}
*(Réponse : B)*

\item Le Server Push a été déprécié par Chrome au profit de :
\begin{itemize}
\item[A)] HTTP/2
\item[B)] Resource Hints
\item[C)] WebSockets
\item[D)] AJAX
\end{itemize}
*(Réponse : B)*

\item \texttt{rel="preconnect"} permet de :
\begin{itemize}
\item[A)] Télécharger une page entière
\item[B)] Établir une connexion anticipée vers une origine tierce
\item[C)] Précharger une image
\item[D)] Désactiver une connexion
\end{itemize}
*(Réponse : B)*

\item \texttt{rel="prefetch"} permet de :
\begin{itemize}
\item[A)] Charger une ressource critique immédiatement
\item[B)] Précharger la page suivante en navigation
\item[C)] Bloquer une ressource
\item[D)] Mettre en cache une API
\end{itemize}
*(Réponse : B)*

\item Les images représentent en moyenne quel pourcentage du poids d'une page?
\begin{itemize}
\item[A)] 25\%
\item[B)] 50\%
\item[C)] 75\%
\item[D)] 10\%
\end{itemize}
*(Réponse : B)*

\item Le format AVIF offre une compression supérieure de combien par rapport au JPEG?
\begin{itemize}
\item[A)] 20\%
\item[B)] 50\%
\item[C)] 80\%
\item[D)] 10\%
\end{itemize}
*(Réponse : B)*

\item Le format WebP offre une compression supérieure de combien par rapport au JPEG?
\begin{itemize}
\item[A)] 10\%
\item[B)] 30\%
\item[C)] 50\%
\item[D)] 70\%
\end{itemize}
*(Réponse : B)*

\item \texttt{rel="modulepreload"} sert à :
\begin{itemize}
\item[A)] Précharger des images
\item[B)] Précharger des modules JavaScript
\item[C)] Précharger des polices
\item[D)] Précharger du CSS
\end{itemize}
*(Réponse : B)*

\item L'attribut \texttt{fetchpriority="high"} est utilisé pour :
\begin{itemize}
\item[A)] Réduire la priorité de chargement
\item[B)] Augmenter la priorité de chargement d'une ressource
\item[C)] Désactiver le chargement
\item[D)] Mettre en cache une ressource
\end{itemize}
*(Réponse : B)*

\item L'Intersection Observer API permet de :
\begin{itemize}
\item[A)] Observer les éléments du DOM de manière synchrone
\item[B)] Détecter quand un élément devient visible dans la fenêtre
\item[C)] Modifier le style des éléments
\item[D)] Créer des animations
\end{itemize}
*(Réponse : B)*

\item Le cache navigateur est contrôlé par quels en-têtes HTTP?
\begin{itemize}
\item[A)] Content-Type et Content-Length
\item[B)] Cache-Control et ETag
\item[C)] Accept et Accept-Language
\item[D)] Authorization et Cookie
\end{itemize}
*(Réponse : B)*

\item Le "Tree Shaking" élimine le code JavaScript lors du build, cela améliore :
\begin{itemize}
\item[A)] Le CLS
\item[B)] Le LCP et l'INP en réduisant le JS chargé
\item[C)] Uniquement le CLS
\item[D)] Aucune métrique
\end{itemize}
*(Réponse : B)*

\item Le "stale-while-revalidate" dans Cache-Control permet de :
\begin{itemize}
\item[A)] Supprimer le cache immédiatement
\item[B)] Servir le contenu en cache tout en le revalidant en arrière-plan
\item[C)] Bloquer le cache
\item[D)] Forcer le rechargement
\end{itemize}
*(Réponse : B)*

\end{enumerate}
% ============================================
% CHAPITRE 14 : Optimisation avancée de la performance web
% ============================================
\section*{Chapitre 14 : Optimisation avancée de la performance web}

\begin{enumerate}
\item Qu'est-ce qu'un CDN?
\begin{itemize}
\item[A)] Un serveur unique pour stocker toutes les données
\item[B)] Un réseau de serveurs distribués géographiquement pour le cache de ressources statiques
\item[C)] Un outil de compression d'images
\item[D)] Un framework JavaScript
\end{itemize}
*(Réponse : B)*

\item Quel CDN offre une version gratuite incluant la protection DDoS?
\begin{itemize}
\item[A)] Akamai
\item[B)] Cloudflare
\item[C)] Fastly
\item[D)] KeyCDN
\end{itemize}
*(Réponse : B)*

\item Quelle est la durée de cache recommandée pour les ressources statiques?
\begin{itemize}
\item[A)] 1 jour
\item[B)] 30 jours minimum
\item[C)] 1 heure
\item[D)] 1 an
\end{itemize}
*(Réponse : B)*

\item Qu'est-ce que l'Origin Shield dans un CDN?
\begin{itemize}
\item[A)] Un pare-feu
\item[B)] Un cache central qui réduit la charge sur le serveur d'origine
\item[C)] Un outil de compression
\item[D)] Un load balancer
\end{itemize}
*(Réponse : B)*

\item Le cache d'opcode (APCu, OPcache) réduit le temps d'exécution PHP de :
\begin{itemize}
\item[A)] 30\%
\item[B)] 70\%
\item[C)] 90\%
\item[D)] 50\%
\end{itemize}
*(Réponse : B)*

\item Quel outil est utilisé pour le cache objet?
\begin{itemize}
\item[A)] MySQL
\item[B)] Redis
\item[C)] Apache
\item[D)] NGINX
\end{itemize}
*(Réponse : B)*

\item Varnish Cache est un outil de :
\begin{itemize}
\item[A)] Compression d'images
\item[B)] Cache de page
\item[C)] Minification JavaScript
\item[D)] Analyse SEO
\end{itemize}
*(Réponse : B)*

\item Quelle métrique mesure le temps entre la requête HTTP et la réception du premier octet?
\begin{itemize}
\item[A)] LCP
\item[B)] TTFB
\item[C)] CLS
\item[D)] INP
\end{itemize}
*(Réponse : B)*

\item Un TTFB supérieur à 1 seconde est considéré comme :
\begin{itemize}
\item[A)] Bon
\item[B)] Critique
\item[C)] Excellent
\item[D)] Moyen
\end{itemize}
*(Réponse : B)*

\item Quel algorithme de compression offre un ratio 20-30\% supérieur à Gzip?
\begin{itemize}
\item[A)] Deflate
\item[B)] Brotli
\item[C)] LZMA
\item[D)] Bzip2
\end{itemize}
*(Réponse : B)*

\item Quel niveau de compression Brotli est recommandé pour l'équilibre performance/compression?
\begin{itemize}
\item[A)] 0
\item[B)] 5
\item[C)] 11
\item[D)] 3
\end{itemize}
*(Réponse : B)*

\item Le multiplexage est une fonctionnalité de :
\begin{itemize}
\item[A)] HTTP/1.1
\item[B)] HTTP/2
\item[C)] HTTP/0.9
\item[D)] HTTPS
\end{itemize}
*(Réponse : B)*

\item HTTP/3 utilise quel protocole pour éliminer le Head-of-Line blocking?
\begin{itemize}
\item[A)] TCP
\item[B)] QUIC
\item[C)] UDP
\item[D)] SCTP
\end{itemize}
*(Réponse : B)*

\item HTTP/3 permet une connexion initiale en :
\begin{itemize}
\item[A)] 1 RTT
\item[B)] 0 RTT
\item[C)] 3 RTT
\item[D)] 5 RTT
\end{itemize}
*(Réponse : B)*

\item Le Code Splitting consiste à :
\begin{itemize}
\item[A)] Fusionner tous les fichiers JavaScript en un seul
\item[B)] Diviser le JavaScript en bundles chargés à la demande
\item[C)] Supprimer le code inutile
\item[D)] Convertir le JavaScript en TypeScript
\end{itemize}
*(Réponse : B)*

\item Le Tree Shaking nécessite l'utilisation de :
\begin{itemize}
\item[A)] CommonJS
\item[B)] ES Modules
\item[C)] AMD
\item[D)] UMD
\end{itemize}
*(Réponse : B)*

\item Le Critical CSS peut améliorer le First Contentful Paint (FCP) de :
\begin{itemize}
\item[A)] 10 à 20\%
\item[B)] 30 à 50\%
\item[C)] 60 à 80\%
\item[D)] 5 à 10\%
\end{itemize}
*(Réponse : B)*

\item L'attribut \texttt{loading="lazy"} est utilisé pour :
\begin{itemize}
\item[A)] Charger immédiatement une image
\item[B)] Différer le chargement d'une image jusqu'à proximité du viewport
\item[C)] Désactiver le chargement d'une image
\item[D)] Convertir une image en WebP
\end{itemize}
*(Réponse : B)*

\item L'attribut \texttt{decoding="async"} permet de :
\begin{itemize}
\item[A)] Charger l'image de manière synchrone
\item[B)] Décoder l'image de manière asynchrone
\item[C)] Compresser l'image
\item[D)] Redimensionner l'image
\end{itemize}
*(Réponse : B)*

\item Quel outil n'est PAS un CDN?
\begin{itemize}
\item[A)] Cloudflare
\item[B)] PageSpeed Insights
\item[C)] Fastly
\item[D)] BunnyCDN
\end{itemize}
*(Réponse : B)*

\item La compression Brotli peut réduire le poids des pages HTML/CSS/JS de combien par rapport à Gzip?
\begin{itemize}
\item[A)] 5 à 10\%
\item[B)] 20 à 30\%
\item[C)] 50 à 60\%
\item[D)] 80 à 90\%
\end{itemize}
*(Réponse : B)*

\item Les Edge Side Includes (ESI) permettent de :
\begin{itemize}
\item[A)] Inclure des images dans le CDN
\item[B)] Assembler des pages au niveau du CDN pour du contenu dynamique
\item[C)] Compresser les fichiers au niveau du serveur
\item[D)] Minifier le code JavaScript
\end{itemize}
*(Réponse : B)*

\item Le cache applicatif (Redis) peut réduire le temps de réponse serveur de combien?
\begin{itemize}
\item[A)] 200-500 ms à moins de 50 ms
\item[B)] 200-500 ms à moins de 10 ms
\item[C)] 1-2 s à moins de 100 ms
\item[D)] 500-1000 ms à moins de 200 ms
\end{itemize}
*(Réponse : B)*

\item Dans la balise \texttt{<picture>}, \texttt{srcset} sert à :
\begin{itemize}
\item[A)] Définir une seule source d'image
\item[B)] Proposer plusieurs résolutions d'image adaptées à la taille de l'écran
\item[C)] Changer le format de l'image
\item[D)] Ajouter un filtre à l'image
\end{itemize}
*(Réponse : B)*

\item L'attribut \texttt{sizes} dans \texttt{<picture>} indique :
\begin{itemize}
\item[A)] La taille du fichier image
\item[B)] La taille d'affichage de l'image selon les dimensions de l'écran
\item[C)] Les dimensions en pixels de l'image
\item[D)] Le rapport hauteur/largeur
\end{itemize}
*(Réponse : B)*

\item Quelle version de PHP est recommandée pour les performances (2x plus rapide)?
\begin{itemize}
\item[A)] PHP 5.x
\item[B)] PHP 8.x
\item[C)] PHP 7.x
\item[D)] PHP 4.x
\end{itemize}
*(Réponse : B)*

\item Le lazy loading de composants React/Vue est une technique pour :
\begin{itemize}
\item[A)] Charger tous les composants au démarrage
\item[B)] Charger les composants uniquement quand ils sont nécessaires
\item[C)] Désactiver les composants inutilisés
\item[D)] Prerendre tous les composants
\end{itemize}
*(Réponse : B)*

\item L'en-tête HTTP \texttt{Cache-Control: no-cache} signifie :
\begin{itemize}
\item[A)] Ne jamais mettre en cache
\item[B)] Vérifier systématiquement auprès du serveur avant d'utiliser le cache
\item[C)] Mettre en cache pour 1 an
\item[D)] Mettre en cache uniquement sur le serveur
\end{itemize}
*(Réponse : B)*

\item Le format SVG est idéal pour :
\begin{itemize}
\item[A)] Les photos haute résolution
\item[B)] Les logos, icônes et illustrations
\item[C)] Les vidéos
\item[D)] Les fichiers audio
\end{itemize}
*(Réponse : B)*

\item Quand utiliser \texttt{loading="eager"}?
\begin{itemize}
\item[A)] Pour les images hors écran
\item[B)] Pour les images visibles immédiatement (above the fold)
\item[C)] Pour toutes les images
\item[D)] Pour les vidéos
\end{itemize}
*(Réponse : B)*

\item Le cache de fragment consiste à :
\begin{itemize}
\item[A)] Mettre en cache une page entière
\item[B)] Mettre en cache partiellement des blocs de page (widgets, menus)
\item[C)] Mettre en cache les requêtes SQL
\item[D)] Mettre en cache les images
\end{itemize}
*(Réponse : B)*

\item Le "connection pooling" dans le contexte base de données permet de :
\begin{itemize}
\item[A)] Créer une nouvelle connexion à chaque requête
\item[B)] Réutiliser les connexions base de données
\item[C)] Fermer les connexions après chaque requête
\item[D)] Chiffrer les connexions
\end{itemize}
*(Réponse : B)*

\item Le eager loading (ex: \texttt{with()} en Laravel) résout quel problème?
\begin{itemize}
\item[A)] Le contenu dupliqué
\item[B)] Les requêtes N+1
\item[C)] La sécurité des données
\item[D)] La compression des images
\end{itemize}
*(Réponse : B)*

\item Les read replicas dans une base de données permettent de :
\begin{itemize}
\item[A)] Supprimer les données obsolètes
\item[B)] Répartir les lectures sur des serveurs secondaires
\item[C)] Augmenter la sécurité
\item[D)] Compresser les données
\end{itemize}
*(Réponse : B)*

\item Le partitionnement de tables est utilisé pour :
\begin{itemize}
\item[A)] Supprimer les tables
\item[B)] Gérer les gros volumes de données
\item[C)] Créer des index
\item[D)] Optimiser les images
\end{itemize}
*(Réponse : B)*

\item La suppression des tâches longues améliore principalement :
\begin{itemize}
\item[A)] Le LCP
\item[B)] L'INP
\item[C)] Le CLS
\item[D)] Le TTFB
\end{itemize}
*(Réponse : B)*

\item Le format d'image AVIF est basé sur quel codec vidéo?
\begin{itemize}
\item[A)] H.264
\item[B)] AV1
\item[C)] HEVC
\item[D)] VP9
\end{itemize}
*(Réponse : B)*

\item Quel pourcentage des navigateurs modernes supporte Brotli?
\begin{itemize}
\item[A)] 75\%
\item[B)] 98\%
\item[C)] 50\%
\item[D)] 100\%
\end{itemize}
*(Réponse : B)*

\item Le module/nomodule pattern permet de :
\begin{itemize}
\item[A)] Supprimer les modules
\item[B)] Fournir un bundle moderne pour navigateurs récents + fallback
\item[C)] Compresser les modules
\item[D)] Chiffrer les modules
\end{itemize}
*(Réponse : B)*

\item Dans une configuration CDN, "Purge sélective" signifie :
\begin{itemize}
\item[A)] Supprimer tout le cache
\item[B)] Invalider le cache uniquement pour les ressources modifiées
\item[C)] Ajouter des ressources au cache
\item[D)] Désactiver le cache
\end{itemize}
*(Réponse : B)*

\item Un en-tête \texttt{Cache-Control: immutable} indique que :
\begin{itemize}
\item[A)] La ressource peut être modifiée à tout moment
\item[B)] La ressource ne changera jamais et peut être mise en cache indéfiniment
\item[C)] La ressource ne doit pas être mise en cache
\item[D)] La ressource est confidentielle
\end{itemize}
*(Réponse : B)*

\item Le Dynamic Import en JavaScript s'écrit avec la syntaxe :
\begin{itemize}
\item[A)] require()
\item[B)] import()
\item[C)] include()
\item[D)] load()
\end{itemize}
*(Réponse : B)*

\item La priorité de chargement \texttt{fetchpriority="low"} est utilisée pour :
\begin{itemize}
\item[A)] Les ressources critiques
\item[B)] Les ressources non essentielles
\item[C)] Les polices web
\item[D)] Le CSS critique
\end{itemize}
*(Réponse : B)*

\item Le "Vendor splitting" dans le code splitting consiste à :
\begin{itemize}
\item[A)] Diviser le code en plusieurs bundles
\item[B)] Séparer les bibliothèques tierces du code applicatif
\item[C)] Supprimer les vendors inutilisés
\item[D)] Fusionner tous les vendors
\end{itemize}
*(Réponse : B)*

\item Les en-têtes HTTP Last-Modified sont utilisés pour :
\begin{itemize}
\item[A)] Chiffrer les données
\item[B)] La validation conditionnelle du cache
\item[C)] Rediriger le client
\item[D)] Compresser les données
\end{itemize}
*(Réponse : B)*

\item Le format JPEG XL offre quel avantage par rapport au JPEG?
\begin{itemize}
\item[A)] Il est plus volumineux
\item[B)] Compression sans perte et avec perte, encodage progressif
\item[C)] Il est moins compatible
\item[D)] Il ne supporte pas les couleurs
\end{itemize}
*(Réponse : B)*

\item Le cache SQL (mise en cache des requêtes fréquentes) utilise généralement :
\begin{itemize}
\item[A)] MySQL
\item[B)] Redis
\item[C)] Apache
\item[D)] NGINX
\end{itemize}
*(Réponse : B)*

\item La priorité de chargement \texttt{fetchpriority="high"} doit être utilisée pour :
\begin{itemize}
\item[A)] Toutes les images
\item[B)] L'élément LCP de la page
\item[C)] Les images en lazy loading
\item[D)] Les vidéos
\end{itemize}
*(Réponse : B)*

\item HTTP/3 offre quelle amélioration pour le changement de réseau?
\begin{itemize}
\item[A)] Aucune amélioration
\item[B)] Migration de connexion sans reconnexion
\item[C)] Connexion plus lente
\item[D)] Plus de latence
\end{itemize}
*(Réponse : B)*

\item L'en-tête Vary est utilisé dans le contexte CDN pour :
\begin{itemize}
\item[A)] La compression
\item[B)] Indiquer que le contenu varie selon un en-tête (ex: User-Agent)
\item[C)] Le chiffrement
\item[D)] Les redirections
\end{itemize}
*(Réponse : B)*

\end{enumerate}

% ============================================
% CHAPITRE 15 : Le Mobile-First Indexing
% ============================================
\section*{Chapitre 15 : Le Mobile-First Indexing}

\begin{enumerate}
\item Depuis quand Google utilise l'indexation mobile-first comme méthode par défaut?
\begin{itemize}
\item[A)] 2017
\item[B)] 2019
\item[C)] 2021
\item[D)] 2015
\end{itemize}
*(Réponse : B)*

\item Que signifie "Mobile-First Indexing"?
\begin{itemize}
\item[A)] Google utilise la version desktop pour l'indexation
\item[B)] Google utilise la version mobile du site pour l'indexation et le classement
\item[C)] Google n'indexe que les sites mobiles
\item[D)] Google pénalise les sites sans application mobile
\end{itemize}
*(Réponse : B)*

\item Quelle est la configuration recommandée par Google pour le mobile-first?
\begin{itemize}
\item[A)] URLs séparées (m.exemple.com)
\item[B)] Design responsive
\item[C)] Dynamic Serving
\item[D)] Application mobile dédiée
\end{itemize}
*(Réponse : B)*

\item Le design responsive utilise :
\begin{itemize}
\item[A)] Des URLs différentes pour mobile et desktop
\item[B)] Le même HTML avec du CSS adaptatif via @media queries
\item[C)] Un serveur dédié pour mobile
\item[D)] Une application native
\end{itemize}
*(Réponse : B)*

\item Qu'est-ce que le Dynamic Serving?
\begin{itemize}
\item[A)] Même URL, HTML identique pour tous les appareils
\item[B)] Même URL, HTML différent selon l'user-agent
\item[C)] URLs différentes pour chaque appareil
\item[D)] Pas de HTML, uniquement du JavaScript
\end{itemize}
*(Réponse : B)*

\item Quelle structure d'URL est déconseillée par Google?
\begin{itemize}
\item[A)] exemple.com
\item[B)] m.exemple.com
\item[C)] exemple.com/produits
\item[D)] www.exemple.com
\end{itemize}
*(Réponse : B)*

\item Dans un contexte mobile-first, les données structurées doivent être :
\begin{itemize}
\item[A)] Uniquement sur la version desktop
\item[B)] Identiques sur les versions mobile et desktop
\item[C)] Uniquement sur la version mobile
\item[D)] Absentes du site
\end{itemize}
*(Réponse : B)*

\item Les meta robots doivent être :
\begin{itemize}
\item[A)] Différents sur mobile et desktop
\item[B)] Identiques sur les deux versions
\item[C)] Présents uniquement sur mobile
\item[D)] Supprimés du site
\end{itemize}
*(Réponse : B)*

\item Faut-il créer un sitemap séparé pour la version mobile?
\begin{itemize}
\item[A)] Oui, obligatoirement
\item[B)] Non, utiliser un sitemap unique
\item[C)] Oui, si le site utilise des URLs séparées
\item[D)] Non, Google ne lit pas les sitemaps
\end{itemize}
*(Réponse : B)*

\item Quel outil Google permet de tester l'optimisation mobile?
\begin{itemize}
\item[A)] Google Analytics
\item[B)] Test d'optimisation mobile
\item[C)] Google Trends
\item[D)] Google Ads
\end{itemize}
*(Réponse : B)*

\item Les Core Web Vitals doivent être optimisés en priorité sur :
\begin{itemize}
\item[A)] Desktop uniquement
\item[B)] Mobile
\item[C)] Tablette
\item[D)] Tous les appareils de manière égale
\end{itemize}
*(Réponse : B)*

\item Google render-t-il le contenu caché dans des tabs ou accordéons?
\begin{itemize}
\item[A)] Non, jamais
\item[B)] Oui, Google render tout maintenant
\item[C)] Uniquement sur desktop
\item[D)] Uniquement avec JavaScript
\end{itemize}
*(Réponse : B)*

\item Que se passe-t-il si la version mobile contient moins de contenu que la desktop?
\begin{itemize}
\item[A)] Google utilise la version desktop
\item[B)] Google peut ne pas indexer le contenu manquant
\item[C)] Google fusionne les deux versions
\item[D)] Google pénalise le site
\end{itemize}
*(Réponse : B)*

\item Le viewport doit être configuré comment pour un site responsive?
\begin{itemize}
\item[A)] <meta name="viewport" content="width=device-width, initial-scale=1">
\item[B)] <meta name="viewport" content="width=1024">
\item[C)] <meta name="viewport" content="user-scalable=no">
\item[D)] Il n'est pas nécessaire
\end{itemize}
*(Réponse : A)*

\item Comment Google détecte-t-il la version mobile d'un site?
\begin{itemize}
\item[A)] Via le user-agent du navigateur
\item[B)] Via la balise viewport
\item[C)] Via le fichier robots.txt
\item[D)] Via le sitemap XML
\end{itemize}
*(Réponse : A)*

\item Sur mobile, les popups intrusifs sont :
\begin{itemize}
\item[A)] Recommandés
\item[B)] Pénalisés par Google
\item[C)] Obligatoires
\item[D)] Sans impact SEO
\end{itemize}
*(Réponse : B)*

\item La vitesse de chargement est-elle plus importante sur mobile ou desktop?
\begin{itemize}
\item[A)] Plus importante sur desktop
\item[B)] Plus importante sur mobile
\item[C)] Aussi importante sur les deux
\item[D)] Sans importance
\end{itemize}
*(Réponse : B)*

\item Le Dynamic Serving utilise quel en-tête pour identifier le type d'appareil?
\begin{itemize}
\item[A)] Accept-Language
\item[B)] User-Agent
\item[C)] Cache-Control
\item[D)] Content-Type
\end{itemize}
*(Réponse : B)*

\item Le "Mobilegeddon" (2015) était une mise à jour qui :
\begin{itemize}
\item[A)] Pénalisait les sites lents
\item[B)] Récompensait les sites compatibles mobiles
\item[C)] Introduisait le machine learning
\item[D)] Modifiait le PageRank
\end{itemize}
*(Réponse : B)*

\item Les feuilles de style CSS doivent-elles être bloquées par robots.txt?
\begin{itemize}
\item[A)] Oui
\item[B)] Non, Google en a besoin pour le rendu mobile
\item[C)] Uniquement sur desktop
\item[D)] Uniquement les fichiers volumineux
\end{itemize}
*(Réponse : B)*

\item Le JavaScript doit-il être accessible par Googlebot sur mobile?
\begin{itemize}
\item[A)] Non, Googlebot ne l'exécute pas
\item[B)] Oui, Googlebot doit pouvoir accéder au JS pour le rendu
\item[C)] Uniquement les fichiers externes
\item[D)] Uniquement les scripts inline
\end{itemize}
*(Réponse : B)*

\item Différence entre responsive design et dynamic serving?
\begin{itemize}
\item[A)] Aucune, c'est la même chose
\item[B)] Le responsive utilise le même HTML, le dynamic serving sert du HTML différent
\item[C)] Le responsive est obsolète
\item[D)] Le dynamic serving nécessite des URLs séparées
\end{itemize}
*(Réponse : B)*

\item Google recommande-t-il le même contenu sur mobile et desktop?
\begin{itemize}
\item[A)] Non, le contenu mobile doit être plus court
\item[B)] Oui, le contenu doit être équivalent
\item[C)] Le contenu mobile doit être plus long
\item[D)] Seul le contenu desktop compte
\end{itemize}
*(Réponse : B)*

\item Quel est le pourcentage de recherches Google ayant une intention locale?
\begin{itemize}
\item[A)] 26\%
\item[B)] 46\%
\item[C)] 66\%
\item[D)] 86\%
\end{itemize}
*(Réponse : B)*

\item L'image "hero" d'une page mobile doit être :
\begin{itemize}
\item[A)] Chargée en lazy loading
\item[B)] Préchargée avec fetchpriority="high"
\item[C)] Supprimée
\item[D)] Redimensionnée en 100x100
\end{itemize}
*(Réponse : B)*

\item Le format d'image recommandé pour le mobile est :
\begin{itemize}
\item[A)] BMP
\item[B)] WebP ou AVIF
\item[C)] TIFF
\item[D)] PSD
\end{itemize}
*(Réponse : B)*

\item Les animations CSS complexes sont-elles recommandées sur mobile?
\begin{itemize}
\item[A)] Oui, sans limite
\item[B)] Non, elles peuvent impacter les performances
\item[C)] Uniquement sur desktop
\item[D)] Uniquement les animations JavaScript
\end{itemize}
*(Réponse : B)*

\item L'indexation mobile-first s'applique-t-elle à tous les sites?
\begin{itemize}
\item[A)] Non, uniquement aux nouveaux sites
\item[B)] Oui, c'est la méthode par défaut pour tous
\item[C)] Non, uniquement aux sites e-commerce
\item[D)] Non, uniquement aux sites avec version mobile
\end{itemize}
*(Réponse : B)*

\item Pour un site responsive, le nombre de sitemaps doit être :
\begin{itemize}
\item[A)] 2 (mobile + desktop)
\item[B)] 1 seul
\item[C)] 3 (mobile, desktop, tablette)
\item[D)] aucun
\end{itemize}
*(Réponse : B)*

\item Quel attribut CSS est utilisé pour adapter le design à la taille de l'écran?
\begin{itemize}
\item[A)] @media screen
\item[B)] @media (max-width: ...)
\item[C)] @media mobile
\item[D)] @media responsive
\end{itemize}
*(Réponse : B)*

\item La largeur du viewport sur mobile doit être définie à :
\begin{itemize}
\item[A)] 1024px
\item[B)] device-width
\item[C)] 100\%
\item[D)] auto
\end{itemize}
*(Réponse : B)*

\item Le texte sur mobile doit avoir une taille de police d'au moins :
\begin{itemize}
\item[A)] 12px
\item[B)] 16px
\item[C)] 20px
\item[D)] 8px
\end{itemize}
*(Réponse : B)*

\item Google Search Console fournit un rapport spécifique pour :
\begin{itemize}
\item[A)] Le trafic desktop
\item[B)] Les problèmes d'utilisabilité mobile
\item[C)] Les backlinks
\item[D)] Les données structurées
\end{itemize}
*(Réponse : B)*

\item Le chargement différé des ressources non critiques sur mobile utilise :
\begin{itemize}
\item[A)] loading="eager"
\item[B)] loading="lazy"
\item[C)] loading="auto"
\item[D)] loading="async"
\end{itemize}
*(Réponse : B)*

\item Les temps de chargement sur mobile doivent être inférieurs à :
\begin{itemize}
\item[A)] 5 secondes
\item[B)] 3 secondes
\item[C)] 10 secondes
\item[D)] 1 seconde
\end{itemize}
*(Réponse : B)*

\item Le contenu "above the fold" sur mobile doit être :
\begin{itemize}
\item[A)] Chargé en lazy loading
\item[B)] Chargé en priorité
\item[C)] Supprimé
\item[D)] Identique au desktop
\end{itemize}
*(Réponse : B)*

\item Les éléments tactiles sur mobile doivent avoir une taille minimale de :
\begin{itemize}
\item[A)] 24x24px
\item[B)] 48x48px
\item[C)] 12x12px
\item[D)] 96x96px
\end{itemize}
*(Réponse : B)*

\item Un site avec m.exemple.com doit utiliser quelle balise de liaison?
\begin{itemize}
\item[A)] <link rel="canonical">
\item[B)] <link rel="alternate" media="only screen and (max-width: 640px)">
\item[C)] <meta name="robots">
\item[D)] <link rel="preload">
\end{itemize}
*(Réponse : B)*

\item Google recommande-t-il AMP pour le mobile-first indexing?
\begin{itemize}
\item[A)] Oui, obligatoire
\item[B)] Non, AMP n'est plus requis
\item[C)] Oui, fortement recommandé
\item[D)] AMP est interdit
\end{itemize}
*(Réponse : B)*

\item Le Vary HTTP header est important en dynamic serving pour :
\begin{itemize}
\item[A)] La compression
\item[B)] La gestion du cache CDN
\item[C)] Le chiffrement
\item[D)] Les redirections
\end{itemize}
*(Réponse : B)*

\item Sur mobile, les images doivent utiliser quel attribut pour l'adaptation responsive?
\begin{itemize}
\item[A)] width="100\%"
\item[B)] srcset et sizes
\item[C)] height="auto"
\item[D)] display="responsive"
\end{itemize}
*(Réponse : B)*

\item Le test "Mobile-Friendly" de Google vérifie :
\begin{itemize}
\item[A)] La vitesse du site
\item[B)] La compatibilité tactile et la lisibilité sur mobile
\item[C)] Le nombre de pages mobiles
\item[D)] Les backlinks mobiles
\end{itemize}
*(Réponse : B)*

\item Googlebot mobile utilise quel user-agent pour crawler?
\begin{itemize}
\item[A)] Googlebot Desktop
\item[B)] Googlebot Mobile / Chrome Mobile
\item[C)] Googlebot Image
\item[D)] Googlebot Video
\end{itemize}
*(Réponse : B)*

\item Les espaces réservés (ad slots) dans les publicités sont importants pour :
\begin{itemize}
\item[A)] Le LCP
\item[B)] Le CLS sur mobile
\item[C)] L'INP
\item[D)] Le TTFB
\end{itemize}
*(Réponse : B)*

\item La navigation au pouce (thumb navigation) est une considération pour :
\begin{itemize}
\item[A)] Le SEO desktop
\item[B)] Le design UX mobile
\item[C)] Le JavaScript SEO
\item[D)] Les données structurées
\end{itemize}
*(Réponse : B)*

\item Un site avec responsive design et dynamic serving peut utiliser :
\begin{itemize}
\item[A)] La même balise canonical pour les deux versions
\item[B)] Des balises canonical adaptées si les contenus diffèrent
\item[C)] Pas de canonical
\item[D)] Un canonical vers la version desktop uniquement
\end{itemize}
*(Réponse : B)*

\item Les balises meta robots peuvent être utilisées pour contrôler l'indexation :
\begin{itemize}
\item[A)] Uniquement sur desktop
\item[B)] De manière identique sur mobile et desktop recommandée
\item[C)] Uniquement sur mobile
\item[D)] Avec des valeurs différentes par appareil
\end{itemize}
*(Réponse : B)*

\item Le texte alternatif des images sur mobile doit être :
\begin{itemize}
\item[A)] Plus court que sur desktop
\item[B)] Présent et descriptif, comme sur desktop
\item[C)] Supprimé pour économiser de la bande passante
\item[D)] Remplacé par des attributs title
\end{itemize}
*(Réponse : B)*

\item Le chargement prioritaire (preload) des ressources est particulièrement important sur mobile car :
\begin{itemize}
\item[A)] Le réseau mobile est plus rapide
\item[B)] La bande passante mobile est limitée
\item[C)] Les mobiles ont plus de mémoire
\item[D)] Le réseau mobile est toujours stable
\end{itemize}
*(Réponse : B)*

\item La conversion des images en WebP sur mobile permet de :
\begin{itemize}
\item[A)] Augmenter le poids des pages
\item[B)] Réduire le temps de chargement et améliorer le LCP
\item[C)] Augmenter la qualité des images
\item[D)] Ajouter des animations
\end{itemize}
*(Réponse : B)*

\end{enumerate}
% ============================================
% CHAPITRE 16 : Les données structurées (Schema Markup)
% ============================================
\section*{Chapitre 16 : Les données structurées (Schema Markup)}

\begin{enumerate}
\item Qu'est-ce que Schema.org?
\begin{itemize}
\item[A)] Un moteur de recherche
\item[B)] Un vocabulaire standardisé de données structurées
\item[C)] Un framework JavaScript
\item[D)] Un outil d'analyse SEO
\end{itemize}
*(Réponse : B)*

\item Quel format d'implémentation des données structurées est recommandé par Google?
\begin{itemize}
\item[A)] Microdata
\item[B)] JSON-LD
\item[C)] RDFa
\item[D)] XML
\end{itemize}
*(Réponse : B)*

\item À quoi servent les données structurées?
\begin{itemize}
\item[A)] À augmenter la vitesse du site
\item[B)] À permettre aux moteurs de recherche de comprendre le contenu précisément
\item[C)] À crypter les données du site
\item[D)] À créer des formulaires
\end{itemize}
*(Réponse : B)*

\item Quel type de schema est utilisé pour un article de blog?
\begin{itemize}
\item[A)] Product
\item[B)] Article / BlogPosting
\item[C)] Event
\item[D)] Recipe
\end{itemize}
*(Réponse : B)*

\item Quel type de schema est utilisé pour une entreprise locale?
\begin{itemize}
\item[A)] Person
\item[B)] LocalBusiness
\item[C)] VideoObject
\item[D)] HowTo
\end{itemize}
*(Réponse : B)*

\item Quel type de schema est utilisé pour un produit?
\begin{itemize}
\item[A)] Article
\item[B)] Product
\item[C)] Event
\item[D)] FAQPage
\end{itemize}
*(Réponse : B)*

\item Quel type de schema est utilisé pour une page de questions fréquentes?
\begin{itemize}
\item[A)] HowTo
\item[B)] FAQPage
\item[C)] Recipe
\item[D)] BreadcrumbList
\end{itemize}
*(Réponse : B)*

\item Les Rich Snippets sont :
\begin{itemize}
\item[A)] Des extraits de code HTML
\item[B)] Des résultats de recherche enrichis avec des informations supplémentaires
\item[C)] Des scripts JavaScript
\item[D)] Des fichiers CSS
\end{itemize}
*(Réponse : B)*

\item Quel schema permet d'afficher des étoiles dans les SERP?
\begin{itemize}
\item[A)] Product
\item[B)] AggregateRating
\item[C)] Review
\item[D)] Article
\end{itemize}
*(Réponse : B)*

\item Quel schema est utilisé pour le fil d'Ariane (breadcrumb)?
\begin{itemize}
\item[A)] Navigation
\item[B)] BreadcrumbList
\item[C)] SiteNavigationElement
\item[D)] WebPage
\end{itemize}
*(Réponse : B)*

\item Les Microdata utilisent quel type d'attributs HTML?
\begin{itemize}
\item[A)] data-*
\item[B)] itemscope, itemtype, itemprop
\item[C)] Des classes CSS
\item[D)] Des ID
\end{itemize}
*(Réponse : B)*

\item Le type de schema pour une vidéo est :
\begin{itemize}
\item[A)] Media
\item[B)] VideoObject
\item[C)] Video
\item[D)] Movie
\end{itemize}
*(Réponse : B)*

\item Le type de schema pour un événement est :
\begin{itemize}
\item[A)] Calendar
\item[B)] Event
\item[C)] Meeting
\item[D)] Appointment
\end{itemize}
*(Réponse : B)*

\item Le type de schema pour une recette est :
\begin{itemize}
\item[A)] Food
\item[B)] Recipe
\item[C)] Cooking
\item[D)] Nutrition
\end{itemize}
*(Réponse : B)*

\item Le type de schema pour un tutoriel étape par étape est :
\begin{itemize}
\item[A)] Tutorial
\item[B)] HowTo
\item[C)] Guide
\item[D)] StepByStep
\end{itemize}
*(Réponse : B)*

\item Le type de schema pour une personne est :
\begin{itemize}
\item[A)] Human
\item[B)] Person
\item[C)] Individual
\item[D)] Profile
\end{itemize}
*(Réponse : B)*

\item Les données structurées JSON-LD sont placées dans :
\begin{itemize}
\item[A)] Un fichier .json séparé
\item[B)] Une balise <script type="application/ld+json"> dans le HTML
\item[C)] Un attribut HTML
\item[D)] Le fichier robots.txt
\end{itemize}
*(Réponse : B)*

\item Google utilise-t-il les données structurées comme facteur de classement direct?
\begin{itemize}
\item[A)] Oui, c'est un facteur majeur
\item[B)] Non, mais elles permettent les Rich Snippets qui améliorent le CTR
\item[C)] Oui, c'est le facteur le plus important
\item[D)] Non, Google les ignore
\end{itemize}
*(Réponse : B)*

\item Quel outil Google permet de tester les données structurées?
\begin{itemize}
\item[A)] Google Analytics
\item[B)] Rich Results Test
\item[C)] Google Trends
\item[D)] Google Ads
\end{itemize}
*(Réponse : B)*

\item Le schema Organization est utilisé pour :
\begin{itemize}
\item[A)] Une personne
\item[B)] Une entreprise ou organisation
\item[C)] Un produit
\item[D)] Un événement
\end{itemize}
*(Réponse : B)*

\item Le schema LocalBusiness hérite de quel type parent?
\begin{itemize}
\item[A)] Product
\item[B)] Organization
\item[C)] Place
\item[D)] Thing
\end{itemize}
*(Réponse : B)*

\item Combien de types de schema sont disponibles sur Schema.org?
\begin{itemize}
\item[A)] Environ 100
\item[B)] Plus de 800
\item[C)] Environ 500
\item[D)] Plus de 2000
\end{itemize}
*(Réponse : B)*

\item Schema.org a été développé par :
\begin{itemize}
\item[A)] Uniquement Google
\item[B)] Google, Microsoft, Yahoo et Yandex
\item[C)] Uniquement Microsoft
\item[D)] Un consortium open source indépendant
\end{itemize}
*(Réponse : B)*

\item Une même page peut-elle avoir plusieurs types de schema?
\begin{itemize}
\item[A)] Non, un seul type par page
\item[B)] Oui, on peut combiner plusieurs types
\item[C)] Non, cela provoque des erreurs
\item[D)] Uniquement avec Microdata
\end{itemize}
*(Réponse : B)*

\item Le schema FAQPage peut générer quel type de résultat enrichi?
\begin{itemize}
\item[A)] Des étoiles
\item[B)] Une liste de questions-réponses déroulante dans les SERP
\item[C)] Un prix
\item[D)] Une vidéo
\end{itemize}
*(Réponse : B)*

\item Le schema HowTo peut générer quel type de résultat enrichi?
\begin{itemize}
\item[A)] Un carrousel d'images
\item[B)] Des étapes avec instructions dans les SERP
\item[C)] Un prix
\item[D)] Un avis
\end{itemize}
*(Réponse : B)*

\item Le schema Product avec AggregateRating peut afficher :
\begin{itemize}
\item[A)] La date de publication
\item[B)] Le prix et les étoiles de notation
\item[C)] L'auteur
\item[D)] La catégorie
\end{itemize}
*(Réponse : B)*

\item Le schema VideoObject permet d'obtenir :
\begin{itemize}
\item[A)] Un classement amélioré
\item[B)] Des Rich Snippets avec vignette vidéo dans les SERP
\item[C)] Des backlinks automatiques
\item[D)] Une indexation plus rapide
\end{itemize}
*(Réponse : B)*

\item Le JSON-LD doit être placé dans quelle partie du HTML?
\begin{itemize}
\item[A)] Uniquement dans le <head>
\item[B)] Dans le <head> ou le <body>
\item[C)] Uniquement dans le <body>
\item[D)] Dans le footer uniquement
\end{itemize}
*(Réponse : B)*

\item Le schema Event peut inclure quel champ pour la localisation?
\begin{itemize}
\item[A)] address
\item[B)] location
\item[C)] place
\item[D)] venue
\end{itemize}
*(Réponse : B)*

\item Le schema Person peut inclure quel champ pour les réseaux sociaux?
\begin{itemize}
\item[A)] social
\item[B)] sameAs
\item[C)] network
\item[D)] profile
\end{itemize}
*(Réponse : B)*

\item Le schema Article peut inclure quel champ pour l'auteur?
\begin{itemize}
\item[A)] writer
\item[B)] author
\item[C)] creator
\item[D)] editor
\end{itemize}
*(Réponse : B)*

\item Le schema Recipe peut inclure quel champ pour le temps de cuisson?
\begin{itemize}
\item[A)] time
\item[B)] cookTime
\item[C)] duration
\item[D)] timer
\end{itemize}
*(Réponse : B)*

\item Les données structurées NewsArticle sont destinées à :
\begin{itemize}
\item[A)] Les articles de blog classiques
\item[B)] Les articles d'actualité
\item[C)] Les fiches produits
\item[D)] Les pages de contact
\end{itemize}
*(Réponse : B)*

\item Le schema JobPosting est utilisé pour :
\begin{itemize}
\item[A)] Les offres d'emploi
\item[B)] Les tâches ménagères
\item[C)] Les projets
\item[D)] Les stages
\end{itemize}
*(Réponse : A)*

\item Le schema Course est utilisé pour :
\begin{itemize}
\item[A)] Les formations et cours en ligne
\item[B)] Les courses sportives
\item[C)] Les courses alimentaires
\item[D)] Les projets
\end{itemize}
*(Réponse : A)*

\item Une erreur dans les données structurées peut-elle impacter le SEO?
\begin{itemize}
\item[A)] Non, Google les ignore si elles sont erronées
\item[B)] Oui, Google peut ne pas afficher les Rich Snippets
\item[C)] Oui, cela pénalise le classement
\item[D)] Non, les erreurs sont automatiquement corrigées
\end{itemize}
*(Réponse : B)*

\item Le schema SoftwareApplication est utilisé pour :
\begin{itemize}
\item[A)] Les applications mobiles
\item[B)] Les logiciels
\item[C)] Les sites web
\item[D)] Les APIs
\end{itemize}
*(Réponse : B)*

\item Le schema Book est utilisé pour :
\begin{itemize}
\item[A)] Les articles
\item[B)] Les livres
\item[C)] Les chapitres
\item[D)] Les documents
\end{itemize}
*(Réponse : B)*

\item Le schema MusicAlbum est utilisé pour :
\begin{itemize}
\item[A)] Les concerts
\item[B)] Les albums musicaux
\item[C)] Les playlists
\item[D)] Les artistes
\end{itemize}
*(Réponse : B)*

\item Le schema Movie est utilisé pour :
\begin{itemize}
\item[A)] Les vidéos YouTube
\item[B)] Les films
\item[C)] Les séries TV
\item[D)] Les documentaires
\end{itemize}
*(Réponse : B)*

\item Le schema Hotel est une spécialisation de quel type?
\begin{itemize}
\item[A)] Organization
\item[B)] LocalBusiness
\item[C)] Place
\item[D)] Accommodation
\end{itemize}
*(Réponse : B)*

\item Le schema Review est généralement associé à quel autre schema?
\begin{itemize}
\item[A)] Product, LocalBusiness, ou CreativeWork
\item[B)] Uniquement Product
\item[C)] Uniquement Article
\item[D)] Uniquement Event
\end{itemize}
*(Réponse : A)*

\item Les données structurées Sitelinks Search Box permettent de :
\begin{itemize}
\item[A)] Ajouter une barre de recherche Google sur le site
\item[B)] Afficher une barre de recherche dans les SERP pour le site
\item[C)] Créer un moteur de recherche interne
\item[D)] Bloquer la recherche Google
\end{itemize}
*(Réponse : B)*

\item Le schema SocialMediaPosting est utilisé pour :
\begin{itemize}
\item[A)] Les réseaux sociaux du site
\item[B)] Les publications sur les réseaux sociaux
\item[C)] Les partages sociaux
\item[D)] Les boutons sociaux
\end{itemize}
*(Réponse : B)*

\item Le schema MedicalCondition fait partie de quel domaine?
\begin{itemize}
\item[A)] Health
\item[B)] Medical
\item[C)] Wellness
\item[D)] Clinical
\end{itemize}
*(Réponse : B)*

\item L'attribut "itemprop" est utilisé dans quel format de données structurées?
\begin{itemize}
\item[A)] JSON-LD
\item[B)] Microdata
\item[C)] RDFa
\item[D)] XML
\end{itemize}
*(Réponse : B)*

\item Le schema BreadcrumbList utilise quel type d'élément de liste?
\begin{itemize}
\item[A)] ListItem
\item[B)] Breadcrumb
\item[C)] Step
\item[D)] Element
\end{itemize}
*(Réponse : A)*

\item Le schema "Offer" est utilisé en conjonction avec :
\begin{itemize}
\item[A)] Article
\item[B)] Product
\item[C)] Event
\item[D)] Person
\end{itemize}
*(Réponse : B)*

\item Le schema "Restaurant" descend de quel type parent?
\begin{itemize}
\item[A)] Bar
\item[B)] LocalBusiness
\item[C)] FoodEstablishment
\item[D)] Place
\end{itemize}
*(Réponse : C)*
\end{enumerate}

% ============================================
% CHAPITRE 17 : Le JavaScript SEO
% ============================================
\section*{Chapitre 17 : Le JavaScript SEO}

\begin{enumerate}
\item Quel est le principal défi du JavaScript SEO?
\begin{itemize}
\item[A)] Le JavaScript est plus rapide que le HTML
\item[B)] Le contenu chargé par JavaScript peut ne pas être vu par Googlebot
\item[C)] Le JavaScript améliore le classement
\item[D)] Le JavaScript est interdit par Google
\end{itemize}
*(Réponse : B)*

\item Combien de temps Googlebot a-t-il environ pour render une page sur mobile?
\begin{itemize}
\item[A)] 10 secondes
\item[B)] 30 secondes
\item[C)] 60 secondes
\item[D)] 5 secondes
\end{itemize}
*(Réponse : B)*

\item Que risque-t-on si des ressources CSS ou JS essentielles sont bloquées par robots.txt?
\begin{itemize}
\item[A)] Rien, le site fonctionne normalement
\item[B)] Le contenu peut ne pas être indexé correctement
\item[C)] Le site est pénalisé
\item[D)] Le site charge plus vite
\end{itemize}
*(Réponse : B)*

\item Combien d'étapes Googlebot traverse-t-il pour indexer une page JavaScript?
\begin{itemize}
\item[A)] 1
\item[B)] 2
\item[C)] 3
\item[D)] 4
\end{itemize}
*(Réponse : C)*

\item Quelle est la première étape de l'indexation d'une page JS par Google?
\begin{itemize}
\item[A)] Google exécute immédiatement le JavaScript
\item[B)] Googlebot télécharge le HTML brut
\item[C)] Google ignore la page
\item[D)] Google met la page en file d'attente
\end{itemize}
*(Réponse : B)*

\item Qu'est-ce que le Server-Side Rendering (SSR)?
\begin{itemize}
\item[A)] Le HTML est généré par le navigateur
\item[B)] Le HTML complet est généré sur le serveur avant d'être envoyé au client
\item[C)] Le JavaScript est supprimé
\item[D)] Les images sont rendues côté serveur
\end{itemize}
*(Réponse : B)*

\item Quel est l'avantage principal du SSR?
\begin{itemize}
\item[A)] Coût serveur réduit
\item[B)] Contenu 100\% visible par Googlebot
\item[C)] Temps de développement réduit
\item[D)] Maintenance simplifiée
\end{itemize}
*(Réponse : B)*

\item Quel framework permet le SSR avec React?
\begin{itemize}
\item[A)] Vue.js
\item[B)] Next.js
\item[C)] Angular
\item[D)] Svelte
\end{itemize}
*(Réponse : B)*

\item Quel framework permet le SSR avec Vue.js?
\begin{itemize}
\item[A)] Next.js
\item[B)] Nuxt.js
\item[C)] Nest.js
\item[D)] Express
\end{itemize}
*(Réponse : B)*

\item Qu'est-ce que la Static Site Generation (SSG)?
\begin{itemize}
\item[A)] Les pages sont générées dynamiquement à chaque requête
\item[B)] Les pages HTML statiques sont générées au moment du build
\item[C)] Les pages sont générées par le client
\item[D)] Les pages sont générées par JavaScript
\end{itemize}
*(Réponse : B)*

\item Quel est l'inconvénient principal du SSG?
\begin{itemize}
\item[A)] Performances médiocres
\item[B)] Contenu statique, régénération nécessaire pour les mises à jour
\item[C)] Coût serveur élevé
\item[D)] Sécurité réduite
\end{itemize}
*(Réponse : B)*

\item Qu'est-ce que l'hydratation (hydration)?
\begin{itemize}
\item[A)] Le processus d'ajout d'eau au serveur
\item[B)] Le processus par lequel une application JS reprend le contrôle d'un HTML rendu par le serveur
\item[C)] Le processus de minification du code
\item[D)] Le processus de compression des images
\end{itemize}
*(Réponse : B)*

\item Qu'est-ce que l'Island Architecture?
\begin{itemize}
\item[A)] Une architecture où tout le JavaScript est chargé au démarrage
\item[B)] Seulement des îlots d'interactivité dans une page autrement statique
\item[C)] Une architecture sans JavaScript
\item[D)] Une architecture avec uniquement du JavaScript
\end{itemize}
*(Réponse : B)*

\item Quel outil permet de visualiser le HTML rendu par Google après exécution JS?
\begin{itemize}
\item[A)] Google Analytics
\item[B)] L'outil d'inspection d'URL dans GSC
\item[C)] Google Trends
\item[D)] PageSpeed Insights
\end{itemize}
*(Réponse : B)*

\item Que vérifier en priorité si le contenu JS n'apparaît pas dans le HTML rendu?
\begin{itemize}
\item[A)] La version de JavaScript
\item[B)] Que le JavaScript n'est pas bloqué par robots.txt et les API sont accessibles
\item[C)] La taille du fichier JavaScript
\item[D)] Le nom des fonctions
\end{itemize}
*(Réponse : B)*

\item Quel navigateur Googlebot utilise-t-il pour le rendu JavaScript?
\begin{itemize}
\item[A)] Firefox
\item[B)] Chrome
\item[C)] Safari
\item[D)] Opera
\end{itemize}
*(Réponse : B)*

\item Le JavaScript SEO est particulièrement important pour quel type d'application?
\begin{itemize}
\item[A)] Les sites statiques HTML
\item[B)] Les Single Page Applications (SPA)
\item[C)] Les sites WordPress
\item[D)] Les blogs simples
\end{itemize}
*(Réponse : B)*

\item L'hydratation progressive consiste à :
\begin{itemize}
\item[A)] Hydrater tous les composants en même temps
\item[B)] Hydrater par priorité les composants visibles
\item[C)] Ne pas hydrater les composants
\item[D)] Hydrater uniquement le header
\end{itemize}
*(Réponse : B)*

\item L'hydratation partielle consiste à :
\begin{itemize}
\item[A)] Hydrater tous les composants
\item[B)] Hydrater seulement les composants interactifs
\item[C)] Ne pas hydrater la page
\item[D)] Hydrater uniquement le footer
\end{itemize}
*(Réponse : B)*

\item Quelle stratégie de rendu offre la meilleure performance pour du contenu statique?
\begin{itemize}
\item[A)] SSR
\item[B)] SSG
\item[C)] CSR (Client-Side Rendering)
\item[D)] Rendre côté client uniquement
\end{itemize}
*(Réponse : B)*

\item Quel est l'inconvénient du SSR?
\begin{itemize}
\item[A)] Le contenu n'est pas visible par Googlebot
\item[B)] Coût serveur plus élevé et TTFB potentiellement plus long
\item[C)] Le site ne fonctionne pas sans JavaScript
\item[D)] La sécurité est réduite
\end{itemize}
*(Réponse : B)*

\item Google Search Console dispose d'un rapport spécifique pour :
\begin{itemize}
\item[A)] Les images
\item[B)] Les pages avec JavaScript
\item[C)] Les vidéos
\item[D)] Les fichiers CSS
\end{itemize}
*(Réponse : B)*

\item Le mode Googlebot dans Chrome DevTools permet de :
\begin{itemize}
\item[A)] Accélérer le site
\item[B)] Tester comment Google voit la page
\item[C)] Modifier le contenu
\item[D)] Supprimer le JavaScript
\end{itemize}
*(Réponse : B)*

\item Les SPA ont quels défis SEO?
\begin{itemize}
\item[A)] Aucun défi spécifique
\item[B)] Gestion d'état complexe et contenu potentiellement non indexable
\item[C)] Elles sont toujours mieux classées
\item[D)] Elles chargent plus vite
\end{itemize}
*(Réponse : B)*

\item Quel outil peut crawler un site en exécutant le JavaScript?
\begin{itemize}
\item[A)] Google Analytics
\item[B)] Screaming Frog SEO Spider (mode JavaScript rendering)
\item[C)] Google Trends
\item[D)] Google Keyword Planner
\end{itemize}
*(Réponse : B)*

\item Qu'est-ce que le Client-Side Rendering (CSR)?
\begin{itemize}
\item[A)] Le HTML est généré sur le serveur
\item[B)] Le navigateur exécute le JavaScript pour générer le contenu
\item[C)] Le contenu est statique
\item[D)] Le rendu est fait par Googlebot
\end{itemize}
*(Réponse : B)*

\item Astro est un framework qui implémente :
\begin{itemize}
\item[A)] SSR uniquement
\item[B)] Island Architecture
\item[C)] CSR uniquement
\item[D)] SSG uniquement
\end{itemize}
*(Réponse : B)*

\item Le pre-rendering génère :
\begin{itemize}
\item[A)] Le HTML au moment de la requête
\item[B)] Des snapshots HTML statiques des pages pour les crawlers
\item[C)] Supprime le JavaScript
\item[D)] Ajoute du JavaScript
\end{itemize}
*(Réponse : B)*

\item Gatsby est un framework de type :
\begin{itemize}
\item[A)] SSR
\item[B)] SSG
\item[C)] CSR
\item[D)] Island Architecture
\end{itemize}
*(Réponse : B)*

\item Hugo est un générateur de site statique écrit en :
\begin{itemize}
\item[A)] JavaScript
\item[B)] Go
\item[C)] Python
\item[D)] Ruby
\end{itemize}
*(Réponse : B)*

\item Angular Universal permet de faire :
\begin{itemize}
\item[A)] SSG
\item[B)] SSR
\item[C)] CSR
\item[D)] Island Architecture
\end{itemize}
*(Réponse : B)*

\item Une page qui nécessite un clic pour charger le contenu est-elle bien indexée?
\begin{itemize}
\item[A)] Oui, parfaitement
\item[B)] Non, Googlebot ne clique pas
\item[C)] Oui, si le clic est simulé
\item[D)] Oui, avec JavaScript
\end{itemize}
*(Réponse : B)*

\item L'architecture Jamstack repose sur quel principe?
\begin{itemize}
\item[A)] JavaScript, API, Markup
\item[B)] Serveur, Base de données, Cache
\item[C)] HTML, CSS, JavaScript
\item[D)] Client, Serveur, Réseau
\end{itemize}
*(Réponse : A)*

\item Fresh est un framework qui implémente :
\begin{itemize}
\item[A)] CSR
\item[B)] Island Architecture
\item[C)] SSR uniquement
\item[D)] Aucun rendu
\end{itemize}
*(Réponse : B)*

\item Le code splitting en JavaScript permet de :
\begin{itemize}
\item[A)] Rendre le code plus complexe
\item[B)] Charger uniquement le code nécessaire pour la page courante
\item[C)] Supprimer tout le code
\item[D)] Doubler la taille du code
\end{itemize}
*(Réponse : B)*

\item Next.js permet quel type de rendu?
\begin{itemize}
\item[A)] SSR uniquement
\item[B)] SSG et SSR
\item[C)] CSR uniquement
\item[D)] Island Architecture uniquement
\end{itemize}
*(Réponse : B)*

\item Gatsby génère quel type de site?
\begin{itemize}
\item[A)] Site dynamique avec rendu côté serveur
\item[B)] Site statique (SSG) avec des données dynamiques via GraphQL
\item[C)] Site avec rendu côté client uniquement
\item[D)] Site sans JavaScript
\end{itemize}
*(Réponse : B)*

\item L'avantage du lazy loading de composants dans une SPA est de :
\begin{itemize}
\item[A)] Rendre l'application plus complexe
\item[B)] Réduire le poids initial du JavaScript téléchargé
\item[C)] Augmenter le nombre de requêtes
\item[D)] Rendre le code moins lisible
\end{itemize}
*(Réponse : B)*

\item Quelle méthode permet de vérifier si Google voit le contenu React?
\begin{itemize}
\item[A)] Regarder le code source dans le navigateur
\item[B)] Utiliser l'outil d'inspection d'URL dans Google Search Console
\item[C)] Utiliser React DevTools
\item[D)] Désactiver JavaScript dans le navigateur
\end{itemize}
*(Réponse : B)*

\item Le rendu Googlebot pour les pages JavaScript utilise :
\begin{itemize}
\item[A)] La version la plus récente de Chrome
\item[B)] Une version Chromium spécifique (généralement stable)
\item[C)] Firefox
\item[D)] Safari
\end{itemize}
*(Réponse : B)*

\item Le pattern PRPL (Push, Render, Pre-cache, Lazy-load) est une stratégie de :
\begin{itemize}
\item[A)] SEO
\item[B)] Performance web
\item[C)] Rédaction de contenu
\item[D)] Link building
\end{itemize}
*(Réponse : B)*

\item Un site utilisant le CSR pur présente quel risque SEO?
\begin{itemize}
\item[A)] Aucun risque
\item[B)] Google peut voir une page blanche ou incomplète
\item[C)] Le site sera plus rapide
\item[D)] Le site sera mieux classé
\end{itemize}
*(Réponse : B)*

\item Le "tree shaking" élimine le code JavaScript :
\begin{itemize}
\item[A)] Chargé dynamiquement
\item[B)] Non utilisé (dead code)
\item[C)] Minifié
\item[D)] Externe
\end{itemize}
*(Réponse : B)*

\item Quelle est la durée maximale approximative du rendu Googlebot sur mobile?
\begin{itemize}
\item[A)] 10 secondes
\item[B)] 30 secondes
\item[C)] 1 minute
\item[D)] 5 minutes
\end{itemize}
*(Réponse : B)*

\item Les ressources JavaScript bloquées par robots.txt peuvent causer :
\begin{itemize}
\item[A)] Une amélioration des performances
\item[B)] Un rendu incomplet et une perte d'indexation
\item[C)] Une augmentation du trafic
\item[D)] Un meilleur classement
\end{itemize}
*(Réponse : B)*

\item L'approche "Progressive Hydration" hydrate les composants :
\begin{itemize}
\item[A)] Tous en même temps
\item[B)] Par ordre de priorité (d'abord les visibles)
\item[C)] Aléatoirement
\item[D)] Uniquement ceux qui ont des bugs
\end{itemize}
*(Réponse : B)*

\item Googlebot utilise la même version de Chrome que les utilisateurs?
\begin{itemize}
\item[A)] Oui, exactement la même
\item[B)] Non, Googlebot utilise une version Chromium parfois légèrement plus ancienne
\item[C)] Oui, mais avec des extensions en moins
\item[D)] Non, Googlebot utilise Firefox
\end{itemize}
*(Réponse : B)*

\item Les applications web progressives (PWA) peuvent améliorer le SEO si :
\begin{itemize}
\item[A)] Elles utilisent uniquement du CSR
\item[B)] Elles sont couplées à du SSR ou du SSG
\item[C)] Elles n'ont pas de contenu HTML
\item[D)] Elles bloquent les crawlers
\end{itemize}
*(Réponse : B)*

\item Le rendu côté serveur (SSR) compare au CSR permet :
\begin{itemize}
\item[A)] Un coût serveur réduit
\item[B)] Une meilleure indexation et un LCP amélioré
\item[C)] Moins de code à écrire
\item[D)] Une sécurité accrue
\end{itemize}
*(Réponse : B)*

\item Le rendu côté client (CSR) peut poser problème pour les moteurs de recherche car :
\begin{itemize}
\item[A)] Le HTML est généré côté serveur
\item[B)] Le contenu est généré dynamiquement par JavaScript et peut ne pas être visible au crawl initial
\item[C)] Le CSS est ignoré
\item[D)] Les images ne se chargent pas
\end{itemize}
*(Réponse : B)*
\end{enumerate}
% ============================================
% CHAPITRE 18 : SEO International et Multilingue
% ============================================
\section*{Chapitre 18 : SEO International et Multilingue}

\begin{enumerate}
\item Qu'est-ce que le SEO international?
\begin{itemize}
\item[A)] Optimiser un site uniquement pour la France
\item[B)] Optimiser un site pour être visible dans différents pays et langues
\item[C)] Créer un site par pays
\item[D)] Traduire automatiquement un site
\end{itemize}
*(Réponse : B)*

\item À quoi servent les balises hreflang?
\begin{itemize}
\item[A)] À indiquer la langue d'une page uniquement
\item[B)] À indiquer à Google quelle version linguistique ou régionale afficher
\item[C)] À traduire automatiquement le contenu
\item[D)] À bloquer certaines langues
\end{itemize}
*(Réponse : B)*

\item Quelle est la syntaxe correcte pour le code de langue français?
\begin{itemize}
\item[A)] FR
\item[B)] fr
\item[C)] french
\item[D)] fra
\end{itemize}
*(Réponse : B)*

\item Que doit contenir la valeur hreflang "x-default"?
\begin{itemize}
\item[A)] La version française
\item[B)] La page par défaut (généralement en anglais)
\item[C)] Aucune page
\item[D)] La version chinoise
\end{itemize}
*(Réponse : B)*

\item Quelle est une erreur fréquente avec hreflang?
\begin{itemize}
\item[A)] Utiliser trop de langues
\item[B)] Absence de lien retour (réciprocité obligatoire)
\item[C)] Utiliser JSON-LD
\item[D)] Placer les balises dans le body
\end{itemize}
*(Réponse : B)*

\item Que signifie ccTLD?
\begin{itemize}
\item[A)] Domaines génériques (.com, .org)
\item[B)] Domaines spécifiques à un pays (.fr, .de, .es)
\item[C)] Sous-domaines dédiés
\item[D)] Domaines internationaux
\end{itemize}
*(Réponse : B)*

\item Quel est l'avantage des ccTLD pour le SEO?
\begin{itemize}
\item[A)] Coût réduit
\item[B)] Signal géographique le plus fort pour Google
\item[C)] Maintenance simplifiée
\item[D)] Autorité consolidée
\end{itemize}
*(Réponse : B)*

\item Quel est l'inconvénient des ccTLD?
\begin{itemize}
\item[A)] Pas de ciblage géographique
\item[B)] Coût élevé et autorité fragmentée
\item[C)] Performance réduite
\item[D)] Blocage par Google
\end{itemize}
*(Réponse : B)*

\item Quelle structure d'URL offre le meilleur équilibre autorité/maintenance/signal géographique?
\begin{itemize}
\item[A)] ccTLD (.fr, .de)
\item[B)] gTLD + sous-répertoire (exemple.com/fr/)
\item[C)] Sous-domaines (fr.exemple.com)
\item[D)] Paramètres d'URL
\end{itemize}
*(Réponse : B)*

\item Dans GSC, comment configurer le ciblage pour un gTLD?
\begin{itemize}
\item[A)] Le ciblage est automatique
\item[B)] Paramètres > Ciblage international > Pays
\item[C)] Ce n'est pas possible
\item[D)] Via le fichier robots.txt
\end{itemize}
*(Réponse : B)*

\item Le contenu dupliqué entre versions linguistiques est-il un problème?
\begin{itemize}
\item[A)] Non, Google comprend les langues
\item[B)] Oui, c'est un défi majeur du SEO international
\item[C)] Oui, mais Google le corrige automatiquement
\item[D)] Non, c'est bénéfique
\end{itemize}
*(Réponse : B)*

\item Comment éviter les problèmes de contenu dupliqué multilingue?
\begin{itemize}
\item[A)] Traduire automatiquement avec Google Translate
\item[B)] Utiliser des balises hreflang et un contenu adapté
\item[C)] Supprimer les versions linguistiques
\item[D)] Utiliser le même contenu pour toutes les langues
\end{itemize}
*(Réponse : B)*

\item Le ciblage géographique (geotargeting) peut être configuré dans :
\begin{itemize}
\item[A)] Le fichier robots.txt
\item[B)] Google Search Console
\item[C)] Le sitemap XML
\item[D)] Les méta descriptions
\end{itemize}
*(Réponse : B)*

\item La détection de la langue du navigateur utilise quel en-tête HTTP?
\begin{itemize}
\item[A)] User-Agent
\item[B)] Accept-Language
\item[C)] Cache-Control
\item[D)] Content-Type
\end{itemize}
*(Réponse : B)*

\item Pour les sites multilingues, la balise canonical doit être :
\begin{itemize}
\item[A)] Pointant vers la version anglaise
\item[B)] Auto-référençante (chaque version pointe vers elle-même)
\item[C)] Absente
\item[D)] Identique pour toutes les langues
\end{itemize}
*(Réponse : B)*

\item L'adaptation culturelle pour le SEO international inclut :
\begin{itemize}
\item[A)] Uniquement la traduction
\item[B)] Mots-clés locaux, monnaie, unités, formats de date, couleurs
\item[C)] Uniquement les images
\item[D)] Uniquement les couleurs
\end{itemize}
*(Réponse : B)*

\item Les balises hreflang peuvent être déclarées dans :
\begin{itemize}
\item[A)] Uniquement dans le <head>
\item[B)] Dans le <head> ou le sitemap XML
\item[C)] Uniquement dans le sitemap XML
\item[D)] Uniquement dans le body
\end{itemize}
*(Réponse : B)*

\item La réciprocité des balises hreflang signifie que :
\begin{itemize}
\item[A)] Chaque page doit pointer vers la version anglaise
\item[B)] Si la page A pointe vers B, B doit pointer vers A
\item[C)] Chaque page doit pointer vers elle-même
\item[D)] Les balises sont optionnelles
\end{itemize}
*(Réponse : B)*

\item Les pages avec hreflang ne doivent pas être :
\begin{itemize}
\item[A)] En HTTPS
\item[B)] En noindex
\item[C)] En HTTP
\item[D)] En cache
\end{itemize}
*(Réponse : B)*

\item Les codes de langue dans hreflang suivent la norme :
\begin{itemize}
\item[A)] ISO 3166-1
\item[B)] ISO 639-1
\item[C)] ISO 8601
\item[D)] ISO 8859-1
\end{itemize}
*(Réponse : B)*

\item Pour pages même langue mais pays différents, on utilise :
\begin{itemize}
\item[A)] Un seul hreflang générique
\item[B)] hreflang avec code régional (ex: fr-FR, fr-BE)
\item[C)] Des URLs identiques
\item[D)] Des redirections
\end{itemize}
*(Réponse : B)*

\item La structure en sous-domaines (fr.exemple.com) peut être traitée par Google comme :
\begin{itemize}
\item[A)] Un seul site unifié
\item[B)] Des sites séparés
\item[C)] Un site mobile
\item[D)] Un site dupliqué
\end{itemize}
*(Réponse : B)*

\item Pour un site multilingue, les cookies de préférence de langue permettent de :
\begin{itemize}
\item[A)] Bloquer certaines langues
\item[B)] Mémoriser le choix de langue de l'utilisateur
\item[C)] Traduire automatiquement
\item[D)] Rediriger vers Google
\end{itemize}
*(Réponse : B)*

\item La traduction automatique non revue peut être :
\begin{itemize}
\item[A)] Bénéfique pour le SEO
\item[B)] Pénalisée par Google comme contenu de faible qualité
\item[C)] Recommandée par Google
\item[D)] Sans impact SEO
\end{itemize}
*(Réponse : B)*

\item Pour un site multilingue avec sous-répertoires, le sitemap doit être :
\begin{itemize}
\item[A)] Un sitemap séparé par langue
\item[B)] Un seul sitemap avec toutes les variantes linguistiques
\item[C)] Pas de sitemap nécessaire
\item[D)] Un sitemap par pays
\end{itemize}
*(Réponse : B)*

\item La pagination doit être dans un contexte multilingue :
\begin{itemize}
\item[A)] Indépendante par langue
\item[B)] Parallèle et cohérente entre les langues
\item[C)] Différente pour chaque langue
\item[D)] Supprimée
\end{itemize}
*(Réponse : B)*

\item Un backlink depuis un site du pays cible est un signal de :
\begin{itemize}
\item[A)] Faible autorité
\item[B)] Pertinence géographique pour le SEO international
\item[C)] Contenu dupliqué
\item[D)] Pénalité potentielle
\end{itemize}
*(Réponse : B)*

\item L'hébergement du site dans le pays cible est-il un signal de géotargeting?
\begin{itemize}
\item[A)] Non, Google ne prend pas en compte la localisation du serveur
\item[B)] Oui, c'est un des signaux de géotargeting
\item[C)] Oui, c'est le seul signal
\item[D)] Non, c'est contre les règles
\end{itemize}
*(Réponse : B)*

\item Le RGPD est une considération importante pour le SEO international dans :
\begin{itemize}
\item[A)] Uniquement en France
\item[B)] L'Union Européenne
\item[C)] Uniquement aux US
\item[D)] En Asie
\end{itemize}
*(Réponse : B)*

\item Le format de date JJ/MM/AAAA est utilisé dans :
\begin{itemize}
\item[A)] Les États-Unis
\item[B)] La France et l'Europe
\item[C)] Le Japon
\item[D)] La Chine
\end{itemize}
*(Réponse : B)*

\item Les URLs descriptives dans chaque langue sont :
\begin{itemize}
\item[A)] Déconseillées
\item[B)] Recommandées pour le SEO international
\item[C)] Obligatoires
\item[D)] Sans impact
\end{itemize}
*(Réponse : B)*

\item La recherche de mots-clés pour chaque marché doit être :
\begin{itemize}
\item[A)] Une simple traduction
\item[B)] Réalisée par des natifs avec une vraie recherche locale
\item[C)] Ignorée
\item[D)] Automatisée
\end{itemize}
*(Réponse : B)*

\item Google Search Console permet de filtrer les performances par :
\begin{itemize}
\item[A)] Uniquement par langue
\item[B)] Par pays
\item[C)] Par navigateur
\item[D)] Par appareil seulement
\end{itemize}
*(Réponse : B)*

\item La meilleure pratique pour les images en contexte international est de :
\begin{itemize}
\item[A)] Utiliser les mêmes images partout
\item[B)] Adapter les images culturellement par marché
\item[C)] Supprimer les images
\item[D)] Utiliser uniquement du texte
\end{itemize}
*(Réponse : B)*

\item La monnaie et les prix doivent être affichés :
\begin{itemize}
\item[A)] Uniquement en dollars
\item[B)] Dans la devise locale de chaque marché
\item[C)] Uniquement en euros
\item[D)] Sans indication de devise
\end{itemize}
*(Réponse : B)*

\item Les balises hreflang sont-elles obligatoires pour les sites multilingues?
\begin{itemize}
\item[A)] Non, uniquement recommandées
\item[B)] Oui, fortement recommandées pour éviter les problèmes de contenu dupliqué
\item[C)] Non, Google les ignore
\item[D)] Oui, obligatoires
\end{itemize}
*(Réponse : B)*

\item Que doit contenir une page avec hreflang="en"?
\begin{itemize}
\item[A)] Du contenu en français
\item[B)] Du contenu en anglais
\item[C)] Du contenu traduit automatiquement
\item[D)] N'importe quel contenu
\end{itemize}
*(Réponse : B)*

\item La redirection basée sur Accept-Language doit toujours proposer :
\begin{itemize}
\item[A)] Aucun choix à l'utilisateur
\item[B)] Une option de changement de langue manuel
\item[C)] Uniquement la langue détectée
\item[D)] Une redirection vers Google
\end{itemize}
*(Réponse : B)*

\item Les en-têtes HTTP Content-Language sont-ils une alternative aux balises hreflang?
\begin{itemize}
\item[A)] Oui, équivalent
\item[B)] Non, ce sont des signaux supplémentaires mais pas équivalents
\item[C)] Oui, meilleurs que hreflang
\item[D)] Non, ignorés par Google
\end{itemize}
*(Réponse : B)*

\item Dans le sitemap XML multilingue, on utilise l'espace de noms :
\begin{itemize}
\item[A)] xmlns:lang
\item[B)] xmlns:xhtml
\item[C)] xmlns:alternate
\item[D)] xmlns:hreflang
\end{itemize}
*(Réponse : B)*

\item La structure en sous-répertoires pour le multilingue permet de :
\begin{itemize}
\item[A)] Fragmenter l'autorité
\item[B)] Consolider l'autorité sur le domaine principal
\item[C)] Réduire la maintenance
\item[D)] Augmenter les coûts
\end{itemize}
*(Réponse : B)*

\item Les backlinks locaux sont importants car ils renforcent :
\begin{itemize}
\item[A)] La vitesse du site
\item[B)] La pertinence géographique locale
\item[C)] Le design
\item[D)] La sécurité
\end{itemize}
*(Réponse : B)*

\item Une erreur hreflang courante est :
\begin{itemize}
\item[A)] Utiliser des codes de langue valides
\item[B)] Avoir des valeurs hreflang qui ne correspondent pas au contenu réel
\item[C)] Déclarer toutes les variantes linguistiques
\item[D)] Utiliser x-default
\end{itemize}
*(Réponse : B)*

\item Le contenu adapté culturellement signifie :
\begin{itemize}
\item[A)] Traduire mot à mot
\item[B)] Aller au-delà de la traduction, adapter exemples et références
\item[C)] Utiliser les mêmes exemples partout
\item[D)] Copier le contenu existant
\end{itemize}
*(Réponse : B)*

\item Un site multilingue bien optimisé doit avoir quel contenu spécifique par marché?
\begin{itemize}
\item[A)] 0\% (100\% traduit)
\item[B)] Un pourcentage de contenu original spécifique à chaque marché
\item[C)] 100\% identique
\item[D)] Uniquement les pages produits traduites
\end{itemize}
*(Réponse : B)*

\item Les codes de pays dans hreflang suivent la norme :
\begin{itemize}
\item[A)] ISO 639-1
\item[B)] ISO 3166-1 Alpha 2
\item[C)] ISO 8601
\item[D)] ISO 8859-15
\end{itemize}
*(Réponse : B)*

\item Le x-default dans hreflang sert à cibler :
\begin{itemize}
\item[A)] Les utilisateurs français
\item[B)] Les utilisateurs sans langue spécifique ou langue non couverte
\item[C)] Les utilisateurs américains
\item[D)] Les utilisateurs allemands
\end{itemize}
*(Réponse : B)*

\item Les ccTLD donnent un signal géographique :
\begin{itemize}
\item[A)] Faible
\item[B)] Très fort
\item[C)] Nul
\item[D)] Négatif
\end{itemize}
*(Réponse : B)*

\item Google Search Console nécessite-t-il une configuration par pays pour les ccTLD?
\begin{itemize}
\item[A)] Oui, obligatoire
\item[B)] Non, le ciblage est automatique
\item[C)] Oui, si le site utilise aussi d'autres langues
\item[D)] Non, les ccTLD ne sont pas supportés
\end{itemize}
*(Réponse : B)*

\item L'ordre de déploiement des marchés dans une stratégie SEO internationale devrait prioriser :
\begin{itemize}
\item[A)] Les marchés les plus difficiles
\item[B)] Les marchés les plus accessibles avec le meilleur potentiel
\item[C)] Les marchés aléatoirement
\item[D)] Tous les marchés en même temps
\end{itemize}
*(Réponse : B)*
\end{enumerate}

% ============================================
% CHAPITRE 19 : Le SEO Sémantique et Entity SEO
% ============================================
\section*{Chapitre 19 : Le SEO Sémantique et Entity SEO}

\begin{enumerate}
\item Qu'est-ce qu'une entité dans le contexte SEO?
\begin{itemize}
\item[A)] Un mot-clé
\item[B)] Un concept ou objet unique identifiable (personne, lieu, organisation, concept)
\item[C)] Un fichier HTML
\item[D)] Un backlink
\end{itemize}
*(Réponse : B)*

\item Qu'est-ce que le Knowledge Graph de Google?
\begin{itemize}
\item[A)] Un outil de graphisme
\item[B)] Une base de connaissances qui comprend les relations entre les entités
\item[C)] Un logiciel de création de graphiques
\item[D)] Un réseau social
\end{itemize}
*(Réponse : B)*

\item Le SEO sémantique consiste à optimiser pour :
\begin{itemize}
\item[A)] Des mots-clés spécifiques uniquement
\item[B)] Des concepts et des thématiques plutôt que des mots-clés isolés
\item[C)] Des balises HTML
\item[D)] Des backlinks uniquement
\end{itemize}
*(Réponse : B)*

\item Qu'est-ce que la Topical Authority?
\begin{itemize}
\item[A)] L'autorité d'un domaine sur tous les sujets
\item[B)] La reconnaissance par Google comme référence complète sur un sujet donné
\item[C)] L'autorité d'un backlink
\item[D)] La popularité sur les réseaux sociaux
\end{itemize}
*(Réponse : B)*

\item Comment démontrer une expertise complète selon le SEO sémantique?
\begin{itemize}
\item[A)] En utilisant beaucoup de mots-clés
\item[B)] En créant un contenu exhaustif qui couvre le sujet dans toutes ses dimensions
\item[C)] En achetant des backlinks
\item[D)] En copiant le contenu des concurrents
\end{itemize}
*(Réponse : B)*

\item Google utilise le Knowledge Graph pour :
\begin{itemize}
\item[A)] Compresser les images
\item[B)] Comprendre les relations entre les entités et fournir des résultats pertinents
\item[C)] Calculer la vitesse du site
\item[D)] Générer des backlinks
\end{itemize}
*(Réponse : B)*

\item Comment implémenter le SEO sémantique?
\begin{itemize}
\item[A)] Ajouter plus de mots-clés
\item[B)] Identifier l'entité principale, utiliser les données structurées, contenu exhaustif
\item[C)] Acheter des liens
\item[D)] Créer des pages dupliquées
\end{itemize}
*(Réponse : B)*

\item Le champ sémantique d'un sujet comprend :
\begin{itemize}
\item[A)] Un seul mot-clé
\item[B)] Un vocabulaire riche et varié autour du sujet principal
\item[C)] Uniquement des synonymes
\item[D)] Des backlinks
\end{itemize}
*(Réponse : B)*

\item L'analyse des "People Also Ask" permet d'identifier :
\begin{itemize}
\item[A)] Les concurrents
\item[B)] Les sous-sujets et questions connexes à couvrir
\item[C)] Les mots-clés les plus difficiles
\item[D)] Les backlinks
\end{itemize}
*(Réponse : B)*

\item Le SEO moderne a évolué d'une approche centrée sur les mots-clés vers :
\begin{itemize}
\item[A)] Les backlinks
\item[B)] Les entités (Entity SEO)
\item[C)] La vitesse
\item[D)] Le design
\end{itemize}
*(Réponse : B)*

\item Le maillage interne dans une stratégie de SEO sémantique doit :
\begin{itemize}
\item[A)] Éviter les liens entre pages connexes
\item[B)] Connecter les pages traitant d'entités connexes
\item[C)] Pointer uniquement vers la page d'accueil
\item[D)] Utiliser des liens nofollow
\end{itemize}
*(Réponse : B)*

\item Le Knowledge Graph de Google contient :
\begin{itemize}
\item[A)] Uniquement des pages web
\item[B)] Des milliards d'entités interconnectées
\item[C)] Uniquement des images
\item[D)] Uniquement des vidéos
\end{itemize}
*(Réponse : B)*

\item Les Topical Clusters sont un concept clé du :
\begin{itemize}
\item[A)] SEO technique
\item[B)] SEO sémantique
\item[C)] Link building
\item[D)] SEO local
\end{itemize}
*(Réponse : B)*

\item Le modèle Hub-and-Spoke est également connu sous le nom de :
\begin{itemize}
\item[A)] Site architecture
\item[B)] Topic Cluster
\item[C)] Sitemap
\item[D)] Funnel
\end{itemize}
*(Réponse : B)*

\item Google BERT améliore la compréhension de quel aspect du langage?
\begin{itemize}
\item[A)] La grammaire
\item[B)] Le contexte des mots dans une phrase (nuances, prépositions)
\item[C)] L'orthographe
\item[D)] La ponctuation
\end{itemize}
*(Réponse : B)*

\item Google MUM est capable de comprendre du contenu dans combien de langues?
\begin{itemize}
\item[A)] 25
\item[B)] 75
\item[C)] 100
\item[D)] 50
\end{itemize}
*(Réponse : B)*

\item MUM est considéré comme combien de fois plus puissant que BERT?
\begin{itemize}
\item[A)] 100 fois
\item[B)] 1000 fois
\item[C)] 10 fois
\item[D)] 10000 fois
\end{itemize}
*(Réponse : B)*

\item Avec BERT, l'écriture du contenu doit être :
\begin{itemize}
\item[A)] Remplie de mots-clés
\item[B)] Naturelle, en langage courant
\item[C)] Technique uniquement
\item[D)] Très courte
\end{itemize}
*(Réponse : B)*

\item Une Pillar Page doit généralement faire :
\begin{itemize}
\item[A)] 500-1000 mots
\item[B)] 3000-7000 mots
\item[C)] 100-200 mots
\item[D)] Plus de 10000 mots
\end{itemize}
*(Réponse : B)*

\item Les recherches quotidiennes jamais vues auparavant représentent environ :
\begin{itemize}
\item[A)] 5\%
\item[B)] 15\%
\item[C)] 25\%
\item[D)] 50\%
\end{itemize}
*(Réponse : B)*

\item MUM peut traiter simultanément :
\begin{itemize}
\item[A)] Uniquement du texte
\item[B)] Texte, images et vidéos
\item[C)] Uniquement des images
\item[D)] Uniquement des vidéos
\end{itemize}
*(Réponse : B)*

\item Les mots-clés de longue traîne sont caractérisés par :
\begin{itemize}
\item[A)] Un volume de recherche élevé
\item[B)] Une intention précise et un volume faible
\item[C)] Une concurrence forte
\item[D)] Un CPC élevé
\end{itemize}
*(Réponse : B)*

\item Le SEO sémantique améliore la pertinence :
\begin{itemize}
\item[A)] Technique
\item[B)] Thématique
\item[C)] Des backlinks
\item[D)] Des images
\end{itemize}
*(Réponse : B)*

\item Google RankBrain a été introduit en :
\begin{itemize}
\item[A)] 2013
\item[B)] 2015
\item[C)] 2017
\item[D)] 2019
\end{itemize}
*(Réponse : B)*

\item Google BERT a été introduit en :
\begin{itemize}
\item[A)] 2015
\item[B)] 2019
\item[C)] 2021
\item[D)] 2023
\end{itemize}
*(Réponse : B)*

\item Google MUM a été introduit en :
\begin{itemize}
\item[A)] 2019
\item[B)] 2021
\item[C)] 2022
\item[D)] 2024
\end{itemize}
*(Réponse : B)*

\item Les entités secondaires dans le SEO sémantique sont :
\begin{itemize}
\item[A)] Des mots-clés secondaires
\item[B)] Des concepts liés à l'entité principale
\item[C)] Des backlinks
\item[D)] Des balises HTML
\end{itemize}
*(Réponse : B)*

\item Le schéma Article avec auteur spécifié aide à :
\begin{itemize}
\item[A)] Améliorer la vitesse
\item[B)] Établir l'entité Person de l'auteur dans le Knowledge Graph
\item[C)] Augmenter le nombre de mots
\item[D)] Créer des backlinks
\end{itemize}
*(Réponse : B)*

\item L'optimisation du champ sémantique consiste à :
\begin{itemize}
\item[A)] Utiliser uniquement des mots-clés exacts
\item[B)] Employer les termes naturellement associés au sujet principal
\item[C)] Supprimer les synonymes
\item[D)] Réduire le vocabulaire
\end{itemize}
*(Réponse : B)*

\item Le Knowledge Panel de Google est alimenté par :
\begin{itemize}
\item[A)] Le contenu du site
\item[B)] Le Knowledge Graph
\item[C)] Les backlinks
\item[D)] Les méta descriptions
\end{itemize}
*(Réponse : B)*

\item Le schema SameAs dans les données structurées relie une entité à :
\begin{itemize}
\item[A)] Des backlinks
\item[B)] Ses profils sur d'autres plateformes (Wikipedia, réseaux sociaux)
\item[C)] Des images
\item[D)] Des vidéos
\end{itemize}
*(Réponse : B)*

\item RankBrain est particulièrement efficace pour comprendre :
\begin{itemize}
\item[A)] Les mots-clés courts
\item[B)] Les requêtes jamais vues auparavant
\item[C)] Les backlinks
\item[D)] Les images
\end{itemize}
*(Réponse : B)*

\item BERT lit le texte :
\begin{itemize}
\item[A)] De gauche à droite uniquement
\item[B)] Dans les deux directions simultanément
\item[C)] De droite à gauche uniquement
\item[D)] De manière aléatoire
\end{itemize}
*(Réponse : B)*

\item La Topical Authority se construit principalement par :
\begin{itemize}
\item[A)] L'achat de backlinks
\item[B)] La publication régulière de contenu exhaustif et interconnecté
\item[C)] Les réseaux sociaux
\item[D)] La publicité payante
\end{itemize}
*(Réponse : B)*

\item L'approche Entity SEO est centrée sur :
\begin{itemize}
\item[A)] Les mots-clés exacts
\item[B)] Les entités et leurs relations
\item[C)] Les backlinks
\item[D)] La technique pure
\end{itemize}
*(Réponse : B)*

\item Pour le SEO sémantique, l'analyse des SERP permet de comprendre :
\begin{itemize}
\item[A)] Le prix des mots-clés
\item[B)] Le type de contenu attendu par Google pour un sujet donné
\item[C)] Les backlinks des concurrents
\item[D)] La vitesse des concurrents
\end{itemize}
*(Réponse : B)*

\item Une entité dans le Knowledge Graph peut être :
\begin{itemize}
\item[A)] Uniquement une personne
\item[B)] Une personne, un lieu, une organisation ou un concept
\item[C)] Uniquement un lieu
\item[D)] Uniquement un mot
\end{itemize}
*(Réponse : B)*

\item L'optimisation pour RankBrain nécessite de se concentrer sur :
\begin{itemize}
\item[A)] Les mots-clés exacts
\item[B)] L'intention de recherche
\item[C)] Les balises title
\item[D)] Les méta descriptions
\end{itemize}
*(Réponse : B)*

\item MUM encourage la création de contenu qui :
\begin{itemize}
\item[A)] Est très court
\item[B)] Est expert et complet, couvrant le sujet en profondeur
\item[C)] Contient beaucoup de mots-clés
\item[D)] Est dupliqué
\end{itemize}
*(Réponse : B)*

\item Le TF-IDF est une métrique qui mesure :
\begin{itemize}
\item[A)] La difficulté d'un mot-clé
\item[B)] L'importance d'un terme dans un document par rapport à un corpus
\item[C)] Le volume de recherche
\item[D)] Le nombre de backlinks
\end{itemize}
*(Réponse : B)*

\item Les termes LSI sont :
\begin{itemize}
\item[A)] Des mots-clés principaux
\item[B)] Des termes associés naturellement présents dans un contenu expert
\item[C)] Des balises HTML
\item[D)] Des données structurées
\end{itemize}
*(Réponse : B)*

\item La couverture thématique complète signifie :
\begin{itemize}
\item[A)] Écrire un seul article très long
\item[B)] Traiter un sujet dans toutes ses dimensions avec plusieurs contenus interconnectés
\item[C)] Acheter des backlinks
\item[D)] Utiliser beaucoup de mots-clés
\end{itemize}
*(Réponse : B)*

\item Les Topical Clusters renforcent :
\begin{itemize}
\item[A)] La vitesse du site
\item[B)] L'expertise thématique aux yeux de Google
\item[C)] Le nombre de pages
\item[D)] Le design
\end{itemize}
*(Réponse : B)*

\item Un des principes du SEO sémantique est :
\begin{itemize}
\item[A)] Répéter le même mot-clé
\item[B)] Utiliser un vocabulaire riche et varié autour du sujet
\item[C)] Éviter les synonymes
\item[D)] Acheter des liens
\end{itemize}
*(Réponse : B)*

\item Les "Related Searches" dans Google aident à :
\begin{itemize}
\item[A)] Trouver des images
\item[B)] Identifier des sujets connexes à couvrir
\item[C)] Analyser la vitesse
\item[D)] Trouver des backlinks
\end{itemize}
*(Réponse : B)*

\item Le Schema.org Article avec author et datePublished aide Google à :
\begin{itemize}
\item[A)] Calculer la vitesse
\item[B)] Comprendre l'entité et la fraîcheur du contenu
\item[C)] Créer des backlinks
\item[D)] Optimiser les images
\end{itemize}
*(Réponse : B)*

\item L'Entity SEO permet de mieux se positionner pour :
\begin{itemize}
\item[A)] Un seul mot-clé
\item[B)] L'ensemble d'un champ sémantique
\item[C)] Les mots-clés de marque uniquement
\item[D)] Les requêtes locales
\end{itemize}
*(Réponse : B)*

\item Le Knowledge Graph de Google a été introduit en :
\begin{itemize}
\item[A)] 2008
\item[B)] 2012
\item[C)] 2015
\item[D)] 2020
\end{itemize}
*(Réponse : B)*

\item L'analyse des "Related Searches" dans Google aide à :
\begin{itemize}
\item[A)] Trouver des images
\item[B)] Découvrir des sujets connexes à couvrir pour enrichir le champ sémantique
\item[C)] Analyser la vitesse du site
\item[D)] Trouver des backlinks
\end{itemize}
*(Réponse : B)*

\item Le schéma "SameAs" dans les données structurées est utilisé pour :
\begin{itemize}
\item[A)] Créer des liens internes
\item[B)] Relier une entité à ses profils sur d'autres plateformes
\item[C)] Définir des auteurs
\item[D)] Indiquer la langue
\end{itemize}
*(Réponse : B)*
\end{enumerate}
% ============================================
% CHAPITRE 20 : Le Netlinking et les Backlinks
% ============================================
\section*{Chapitre 20 : Le Netlinking et les Backlinks}

\begin{enumerate}
\item Que sont les backlinks en SEO?
\begin{itemize}
\item[A)] Des liens internes entre pages d'un même site
\item[B)] Des liens externes pointant d'un site A vers un site B
\item[C)] Des liens morts
\item[D)] Des liens vers des images
\end{itemize}
*(Réponse : B)*

\item Les backlinks sont considérés par Google comme :
\begin{itemize}
\item[A)] Un facteur mineur de classement
\item[B)] Un des facteurs de classement les plus importants
\item[C)] Sans importance
\item[D)] Uniquement pour les sites e-commerce
\end{itemize}
*(Réponse : B)*

\item Un lien dofollow transmet :
\begin{itemize}
\item[A)] Aucune autorité
\item[B)] L'autorité (PageRank)
\item[C)] Uniquement du trafic
\item[D)] Des métadonnées
\end{itemize}
*(Réponse : B)*

\item L'attribut rel="nofollow" sur un lien indique à Google de :
\begin{itemize}
\item[A)] Suivre et indexer le lien
\item[B)] Ignorer le lien pour le PageRank
\item[C)] Accélérer le crawling
\item[D)] Bloquer la page cible
\end{itemize}
*(Réponse : B)*

\item L'attribut rel="ugc" est destiné aux liens provenant de :
\begin{itemize}
\item[A)] Sites publicitaires
\item[B)] Contenus générés par les utilisateurs (commentaires, forums)
\item[C)] Sites gouvernementaux
\item[D)] Réseaux sociaux
\end{itemize}
*(Réponse : B)*

\item L'attribut rel="sponsored" est destiné aux liens :
\begin{itemize}
\item[A)] Naturels
\item[B)] Publicitaires ou sponsorisés
\item[C)] Internes
\item[D)] Vers des images
\end{itemize}
*(Réponse : B)*

\item Parmi les critères de qualité d'un backlink, lequel est important?
\begin{itemize}
\item[A)] La couleur du lien
\item[B)] La pertinence thématique entre les deux sites
\item[C)] La longueur de l'URL
\item[D)] Le nombre de caractères du lien
\end{itemize}
*(Réponse : B)*

\item Où un backlink a-t-il le plus de valeur?
\begin{itemize}
\item[A)] Dans le footer
\item[B)] Dans le contenu principal
\item[C)] Dans la sidebar
\item[D)] Dans les commentaires
\end{itemize}
*(Réponse : B)*

\item L'anchor text (texte du lien) doit idéalement être :
\begin{itemize}
\item[A)] Identique pour tous les liens
\item[B)] Naturel et descriptif
\item[C)] Uniquement le mot-clé exact
\item[D)] Très long
\end{itemize}
*(Réponse : B)*

\item Un profil de backlinks naturel inclut :
\begin{itemize}
\item[A)] Uniquement des dofollow
\item[B)] Un mélange de dofollow et nofollow
\item[C)] Uniquement des nofollow
\item[D)] Aucun lien
\end{itemize}
*(Réponse : B)*

\item La qualité des backlinks est plus importante que :
\begin{itemize}
\item[A)] La pertinence
\item[B)] La quantité
\item[C)] L'emplacement
\item[D)] L'ancrage
\end{itemize}
*(Réponse : B)*

\item Un backlink depuis un site avec une forte autorité de domaine a :
\begin{itemize}
\item[A)] Peu de valeur
\item[B)] Une valeur élevée
\item[C)] Aucune valeur
\item[D)] Un impact négatif
\end{itemize}
*(Réponse : B)*

\item Quel type de lien est le lien par défaut (sans attribut)?
\begin{itemize}
\item[A)] Nofollow
\item[B)] Dofollow
\item[C)] Sponsored
\item[D)] UGC
\end{itemize}
*(Réponse : B)*

\item Les liens provenant de fermes de liens sont considérés comme :
\begin{itemize}
\item[A)] Bénéfiques
\item[B)] Potentiellement dangereux et contraires aux guidelines Google
\item[C)] Neutres
\item[D)] Recommandés
\end{itemize}
*(Réponse : B)*

\item Un lien dans un article invité de qualité peut être bénéfique car :
\begin{itemize}
\item[A)] Il est toujours nofollow
\item[B)] Il provient d'un contexte éditorial pertinent
\item[C)] Il est payant
\item[D)] Il est automatique
\end{itemize}
*(Réponse : B)*

\item Google peut pénaliser un site pour :
\begin{itemize}
\item[A)] Avoir trop de backlinks naturels
\item[B]) Avoir des backlinks artificiels ou manipulés
\item[C)] Avoir des backlinks de qualité
\item[D)] Ne pas avoir assez de nofollow
\end{itemize}
*(Réponse : B)*

\item Le taux de croissance des backlinks doit idéalement être :
\begin{itemize}
\item[A)] Très rapide
\item[B)] Progressif et naturel
\item[C)] Nul
\item[D)] Exponentiel
\end{itemize}
*(Réponse : B)*

\item Une diversité des sources de backlinks est :
\begin{itemize}
\item[A)] Négative pour le SEO
\item[B)] Positive et recommandée
\item[C)] Sans importance
\item[D]) Uniquement pour les sites locaux
\end{itemize}
*(Réponse : B)*

\item Les backlinks provenant de sites .edu et .gov sont généralement considérés comme :
\begin{itemize}
\item[A)] De faible qualité
\item[B)] De haute qualité et précieux
\item[C)] Interdits par Google
\item[D)] Identiques aux autres
\end{itemize}
*(Réponse : B)*

\item Le PageRank est un algorithme qui mesure :
\begin{itemize}
\item[A)] La vitesse de la page
\item[B)] L'importance d'une page via les liens qui pointent vers elle
\item[C)] La qualité du contenu
\item[D)] La sécurité du site
\end{itemize}
*(Réponse : B)*

\item Un lien nofollow peut-il quand même apporter de la valeur?
\begin{itemize}
\item[A)] Non, jamais
\item[B)] Oui, il peut générer du trafic et de la visibilité
\item[C)] Oui, autant qu'un dofollow
\item[D)] Non, il est ignoré totalement
\end{itemize}
*(Réponse : B)*

\item La désaveu (disavow) de backlinks est une fonctionnalité de :
\begin{itemize}
\item[A)] Google Analytics
\item[B)] Google Search Console
\item[C)] Google Ads
\item[D)] Google Tag Manager
\end{itemize}
*(Réponse : B)*

\item Il est recommandé d'utiliser l'outil de désaveu lorsqu'on a :
\begin{itemize}
\item[A)] Trop de backlinks de qualité
\item[B)] Des backlinks toxiques ou artificiels qu'on ne peut pas supprimer
\item[C)] Pas assez de backlinks
\item[D)] Uniquement des nofollow
\end{itemize}
*(Réponse : B)*

\item La pertinence thématique entre sites liés est un signal de qualité car elle indique :
\begin{itemize}
\item[A)] Un échange de liens automatique
\item[B)] Un lien naturel et contextuel
\item[C)] Un achat de lien
\item[D)] Un spam
\end{itemize}
*(Réponse : B)*

\item Un backlink récent a généralement :
\begin{itemize}
\item[A)] Moins de poids qu'un ancien
\item[B)] Plus de poids qu'un ancien
\item[C)] Le même poids
\item[D)] Aucun poids
\end{itemize}
*(Réponse : B)*

\item L'optimisation excessive de l'anchor text avec des mots-clés exacts peut être :
\begin{itemize}
\item[A)] Toujours bénéfique
\item[B)] Risquée (pénalité Google Penguin)
\item[C)] Sans effet
\item[D)] Obligatoire
\end{itemize}
*(Réponse : B)*

\item Google Penguin cible spécifiquement :
\begin{itemize}
\item[A)] Le contenu de faible qualité
\item[B)] Les profils de backlinks artificiels et le keyword stuffing
\item[C)] La lenteur du site
\item[D)] Les problèmes mobiles
\end{itemize}
*(Réponse : B)*

\item Quel est le ratio idéal entre liens dofollow et nofollow?
\begin{itemize}
\item[A)] 100\% dofollow
\item[B)] Un ratio naturel, variable selon le secteur
\item[C)] 50/50 exactement
\item[D)] 100\% nofollow
\end{itemize}
*(Réponse : B)*

\item Un backlink depuis un site qui n'a aucun rapport avec votre secteur est :
\begin{itemize}
\item[A)] Très précieux
\item[B)] Généralement moins pertinent
\item[C)] Aussi précieux qu'un lien pertinent
\item[D)] Toujours toxique
\end{itemize}
*(Réponse : B)*

\item La provenance géographique des backlinks est importante pour :
\begin{itemize}
\item[A)] Le SEO technique
\item[B)] Le SEO local
\item[C)] La vitesse du site
\item[D)] Les données structurées
\end{itemize}
*(Réponse : B)*

\item Le nombre de domaines référents uniques est une métrique qui mesure :
\begin{itemize}
\item[A)] Le nombre total de backlinks
\item[B)] Le nombre de sites différents qui linkent vers vous
\item[C)] Le nombre de pages indexées
\item[D)] Le trafic mensuel
\end{itemize}
*(Réponse : B)*

\item Un pic soudain de nouveaux backlinks peut être interprété par Google comme :
\begin{itemize}
\item[A)] Un signe de popularité naturelle
\item[B)] Une possible tentative de manipulation
\item[C)] Un changement d'algorithme
\item[D)] Une mise à jour de contenu
\end{itemize}
*(Réponse : B)*

\item Les backlinks dans les commentaires de blog avec rel="nofollow" sont généralement :
\begin{itemize}
\item[A)] Très puissants pour le SEO
\item[B)] Faibles mais peuvent apporter du trafic
\item[C)] Interdits
\item[D)] Aussi efficaces que les liens éditoriaux
\end{itemize}
*(Réponse : B)*

\item L'autorité du domaine (Domain Authority) est une métrique créée par :
\begin{itemize}
\item[A)] Google
\item[B)] Moz
\item[C)] Yahoo
\item[D)] Bing
\end{itemize}
*(Réponse : B)*

\item Une page avec beaucoup de backlinks de qualité sera considérée par Google comme :
\begin{itemize}
\item[A)] Suspecte
\item[B)] Une page de référence et d'autorité
\item[C)] Dupliquée
\item[D)] Sans intérêt
\end{itemize}
*(Réponse : B)*

\item Les backlinks sont un signal off-page qui indique à Google :
\begin{itemize}
\item[A)] La vitesse du site
\item[B)] La popularité et la confiance accordée par d'autres sites
\item[C)] La qualité du code HTML
\item[D)] La compatibilité mobile
\end{itemize}
*(Réponse : B)*

\item Les annuaires de faible qualité sont déconseillés car ils sont considérés comme :
\begin{itemize}
\item[A)] Des sources de trafic
\item[B)] Des schémas de liens non naturels
\item[C)] Des sites utiles
\item[D)] Des partenaires recommandés
\end{itemize}
*(Réponse : B)*

\item Les échanges de liens réciproques excessifs (link exchange) sont considérés par Google comme :
\begin{itemize}
\item[A)] Une bonne pratique
\item[B)] Une manipulation potentielle si elle est abusive
\item[C)] Sans effet
\item[D)] Recommandés
\end{itemize}
*(Réponse : B)*

\item Un lien dofollow provenant d'un site d'actualités reconnu est :
\begin{itemize}
\item[A)] Sans valeur
\item[B)] Très précieux (autorité + trafic)
\item[C)] Toujours nofollow
\item[D)] Interdit
\end{itemize}
*(Réponse : B)*

\item Les liens internes sont différents des backlinks car ils relient :
\begin{itemize}
\item[A)] Deux sites différents
\item[B)] Deux pages du même site
\item[C)] Un site à un moteur de recherche
\item[D)] Un site à un réseau social
\end{itemize}
*(Réponse : B)*

\item Un backlink positionné en haut du contenu est généralement :
\begin{itemize}
\item[A)] Moins efficace
\item[B)] Plus efficace qu'en bas de page
\item[C)] Identique
\item[D)] Ignoré par Google
\end{itemize}
*(Réponse : B)*

\item Pour vérifier son profil de backlinks, on peut utiliser :
\begin{itemize}
\item[A)] Google Analytics uniquement
\item[B)] Google Search Console, Ahrefs, Majestic, SEMrush
\item[C)] Uniquement Google Ads
\item[D)] Le fichier robots.txt
\end{itemize}
*(Réponse : B)*

\item La perte soudaine de backlinks peut indiquer :
\begin{itemize}
\item[A)] Une amélioration du SEO
\item[B)] La suppression de contenu ou la fermeture de sites partenaires
\item[C)] Une augmentation de trafic
\item[D)] Un changement d'algorithme positif
\end{itemize}
*(Réponse : B)*

\item Les backlinks sont mesurés en termes de "jus" SEO qui correspond au :
\begin{itemize}
\item[A)] Trafic direct
\item[B)] PageRank transmis
\item[C)] Nombre de mots
\item[D)] Taux de conversion
\end{itemize}
*(Réponse : B)*

\item Les réseaux de blogs privés (PBN) sont considérés par Google comme :
\begin{itemize}
\item[A)] Une excellente stratégie
\item[B)] Une violation des guidelines (risque de pénalité)
\item[C)] Une pratique neutre
\item[D]) Un outil recommandé
\end{itemize}
*(Réponse : B)*

\item La création de liens via des articles invités est efficace si :
\begin{itemize}
\item[A)] On utilise le même article partout
\item[B)] On cible des sites pertinents avec un contenu original de valeur
\item[C)] On paie pour chaque publication
\item[D)] On utilise des fermes de contenu
\end{itemize}
*(Réponse : B)*

\item Le Trust Flow est une métrique développée par :
\begin{itemize}
\item[A)] Google
\item[B)] Majestic
\item[C)] Moz
\item[D)] SEMrush
\end{itemize}
*(Réponse : B)*

\item Le Citation Flow mesure :
\begin{itemize}
\item[A)] La qualité des backlinks
\item[B)] La quantité de backlinks
\item[C)] La vitesse du site
\item[D)] Le trafic organique
\end{itemize}
*(Réponse : B)*

\item Le ratio Citation Flow / Trust Flow permet d'identifier :
\begin{itemize}
\item[A)] Les sites rapides
\item[B)] Les profils de liens potentiellement artificiels
\item[C)] Les pages bien structurées
\item[D)] Les sites sécurisés
\end{itemize}
*(Réponse : B)*

\item Un profil de backlinks sain doit idéalement montrer :
\begin{itemize}
\item[A)] Un seul type de source
\item[B)] Une diversité de domaines, d'ancres et de types de liens
\item[C)] Uniquement des ancres exactes
\item[D)] Uniquement des liens depuis la même thématique
\end{itemize}
*(Réponse : B)*

\end{enumerate}

% ============================================
% CHAPITRE 21 : Les stratégies avancées de Link Building
% ============================================
\section*{Chapitre 21 : Les stratégies avancées de Link Building}

\begin{enumerate}
\item Qu'est-ce que le Digital PR dans le contexte du link building?
\begin{itemize}
\item[A)] La publicité payante sur Google
\item[B)] Les relations presse numériques pour obtenir des mentions et backlinks naturels
\item[C)] Le design de sites web
\item[D)] La gestion de réseaux sociaux
\end{itemize}
*(Réponse : B)*

\item Qu'est-ce qu'un "Linkable Asset"?
\begin{itemize}
\item[A)] Un outil payant de création de liens
\item[B)] Une ressource de valeur qui attire naturellement des backlinks
\item[C)] Un plugin WordPress
\item[D)] Un type de serveur
\end{itemize}
*(Réponse : B)*

\item Les études originales sont efficaces pour le link building car elles :
\begin{itemize}
\item[A)] Sont faciles à copier
\item[B)] Fournissent des données uniques que d'autres voudront citer
\item[C)] Sont gratuites
\item[D)] Sont courtes
\end{itemize}
*(Réponse : B)*

\item Les infographies attirent des backlinks car :
\begin{itemize}
\item[A)] Elles sont faciles à créer
\item[B)] Elles permettent de visualiser des données complexes de façon engageante
\item[C)] Elles sont obligatoires
\item[D)] Elles sont longues
\end{itemize}
*(Réponse : B)*

\item Le guest blogging (publication invitée) reste efficace si :
\begin{itemize}
\item[A)] On écrit sur n'importe quel site
\item[B)] On apporte une valeur réelle sur un site pertinent avec une vraie audience
\item[C)] On utilise toujours le même contenu
\item[D)] On paie pour être publié
\end{itemize}
*(Réponse : B)*

\item Le broken link building consiste à :
\begin{itemize}
\item[A)] Créer des liens cassés volontairement
\item[B)] Identifier des liens brisés et proposer son contenu comme remplacement
\item[C)] Supprimer des backlinks
\item[D)] Ignorer les erreurs 404
\end{itemize}
*(Réponse : B)*

\item Quelle est la première étape du broken link building?
\begin{itemize}
\item[A)] Créer du contenu
\item[B)] Identifier des pages pertinentes avec des liens brisés (404)
\item[C)] Contacter les webmasters
\item[D)] Analyser les concurrents
\end{itemize}
*(Réponse : B)*

\item La Skyscraper Technique popularisée par Brian Dean consiste à :
\begin{itemize}
\item[A)] Copier le contenu des concurrents
\item[B)] Trouver un contenu performant, créer une version améliorée, et contacter les sites qui le linkent
\item[C)] Acheter des backlinks
\item[D)] Créer des pages très courtes
\end{itemize}
*(Réponse : B)*

\item La Skyscraper Technique peut générer combien de fois plus de backlinks que le contenu original?
\begin{itemize}
\item[A)] 1 à 2 fois
\item[B)] 2 à 3 fois
\item[C)] 5 à 10 fois
\item[D)] 10 à 20 fois
\end{itemize}
*(Réponse : B)*

\item Les brand mentions sans lien sont désormais reconnues par Google comme :
\begin{itemize}
\item[A)] Sans importance
\item[B)] Des signaux d'autorité
\item[C)] Du spam
\item[D)] Des erreurs
\end{itemize}
*(Réponse : B)*

\item La Topical Authority signifie :
\begin{itemize}
\item[A)] Avoir des backlinks de tous les sujets
\item[B)] Être reconnu comme une référence dans un domaine spécifique
\item[C)] Avoir le plus de pages possible
\item[D)] Utiliser tous les mots-clés possibles
\end{itemize}
*(Réponse : B)*

\item Le brand building (construction de marque) est important pour le link building car :
\begin{itemize}
\item[A)] Il remplace le SEO
\item[B)] Une marque forte attire naturellement des liens et des mentions
\item[C)] Il est moins cher
\item[D)] Il est plus rapide
\end{itemize}
*(Réponse : B)*

\item La content syndication consiste à :
\begin{itemize}
\item[A)] Copier le contenu d'autres sites
\item[B)] Distribuer votre contenu sur d'autres plateformes autorisées (Medium, LinkedIn)
\item[C)] Supprimer du contenu
\item[D)] Créer du contenu en plusieurs langues
\end{itemize}
*(Réponse : B)*

\item Dans la content syndication, la balise canonique doit :
\begin{itemize}
\item[A)] Pointer vers la plateforme de syndication
\item[B)] Pointer vers l'article original sur votre site
\item[C)] Être absente
\item[D)] Être en nofollow
\end{itemize}
*(Réponse : B)*

\item Haro (Help A Reporter Out) est une plateforme utile pour :
\begin{itemize}
\item[A)] Acheter des liens
\item[B)] Obtenir des mentions dans des articles de presse et des backlinks
\item[C)] Créer des infographies
\item[D)] Analyser les concurrents
\end{itemize}
*(Réponse : B)*

\item Les outils gratuits (calculateurs, générateurs) sont d'excellents linkable assets car :
\begin{itemize}
\item[A)] Ils sont faciles à créer
\item[B)] Ils attirent des backlinks car ils sont utiles et partagés
\item[C)] Ils sont obligatoires pour le SEO
\item[D)] Ils sont rapides à développer
\end{itemize}
*(Réponse : B)*

\item Le link building moderne se concentre sur :
\begin{itemize}
\item[A)] Les échanges de liens automatiques
\item[B)] La création de valeur réelle et les relations publiques numériques
\item[C)] L'achat massif de liens
\item[D)] Les annuaires
\end{itemize}
*(Réponse : B)*

\item Les benchmarks sectoriels sont des linkable assets car :
\begin{itemize}
\item[A)] Ils sont secrets
\item[B)] Ils fournissent des données comparatives que le secteur veut citer
\item[C)] Ils sont payants
\item[D)] Ils sont interdits
\end{itemize}
*(Réponse : B)*

\item Le outreach (prospection) est une étape cruciale du link building qui consiste à :
\begin{itemize}
\item[A)] Acheter des liens
\item[B)] Contacter les webmasters pour proposer votre contenu
\item[C)] Créer du contenu
\item[D)] Analyser les mots-clés
\end{itemize}
*(Réponse : B)*

\item Un email d'outreach efficace doit être :
\begin{itemize}
\item[A)] Très long et détaillé
\item[B)] Personnalisé, concis et value-driven
\item[C)] Automatisé et impersonnel
\item[D]) Uniquement en anglais
\end{itemize}
*(Réponse : B)*

\item Les "Roundups" (articles récapitulatifs) sont une bonne opportunité de link building car :
\begin{itemize}
\item[A)] Ils sont faciles à ignorer
\item[B)] Les auteurs citent souvent plusieurs sources et experts
\item[C)] Ils sont payants
\item[D)] Ils sont réservés aux grands médias
\end{itemize}
*(Réponse : B)*

\item Les témoignages et avis clients sur d'autres sites peuvent générer des backlinks :
\begin{itemize}
\item[A)] Toujours en nofollow
\item[B)] Parfois en dofollow depuis des sites partenaires
\item[C)] Jamais
\item[D)] Uniquement sur les réseaux sociaux
\end{itemize}
*(Réponse : B)*

\item Les ressources "Best of" (meilleures listes, meilleurs outils) sont recherchées car :
\begin{itemize}
\item[A)] Elles sont courtes
\item[B)] Les sites aiment recommander des ressources utiles à leur audience
\item[C)] Elles sont payantes
\item[D)] Elles sont interdites
\end{itemize}
*(Réponse : B)*

\item Le link reclamation consiste à :
\begin{itemize}
\item[A)] Supprimer des liens
\item[B)] Identifier les mentions de marque sans lien et demander l'ajout du lien
\item[C)] Créer des liens internes
\item[D)] Analyser les concurrents
\end{itemize}
*(Réponse : B)*

\item Le study outreach (prospection basée sur des études) consiste à :
\begin{itemize}
\item[A)] Copier les études des concurrents
\item[B)] Créer une étude originale et contacter les journalistes/couvreurs du secteur
\item[C)] Acheter des études
\item[D)] Ignorer les données
\end{itemize}
*(Réponse : B)*

\item Les backlinks depuis des sites d'actualités sont précieux car ils apportent :
\begin{itemize}
\item[A)] Uniquement du trafic direct
\item[B)] Autorité, trafic référent et signal de confiance fort
\item[C)] Uniquement de la visibilité sociale
\item[D)] Des conversions garanties
\end{itemize}
*(Réponse : B)*

\item L'Entity SEO appliqué au netlinking signifie :
\begin{itemize}
\item[A)] Être linké par n'importe quel site
\item[B)] Être associé à des entités reconnues par Google dans votre secteur
\item[C)] Acheter des entités
\item[D)] Créer des entités fictives
\end{itemize}
*(Réponse : B)*

\item Les guides exhaustifs (definitive guides) attirent des backlinks car :
\begin{itemize}
\item[A)] Ils sont courts
\item[B)] Ils deviennent des ressources de référence citées par d'autres
\item[C)] Ils sont faciles à écrire
\item[D)] Ils sont payants
\end{itemize}
*(Réponse : B)*

\item Une stratégie de link building doit être alignée sur :
\begin{itemize}
\item[A)] Le budget disponible uniquement
\item[B)] Les objectifs business et la stratégie de contenu
\item[C)] Le nombre de concurrents
\item[D)] La couleur du site
\end{itemize}
*(Réponse : B)*

\item Les enquêtes et statistiques originales sont efficaces car :
\begin{itemize}
\item[A)] Elles sont faciles à fabriquer
\item[B)] Elles fournissent des données fraîches que les médias veulent citer
\item[C)] Elles sont toujours gratuites
\item[D)] Elles sont courtes
\end{itemize}
*(Réponse : B)*

\item La construction de liens via Digital PR nécessite :
\begin{itemize}
\item[A)] Un gros budget publicitaire
\item[B)] Des relations avec les journalistes et influenceurs du secteur
\item[C)] Un site en plusieurs langues
\item[D)] Un design particulier
\end{itemize}
*(Réponse : B)*

\item Le skyscraper technique nécessite d'identifier d'abord :
\begin{itemize}
\item[A)] Les mots-clés les plus courts
\item[B)] Un contenu performant avec beaucoup de backlinks dans votre niche
\item[C)] Les sites les moins chers
\item[D)] Les images les plus grandes
\end{itemize}
*(Réponse : B)*

\item L'analyse des backlinks des concurrents permet de :
\begin{itemize}
\item[A)] Copier leur site
\item[B)] Identifier des opportunités de sources de liens pertinentes
\item[C)] Les signaler à Google
\item[D)] Ignorer leur stratégie
\end{itemize}
*(Réponse : B)*

\item Le "ego bait" est une technique qui consiste à :
\begin{itemize}
\item[A)] Critiquer les concurrents
\item[B)] Mentionner des influenceurs dans votre contenu pour qu'ils le partagent
\item[C)] Créer du contenu sur vous-même
\item[D)] Acheter des liens
\end{itemize}
*(Réponse : B)*

\item Les ressources de type "ultimate guide" sont efficaces car elles :
\begin{itemize}
\item[A)] Sont rapides à écrire
\item[B)] Couvrent un sujet de façon exhaustive, devenant une référence
\item[C)] Sont toujours courtes
\item[D)] Contiennent beaucoup d'images
\end{itemize}
*(Réponse : B)*

\item Une campagne de Digital PR réussie peut générer :
\begin{itemize}
\item[A)] Uniquement des liens nofollow
\item[B)] Des backlinks de haute qualité depuis des sites médiatiques
\item[C)] Uniquement du trafic direct
\item[D)] Aucun impact SEO
\end{itemize}
*(Réponse : B)*

\item Le "resource page link building" consiste à :
\begin{itemize}
\item[A)] Créer des pages de ressources sur votre site
\item[B)] Trouver des pages de ressources dans votre niche et proposer d'y figurer
\item[C)] Copier les ressources des concurrents
\item[D)] Supprimer des pages de ressources
\end{itemize}
*(Réponse : B)*

\item Les données originales (recherche first-party) sont les plus efficaces car :
\begin{itemize}
\item[A)] Elles sont les moins chères
\item[B)] Elles sont uniques et ne peuvent pas être trouvées ailleurs
\item[C)] Elles sont faciles à produire
\item[D)] Elles sont toujours gratuites
\end{itemize}
*(Réponse : B)*

\item Les liens contextuels éditoriaux (dans le corps d'un article) sont :
\begin{itemize}
\item[A)] Moins efficaces que les liens dans le footer
\item[B)] Les plus efficaces car en contexte pertinent
\item[C)] Toujours nofollow
\item[D)] Interdits
\end{itemize}
*(Réponse : B)*

\item Le "niche edit" consiste à :
\begin{itemize}
\item[A)] Créer un nouvel article avec un lien
\item[B)] Ajouter un lien dans un article existant et pertinent
\item[C)] Modifier le design du site
\item[D)] Supprimer des pages
\end{itemize}
*(Réponse : B)*

\item La création de contenu visuel (infographies, vidéos) aide au link building car :
\begin{itemize}
\item[A)] C'est moins cher
\item[B)] Le contenu visuel est plus partagé et linké que le texte seul
\item[C)] C'est plus rapide
\item[D)] C'est obligatoire
\end{itemize}
*(Réponse : B)*

\item Les communautés en ligne (Reddit, Quora, groupes Facebook) peuvent servir au link building en :
\begin{itemize}
\item[A)] Spammant des liens partout
\item[B)] Partagent du contenu utile que les membres voudront citer
\item[C)] Achetant des upvotes
\item[D)] Créant des faux profils
\end{itemize}
*(Réponse : B)*

\item La technique du "moving man" consiste à :
\begin{itemize}
\item[A)] Changer de site régulièrement
\item[B)] Trouver des entreprises qui ont fermé ou changé de nom pour récupérer leurs liens
\item[C)] Déplacer du contenu
\item[D)] Changer de serveur
\end{itemize}
*(Réponse : B)*

\item Les partenariats avec des influenceurs peuvent générer :
\begin{itemize}
\item[A)] Uniquement du trafic social
\item[B)] Des backlinks de qualité si les influenceurs mentionnent votre site
\item[C)] Uniquement des ventes
\item[D)] Aucun bénéfice SEO
\end{itemize}
*(Réponse : B)*

\item Le monitoring des mentions de marque permet de :
\begin{itemize}
\item[A)] Ignorer ce qui se dit
\item[B)] Identifier les sites qui parlent de vous sans vous linkers
\item[C)] Supprimer les mauvais avis
\item[D)] Créer des faux comptes
\end{itemize}
*(Réponse : B)*

\item L'outil Ahrefs est particulièrement utile pour le link building car il permet de :
\begin{itemize}
\item[A)] Créer des liens automatiquement
\item[B)] Analyser les backlinks des concurrents et trouver des opportunités
\item[C)] Acheter des domaines
\item[D)] Héberger des sites
\end{itemize}
*(Réponse : B)*

\item Un bon profil de backlinks doit idéalement montrer une croissance :
\begin{itemize}
\item[A)] Très rapide et soudaine
\item[B)] Organique, progressive et diversifiée
\item[C)] Nulle
\item[D]) Aléatoire
\end{itemize}
*(Réponse : B)*

\item La diversification des ancres (anchor text) est importante pour éviter :
\begin{itemize}
\item[A)] Les clics
\item[B)] Les pénalités pour sur-optimisation
\item[C)] Le trafic
\item[D)] Les backlinks
\end{itemize}
*(Réponse : B)*

\item Le HARO (Help A Reporter Out) permet aux experts de :
\begin{itemize}
\item[A)] Acheter de la visibilité
\item[B)] Répondre à des demandes de journalistes et obtenir des citations avec backlinks
\item[C)] Vendre des produits
\item[D)] Créer des sites
\end{itemize}
*(Réponse : B)*

\item Le "link reclamation" dans le link building consiste à :
\begin{itemize}
\item[A)] Supprimer des liens de mauvaise qualité
\item[B)] Identifier les mentions de marque sans lien et demander l'ajout du lien
\item[C)] Créer des liens internes supplémentaires
\item[D)] Analyser les backlinks des concurrents
\end{itemize}
*(Réponse : B)*
\end{enumerate}
% ============================================
% CHAPITRE 22 : E-E-A-T (Expérience, Expertise, Autorité, Confiance)
% ============================================
\section*{Chapitre 22 : E-E-A-T : Expérience, Expertise, Autorité, Confiance}

\begin{enumerate}
\item Que signifie l'acronyme E-E-A-T?
\begin{itemize}
\item[A)] Efficiency, Effectiveness, Accuracy, Trust
\item[B)] Experience, Expertise, Authoritativeness, Trustworthiness
\item[C)] Evaluation, Engagement, Authority, Testing
\item[D)] Email, Engagement, Analytics, Traffic
\end{itemize}
*(Réponse : B)*

\item E-E-A-T a été introduit par Google dans :
\begin{itemize}
\item[A)] Les guidelines Google Ads
\item[B)] Les Search Quality Rater Guidelines
\item[C)] Les conditions d'utilisation de GSC
\item[D)] Le fichier robots.txt
\end{itemize}
*(Réponse : B)*

\item E-E-A-T est :
\begin{itemize}
\item[A)] Un facteur de classement direct
\item[B)] Un cadre d'évaluation de la qualité du contenu
\item[C)] Un algorithme de ranking
\item[D)] Une balise HTML
\end{itemize}
*(Réponse : B)*

\item La composante "Experience" d'E-E-A-T évalue :
\begin{itemize}
\item[A)] Le nombre d'années du site
\item[B)] Si l'auteur a une expérience directe du sujet traité
\item[C)] La vitesse du site
\item[D)] Le nombre de backlinks
\end{itemize}
*(Réponse : B)*

\item La composante "Expertise" d'E-E-A-T évalue :
\begin{itemize}
\item[A)] Le design du site
\item[B)] Les qualifications et la légitimité de l'auteur sur le sujet
\item[C)] Le nombre de pages
\item[D)] La popularité sur les réseaux sociaux
\end{itemize}
*(Réponse : B)*

\item La composante "Authoritativeness" (Autorité) évalue :
\begin{itemize}
\item[A)] Le nombre de visiteurs
\item[B)] Si le site est reconnu comme une source de référence
\item[C)] La qualité du code HTML
\item[D)] Le nom de domaine
\end{itemize}
*(Réponse : B)*

\item La composante "Trustworthiness" (Confiance) évalue :
\begin{itemize}
\item[A)] La vitesse de chargement
\item[B)] La sécurité, la transparence et la fiabilité du site
\item[C)] Le nombre de mots
\item[D)] La couleur du site
\end{itemize}
*(Réponse : B)*

\item Que signifie l'acronyme YMYL?
\begin{itemize}
\item[A)] Your Money Your Life
\item[B)] Your Money or Your Life
\item[C)] Young Marketing Youth League
\item[D)] Your Marketing Yearly List
\end{itemize}
*(Réponse : B)*

\item Les pages YMYL sont soumises à des critères E-E-A-T :
\begin{itemize}
\item[A)] Moins stricts
\item[B)] Beaucoup plus stricts
\item[C)] Identiques
\item[D)] Aucun critère particulier
\end{itemize}
*(Réponse : B)*

\item Quel domaine est considéré comme YMYL?
\begin{itemize}
\item[A)] Le divertissement
\item[B)] La santé et la médecine
\item[C)] Les loisirs
\item[D)] Le sport
\end{itemize}
*(Réponse : B)*

\item Quel domaine est considéré comme YMYL?
\begin{itemize}
\item[A)] La mode
\item[B)] La finance et les conseils financiers
\item[C)] La cuisine
\item[D)] Les voyages
\end{itemize}
*(Réponse : B)*

\item Quel domaine est considéré comme YMYL?
\begin{itemize}
\item[A)] Le cinéma
\item[B)] Le droit et les conseils juridiques
\item[C)] La musique
\item[D)] Les jeux vidéo
\end{itemize}
*(Réponse : B)*

\item Quel domaine est considéré comme YMYL?
\begin{itemize}
\item[A)] Les réseaux sociaux
\item[B)] La sécurité des produits et la sécurité en ligne
\item[C)] La photographie
\item[D)] Le jardinage
\end{itemize}
*(Réponse : B)*

\item L'actualité d'importance publique est-elle considérée comme YMYL?
\begin{itemize}
\item[A)] Non
\item[B)] Oui
\item[C)] Uniquement en période électorale
\item[D)] Jamais
\end{itemize}
*(Réponse : B)*

\item L'éducation (formation, certification) est-elle considérée comme YMYL?
\begin{itemize}
\item[A)] Non
\item[B)] Oui
\item[C)] Uniquement pour les universités
\item[D)] Jamais
\end{itemize}
*(Réponse : B)*

\item Pour améliorer l'E-E-A-T, il est recommandé d'afficher :
\begin{itemize}
\item[A)] Uniquement le nom de l'auteur
\item[B)] Le nom, photo, biographie et qualifications de l'auteur
\item[C)] Aucune information sur l'auteur
\item[D)] Uniquement la photo
\end{itemize}
*(Réponse : B)*

\item Citer ses sources et références contribue à améliorer :
\begin{itemize}
\item[A)] La vitesse du site
\item[B)] La confiance (Trustworthiness) et l'expertise
\item[C)] Le design
\item[D)] Le nombre de pages
\end{itemize}
*(Réponse : B)*

\item La mise à jour régulière du contenu est importante pour E-E-A-T car elle démontre :
\begin{itemize}
\item[A)] Une activité constante
\item[B)] La fraîcheur et la précision des informations
\item[C)] Un grand budget
\item[D)] Une équipe nombreuse
\end{itemize}
*(Réponse : B)*

\item Les mentions légales (CGU, politique de confidentialité) contribuent à :
\begin{itemize}
\item[A)] L'expérience
\item[B)] La confiance (Trustworthiness)
\item[C)] La vitesse
\item[D)] Le design
\end{itemize}
*(Réponse : B)*

\item L'utilisation de HTTPS est importante pour E-E-A-T car elle renforce :
\begin{itemize}
\item[A)] L'expertise
\item[B)] La sécurité et la confiance
\item[C)] L'expérience
\item[D)] L'autorité
\end{itemize}
*(Réponse : B)*

\item Les avis vérifiés (Google Reviews, Trustpilot) contribuent à :
\begin{itemize}
\item[A)] L'expertise
\item[B)] La confiance et l'autorité
\item[C)] La vitesse
\item[D)] Le design
\end{itemize}
*(Réponse : B)*

\item La présence hors ligne (conférences, publications, médias) renforce :
\begin{itemize}
\item[A)] Uniquement l'expérience
\item[B)] L'autorité et la crédibilité globale
\item[C)] Uniquement la vitesse
\item[D)] Uniquement le design
\end{itemize}
*(Réponse : B)*

\item E-E-A-T est particulièrement important pour les sites :
\begin{itemize}
\item[A)] De divertissement pur
\item[B)] Qui traitent de sujets sensibles (santé, finance, droit)
\item[C)] De musique
\item[D)] De sport
\end{itemize}
*(Réponse : B)*

\item Une page sur un traitement médical rédigée par un médecin aura un meilleur E-E-A-T qu'une page rédigée par :
\begin{itemize}
\item[A)] Un autre médecin
\item[B)] Un non-professionnel sans expérience médicale
\item[C)] Un journaliste scientifique
\item[D)] Un patient
\end{itemize}
*(Réponse : B)*

\item Pour les pages YMYL, Google exige un niveau d'expertise :
\begin{itemize}
\item[A)] Faible
\item[B)] Élevé et démontrable
\item[C)] Optionnel
\item[D)] Nul
\end{itemize}
*(Réponse : B)*

\item Les backlinks de qualité contribuent à quelle composante E-E-A-T?
\begin{itemize}
\item[A)] Experience
\item[B)] Authoritativeness (Autorité)
\item[C)] Trustworthiness
\item[D)] Expertise
\end{itemize}
*(Réponse : B)*

\item Les "About Us" et "Contact" pages complètes renforcent :
\begin{itemize}
\item[A)] Uniquement l'expertise
\item[B)] La transparence et donc la confiance
\item[C)] Uniquement la vitesse
\item[D)] Uniquement le design
\end{itemize}
*(Réponse : B)*

\item E-E-A-T inclut une quatrième composante ajoutée récemment, laquelle?
\begin{itemize}
\item[A)] Efficiency
\item[B)] Experience (Expérience)
\item[C)] Engagement
\item[D)] Evaluation
\end{itemize}
*(Réponse : B)*

\item Les critères E-E-A-T sont évalués par :
\begin{itemize}
\item[A)] Uniquement des algorithmes
\item[B)] Des évaluateurs humains (Quality Raters) et des algorithmes
\item[C)] Uniquement des robots
\item[D)] Des annonceurs
\end{itemize}
*(Réponse : B)*

\item Google a ajouté le "E" (Experience) à E-A-T en :
\begin{itemize}
\item[A)] 2014
\item[B)] 2022
\item[C)] 2020
\item[D)] 2024
\end{itemize}
*(Réponse : B)*

\item Un contenu sur la randonnée écrit par quelqu'un qui a parcouru le GR20 illustre la composante :
\begin{itemize}
\item[A)] Expertise
\item[B)] Experience (Expérience)
\item[C)] Authoritativeness
\item[D)] Trustworthiness
\end{itemize}
*(Réponse : B)*

\item La transparence sur la politique de confidentialité et les CGU renforce surtout :
\begin{itemize}
\item[A)] L'Experience
\item[B)] La Trustworthiness (Confiance)
\item[C)] L'Expertise
\item[D)] L'Authoritativeness
\end{itemize}
*(Réponse : B)*

\item Les sites YMYL de faible qualité peuvent être :
\begin{itemize}
\item[A)] Recommandés par Google
\item[B)] Déclassés ou exclus des résultats de recherche
\item[C)] Mis en avant
\item[D)] Subventionnés
\end{itemize}
*(Réponse : B)*

\item Pour démontrer l'expertise, il est utile de :
\begin{itemize}
\item[A)] Cacher ses qualifications
\item[B)] Mettre en avant les diplômes, certifications et reconnaissances
\item[C)] Utiliser un pseudonyme
\item[D]) Éviter de parler de soi
\end{itemize}
*(Réponse : B)*

\item L'autorité (Authoritativeness) se construit principalement par :
\begin{itemize}
\item[A)] Le nombre de pages
\item[B)] Les backlinks, mentions et citations provenant de sources reconnues
\item[C)] La vitesse du site
\item[D)] Le design responsive
\end{itemize}
*(Réponse : B)*

\item Une page sans auteur identifiable aura du mal à démontrer :
\begin{itemize}
\item[A)] La vitesse
\item[B)] L'expertise et l'expérience
\item[C)] Le responsive design
\item[D)] La compatibilité mobile
\end{itemize}
*(Réponse : B)*

\item Les pages "A propos" détaillées contribuent à :
\begin{itemize}
\item[A)] L'indexation
\item[B)] La transparence et la confiance
\item[C)] La vitesse
\item[D)] Les backlinks
\end{itemize}
*(Réponse : B)*

\item Les avis négatifs non traités peuvent impacter négativement :
\begin{itemize}
\item[A)] La vitesse
\item[B)] La confiance (Trustworthiness)
\item[C)] Le responsive
\item[D)] Le maillage interne
\end{itemize}
*(Réponse : B)*

\item E-E-A-T est un concept qui s'applique :
\begin{itemize}
\item[A)] Uniquement aux sites YMYL
\item[B)] À tous les sites, mais avec plus de rigueur pour les sites YMYL
\item[C)] Uniquement aux sites e-commerce
\item[D)] Uniquement aux blogs personnels
\end{itemize}
*(Réponse : B)*

\item Les Quality Raters de Google utilisent E-E-A-T pour :
\begin{itemize}
\item[A)] Classer directement les sites
\item[B)] Évaluer la qualité des résultats et améliorer les algorithmes
\item[C)] Supprimer des pages
\item[D)] Créer du contenu
\end{itemize}
*(Réponse : B)*

\item Le contenu généré par IA de faible qualité peut impacter négativement :
\begin{itemize}
\item[A)] La vitesse
\item[B)] La perception d'expertise et de confiance (E-E-A-T)
\item[C)] Le responsive
\item[D)] Le maillage interne
\end{itemize}
*(Réponse : B)*

\item Pour améliorer l'Authoritativeness, il faut chercher à être cité par :
\begin{itemize}
\item[A)] N'importe quel site
\item[B)] Des sites reconnus et pertinents dans votre secteur
\item[C)] Uniquement des sites .com
\item[D)] Des sites anonymes
\end{itemize}
*(Réponse : B)*

\item La qualité rédactionnelle (orthographe, grammaire) contribue à :
\begin{itemize}
\item[A)] La vitesse
\item[B)] La crédibilité et donc l'expertie et la confiance
\item[C)] L'indexation
\item[D)] Le responsive
\end{itemize}
*(Réponse : B)*

\item Les pages de témoignages clients avec des cas réels renforcent :
\begin{itemize}
\item[A)] Uniquement l'expertise
\item[B)] La preuve sociale et la confiance
\item[C)] Uniquement la vitesse
\item[D)] Uniquement le design
\end{itemize}
*(Réponse : B)*

\item E-E-A-T est un facteur de classement ? (vrai ou faux)
\begin{itemize}
\item[A)] Vrai, c'est un facteur direct
\item[B)] Faux, c'est un cadre d'évaluation qui influence les algorithmes
\item[C)] Vrai, c'est le facteur le plus important
\item[D)] Faux, Google ne l'utilise pas
\end{itemize}
*(Réponse : B)*

\item Les biographies d'auteurs doivent idéalement inclure :
\begin{itemize}
\item[A)] Uniquement le nom
\item[B)] Les qualifications, l'expérience et les réalisations pertinentes
\item[C)] Uniquement la photo
\item[D)] Uniquement les loisirs
\end{itemize}
*(Réponse : B)*

\item Un site de conseils financiers sans indications sur les qualifications de l'auteur aura :
\begin{itemize}
\item[A)] Un bon E-E-A-T
\item[B)] Un mauvais E-E-A-T (risque de déclassement)
\item[C)] Aucun impact
\item[D)] Un classement amélioré
\end{itemize}
*(Réponse : B)*

\item Les données structurées (schema.org) comme "Author" et "Person" aident à :
\begin{itemize}
\item[A)] Accélérer le site
\item[B)] Communiquer explicitement les informations d'auteur à Google
\item[C)] Réduire le poids des pages
\item[D)] Améliorer le design
\end{itemize}
*(Réponse : B)*

\item Le content marketing de qualité aligné sur E-E-A-T doit :
\begin{itemize}
\item[A)] Être produit rapidement et en grande quantité
\item[B)] Être créé par des experts, sourcé et régulièrement mis à jour
\item[C)] Être copié des concurrents
\item[D)] Être rédigé par IA sans relecture
\end{itemize}
*(Réponse : B)*

\item L'évaluation E-E-A-T d'un site évolue :
\begin{itemize}
\item[A)] Très lentement et rarement
\item[B)] Avec le temps, en fonction des actions d'amélioration continue
\item[C)] Uniquement lors du lancement
\item[D)] De façon aléatoire
\end{itemize}
*(Réponse : B)*

\end{enumerate}

% ============================================
% CHAPITRE 23 : Stratégie de contenu et Content Marketing SEO
% ============================================
\section*{Chapitre 23 : Stratégie de contenu et Content Marketing SEO}

\begin{enumerate}
\item Qu'est-ce que le Content Marketing SEO?
\begin{itemize}
\item[A)] Publier le plus d'articles possible
\item[B)] Créer, distribuer et promouvoir du contenu pour attirer un trafic organique qualifié
\item[C)] Copier le contenu des concurrents
\item[D)] Uniquement écrire des articles de blog
\end{itemize}
*(Réponse : B)*

\item Combien de piliers compte le Content Marketing SEO selon le cours?
\begin{itemize}
\item[A)] 3
\item[B)] 4
\item[C)] 5
\item[D)] 6
\end{itemize}
*(Réponse : B)*

\item Quel est le premier pilier du Content Marketing SEO?
\begin{itemize}
\item[A)] Distribution
\item[B)] Recherche (data-driven)
\item[C)] Création
\item[D)] Mesure
\end{itemize}
*(Réponse : B)*

\item Le calendrier éditorial est un outil qui permet de :
\begin{itemize}
\item[A)] Supprimer du contenu
\item[B)] Structurer la production de contenu dans le temps
\item[C)] Analyser les concurrents
\item[D)] Acheter des backlinks
\end{itemize}
*(Réponse : B)*

\item Le contenu TOFU (Top of Funnel) vise quel type d'audience?
\begin{itemize}
\item[A)] Des clients prêts à acheter
\item[B)] Des prospects en début de parcours d'achat, besoin d'information
\item[C)] Des clients fidèles
\item[D)] Des anciens clients
\end{itemize}
*(Réponse : B)*

\item Le contenu MOFU (Middle of Funnel) vise à :
\begin{itemize}
\item[A)] Attirer de nouveaux visiteurs
\item[B)] Éduquer et faire considérer la solution
\item[C)] Finaliser la vente
\item[D)] Fidéliser
\end{itemize}
*(Réponse : B)*

\item Le contenu BOFU (Bottom of Funnel) a pour objectif de :
\begin{itemize}
\item[A)] Attirer du trafic
\item[B)] Convertir le prospect en client
\item[C)] Informer
\item[D)] Éduquer
\end{itemize}
*(Réponse : B)*

\item Le modèle Hub-and-Spoke est aussi appelé :
\begin{itemize}
\item[A)] Site structure
\item[B)] Topic Cluster
\item[C)] Funnel model
\item[D)] Star model
\end{itemize}
*(Réponse : B)*

\item Dans le modèle Hub-and-Spoke, la page pilier doit faire environ :
\begin{itemize}
\item[A)] 500-1000 mots
\item[B)] 3000-5000 mots
\item[C)] 100-200 mots
\item[D)] Plus de 10000 mots
\end{itemize}
*(Réponse : B)*

\item Les pages clusters dans le modèle Hub-and-Spoke font environ :
\begin{itemize}
\item[A)] 300-500 mots
\item[B)] 1500-2500 mots
\item[C)] 5000-7000 mots
\item[D)] 10000 mots
\end{itemize}
*(Réponse : B)*

\item Dans le modèle Hub-and-Spoke, les liens internes doivent :
\begin{itemize}
\item[A)] Pointer uniquement du cluster vers le pilier
\item[B)] Connecter tous les clusters vers le pilier et vice-versa
\item[C)] Ne pas exister
\item[D)] Pointer uniquement vers la page d'accueil
\end{itemize}
*(Réponse : B)*

\item La Skyscraper Technique a été popularisée par :
\begin{itemize}
\item[A)] Neil Patel
\item[B)] Brian Dean (Backlinko)
\item[C)] Rand Fishkin
\item[D)] John Mueller
\end{itemize}
*(Réponse : B)*

\item La première étape de la Skyscraper Technique est :
\begin{itemize}
\item[A)] Créer un contenu amélioré
\item[B)] Trouver un contenu performant dans votre niche
\item[C)] Promouvoir le contenu
\item[D)] Analyser les mots-clés
\end{itemize}
*(Réponse : B)*

\item Les Content Pillars sont :
\begin{itemize}
\item[A)] Tous les sujets possibles
\item[B)] Les 3 à 5 sujets fondamentaux autour desquels s'articule la stratégie
\item[C)] Les articles les plus longs
\item[D)] Les pages les plus visitées
\end{itemize}
*(Réponse : B)*

\item Quel outil permet de trouver les questions posées par les utilisateurs?
\begin{itemize}
\item[A)] Google Analytics
\item[B)] AnswerThePublic
\item[C)] Google Search Console
\item[D)] Screaming Frog
\end{itemize}
*(Réponse : B)*

\item Le Content Gap Analysis permet d'identifier :
\begin{itemize}
\item[A)] Les pages cassées
\item[B)] Les sujets que vos concurrents couvrent mais pas vous
\item[C)] Les mots-clés trop chers
\item[D)] Les backlinks toxiques
\end{itemize}
*(Réponse : B)*

\item Le contenu long-form (plus de 2000 mots) tend à :
\begin{itemize}
\item[A)] Moins bien performer
\item[B)] Mieux performer dans les SERP
\item[C)] Être ignoré
\item[D)] Être pénalisé
\end{itemize}
*(Réponse : B)*

\item Le repurposing de contenu consiste à :
\begin{itemize}
\item[A)] Supprimer du contenu
\item[B)] Adapter un même contenu sous différents formats
\item[C)] Copier du contenu
\item[D)] Acheter du contenu
\end{itemize}
*(Réponse : B)*

\item La règle des 7 contacts indique qu'un prospect a besoin de voir votre contenu en moyenne combien de fois avant d'agir?
\begin{itemize}
\item[A)] 3 fois
\item[B)] 7 fois
\item[C)] 10 fois
\item[D)] 1 fois
\end{itemize}
*(Réponse : B)*

\item Le modèle d'attribution "first-click" attribue la conversion :
\begin{itemize}
\item[A)] Au dernier point de contact
\item[B)] Au premier point de contact
\item[C)] À tous les points de contact
\item[D)] À aucun point de contact
\end{itemize}
*(Réponse : B)*

\item Le modèle d'attribution "last-click" attribue la conversion :
\begin{itemize}
\item[A)] Au premier point de contact
\item[B)] Au dernier point de contact
\item[C)] À tous les points de contact
\item[D)] De façon algorithmique
\end{itemize}
*(Réponse : B)*

\item Quel KPI mesure le nombre de sessions provenant de Google?
\begin{itemize}
\item[A)] Taux de conversion
\item[B)] Trafic organique
\item[C)] Engagement
\item[D)] Brand awareness
\end{itemize}
*(Réponse : B)*

\item Le "Share of Voice" mesure :
\begin{itemize}
\item[A)] Le nombre de backlinks
\item[B)] La visibilité par rapport aux concurrents
\item[C)] La vitesse du site
\item[D)] Le nombre de pages
\end{itemize}
*(Réponse : B)*

\item Le contenu intermédiaire (1000-2000 mots) comprend notamment :
\begin{itemize}
\item[A)] Les fiches produits
\item[B)] Les articles de blog standard, tutoriels et listicles
\item[C)] Les livres blancs
\item[D)] Les glossaires
\end{itemize}
*(Réponse : B)*

\item Le contenu court (moins de 1000 mots) comprend notamment :
\begin{itemize}
\item[A)] Les guides complets
\item[B)] Les fiches produits optimisées et FAQ structurées
\item[C)] Les études de cas
\item[D)] Les whitepapers
\end{itemize}
*(Réponse : B)*

\item La distribution de contenu par email marketing est un canal de :
\begin{itemize}
\item[A)] Distribution payante
\item[B)] Distribution directe vers votre base d'abonnés
\item[C)] Distribution organique uniquement
\item[D)] Distribution sociale
\end{itemize}
*(Réponse : B)*

\item Le repurposing d'un article de blog peut donner :
\begin{itemize}
\item[A)] Un seul format
\item[B)] Vidéo YouTube + Podcast + Infographie + LinkedIn carousel
\item[C)] Rien d'autre
\item[D)] Uniquement une infographie
\end{itemize}
*(Réponse : B)*

\item L'analyse des SERP avant de créer du contenu permet de comprendre :
\begin{itemize}
\item[A)] Le prix des mots-clés
\item[B)] Le type de contenu attendu par Google pour cette requête
\item[C)] Les backlinks des concurrents
\item[D)] La vitesse des concurrents
\end{itemize}
*(Réponse : B)*

\item Google Trends est un outil utile pour analyser :
\begin{itemize}
\item[A)] Les backlinks
\item[B)] Les tendances et la saisonnalité des sujets
\item[C)] La vitesse du site
\item[D)] Les données structurées
\end{itemize}
*(Réponse : B)*

\item Ahrefs Content Explorer permet de :
\begin{itemize}
\item[A)] Analyser les mots-clés
\item[B)] Trouver le contenu le plus partagé par thématique
\item[C)] Auditer la vitesse
\item[D)] Vérifier le mobile-first
\end{itemize}
*(Réponse : B)*

\item La création de contenu doit être alignée sur :
\begin{itemize}
\item[A)] Les préférences personnelles
\item[B)] Les objectifs business et l'intention de recherche
\item[C)] Le budget disponible
\item[D)] La couleur du site
\end{itemize}
*(Réponse : B)*

\item Les "People Also Ask" dans Google SERP aident à :
\begin{itemize}
\item[A)] Trouver des images
\item[B)] Identifier les questions connexes pour enrichir le contenu
\item[C)] Analyser la vitesse
\item[D)] Vérifier le responsive
\end{itemize}
*(Réponse : B)*

\item Un calendrier éditorial doit inclure :
\begin{itemize}
\item[A)] Uniquement les titres
\item[B)] Périodicité, saisonnalité, événements, étapes du funnel et piliers thématiques
\item[C)] Uniquement les dates
\item[D)] Uniquement les auteurs
\end{itemize}
*(Réponse : B)*

\item L'outil Notion est utilisé en content marketing pour :
\begin{itemize}
\item[A)] Analyser les mots-clés
\item[B)] Créer une base de données éditoriale complète
\item[C)] Auditer le site
\item[D)] Suivre les backlinks
\end{itemize}
*(Réponse : B)*

\item La mesure du ROI du contenu inclut des KPIs comme :
\begin{itemize}
\item[A)] La météo
\item[B)] Le trafic organique, les backlinks générés, le taux de conversion
\item[C)] La température du serveur
\item[D)] La couleur du logo
\end{itemize}
*(Réponse : B)*

\item L'analyse concurrentielle de contenu commence par :
\begin{itemize}
\item[A)] Copier les concurrents
\item[B)] Identifier ses vrais concurrents SEO (pas seulement business)
\item[C)] Ignorer les concurrents
\item[D)] Acheter leurs backlinks
\end{itemize}
*(Réponse : B)*

\item Un article de blog standard fait partie du contenu :
\begin{itemize}
\item[A)] Long-form (>2000 mots)
\item[B)] Intermédiaire (1000-2000 mots)
\item[C)] Court (<1000 mots)
\item[D)] Micro-contenu
\end{itemize}
*(Réponse : B)*

\item Les whitepapers sont un format de contenu :
\begin{itemize}
\item[A)] Court
\item[B)] Long-form, premium, pour génération de leads
\item[C)] Intermédiaire
\item[D)] Micro
\end{itemize}
*(Réponse : B)*

\item La promotion multi-canal du contenu fait partie du pilier :
\begin{itemize}
\item[A)] Recherche
\item[B)] Distribution
\item[C)] Création
\item[D)] Mesure
\end{itemize}
*(Réponse : B)*

\item Le modèle d'attribution multi-touch distribue le crédit de conversion :
\begin{itemize}
\item[A)] Au premier contact uniquement
\item[B)] Entre tous les points de contact
\item[C)] Au dernier contact uniquement
\item[D)] De façon aléatoire
\end{itemize}
*(Réponse : B)*

\item Une stratégie de contenu efficace doit inclure du contenu pour chaque étape :
\begin{itemize}
\item[A)] Uniquement BOFU
\item[B)] TOFU, MOFU et BOFU
\item[C)] Uniquement TOFU
\item[D)] Uniquement MOFU
\end{itemize}
*(Réponse : B)*

\item Les études de cas sont un format de contenu :
\begin{itemize}
\item[A)] Court et rapide
\item[B)] Long-form, data-driven, qui démontre des résultats concrets
\item[C)] Uniquement pour les réseaux sociaux
\item[D]) Micro-contenu
\end{itemize}
*(Réponse : B)*

\item Le rythme de publication dans un calendrier éditorial doit être :
\begin{itemize}
\item[A)] Aléatoire
\item[B)] Régulier (hebdomadaire, bi-mensuel)
\item[C)] Quotidien
\item[D)] Annuel
\end{itemize}
*(Réponse : B)*

\item Google Autocomplete est utile pour la recherche de contenu car il montre :
\begin{itemize}
\item[A)] Les prix des mots-clés
\item[B)] Les suggestions de recherche en temps réel des utilisateurs
\item[C)] Les backlinks
\item[D)] La vitesse du site
\end{itemize}
*(Réponse : B)*

\item Le Content Marketing SEO vise à attirer du trafic :
\begin{itemize}
\item[A)] Payant
\item[B)] Organique qualifié
\item[C)] Social uniquement
\item[D)] Direct
\end{itemize}
*(Réponse : B)*

\item La création de contenu sans distribution est comparée à :
\begin{itemize}
\item[A)] Un livre sans couverture
\item[B)] Une fête sans invités
\item[C)] Une voiture sans roues
\item[D)] Une maison sans toit
\end{itemize}
*(Réponse : B)*

\item Les benchmarks sectoriels sont des contenus qui :
\begin{itemize}
\item[A)] Comparent les prix
\item[B)] Fournissent des données comparatives utiles pour le secteur
\item[C)] Copient les concurrents
\item[D)] Sont confidentiels
\end{itemize}
*(Réponse : B)*

\item La période idéale pour actualiser un contenu existant dépend de :
\begin{itemize}
\item[A)] La couleur du site
\item[B)] La saisonnalité et l'évolution du sujet
\item[C)] Le nombre de mots
\item[D)] Le nombre d'images
\end{itemize}
*(Réponse : B)*

\item Un Topic Cluster (Hub-and-Spoke) permet de :
\begin{itemize}
\item[A)] Réduire le nombre de pages
\item[B)] Couvrir un sujet de façon exhaustive et interconnectée
\item[C)] Supprimer du contenu
\item[D)] Acheter des backlinks
\end{itemize}
*(Réponse : B)*

\item Le contenu saisonnier dans un calendrier éditorial est important pour :
\begin{itemize}
\item[A)] Réduire les coûts
\item[B)] Capitaliser sur les cycles saisonniers du marché
\item[C)] Simplifier la production
\item[D]) Remplir le calendrier
\end{itemize}
*(Réponse : B)*

\end{enumerate}

% ============================================
% CHAPITRE 24 : Le SEO Local
% ============================================
\section*{Chapitre 24 : Le SEO Local}

\begin{enumerate}
\item Quel pourcentage des recherches Google a une intention locale?
\begin{itemize}
\item[A)] 25\%
\item[B)] 46\%
\item[C)] 76\%
\item[D)] 90\%
\end{itemize}
*(Réponse : B)*

\item Quel pourcentage des personnes qui effectuent une recherche locale visitent un magasin dans la journée?
\begin{itemize}
\item[A)] 46\%
\item[B)] 76\%
\item[C)] 90\%
\item[D)] 50\%
\end{itemize}
*(Réponse : B)*

\item Quel est l'outil le plus important pour le SEO local?
\begin{itemize}
\item[A)] Google Analytics
\item[B)] Google Business Profile (GBP)
\item[C)] Google Ads
\item[D)] Google Tag Manager
\end{itemize}
*(Réponse : B)*

\item Google Business Profile était anciennement connu sous le nom de :
\begin{itemize}
\item[A)] Google Places
\item[B)] Google My Business
\item[C)] Google Local
\item[D)] Google Maps
\end{itemize}
*(Réponse : B)*

\item Que signifie NAP dans le contexte du SEO local?
\begin{itemize}
\item[A)] Network Access Protocol
\item[B)] Name, Address, Phone
\item[C)] News and Press
\item[D)] Navigation and Pages
\end{itemize}
*(Réponse : B)*

\item La consistance NAP signifie que :
\begin{itemize}
\item[A)] Les informations peuvent varier selon les annuaires
\item[B)] Le nom, adresse et téléphone doivent être identiques sur tous les annuaires
\item[C)] Seul le nom doit être cohérent
\item[D)] L'adresse peut changer
\end{itemize}
*(Réponse : B)*

\item Dans Google Business Profile, la catégorie d'activité doit être :
\begin{itemize}
\item[A)] Générique
\item[B)] La plus précise possible (principale et secondaires)
\item[C)] Changée chaque semaine
\item[D)] Absente
\end{itemize}
*(Réponse : B)*

\item Les horaires d'ouverture dans GBP doivent inclure :
\begin{itemize}
\item[A)] Uniquement les heures d'ouverture standard
\item[B)] Les heures précises incluant les jours fériés
\item[C)] Aucune indication
\item[D)] Uniquement la fermeture
\end{itemize}
*(Réponse : B)*

\item Les photos et vidéos dans GBP doivent être :
\begin{itemize}
\item[A)] Peu nombreuses
\item[B)] De haute qualité
\item[C)] Uniquement des photos
\item[D)] Uniquement des vidéos
\end{itemize}
*(Réponse : B)*

\item La gestion des avis clients dans GBP inclut :
\begin{itemize}
\item[A)] Supprimer les avis négatifs
\item[B)] Répondre professionnellement aux avis (positifs et négatifs)
\item[C)] Ignorer les avis
\item[D]) Acheter des avis positifs
\end{itemize}
*(Réponse : B)*

\item Les posts réguliers dans GBP peuvent contenir :
\begin{itemize}
\item[A)] Uniquement du texte
\item[B)] Actualités, offres et événements
\item[C)] Uniquement des offres
\item[D)] Uniquement des liens
\end{itemize}
*(Réponse : B)*

\item La section Questions/Réponses dans GBP doit être :
\begin{itemize}
\item[A)] Ignorée
\item[B)] Gérée activement
\item[C)] Supprimée
\item[D)] Automatisée
\end{itemize}
*(Réponse : B)*

\item Parmi ces annuaires, lequel est un annuaire clé pour le SEO local en France?
\begin{itemize}
\item[A)] Google Ads
\item[B)] PagesJaunes
\item[C)] Facebook
\item[D)] Instagram
\end{itemize}
*(Réponse : B)*

\item Bing Places est un annuaire pour quel moteur de recherche?
\begin{itemize}
\item[A)] Google
\item[B)] Bing
\item[C)] Yahoo
\item[D)] DuckDuckGo
\end{itemize}
*(Réponse : B)*

\item Yelp et TripAdvisor sont des annuaires importants pour quels types d'activités?
\begin{itemize}
\item[A)] Les industries
\item[B)] La restauration, l'hôtellerie et les services
\item[C)] La finance
\item[D)] La technologie
\end{itemize}
*(Réponse : B)*

\item Les avis Google Reviews impactent directement :
\begin{itemize}
\item[A)] La vitesse du site
\item[B)] Le classement local et la confiance des utilisateurs
\item[C)] Le design
\item[D)] Les backlinks
\end{itemize}
*(Réponse : B)*

\item Une fiche GBP complète et vérifiée est un signal de :
\begin{itemize}
\item[A)] Faible qualité
\item[B)] Confiance et légitimité pour Google
\item[C)] Design moderne
\item[D)] Grande taille d'entreprise
\end{itemize}
*(Réponse : B)*

\item Le SEO local est essentiel pour :
\begin{itemize}
\item[A)] Les sites purement e-commerce sans adresse physique
\item[B)] Les entreprises ayant une présence physique (magasin, cabinet, restaurant)
\item[C)] Les sites internationaux
\item[D)] Les blogs seulement
\end{itemize}
*(Réponse : B)*

\item La proximité géographique est un facteur de classement dans le :
\begin{itemize}
\item[A)] SEO technique
\item[B)] Local Pack (les 3 résultats locaux de Google)
\item[C)] SEO on-page
\item[D)] Netlinking
\end{itemize}
*(Réponse : B)*

\item Les avis positifs et nombreux sur GBP améliorent :
\begin{itemize}
\item[A)] Uniquement le taux de clic
\item[B)] Le classement local et le taux de conversion
\item[C)] Uniquement la vitesse
\item[D)] Uniquement le design
\end{itemize}
*(Réponse : B)*

\item Le Local Pack de Google affiche généralement combien de résultats?
\begin{itemize}
\item[A)] 1
\item[B)] 3
\item[C)] 5
\item[D)] 10
\end{itemize}
*(Réponse : B)*

\item Les backlinks provenant de sites locaux (journal local, chambre de commerce) sont :
\begin{itemize}
\item[A)] Sans intérêt
\item[B)] Particulièrement pertinents pour le SEO local
\item[C)] Interdits
\item[D)] Identiques aux autres backlinks
\end{itemize}
*(Réponse : B)*

\item Le NAP doit être cohérent sur :
\begin{itemize}
\item[A)] Uniquement Google Business Profile
\item[B)] Tous les annuaires et mentions en ligne
\item[C)] Uniquement le site web
\item[D)] Uniquement les réseaux sociaux
\end{itemize}
*(Réponse : B)*

\item Les citations NAP sont des mentions de votre entreprise sur :
\begin{itemize}
\item[A)] Votre propre site
\item[B)] Des annuaires et sites tiers
\item[C)] Google Analytics
\item[D)] Votre serveur
\end{itemize}
*(Réponse : B)*

\item La catégorie secondaire dans GBP permet de :
\begin{itemize}
\item[A)] Ajouter des services supplémentaires que vous proposez
\item[B)] Remplacer la catégorie principale
\item[C)] N'a aucun impact
\item[D]) Être supprimée
\end{itemize}
*(Réponse : B)*

\item Google Maps est étroitement lié à quel outil de SEO local?
\begin{itemize}
\item[A)] Google Analytics
\item[B)] Google Business Profile
\item[C)] Google Ads
\item[D)] Google Drive
\end{itemize}
*(Réponse : B)*

\item Les avis Google avec photos sont considérés comme :
\begin{itemize}
\item[A)] Moins fiables
\item[B)] Plus engageants et dignes de confiance
\item[C)] Identiques aux avis sans photos
\item[D)] Suspects
\end{itemize}
*(Réponse : B)*

\item Pour un restaurant, quel type d'annuaire est particulièrement important?
\begin{itemize}
\item[A)] LinkedIn
\item[B)] TripAdvisor
\item[C)] GitHub
\item[D)] Stack Overflow
\end{itemize}
*(Réponse : B)*

\item Les horaires spéciaux (jours fériés) dans GBP doivent être :
\begin{itemize}
\item[A)] Ignorés
\item[B)] Mis à jour dans GBP
\item[C)] Uniquement mentionnés sur le site
\item[D)] Laissés vides
\end{itemize}
*(Réponse : B)*

\item Une incohérence NAP (téléphone différent sur deux annuaires) peut :
\begin{itemize}
\item[A)] Améliorer le classement
\item[B)] Nuire à la confiance de Google et au classement local
\item[C)] Être ignorée
\item[D]) Être bénéfique
\end{itemize}
*(Réponse : B)*

\item Les questions fréquentes dans la fiche GBP doivent être :
\begin{itemize}
\item[A)] Ignorées
\item[B)] Renseignées avec des réponses précises
\item[C)] Supprimées
\item[D)] Automatisées
\end{itemize}
*(Réponse : B)*

\item Le SEO local est particulièrement pertinent pour les recherches avec :
\begin{itemize}
\item[A)] "Comment" et "Pourquoi"
\item[B)] "Près de moi" ou "à proximité"
\item[C)] "Définition"
\item[D)] "Histoire de"
\end{itemize}
*(Réponse : B)*

\item L'utilisation de photos récentes dans GBP est importante pour montrer :
\begin{itemize}
\item[A)] Le design du site
\item[B)] L'activité réelle et actuelle de l'établissement
\item[C)] Les backlinks
\item[D)] La vitesse
\end{itemize}
*(Réponse : B)*

\item Les entreprises avec une adresse physique vérifiée ont quel avantage en SEO local?
\begin{itemize}
\item[A)] Aucun avantage
\item[B)] Une meilleure crédibilité et un meilleur classement dans le Local Pack
\item[C)] Un trafic réduit
\item[D)] Des backlinks automatiques
\end{itemize}
*(Réponse : B)*

\item Le fait de répondre aux avis Google montre que l'entreprise :
\begin{itemize}
\item[A)] Est grande
\item[B)] Est à l'écoute et engagée envers ses clients
\item[C)] A beaucoup de temps libre
\item[D)] Paye pour des avis
\end{itemize}
*(Réponse : B)*

\item Les annuaires professionnels par secteur sont importants pour :
\begin{itemize}
\item[A)] Le SEO général
\item[B)] Renforcer les citations NAP dans un contexte pertinent
\item[C)] La vitesse du site
\item[D)] Le design
\end{itemize}
*(Réponse : B)*

\item Le champ "Questions et réponses" dans GBP est alimenté par :
\begin{itemize}
\item[A)] Uniquement le propriétaire
\item[B)] Les utilisateurs et le propriétaire
\item[C)] Uniquement Google
\item[D)] Uniquement les concurrents
\end{itemize}
*(Réponse : B)*

\item La section "Avis" dans GBP influence quel aspect du SEO local?
\begin{itemize}
\item[A)] Uniquement le trafic direct
\item[B)] La réputation et le classement local
\item[C)] Uniquement la vitesse
\item[D)] Uniquement le design
\end{itemize}
*(Réponse : B)*

\item Une entreprise avec plusieurs établissements doit :
\begin{itemize}
\item[A)] Utiliser une seule fiche GBP
\item[B)] Créer une fiche GBP distincte pour chaque établissement
\item[C)] Ne pas utiliser GBP
\item[D)] Utiliser une fiche par ville uniquement
\end{itemize}
*(Réponse : B)*

\item La recherche "près de moi" est un exemple typique de recherche :
\begin{itemize}
\item[A)] Transactionnelle
\item[B)] Locale
\item[C)] Informationnelle
\item[D)] Navigationnelle
\end{itemize}
*(Réponse : B)*

\item Les photos de l'établissement dans GBP aident les utilisateurs à :
\begin{itemize}
\item[A)] Voir le site web
\item[B)] Se faire une idée réaliste de l'établissement avant de visiter
\item[C)] Lire les avis
\item[D)] Contacter l'entreprise
\end{itemize}
*(Réponse : B)*

\item GBP permet de publier des offres spéciales et des événements sous forme de :
\begin{itemize}
\item[A)] Emails
\item[B)] Posts
\item[C)] Publicités
\item[D]) Notifications push
\end{itemize}
*(Réponse : B)*

\item Les backlinks locaux proviennent idéalement de :
\begin{itemize}
\item[A)] Sites internationaux
\item[B)] Sites locaux pertinents (presse locale, partenaires locaux)
\item[C)] Sites étrangers
\item[D)] Annuaires génériques
\end{itemize}
*(Réponse : B)*

\item La vérification de la fiche GBP est :
\begin{itemize}
\item[A)] Optionnelle
\item[B)] Obligatoire pour apparaître dans Google Maps et le Local Pack
\item[C)] Impossible
\item[D)] Payante
\end{itemize}
*(Réponse : B)*

\item Le SEO local utilise les signaux de géolocalisation pour :
\begin{itemize}
\item[A)] Déterminer la langue du site
\item[B)] Classer les entreprises pertinentes à proximité de l'utilisateur
\item[C)] Analyser les mots-clés
\item[D)] Créer des backlinks
\end{itemize}
*(Réponse : B)*

\item Les catégories d'activité dans GBP doivent représenter :
\begin{itemize}
\item[A)] Uniquement l'activité principale
\item[B)] L'activité principale et les services spécifiques proposés
\item[C)] Des activités fictives
\item[D)] Des activités de concurrents
\end{itemize}
*(Réponse : B)*

\item La fréquence de publication des posts GBP recommandée est :
\begin{itemize}
\item[A)] Une fois par an
\item[B)] Régulière (au moins une fois par semaine)
\item[C)] Une seule fois
\item[D]) Jamais
\end{itemize}
*(Réponse : B)*

\item Une adresse complète dans GBP doit inclure :
\begin{itemize}
\item[A)] Uniquement la rue
\item[B)] Rue, code postal, ville, pays (boîte postale déconseillée)
\item[C)] Uniquement la ville
\item[D)] Uniquement le code postal
\end{itemize}
*(Réponse : B)*

\item L'optimisation du SEO local permet d'apparaître dans :
\begin{itemize}
\item[A)] Google Shopping
\item[B)] Le Local Pack, Google Maps et les résultats locaux organiques
\item[C)] Uniquement Google Ads
\item[D)] Uniquement les résultats payants
\end{itemize}
*(Réponse : B)*

\item Le numéro de téléphone dans GBP doit être :
\begin{itemize}
\item[A)] Un numéro surtaxé
\item[B)] Un numéro local qui répond directement à l'entreprise
\item[C)] Un numéro international
\item[D)] Un numéro sans indicatif
\end{itemize}
*(Réponse : B)*
\end{enumerate}

% ======================================================================
% CHAPITRE 25 : MIGRATION SEO ET REDIRECTION
% ======================================================================
\section*{Chapitre 25 : Migration SEO et Redirection}

\begin{enumerate}
\item Qu'est-ce qu'une migration SEO ?
\begin{itemize}
\item[A)] Un changement de nom de domaine uniquement
\item[B)] Tout changement significatif apporté à un site qui peut impacter sa visibilité
\item[C)] Une simple mise à jour de contenu
\item[D)] Le passage de HTTP à HTTPS uniquement
\end{itemize}
*(Réponse : B)*

\item Quel code HTTP est recommandé pour une redirection permanente ?
\begin{itemize}
\item[A)] 302
\item[B)] 307
\item[C)] 301
\item[D)] 404
\end{itemize}
*(Réponse : C)*

\item Quel pourcentage du PageRank un code 301 transmet-il selon Google ?
\begin{itemize}
\item[A)] 50-60\%
\item[B)] 75-80\%
\item[C)] 90-99\%
\item[D)] 100\%
\end{itemize}
*(Réponse : C)*

\item Dans quel contexte utilise-t-on une redirection 302 ?
\begin{itemize}
\item[A)] Changement permanent d'URL
\item[B)] Migration de domaine définitive
\item[C)] Test A/B, page en maintenance, contenu saisonnier
\item[D)] Suppression définitive d'une page
\end{itemize}
*(Réponse : C)*

\item Qu'est-ce que le mapping URL dans le cadre d'une migration ?
\begin{itemize}
\item[A)] Un plan du site web
\item[B)] La table de correspondance entre anciennes et nouvelles URLs
\item[C)] Une carte des serveurs
\item[D)] Un fichier de configuration DNS
\end{itemize}
*(Réponse : B)*

\item Quel est le premier élément à réaliser avant une migration SEO ?
\begin{itemize}
\item[A)] Changer les DNS immédiatement
\item[B)] Effectuer un backup complet du site
\item[C)] Supprimer l'ancien site
\item[D)] Désactiver tous les plugins
\end{itemize}
*(Réponse : B)*

\item Combien de temps peut prendre la transmission complète du PageRank ?
\begin{itemize}
\item[A)] 1 à 3 jours
\item[B)] 1 à 3 mois
\item[C)] 6 à 12 mois
\item[D)] 24 heures
\end{itemize}
*(Réponse : B)*

\item Quel code HTTP est le remplacement moderne du 301 ?
\begin{itemize}
\item[A)] 302
\item[B)] 303
\item[C)] 307
\item[D)] 308
\end{itemize}
*(Réponse : D)*

\item Que faut-il faire du TTL DNS avant une migration ?
\begin{itemize}
\item[A)] L'augmenter à 86400 secondes
\item[B)] Le réduire à 300 secondes (5 min) au moins 48h avant
\item[C)] Le supprimer complètement
\item[D)] Ne pas y toucher
\end{itemize}
*(Réponse : B)*

\item Quel est le risque d'une chaîne de redirection ?
\begin{itemize}
\item[A)] Amélioration du SEO
\item[B)] Augmentation de la vitesse
\item[C)] Dilution du PageRank et perte d'autorité
\item[D)] Meilleure indexation
\end{itemize}
*(Réponse : C)*

\item Quel outil est recommandé pour le crawl complet avant migration ?
\begin{itemize}
\item[A)] Google Analytics
\item[B)] Screaming Frog SEO Spider
\item[C)] Google Trends
\item[D)] PageSpeed Insights
\end{itemize}
*(Réponse : B)*

\item Que faire des URLs sans équivalent dans la nouvelle structure ?
\begin{itemize}
\item[A)] Les supprimer sans redirection
\item[B)] Les rediriger en 301 vers la page parente la plus proche
\item[C)] Les laisser en 404
\item[D)] Les rediriger toutes vers la home page
\end{itemize}
*(Réponse : B)*

\item Pourquoi éviter de rediriger toutes les anciennes URLs vers la home ?
\begin{itemize}
\item[A)] Cela ralentit le serveur
\item[B)] Google considère cela comme des soft 404
\item[C)] Cela améliore le PageRank de la home
\item[D)] Cela augmente le trafic
\end{itemize}
*(Réponse : B)*

\item Combien de temps dure la propagation DNS complète ?
\begin{itemize}
\item[A)] 1 à 2 heures
\item[B)] 24 à 72 heures
\item[C)] 1 à 2 semaines
\item[D)] Immédiate
\end{itemize}
*(Réponse : B)*

\item Quel rapport GSC est crucial à surveiller après une migration ?
\begin{itemize}
\item[A)] Performance
\item[B)] Couverture (indexation)
\item[C)] Core Web Vitals
\item[D)] Vitesse
\end{itemize}
*(Réponse : B)*

\item Quel est le délai recommandé pour le suivi post-migration ?
\begin{itemize}
\item[A)] 1 semaine
\item[B)] 1 mois
\item[C)] 3 mois
\item[D)] 1 an
\end{itemize}
*(Réponse : C)*

\item Qu'est-ce qu'une chaîne de redirection ?
\begin{itemize}
\item[A)] A redirige vers B, B redirige vers C
\item[B)] A redirige vers B, B redirige vers A
\item[C)] Redirection vers la même page
\item[D)] Absence de redirection
\end{itemize}
*(Réponse : A)*

\item Que signifie une baisse du crawl Googlebot après migration ?
\begin{itemize}
\item[A)] Le site est mieux indexé
\item[B)] Un problème technique empêche Google d'accéder au site
\item[C)] Le contenu est de meilleure qualité
\item[D)] Les backlinks ont augmenté
\end{itemize}
*(Réponse : B)*

\item Quel élément ne doit jamais être bloqué dans robots.txt ?
\begin{itemize}
\item[A)] Les images
\item[B)] Les ressources CSS et JavaScript
\item[C)] Les fichiers PDF
\item[D)] Les archives
\end{itemize}
*(Réponse : B)*

\item À quel moment activer les redirections 301 ?
\begin{itemize}
\item[A)] Une semaine avant le basculement DNS
\item[B)] Dès le basculement DNS
\item[C)] Un mois après le basculement
\item[D)] Avant même de préparer le nouveau site
\end{itemize}
*(Réponse : B)*

\item Qu'est-ce qu'un plan de rollback ?
\begin{itemize}
\item[A)] Un plan pour accélérer la migration
\item[B)] Un plan de retour arrière en cas d'échec
\item[C)] Un plan de référencement
\item[D)] Un plan de sauvegarde des emails
\end{itemize}
*(Réponse : B)*

\item Quel est le code HTTP 307 ?
\begin{itemize}
\item[A)] Redirection permanente
\item[B)] Redirection temporaire avec conservation de la méthode HTTP
\item[C)] Page non trouvée
\item[D)] Erreur serveur
\end{itemize}
*(Réponse : B)*

\item Conséquence d'un noindex accidentel sur le nouveau site ?
\begin{itemize}
\item[A)] Meilleur classement
\item[B)] Les pages ne seront pas indexées par Google
\item[C)] Augmentation du trafic
\item[D)] Amélioration du CTR
\end{itemize}
*(Réponse : B)*

\item Comment informer Google d'un changement d'adresse ?
\begin{itemize}
\item[A)] Par email
\item[B)] Via l'outil "Changement d'adresse" dans GSC
\item[C)] Via un tweet
\item[D)] En modifiant le fichier .htaccess
\end{itemize}
*(Réponse : B)*

\item Impact d'une migration HTTPS sur le CTR ?
\begin{itemize}
\item[A)] Aucun impact
\item[B)] Diminution du CTR
\item[C)] Amélioration du CTR grâce au cadenas de sécurité
\item[D)] Le CTR est divisé par deux
\end{itemize}
*(Réponse : C)*

\item Que doit contenir un inventaire pré-migration complet ?
\begin{itemize}
\item[A)] Uniquement les URLs principales
\item[B)] URLs, backlinks, trafic, positions et conversions
\item[C)] Uniquement les backlinks
\item[D)] Uniquement les fichiers CSS
\end{itemize}
*(Réponse : B)*

\item Différence entre redirection 301 et 308 ?
\begin{itemize}
\item[A)] La 308 conserve la méthode HTTP
\item[B)] La 301 est temporaire
\item[C)] La 308 est plus rapide
\item[D)] Aucune différence
\end{itemize}
*(Réponse : A)*

\item Premier indicateur à surveiller le jour de la migration ?
\begin{itemize}
\item[A)] Les positions Google
\item[B)] Les erreurs 404 et 500 dans les logs serveur
\item[C)] Le taux de conversion
\item[D)] Les backlinks
\end{itemize}
*(Réponse : B)*

\item Qu'est-ce qu'une boucle de redirection ?
\begin{itemize}
\item[A)] A vers B, B vers C
\item[B)] A vers B, B vers A
\item[C)] A vers B en 301
\item[D)] B vers A en 302
\end{itemize}
*(Réponse : B)*

\item Que faire si le trafic n'est pas revenu après 3 mois ?
\begin{itemize}
\item[A)] Attendre encore 3 mois
\item[B)] Effectuer un audit des causes possibles
\item[C)] Revenir à l'ancien site
\item[D)] Tout laisser en l'état
\end{itemize}
*(Réponse : B)*

\item Quel pourcentage de trafic pré-migration à J+7 ?
\begin{itemize}
\item[A)] 50\%
\item[B)] 80-100\%
\item[C)] 30\%
\item[D)] 150\%
\end{itemize}
*(Réponse : B)*

\item Migration qui consiste à changer de CMS ?
\begin{itemize}
\item[A)] Migration de domaine
\item[B)] Migration de CMS
\item[C)] Migration de protocole
\item[D)] Migration de serveur
\end{itemize}
*(Réponse : B)*

\item Comment Google interprète-t-il les 302 répétés ?
\begin{itemize}
\item[A)] Comme des 301
\item[B)] Comme des 404
\item[C)] Comme des 500
\item[D)] Comme des 200
\end{itemize}
*(Réponse : A)*

\item Pourquoi conserver l'ancien site intact après migration ?
\begin{itemize}
\item[A)] Pour le backup et le rollback
\item[B)] Pour continuer à recevoir du trafic
\item[C)] Pour des raisons fiscales
\item[D)] Pour héberger des fichiers inutiles
\end{itemize}
*(Réponse : A)*

\item Outil pour vérifier la propagation DNS ?
\begin{itemize}
\item[A)] Google Analytics
\item[B)] whatsmydns.net
\item[C)] PageSpeed Insights
\item[D)] GTmetrix
\end{itemize}
*(Réponse : B)*

\item Quelle métrique indique une baisse du CTR après migration ?
\begin{itemize}
\item[A)] Titres et descriptions mal migrés
\item[B)] Augmentation du budget de crawl
\item[C)] Amélioration du LCP
\item[D)] Augmentation des backlinks
\end{itemize}
*(Réponse : A)*

\item Code HTTP qui conserve la méthode HTTP ?
\begin{itemize}
\item[A)] 301
\item[B)] 302
\item[C)] 307
\item[D)] 308
\end{itemize}
*(Réponse : D)*

\item Que faire des backlinks avant une migration ?
\begin{itemize}
\item[A)] Les supprimer tous
\item[B)] Les exporter et documenter pour vérifier la transmission
\item[C)] Les ignorer
\item[D)] Les recopier manuellement
\end{itemize}
*(Réponse : B)*

\item Risque d'une migration sans test mobile ?
\begin{itemize}
\item[A)] Le site sera plus rapide
\item[B)] Le nouveau site pourrait ne pas être responsive
\item[C)] Le site sera mieux classé
\item[D)] Aucun risque
\end{itemize}
*(Réponse : B)*

\item Quand réduire le TTL DNS avant migration ?
\begin{itemize}
\item[A)] 1 heure avant
\item[B)] Au moins 48 heures avant
\item[C)] 1 semaine avant
\item[D)] 1 mois avant
\end{itemize}
*(Réponse : B)*

\item Première action après le basculement DNS le jour J ?
\begin{itemize}
\item[A)] Surveiller les logs serveur en temps réel
\item[B)] Partir en weekend
\item[C)] Supprimer l'ancien site
\item[D)] Désactiver les redirections
\end{itemize}
*(Réponse : A)*

\item Code HTTP pour contenu saisonnier indisponible ?
\begin{itemize}
\item[A)] 301
\item[B)] 302
\item[C)] 404
\item[D)] 500
\end{itemize}
*(Réponse : B)*

\item Contenu du fichier de mapping URL complet ?
\begin{itemize}
\item[A)] Ancienne URL et nouvelle URL
\item[B)] Ancienne URL, nouvelle URL, type, statut
\item[C)] Uniquement les nouvelles URLs
\item[D)] Les adresses IP des serveurs
\end{itemize}
*(Réponse : B)*

\item Erreur la plus fréquente en migration SEO ?
\begin{itemize}
\item[A)] Trop de backlinks
\item[B)] Absence de mapping URL et redirections génériques
\item[C)] Trop de contenu
\item[D)] Vitesse excessive
\end{itemize}
*(Réponse : B)*

\item Comment protéger le staging de l'indexation ?
\begin{itemize}
\item[A)] Mot de passe ou restriction IP
\item[B)] Accès libre à tous
\item[C)] Supprimer les fichiers
\item[D)] Augmenter le TTL
\end{itemize}
*(Réponse : A)*

\item Effet d'un canonical incorrect vers l'ancienne version ?
\begin{itemize}
\item[A)] Amélioration du classement
\item[B)] Google peut ignorer le nouveau site
\item[C)] Augmentation du trafic
\item[D)] Aucun impact
\end{itemize}
*(Réponse : B)*

\item Période déconseillée pour effectuer une migration ?
\begin{itemize}
\item[A)] Période de faible trafic
\item[B)] Noël, Black Friday, soldes
\item[C)] Le weekend
\item[D)] En début de mois
\end{itemize}
*(Réponse : B)*

\item Temps de propagation DNS complet ?
\begin{itemize}
\item[A)] 1 heure maximum
\item[B)] 24 à 72 heures
\item[C)] 5 minutes
\item[D)] 1 mois
\end{itemize}
*(Réponse : B)*

\item Outil Google pour demander l'indexation post-migration ?
\begin{itemize}
\item[A)] Google Trends
\item[B)] Inspection d'URL dans GSC
\item[C)] Google Ads
\item[D)] Google Tag Manager
\end{itemize}
*(Réponse : B)*

\item Que faire si des 404 apparaissent après migration ?
\begin{itemize}
\item[A)] Les ignorer
\item[B)] Créer les redirections 301 manquantes
\item[C)] Les transformer en 500
\item[D)] Supprimer les pages concernées
\end{itemize}
*(Réponse : B)*
\end{enumerate}

% ======================================================================
% CHAPITRE 26 : LE SEO POUR L'ECOMMERCE
% ======================================================================
\section*{Chapitre 26 : Le SEO pour l'E-commerce}

\begin{enumerate}
\item Principal défi SEO des sites e-commerce ?
\begin{itemize}
\item[A)] Le manque de contenu
\item[B)] Le contenu dupliqué et la grande quantité de pages
\item[C)] L'absence de backlinks
\item[D)] La lenteur des transactions
\end{itemize}
*(Réponse : B)*

\item Profondeur de clics recommandée pour pages importantes ?
\begin{itemize}
\item[A)] 5 à 6 clics
\item[B)] 1 à 3 clics depuis la racine
\item[C)] Plus de 10 clics
\item[D)] 7 à 8 clics
\end{itemize}
*(Réponse : B)*

\item Comment gérer les catégories sans contenu ?
\begin{itemize}
\item[A)] Les supprimer
\item[B)] Ajouter des descriptions uniques
\item[C)] Les masquer avec noindex
\item[D)] Les rediriger vers la home
\end{itemize}
*(Réponse : B)*

\item Problème SEO des filtres et facettes ?
\begin{itemize}
\item[A)] Ils améliorent le référencement
\item[B)] URLs infinies et contenu dupliqué
\item[C)] Ils augmentent le PageRank
\item[D)] Ils sont ignorés
\end{itemize}
*(Réponse : B)*

\item Produits hors stock en SEO ?
\begin{itemize}
\item[A)] Supprimer les pages
\item[B)] Garder indexés avec mention de rupture
\item[C)] Rediriger en 302
\item[D)] Noindex définitif
\end{itemize}
*(Réponse : B)*

\item Meilleure pratique pour descriptions produits ?
\begin{itemize}
\item[A)] Utiliser les descriptions fournisseur
\item[B)] Rédiger des descriptions uniques
\item[C)] Ne pas mettre de description
\item[D)] Copier les concurrents
\end{itemize}
*(Réponse : B)*

\item Comment optimiser la pagination e-commerce ?
\begin{itemize}
\item[A)] Tout mettre en noindex
\item[B)] Utiliser rel next/prev ou canonical
\item[C)] Supprimer la pagination
\item[D)] Mettre en nofollow
\end{itemize}
*(Réponse : B)*

\item Architecture URL recommandée pour e-commerce ?
\begin{itemize}
\item[A)] exemple.com/p?id=123
\item[B)] exemple.com/categorie/sous-cat/produit
\item[C)] exemple.com/page123
\item[D)] exemple.com/produit.php?id=
\end{itemize}
*(Réponse : B)*

\item Variations produits (tailles, couleurs) en SEO ?
\begin{itemize}
\item[A)] URL par variation
\item[B)] Canonical sur la page principale
\item[C)] Tout noindex
\item[D)] Ignorer
\end{itemize}
*(Réponse : B)*

\item Données structurées essentielles e-commerce ?
\begin{itemize}
\item[A)] Event
\item[B)] Product + AggregateRating + Offer
\item[C)] Recipe
\item[D)] FAQPage
\end{itemize}
*(Réponse : B)*

\item Avis clients en SEO e-commerce ?
\begin{itemize}
\item[A)] Les ignorer
\item[B)] Review schema pour les étoiles SERP
\item[C)] Les supprimer
\item[D)] Les cacher
\end{itemize}
*(Réponse : B)*

\item Structure de site idéale ?
\begin{itemize}
\item[A)] Plate (tout à la racine)
\item[B)] Hiérarchique catégories/sous-catégories/produits
\item[C)] Aléatoire
\item[D)] Blog uniquement
\end{itemize}
*(Réponse : B)*

\item Impact du contenu dupliqué ?
\begin{itemize}
\item[A)] Aucun impact
\item[B)] Pénalité ou mauvaise page classée
\item[C)] Amélioration du classement
\item[D)] Augmentation du trafic
\end{itemize}
*(Réponse : B)*

\item Optimisation images produits ?
\begin{itemize}
\item[A)] Basse résolution
\item[B)] Alt text + noms + WebP
\item[C)] Pas d'images
\item[D)] Images d'autres sites
\end{itemize}
*(Réponse : B)*

\item Problème sessions PHP dans URLs ?
\begin{itemize}
\item[A)] Aucun problème
\item[B)] Création d'URLs dupliquées
\item[C)] Améliore le SEO
\item[D)] Augmente la vitesse
\end{itemize}
*(Réponse : B)*

\item Optimisation pages catégories ?
\begin{itemize}
\item[A)] Images seulement
\item[B)] Texte unique + title + meta description
\item[C)] Laisser vides
\item[D)] Copier les fournisseurs
\end{itemize}
*(Réponse : B)*

\item Avantage d'un blog e-commerce ?
\begin{itemize}
\item[A)] Aucun
\item[B)] Trafic haut de funnel
\item[C)] Remplace les fiches produits
\item[D)] Ralentit le site
\end{itemize}
*(Réponse : B)*

\item Pages de marques ?
\begin{itemize}
\item[A)] Supprimer
\item[B)] Pages optimisées contenu unique
\item[C)] Rediriger vers catégories
\item[D)] Noindex
\end{itemize}
*(Réponse : B)*

\item Maillage interne e-commerce ?
\begin{itemize}
\item[A)] Aucun lien
\item[B)] Liens produits complémentaires
\item[C)] Tout vers la home
\item[D)] Footer seulement
\end{itemize}
*(Réponse : B)*

\item Format image recommandé photos produits ?
\begin{itemize}
\item[A)] BMP
\item[B)] WebP ou AVIF
\item[C)] TIFF
\item[D)] GIF
\end{itemize}
*(Réponse : B)*

\item Contenu dupliqué fiches similaires ?
\begin{itemize}
\item[A)] Supprimer
\item[B)] Descriptions uniques + canonique
\item[C)] Fusionner
\item[D)] Noindex
\end{itemize}
*(Réponse : B)*

\item URLs produits recommandées ?
\begin{itemize}
\item[A)] Numériques
\item[B)] Descriptives
\item[C)] Aléatoires
\item[D)] ID seul
\end{itemize}
*(Réponse : B)*

\item Catégories vides ?
\begin{itemize}
\item[A)] Garder en l'état
\item[B)] Contenu unique ou redirection
\item[C)] Nofollow
\item[D)] Dupliquer
\end{itemize}
*(Réponse : B)*

\item Impact schema Product ?
\begin{itemize}
\item[A)] Aucun
\item[B)] Rich snippets prix/disponibilité
\item[C)] Ralentit le site
\item[D)] Remplace meta desc
\end{itemize}
*(Réponse : B)*

\item Liens catégories optimaux ?
\begin{itemize}
\item[A)] Vers la home
\item[B)] Hiérarchiques + breadcrumbs
\item[C)] Pas de liens
\item[D)] Aléatoires
\end{itemize}
*(Réponse : B)*

\item Produits en fin de vie ?
\begin{itemize}
\item[A)] Laisser en l'état
\item[B)] 301 vers le remplacement
\item[C)] Supprimer sans redirection
\item[D)] 404
\end{itemize}
*(Réponse : B)*

\item Landing pages saisonnières ?
\begin{itemize}
\item[A)] Supprimer hors saison
\item[B)] Mise à jour régulière
\item[C)] Noindex hors saison
\item[D)] Rediriger
\end{itemize}
*(Réponse : B)*

\item Profondeur max recommandée ?
\begin{itemize}
\item[A)] 5 clics
\item[B)] 3 clics depuis l'accueil
\item[C)] 10 clics
\item[D)] Illimitée
\end{itemize}
*(Réponse : B)*

\item Breadcrumbs en e-commerce ?
\begin{itemize}
\item[A)] Ignorer
\item[B)] BreadcrumbList schema
\item[C)] Cacher
\item[D)] Nofollow
\end{itemize}
*(Réponse : B)*

\item Impact temps de chargement ?
\begin{itemize}
\item[A)] Aucun
\item[B)] Perte trafic + déclassement
\item[C)] Amélioration
\item[D)] Mobile seulement
\end{itemize}
*(Réponse : B)*

\item Filtres créant des milliers d'URLs ?
\begin{itemize}
\item[A)] Tout indexer
\item[B)] Paramètres GSC + noindex
\item[C)] Supprimer les filtres
\item[D)] Nofollow
\end{itemize}
*(Réponse : B)*

\item Recherches longue traîne en e-commerce ?
\begin{itemize}
\item[A)] Nulles
\item[B)] Forte intention d'achat
\item[C)] Faibles
\item[D)] Pénalisées
\end{itemize}
*(Réponse : B)*

\item Page d'accueil optimale ?
\begin{itemize}
\item[A)] Images seulement
\item[B)] Title + meta desc + catégories + contenu
\item[C)] Pas de texte
\item[D)] Redirection
\end{itemize}
*(Réponse : B)*

\item Schema pour avis produits ?
\begin{itemize}
\item[A)] FAQPage
\item[B)] AggregateRating + Review
\item[C)] Event
\item[D)] HowTo
\end{itemize}
*(Réponse : B)*

\item Pages de recherche interne ?
\begin{itemize}
\item[A)] Indexer
\item[B)] Bloquer robots.txt/noindex
\item[C)] Rendre publiques
\item[D)] Supprimer
\end{itemize}
*(Réponse : B)*

\item Catégories parentes ?
\begin{itemize}
\item[A)] Vides
\item[B)] Contenu éditorial + liens
\item[C)] Rediriger
\item[D)] Supprimer
\end{itemize}
*(Réponse : B)*

\item Tags en e-commerce ?
\begin{itemize}
\item[A)] Indexer
\item[B)] Noindex
\item[C)] Supprimer
\item[D)] En catégories
\end{itemize}
*(Réponse : B)*

\item Impact délais livraison SEO ?
\begin{itemize}
\item[A)] Aucun
\item[B)] Indirect via avis et taux conversion
\item[C)] Direct classement
\item[D)] Pénalité
\end{itemize}
*(Réponse : B)*

\item Pages panier/checkout ?
\begin{itemize}
\item[A)] Indexer
\item[B)] Noindex (entonnoir conversion)
\item[C)] Mettre en avant
\item[D)] Supprimer
\end{itemize}
*(Réponse : B)*

\item Taille catalogue idéale ?
\begin{itemize}
\item[A)] Moins de 100 produits
\item[B)] Pas de taille, architecture clé
\item[C)] Plus de 10000
\item[D)] 500 exactement
\end{itemize}
*(Réponse : B)*

\item Sous-catégories ?
\begin{itemize}
\item[A)] Supprimer
\item[B)] Hiérarchie logique URLs descr.
\item[C)] Plates
\item[D)] Sous-domaines
\end{itemize}
*(Réponse : B)*

\item Effet dropshipping SEO ?
\begin{itemize}
\item[A)] Aucun
\item[B)] Risque contenu dupliqué élevé
\item[C)] Amélioration
\item[D)] Pénalité auto
\end{itemize}
*(Réponse : B)*

\item Variantes de prix multiples ?
\begin{itemize}
\item[A)] Autant de pages
\item[B)] Schema Product options prix
\item[C)] Masquer les prix
\item[D)] Prix minimum
\end{itemize}
*(Réponse : B)*

\item Pages FAQ e-commerce ?
\begin{itemize}
\item[A)] Ignorer
\item[B)] FAQPage schema
\item[C)] Supprimer
\item[D)] Noindex
\end{itemize}
*(Réponse : B)*

\item Recherche produits pour Google ?
\begin{itemize}
\item[A)] Accès bots
\item[B)] Noindex + optimisation
\item[C)] Supprimer
\item[D)] Bloquer
\end{itemize}
*(Réponse : B)*

\item Rôle URL images SEO e-commerce ?
\begin{itemize}
\item[A)] Aucun
\item[B)] Noms descriptifs = Google Images
\item[C)] Ignorés
\item[D)] Pénalisent
\end{itemize}
*(Réponse : B)*

\item Ventes croisées SEO ?
\begin{itemize}
\item[A)] Supprimer
\item[B)] Liens internes produits
\item[C)] Cacher
\item[D)] Nofollow
\end{itemize}
*(Réponse : B)*

\item Métrique prioritaire e-commerce SEO ?
\begin{itemize}
\item[A)] Nombre de pages
\item[B)] Trafic organique + taux conversion
\item[C)] Taille panier
\item[D)] Commandes téléphone
\end{itemize}
*(Réponse : B)*

\item Meilleur KPI e-commerce SEO ?
\begin{itemize}
\item[A)] Pages vues
\item[B)] Revenus organiques
\item[C)] Taux rebond
\item[D)] Articles blog
\end{itemize}
*(Réponse : B)*

    \item Quel format de données structurées est essentiel pour un site e-commerce ?
    \begin{itemize}
        \item[A)] Event schema
        \item[B)] Product schema
        \item[C)] VideoObject schema
        \item[D)] Person schema
    \end{itemize}
    *(Réponse : B)*
\end{enumerate}

% ======================================================================
% CHAPITRE 27 : SEO POUR LES CMS POPULAIRES
% ======================================================================
\section*{Chapitre 27 : SEO pour les CMS populaires (WordPress, Shopify, Webflow)}

\begin{enumerate}
\item CMS le plus utilisé au monde ?
\begin{itemize}
\item[A)] Shopify
\item[B)] WordPress
\item[C)] Webflow
\item[D)] Wix
\end{itemize}
*(Réponse : B)*

\item Structure permaliens WordPress recommandée ?
\begin{itemize}
\item[A)] /?p=123
\item[B)] /\%postname\%/
\item[C)] /category/id/
\item[D)] /2024/05/nom/
\end{itemize}
*(Réponse : B)*

\item Plugin SEO le plus répandu WordPress ?
\begin{itemize}
\item[A)] Rank Math
\item[B)] Yoast SEO
\item[C)] SEOPress
\item[D)] AIOSEO
\end{itemize}
*(Réponse : B)*

\item Force SEO principale Shopify ?
\begin{itemize}
\item[A)] robots.txt illimité
\item[B)] Titre/meta auto + CDN + HTTPS
\item[C)] Plugins gratuits
\item[D)] URL rewriting
\end{itemize}
*(Réponse : B)*

\item Faiblesse SEO Shopify pour blogs ?
\begin{itemize}
\item[A)] Pas de blog
\item[B)] URL rigide /blogs/ + limitations
\item[C)] Blog non indexable
\item[D)] Pas de contenu
\end{itemize}
*(Réponse : B)*

\item Alternative française respectueuse données ?
\begin{itemize}
\item[A)] Yoast
\item[B)] SEOPress
\item[C)] Rank Math
\item[D)] AIOSEO
\end{itemize}
*(Réponse : B)*

\item Optimisation images WordPress ?
\begin{itemize}
\item[A)] Pas d'images
\item[B)] Légendes + alt text + WebP auto
\item[C)] HD non compressées
\item[D)] Ignorer
\end{itemize}
*(Réponse : B)*

\item Avantage Webflow SEO ?
\begin{itemize}
\item[A)] Gratuit
\item[B)] CMS visuel + URLs custom + sitemap + perf
\item[C)] 1000 plugins
\item[D)] Accès serveur
\end{itemize}
*(Réponse : B)*

\item Limite Webflow pour redirections ?
\begin{itemize}
\item[A)] Pas possible
\item[B)] 100-500 selon plan
\item[C)] Illimité
\item[D)] 302 seulement
\end{itemize}
*(Réponse : B)*

\item Approche recommandée tags WordPress ?
\begin{itemize}
\item[A)] Tout indexer
\item[B)] Noindex
\item[C)] Supprimer
\item[D)] En catégories
\end{itemize}
*(Réponse : B)*

\item Plugin Schema gratuit WordPress ?
\begin{itemize}
\item[A)] Yoast gratuit
\item[B)] Rank Math
\item[C)] SEOPress gratuit
\item[D)] The SEO Framework
\end{itemize}
*(Réponse : B)*

\item Génération sitemap Shopify ?
\begin{itemize}
\item[A)] Manuelle
\item[B)] Automatique
\item[C)] Plugin payant
\item[D)] Pas de sitemap
\end{itemize}
*(Réponse : B)*

\item Limite robots.txt Shopify ?
\begin{itemize}
\item[A)] Pas de fichier
\item[B)] Personnalisation limitée
\item[C)] Complet
\item[D)] Bloqué
\end{itemize}
*(Réponse : B)*

\item Collections CMS Webflow SEO ?
\begin{itemize}
\item[A)] Pas de collections
\item[B)] Pages dynamiques URLs customisables
\item[C)] Non indexables
\item[D)] Identiques
\end{itemize}
*(Réponse : B)*

\item Meilleure flexibilité SEO ?
\begin{itemize}
\item[A)] Shopify
\item[B)] WordPress
\item[C)] Wix
\item[D)] Squarespace
\end{itemize}
*(Réponse : B)*

\item Hébergement recommandé WordPress SEO ?
\begin{itemize}
\item[A)] Mutualisé
\item[B)] WP Engine, Kinsta, SiteGround
\item[C)] Gratuit
\item[D)] Serveur perso
\end{itemize}
*(Réponse : B)*

\item Données structurées produits Shopify ?
\begin{itemize}
\item[A)] Pas de données
\item[B)] Génération auto Product schema
\item[C)] Manuel
\item[D)] Plugin
\end{itemize}
*(Réponse : B)*

\item Wix SEO Wiz ?
\begin{itemize}
\item[A)] Inutile
\item[B)] Assistant SEO guidé
\item[C)] Payant
\item[D)] Suppression contenu
\end{itemize}
*(Réponse : B)*

\item Performances WordPress SEO ?
\begin{itemize}
\item[A)] 50 plugins
\item[B)] Redis + CDN + compression + cache
\item[C)] 0 plugin
\item[D)] Thème lourd
\end{itemize}
*(Réponse : B)*

\item Catégories WordPress recommandation ?
\begin{itemize}
\item[A)] Supprimer
\item[B)] Optimiser + noindex tags
\item[C)] Tags oui catégories non
\item[D)] Tout noindex
\end{itemize}
*(Réponse : B)*

\item Inconvénient Shopify SEO avancé ?
\begin{itemize}
\item[A)] Pas SSL
\item[B)] robots.txt et URLs limités
\item[C)] Site inaccessible
\item[D)] Pas de contenu
\end{itemize}
*(Réponse : B)*

\item Plugin WordPress léger sans pub ?
\begin{itemize}
\item[A)] Yoast
\item[B)] The SEO Framework
\item[C)] Rank Math
\item[D)] AIOSEO
\end{itemize}
*(Réponse : B)*

\item Noindex dans Webflow ?
\begin{itemize}
\item[A)] Pas possible
\item[B)] Disponible par page
\item[C)] Par collection
\item[D)] Payant
\end{itemize}
*(Réponse : B)*

\item Pagination CMS sur mesure ?
\begin{itemize}
\item[A)] Sans balises
\item[B)] rel next/prev
\item[C)] Pas de pagination
\item[D)] Infinie
\end{itemize}
*(Réponse : B)*

\item Contenu dupliqué WordPress ?
\begin{itemize}
\item[A)] Ignorer
\item[B)] Canonique par plugin SEO
\item[C)] Supprimer
\item[D)] Noindex
\end{itemize}
*(Réponse : B)*

\item Différence Shopify vs WordPress blogging ?
\begin{itemize}
\item[A)] Aucune
\item[B)] WordPress plus flexible
\item[C)] Shopify meilleur
\item[D)] WP sans blog
\end{itemize}
*(Réponse : B)*

\item Sitemap Webflow ?
\begin{itemize}
\item[A)] Manuel
\item[B)] Auto avec exclusion possible
\item[C)] Pas de sitemap
\item[D)] Plugin
\end{itemize}
*(Réponse : B)*

\item Wix pour débutant SEO ?
\begin{itemize}
\item[A)] Aucun
\item[B)] Wix SEO Wiz guide pas à pas
\item[C)] Meilleur CMS
\item[D)] Gratuit
\end{itemize}
*(Réponse : B)*

\item Breadcrumbs dans CMS ?
\begin{itemize}
\item[A)] Ignorer
\item[B)] BreadcrumbList schema
\item[C)] Cacher
\item[D)] Nofollow
\end{itemize}
*(Réponse : B)*

\item Permaliens CMS sur mesure ?
\begin{itemize}
\item[A)] IDs numériques
\item[B)] Slugs custom sans ID
\item[C)] Aléatoires
\item[D)] Paramètres GET
\end{itemize}
*(Réponse : B)*

\item Cache optimisé pour WordPress ?
\begin{itemize}
\item[A)] Varnish
\item[B)] Cloudflare APO
\item[C)] Redis seul
\item[D)] Memcached
\end{itemize}
*(Réponse : B)*

\item Multilingue WordPress SEO ?
\begin{itemize}
\item[A)] Impossible
\item[B)] WPML/Polylang + hreflang
\item[C)] Traduction manuelle
\item[D)] Google Translate
\end{itemize}
*(Réponse : B)*

\item CMS sans accès serveur ?
\begin{itemize}
\item[A)] WordPress
\item[B)] Shopify et Wix (SaaS)
\item[C)] Webflow
\item[D)] Drupal
\end{itemize}
*(Réponse : B)*

\item Performances Shopify SEO ?
\begin{itemize}
\item[A)] Plugins cache
\item[B)] Images optimisées + apps limitées + thème léger
\item[C)] Accès serveur
\item[D)] .htaccess
\end{itemize}
*(Réponse : B)*

\item Force Shopify dropshipping SEO ?
\begin{itemize}
\item[A)] Aucune
\item[B)] Apps natives + sitemap auto
\item[C)] Meilleur que WooCommerce
\item[D)] Parfait
\end{itemize}
*(Réponse : B)*

\item Sitemap dynamique CMS sur mesure ?
\begin{itemize}
\item[A)] Manuel
\item[B)] XML auto à chaque publication
\item[C)] Pas de sitemap
\item[D)] Statique
\end{itemize}
*(Réponse : B)*

\item CMS le plus performant ?
\begin{itemize}
\item[A)] WordPress
\item[B)] Dépend de la config, Webflow bon
\item[C)] Shopify
\item[D)] Wix
\end{itemize}
*(Réponse : B)*

\item Redirections 301 Webflow ?
\begin{itemize}
\item[A)] Pas possible
\item[B)] Configurables dans paramètres site
\item[C)] Plugin
\item[D)] .htaccess
\end{itemize}
*(Réponse : B)*

\item Cache WordPress recommandé ?
\begin{itemize}
\item[A)] Aucun
\item[B)] WP Rocket / W3 Total / LiteSpeed
\item[C)] Yoast Cache
\item[D)] Rank Math Cache
\end{itemize}
*(Réponse : B)*

\item Images produits Shopify SEO ?
\begin{itemize}
\item[A)] Pas d'images
\item[B)] Alt text + noms + compression apps
\item[C)] Automatique
\item[D)] Non optimisables
\end{itemize}
*(Réponse : B)*

\item Limite Wix SEO avancé ?
\begin{itemize}
\item[A)] Aucune
\item[B)] Pas d'accès serveur/plugins/internationalisation
\item[C)] Meilleur que WP
\item[D)] Tout possible
\end{itemize}
*(Réponse : B)*

\item Balises canoniques WordPress ?
\begin{itemize}
\item[A)] Manuel
\item[B)] Auto par plugin SEO
\item[C)] Pas de canonique
\item[D)] .htaccess
\end{itemize}
*(Réponse : B)*

\item URL impossible à modifier Shopify ?
\begin{itemize}
\item[A)] Produits
\item[B)] Structure /blogs/
\item[C)] Pages
\item[D)] Collections
\end{itemize}
*(Réponse : B)*

\item Sécurité WordPress SEO ?
\begin{itemize}
\item[A)] Inutile
\item[B)] MAJ régulières + plugins limités + pare-feu
\item[C)] Désactiver MAJ
\item[D)] 100 plugins
\end{itemize}
*(Réponse : B)*

\item Avantage SEOPress français ?
\begin{itemize}
\item[A)] Moins performant
\item[B)] Respect données + sans pub + complet
\item[C)] Anglais seulement
\item[D)] Payant
\end{itemize}
*(Réponse : B)*

\item Contenu dupliqué Shopify ?
\begin{itemize}
\item[A)] Impossible
\item[B)] Canoniques auto + descriptions uniques
\item[C)] Supprimer
\item[D)] Ignorer
\end{itemize}
*(Réponse : B)*

\item Nettoyage base WordPress ?
\begin{itemize}
\item[A)] Quotidien
\item[B)] Régulier (révisions, spams, transients)
\item[C)] Jamais
\item[D)] Annuel
\end{itemize}
*(Réponse : B)*

\item Contrôle données structurées max ?
\begin{itemize}
\item[A)] Shopify
\item[B)] WordPress + Rank Math/Yoast
\item[C)] Wix
\item[D)] Squarespace
\end{itemize}
*(Réponse : B)*

\item Paramètres URL Wix SEO ?
\begin{itemize}
\item[A)] Non modifiables
\item[B)] Réécriture disponible
\item[C)] Fixes
\item[D)] Pas d'URL
\end{itemize}
*(Réponse : B)*

\item Hébergement optimisation auto SEO ?
\begin{itemize}
\item[A)] Partagé basique
\item[B)] WP Engine/Kinsta + Redis + CDN
\item[C)] Gratuit
\item[D)] Serveur non configuré
\end{itemize}
*(Réponse : B)*
\end{enumerate}
\section*{Chapitre 28 : Le SEO pour la Vidéo et le Contenu Multimédia}
\begin{enumerate}
    \item Quel pourcentage du trafic internet la vidéo représente-t-elle aujourd'hui ?
    \begin{itemize}
        \item[A)] Plus de 50\%
        \item[B)] Plus de 65\%
        \item[C)] Plus de 82\%
        \item[D)] Plus de 90\%
    \end{itemize}
    *(Réponse : C)*

    \item Quel est le deuxième moteur de recherche mondial ?
    \begin{itemize}
        \item[A)] Bing
        \item[B)] YouTube
        \item[C)] Yahoo
        \item[D)] DuckDuckGo
    \end{itemize}
    *(Réponse : B)*

    \item Où doit-on idéalement inclure le mot-clé principal dans une vidéo YouTube ?
    \begin{itemize}
        \item[A)] Dans les commentaires
        \item[B)] Au début du titre
        \item[C)] À la fin de la description
        \item[D)] Uniquement dans les tags
    \end{itemize}
    *(Réponse : B)*

    \item Quelle longueur de description est recommandée pour une vidéo YouTube ?
    \begin{itemize}
        \item[A)] 50-100 mots
        \item[B)] 100-150 mots
        \item[C)] 200-300 mots
        \item[D)] 500-600 mots
    \end{itemize}
    *(Réponse : C)*

    \item Où doit apparaître le mot-clé dans la description YouTube ?
    \begin{itemize}
        \item[A)] Dans les 50 premiers caractères
        \item[B)] Dans les 150 premiers caractères
        \item[C)] À la fin de la description
        \item[D)] Uniquement dans les liens
    \end{itemize}
    *(Réponse : B)*

    \item Quel élément vidéo améliore les apparitions dans les featured snippets ?
    \begin{itemize}
        \item[A)] Les tags
        \item[B)] Les playlists
        \item[C)] Les chapitres vidéo
        \item[D)] Le nombre de likes
    \end{itemize}
    *(Réponse : C)*

    \item Quel format de données structurées est utilisé pour les vidéos ?
    \begin{itemize}
        \item[A)] Product schema
        \item[B)] VideoObject schema
        \item[C)] Article schema
        \item[D)] FAQ schema
    \end{itemize}
    *(Réponse : B)*

    \item Quel élément visuel est crucial pour le taux de clic d'une vidéo ?
    \begin{itemize}
        \item[A)] La durée de la vidéo
        \item[B)] Le thumbnail personnalisé
        \item[C)] Le nombre de vues
        \item[D)] La date de publication
    \end{itemize}
    *(Réponse : B)*

    \item Comment s'appelle le texte complet de la vidéo utilisé pour l'indexation ?
    \begin{itemize}
        \item[A)] Le script
        \item[B)] La transcription
        \item[C)] Le sous-titrage
        \item[D)] Le résumé
    \end{itemize}
    *(Réponse : B)*

    \item Quel est l'avantage des playlists YouTube pour le SEO ?
    \begin{itemize}
        \item[A)] Elles augmentent le nombre de likes
        \item[B)] Elles permettent une organisation thématique
        \item[C)] Elles réduisent le temps de chargement
        \item[D)] Elles suppriment les publicités
    \end{itemize}
    *(Réponse : B)*

    \item Quel type de contenu multimédia nécessite une optimisation spécifique pour le référencement ?
    \begin{itemize}
        \item[A)] Uniquement les images
        \item[B)] Uniquement le texte
        \item[C)] La vidéo et le contenu multimédia
        \item[D)] Uniquement les fichiers PDF
    \end{itemize}
    *(Réponse : C)*

    \item Les données structurées vidéo permettent d'obtenir :
    \begin{itemize}
        \item[A)] Un meilleur PageRank
        \item[B)] Des rich snippets
        \item[C)] Plus de backlinks
        \item[D)] Une vitesse de chargement améliorée
    \end{itemize}
    *(Réponse : B)*

    \item Quelle est la première étape pour optimiser une vidéo sur YouTube ?
    \begin{itemize}
        \item[A)] Choisir une musique de fond
        \item[B)] Optimiser le titre avec le mot-clé principal
        \item[C)] Ajouter des annotations
        \item[D)] Paramétrer la confidentialité
    \end{itemize}
    *(Réponse : B)*

    \item Les chapitres vidéo YouTube apparaissent sous quelle forme dans les résultats de recherche ?
    \begin{itemize}
        \item[A)] Encadré promotionnel
        \item[B)] Lien sponsorisé
        \item[C)] Featured snippets
        \item[D)] Résultat enrichi
    \end{itemize}
    *(Réponse : C)*

    \item Quel pourcentage des recherches sur YouTube sont effectuées sur mobile ?
    \begin{itemize}
        \item[A)] Environ 40\%
        \item[B)] Environ 50\%
        \item[C)] Plus de 70\%
        \item[D)] Environ 30\%
    \end{itemize}
    *(Réponse : C)*

    \item Quelle est la recommandation concernant la miniature d'une vidéo ?
    \begin{itemize}
        \item[A)] Utiliser une image générée automatiquement
        \item[B)] Créer une image personnalisée avec texte
        \item[C)] Ne pas utiliser de miniature
        \item[D)] Utiliser une image floue
    \end{itemize}
    *(Réponse : B)*

    \item Le référencement vidéo permet d'apparaître dans quel type de résultat Google ?
    \begin{itemize}
        \item[A)] Google Shopping
        \item[B)] Google Images
        \item[C)] Google Vidéo et SERP enrichies
        \item[D)] Google News
    \end{itemize}
    *(Réponse : C)*

    \item Quelle information est essentielle dans le fichier de transcription vidéo ?
    \begin{itemize}
        \item[A)] Les noms des acteurs
        \item[B)] Le texte complet de la vidéo
        \item[C)] Le budget de production
        \item[D)] Les remerciements
    \end{itemize}
    *(Réponse : B)*

    \item Google utilise la transcription vidéo pour :
    \begin{itemize}
        \item[A)] Générer des sous-titres automatiques
        \item[B)] Indexer le contenu de la vidéo
        \item[C)] Créer des playlists
        \item[D)] Analyser les visages
    \end{itemize}
    *(Réponse : B)*

    \item Quel type de données structurées permet d'enrichir l'affichage des vidéos dans les SERP ?
    \begin{itemize}
        \item[A)] Event schema
        \item[B)] VideoObject schema
        \item[C)] Recipe schema
        \item[D)] Person schema
    \end{itemize}
    *(Réponse : B)*

    \item À quoi servent les playlists dans une stratégie SEO vidéo ?
    \begin{itemize}
        \item[A)] À augmenter le nombre d'abonnés
        \item[B)] À organiser thématiquement les vidéos
        \item[C)] À monétiser les vidéos
        \item[D)] À supprimer les commentaires négatifs
    \end{itemize}
    *(Réponse : B)*

    \item Le thumbnail personnalisé doit contenir :
    \begin{itemize}
        \item[A)] Uniquement du texte
        \item[B)] Uniquement une image
        \item[C)] Une image avec du texte explicite
        \item[D)] Uniquement un logo
    \end{itemize}
    *(Réponse : C)*

    \item Quelle est la taille recommandée pour le titre d'une vidéo SEO ?
    \begin{itemize}
        \item[A)] Maximum 30 caractères
        \item[B)] Maximum 60 caractères
        \item[C)] Maximum 100 caractères
        \item[D)] Maximum 200 caractères
    \end{itemize}
    *(Réponse : C)*

    \item Les tags YouTube permettent :
    \begin{itemize}
        \item[A)] D'augmenter la vitesse de chargement
        \item[B)] De préciser le sujet de la vidéo aux moteurs de recherche
        \item[C)] De modifier la miniature
        \item[D)] De bloquer certains utilisateurs
    \end{itemize}
    *(Réponse : B)*

    \item Quel élément améliore le SEO des vidéos intégrées dans un site web ?
    \begin{itemize}
        \item[A)] L'autoplay
        \item[B)] Le lazy loading
        \item[C)] Les données structurées VideoObject
        \item[D)] La boucle de lecture
    \end{itemize}
    *(Réponse : C)*

    \item La première étape de l'optimisation d'une chaîne YouTube est :
    \begin{itemize}
        \item[A)] Publier chaque jour
        \item[B)] Définir des mots-clés cibles
        \item[C)] Acheter des vues
        \item[D)] Créer des playlists aléatoires
    \end{itemize}
    *(Réponse : B)*

    \item Les vidéos courtes (Shorts) sur YouTube nécessitent-elles une optimisation SEO ?
    \begin{itemize}
        \item[A)] Non, elles sont automatiquement référencées
        \item[B)] Oui, elles bénéficient aussi d'une optimisation SEO
        \item[C)] Uniquement si elles durent plus de 30 secondes
        \item[D)] Uniquement si elles sont en 4K
    \end{itemize}
    *(Réponse : B)*

    \item Quel réseau social vidéo est également un moteur de recherche important ?
    \begin{itemize}
        \item[A)] Instagram
        \item[B)] TikTok
        \item[C)] YouTube
        \item[D)] Facebook
    \end{itemize}
    *(Réponse : C)*

    \item L'optimisation des descriptions YouTube doit inclure :
    \begin{itemize}
        \item[A)] Uniquement des emojis
        \item[B)] Des mots-clés et des liens pertinents
        \item[C)] Uniquement le titre de la vidéo
        \item[D)] Des informations de copyright
    \end{itemize}
    *(Réponse : B)*

    \item Google peut afficher des vidéos dans les featured snippets. Cette affirmation est-elle vraie ?
    \begin{itemize}
        \item[A)] Oui
        \item[B)] Non
        \item[C)] Uniquement pour les vidéos YouTube Premium
        \item[D)] Uniquement pour les vidéos en direct
    \end{itemize}
    *(Réponse : A)*

    \item Quel élément n'est PAS un facteur de classement pour les vidéos YouTube ?
    \begin{itemize}
        \item[A)] Le titre
        \item[B)] La description
        \item[C)] La couleur de la chaîne
        \item[D)] Les tags
    \end{itemize}
    *(Réponse : C)*

    \item Les données structurées VideoObject doivent être placées :
    \begin{itemize}
        \item[A)] Dans le fichier robots.txt
        \item[B)] Dans le sitemap
        \item[C)] Dans le code HTML de la page (JSON-LD)
        \item[D)] Dans le CSS
    \end{itemize}
    *(Réponse : C)*

    \item Quelle est la durée idéale d'une vidéo pour le référencement YouTube ?
    \begin{itemize}
        \item[A)] Moins de 1 minute
        \item[B)] Entre 5 et 15 minutes
        \item[C)] Plus de 30 minutes
        \item[D)] Peu importe, la qualité prime
    \end{itemize}
    *(Réponse : D)*

    \item Le SEO vidéo concerne également quel type de contenu ?
    \begin{itemize}
        \item[A)] Uniquement les vidéos YouTube
        \item[B)] Les vidéos hébergées sur n'importe quelle plateforme
        \item[C)] Uniquement les vidéos en direct
        \item[D)] Uniquement les tutoriels
    \end{itemize}
    *(Réponse : B)*

    \item Pour optimiser une vidéo auto-hébergée, il faut :
    \begin{itemize}
        \item[A)] La compresser et ajouter des métadonnées
        \item[B)] La mettre uniquement en 4K
        \item[C)] Supprimer le son
        \item[D)] La convertir en GIF
    \end{itemize}
    *(Réponse : A)*

    \item Quel format de fichier vidéo est recommandé pour le web ?
    \begin{itemize}
        \item[A)] AVI
        \item[B)] MP4
        \item[C)] MOV
        \item[D)] WMV
    \end{itemize}
    *(Réponse : B)*

    \item Quel est le rôle des sous-titres dans le SEO vidéo ?
    \begin{itemize}
        \item[A)] Améliorer l'accessibilité et l'indexation
        \item[B)] Uniquement esthétique
        \item[C)] Réduire le poids du fichier
        \item[D)] Augmenter la luminosité
    \end{itemize}
    *(Réponse : A)*

    \item Google utilise quel signal pour comprendre le contenu d'une vidéo ?
    \begin{itemize}
        \item[A)] La couleur dominante
        \item[B)] Le nombre de commentaires
        \item[C)] La transcription et les métadonnées
        \item[D)] La résolution
    \end{itemize}
    *(Réponse : C)*

    \item Le référencement vidéo est important car :
    \begin{itemize}
        \item[A)] YouTube est le deuxième moteur de recherche mondial
        \item[B)] Les vidéos ne se référencent pas
        \item[C)] Google ignore les vidéos
        \item[D)] Les vidéos sont toujours payantes
    \end{itemize}
    *(Réponse : A)*

    \item Une bonne stratégie de mots-clés pour vidéo doit cibler :
    \begin{itemize}
        \item[A)] Les mots-clés les plus courts possible
        \item[B)] Des mots-clés en rapport avec le contenu vidéo
        \item[C)] Uniquement des marques
        \item[D)] Des mots-clés en anglais uniquement
    \end{itemize}
    *(Réponse : B)*

    \item Les vidéos apparaissent dans quel type de résultat Google ?
    \begin{itemize}
        \item[A)] Google Shopping uniquement
        \item[B)] Google Vidéo et résultats enrichis
        \item[C)] Google Maps uniquement
        \item[D)] Google Flights uniquement
    \end{itemize}
    *(Réponse : B)*

    \item Quelle est une bonne pratique pour le référencement d'une playlist YouTube ?
    \begin{itemize}
        \item[A)] Mélanger des sujets différents
        \item[B)] Créer des playlists thématiques cohérentes
        \item[C)] Mettre toutes les vidéos dans une seule playlist
        \item[D)] Ne pas utiliser de playlists
    \end{itemize}
    *(Réponse : B)*

    \item Quel élément est déconseillé pour le thumbnail d'une vidéo ?
    \begin{itemize}
        \item[A)] Une image de haute qualité
        \item[B)] Du texte lisible
        \item[C)] Une image trompeuse (clickbait)
        \item[D)] Des couleurs contrastées
    \end{itemize}
    *(Réponse : C)*

    \item La transcription d'une vidéo doit être placée :
    \begin{itemize}
        \item[A)] Dans les commentaires
        \item[B)] Dans la description ou à côté de la vidéo
        \item[C)] Uniquement dans le titre
        \item[D)] Dans les tags
    \end{itemize}
    *(Réponse : B)*

    \item Les données structurées vidéo aident Google à :
    \begin{itemize}
        \item[A)] Connaître la date de création de la vidéo
        \item[B)] Comprendre le sujet et le contexte de la vidéo
        \item[C)] Trouver le créateur
        \item[D)] Estimer le budget
    \end{itemize}
    *(Réponse : B)*

    \item Google peut indexer le contenu audio d'une vidéo. Cette affirmation est-elle vraie ?
    \begin{itemize}
        \item[A)] Oui, grâce à la transcription
        \item[B)] Non, Google n'indexe que les images
        \item[C)] Uniquement si la vidéo est en anglais
        \item[D)] Uniquement pour les vidéos en direct
    \end{itemize}
    *(Réponse : A)*

    \item Combien de caractères maximum sont recommandés pour la description YouTube ?
    \begin{itemize}
        \item[A)] 200 caractères
        \item[B)] 500 caractères
        \item[C)] Environ 5000 caractères
        \item[D)] Illimité
    \end{itemize}
    *(Réponse : C)*

    \item Le SEO vidéo inclut l'optimisation de quel élément dans le code HTML ?
    \begin{itemize}
        \item[A)] La balise <video> avec ses attributs
        \item[B)] Uniquement la balise <title>
        \item[C)] Uniquement les balises <h1>
        \item[D)] Uniquement les liens
    \end{itemize}
    *(Réponse : A)*

    \item Quelle est la meilleure pratique pour nommer un fichier vidéo ?
    \begin{itemize}
        \item[A)] video1.mp4
        \item[B)] Un nom descriptif avec mots-clés (ex: tutoriel-seo-video.mp4)
        \item[C)] Uniquement des chiffres
        \item[D)] majuscules-sans-separateur.mp4
    \end{itemize}
    *(Réponse : B)*

    \item Les vidéos en direct (live) peuvent-elles être optimisées pour le SEO ?
    \begin{itemize}
        \item[A)] Oui, avec un titre, une description et des tags optimisés
        \item[B)] Non, les lives ne sont jamais indexés
        \item[C)] Uniquement après la fin du live
        \item[D)] Uniquement si elles durent plus d'une heure
    \end{itemize}
    *(Réponse : A)*
\end{enumerate}
\section*{Chapitre 29 : Le SEO Vocal (Voice Search)}
\begin{enumerate}
    \item Quels sont les principaux assistants vocaux mentionnés dans le chapitre ?
    \begin{itemize}
        \item[A)] Cortana, Bixby, Alexa
        \item[B)] Google Assistant, Siri, Alexa
        \item[C)] Google Assistant, Cortana, Bixby
        \item[D)] Siri, Bixby, Cortana
    \end{itemize}
    *(Réponse : B)*

    \item Comment sont les requêtes vocales comparées aux requêtes textuelles ?
    \begin{itemize}
        \item[A)] Plus courtes
        \item[B)] Plus longues (phrases complètes)
        \item[C)] Identiques
        \item[D)] Moins précises
    \end{itemize}
    *(Réponse : B)*

    \item Quel est le pourcentage de recherches vocales concernant des entreprises locales ?
    \begin{itemize}
        \item[A)] 25\%
        \item[B)] 42\%
        \item[C)] 58\%
        \item[D)] 75\%
    \end{itemize}
    *(Réponse : C)*

    \item Quel élément de recherche est la source principale des réponses vocales ?
    \begin{itemize}
        \item[A)] Les publicités
        \item[B)] Les featured snippets
        \item[C)] Les résultats payants
        \item[D)] Les réseaux sociaux
    \end{itemize}
    *(Réponse : B)*

    \item Quel format de données structurées est recommandé pour la recherche vocale ?
    \begin{itemize}
        \item[A)] Product schema
        \item[B)] FAQ schema
        \item[C)] Event schema
        \item[D)] VideoObject schema
    \end{itemize}
    *(Réponse : B)*

    \item Quel type de mots-clés faut-il viser pour la recherche vocale ?
    \begin{itemize}
        \item[A)] Mots-clés courts et génériques
        \item[B)] Mots-clés long-tail en format question
        \item[C)] Mots-clés à forte concurrence
        \item[D)] Mots-clés en anglais uniquement
    \end{itemize}
    *(Réponse : B)*

    \item Quel aspect technique est critique pour le SEO vocal ?
    \begin{itemize}
        \item[A)] La taille du logo
        \item[B)] La vitesse de chargement
        \item[C)] La couleur du site
        \item[D)] Le nombre de pages
    \end{itemize}
    *(Réponse : B)*

    \item La recherche vocale privilégie quel type de réponses ?
    \begin{itemize}
        \item[A)] Les réponses longues et détaillées
        \item[B)] Les réponses concises
        \item[C)] Les réponses en vidéo
        \item[D)] Les réponses en images
    \end{itemize}
    *(Réponse : B)*

    \item Quel langage doit être utilisé pour optimiser le contenu pour la recherche vocale ?
    \begin{itemize}
        \item[A)] Un langage technique et complexe
        \item[B)] Un langage naturel et conversationnel
        \item[C)] Un langage juridique
        \item[D)] Un langage poétique
    \end{itemize}
    *(Réponse : B)*

    \item Les assistants vocaux transforment la façon dont les utilisateurs interagissent avec :
    \begin{itemize}
        \item[A)] Les sites web uniquement
        \item[B)] Les moteurs de recherche
        \item[C)] Les emails
        \item[D)] Les applications mobiles uniquement
    \end{itemize}
    *(Réponse : B)*

    \item Quelle est une caractéristique des requêtes vocales ?
    \begin{itemize}
        \item[A)] Intention floue
        \item[B)] Intention claire (acheter, réserver, trouver)
        \item[C)] Aucune intention
        \item[D)] Intention de divertissement uniquement
    \end{itemize}
    *(Réponse : B)*

    \item Les requêtes vocales sont souvent formulées comment ?
    \begin{itemize}
        \item[A)] En mots-clés isolés
        \item[B)] En format question
        \item[C)] En langage codé
        \item[D)] En abréviations
    \end{itemize}
    *(Réponse : B)*

    \item La recherche vocale est particulièrement importante pour quel type de SEO ?
    \begin{itemize}
        \item[A)] SEO e-commerce uniquement
        \item[B)] SEO local
        \item[C)] SEO international uniquement
        \item[D)] SEO vidéo uniquement
    \end{itemize}
    *(Réponse : B)*

    \item Quelle position dans les SERP est particulièrement visée par le SEO vocal ?
    \begin{itemize}
        \item[A)] La position 1
        \item[B)] La position zéro (featured snippet)
        \item[C)] La position 5
        \item[D)] La position 10
    \end{itemize}
    *(Réponse : B)*

    \item Quel type de contenu est le plus adapté à la recherche vocale ?
    \begin{itemize}
        \item[A)] Des pages produits
        \item[B)] Des questions-réponses (FAQ)
        \item[C)] Des pages d'accueil
        \item[D)] Des pages de contact
    \end{itemize}
    *(Réponse : B)*

    \item L'optimisation pour la recherche vocale doit inclure :
    \begin{itemize}
        \item[A)] Des phrases complexes
        \item[B)] Du langage naturel
        \item[C)] Du jargon technique
        \item[D)] Des acronymes
    \end{itemize}
    *(Réponse : B)*

    \item Le SEO vocal est en croissance grâce à :
    \begin{itemize}
        \item[A)] La baisse des recherches textuelles
        \item[B)] L'adoption massive des assistants vocaux
        \item[C)] La disparition de Google
        \item[D)] La baisse de l'utilisation d'Internet
    \end{itemize}
    *(Réponse : B)*

    \item Pour optimiser le contenu pour la recherche vocale, il faut utiliser :
    \begin{itemize}
        \item[A)] Uniquement des listes à puces
        \item[B)] Un format de questions-réponses
        \item[C)] Uniquement des tableaux
        \item[D)] Uniquement des images
    \end{itemize}
    *(Réponse : B)*

    \item Les featured snippets sont également appelés :
    \begin{itemize}
        \item[A)] Position un
        \item[B)] Position zéro
        \item[C)] Position négative
        \item[D)] Position premium
    \end{itemize}
    *(Réponse : B)*

    \item Quel pourcentage des recherches vocales sont effectuées sur mobile ?
    \begin{itemize}
        \item[A)] Environ 20\%
        \item[B)] Environ 50\%
        \item[C)] Environ 80\%
        \item[D)] Environ 95\%
    \end{itemize}
    *(Réponse : C)*

    \item Les recherches vocales locales sont souvent formulées ainsi :
    \begin{itemize}
        \item[A)] "Pizza"
        \item[B)] "Où trouver une pizza près de chez moi ?"
        \item[C)] "Recette pizza"
        \item[D)] "Histoire de la pizza"
    \end{itemize}
    *(Réponse : B)*

    \item Google Assistant, Siri et Alexa sont des exemples de :
    \begin{itemize}
        \item[A)] Navigateurs web
        \item[B)] Assistants vocaux
        \item[C)] Réseaux sociaux
        \item[D)] Antivirus
    \end{itemize}
    *(Réponse : B)*

    \item La recherche vocale utilise principalement quel type d'intention ?
    \begin{itemize}
        \item[A)] Intention informationnelle et locale
        \item[B)] Intention de divertissement uniquement
        \item[C)] Intention de téléchargement
        \item[D)] Intention sociale
    \end{itemize}
    *(Réponse : A)*

    \item Pour apparaître dans les résultats vocaux, le contenu doit :
    \begin{itemize}
        \item[A)] Être très long
        \item[B)] Répondre directement et clairement aux questions
        \item[C)] Être uniquement en vidéo
        \item[D)] Contenir beaucoup de publicités
    \end{itemize}
    *(Réponse : B)*

    \item Quelle stratégie de mots-clés est recommandée pour le vocal ?
    \begin{itemize}
        \item[A)] Mots-clés génériques courts
        \item[B)] Expressions de mots-clés en langage naturel
        \item[C)] Mots-clés en majuscules
        \item[D)] Mots-clés sans rapport avec le contenu
    \end{itemize}
    *(Réponse : B)*

    \item Quel type de schéma est recommandé pour le SEO vocal ?
    \begin{itemize}
        \item[A)] Product schema
        \item[B)] FAQ schema et HowTo schema
        \item[C)] VideoObject schema
        \item[D)] Event schema
    \end{itemize}
    *(Réponse : B)*

    \item La recherche vocale privilégie les résultats :
    \begin{itemize}
        \item[A)] Payants
        \item[B)] Locaux et pertinents
        \item[C)] Les plus longs
        \item[D)] Les plus récents uniquement
    \end{itemize}
    *(Réponse : B)*

    \item Le SEO vocal nécessite d'optimiser pour :
    \begin{itemize}
        \item[A)] Le format PDF
        \item[B)] Les featured snippets
        \item[C)] Les pop-ups
        \item[D)] Les vidéos uniquement
    \end{itemize}
    *(Réponse : B)*

    \item Quel type d'appareil utilise principalement la recherche vocale ?
    \begin{itemize}
        \item[A)] Les ordinateurs de bureau
        \item[B)] Les smartphones et enceintes connectées
        \item[C)] Les imprimantes
        \item[D)] Les téléviseurs anciens
    \end{itemize}
    *(Réponse : B)*

    \item La croissance de la recherche vocale est due à :
    \begin{itemize}
        \item[A)] La diminution de la population
        \item[B)] L'amélioration de la reconnaissance vocale
        \item[C)] La baisse du niveau d'éducation
        \item[D)] La disparition des écrans
    \end{itemize}
    *(Réponse : B)*

    \item Le contenu optimisé pour le vocal doit être :
    \begin{itemize}
        \item[A)] Complexe et technique
        \item[B)] Clair, concis et conversationnel
        \item[C)] Long et détaillé
        \item[D)] En langage soutenu
    \end{itemize}
    *(Réponse : B)*

    \item Google utilise quel algorithme pour comprendre les requêtes vocales ?
    \begin{itemize}
        \item[A)] Panda
        \item[B)] BERT
        \item[C)] Penguin
        \item[D)] Hummingbird
    \end{itemize}
    *(Réponse : B)*

    \item La recherche vocale est particulièrement utilisée pour :
    \begin{itemize}
        \item[A)] La recherche académique
        \item[B)] Les recherches locales et les directions
        \item[C)] La programmation
        \item[D)] Les calculs scientifiques
    \end{itemize}
    *(Réponse : B)*

    \item Quel est l'avantage des données structurées FAQ pour le vocal ?
    \begin{itemize}
        \item[A)] Elles améliorent le design du site
        \item[B)] Elles permettent aux moteurs de recherche d'identifier facilement les réponses
        \item[C)] Elles réduisent le temps de chargement
        \item[D)] Elles augmentent le nombre de pages
    \end{itemize}
    *(Réponse : B)*

    \item Les featured snippets sont importants pour le vocal car :
    \begin{itemize}
        \item[A)] Google lit automatiquement le premier résultat
        \item[B)] Les assistants vocaux lisent souvent le featured snippet
        \item[C)] Ils sont plus beaux visuellement
        \item[D)] Ils contiennent plus de mots-clés
    \end{itemize}
    *(Réponse : B)*

    \item La recherche vocale utilise des phrases complètes plutôt que :
    \begin{itemize}
        \item[A)] Des images
        \item[B)] Des mots-clés isolés
        \item[C)] Des vidéos
        \item[D)] Des fichiers audio
    \end{itemize}
    *(Réponse : B)*

    \item Le SEO vocal est une composante importante du SEO moderne car :
    \begin{itemize}
        \item[A)] Il remplace totalement le SEO traditionnel
        \item[B)] De plus en plus de recherches sont effectuées vocalement
        \item[C)] Il ne nécessite aucune optimisation
        \item[D)] Il est moins efficace
    \end{itemize}
    *(Réponse : B)*

    \item Quel réseau social est aussi utilisé pour la recherche vocale ?
    \begin{itemize}
        \item[A)] Facebook
        \item[B)] YouTube (via Google Assistant)
        \item[C)] LinkedIn
        \item[D)] Pinterest
    \end{itemize}
    *(Réponse : B)*

    \item Comment s'appelle la position dans les SERP occupée par les featured snippets ?
    \begin{itemize}
        \item[A)] Position premium
        \item[B)] Position zéro
        \item[C)] Position première
        \item[D)] Position vedette
    \end{itemize}
    *(Réponse : B)*

    \item Le langage naturel et conversationnel est important pour :
    \begin{itemize}
        \item[A)] Le SEO technique uniquement
        \item[B)] Le SEO vocal
        \item[C)] Le netlinking uniquement
        \item[D)] Le SEO vidéo uniquement
    \end{itemize}
    *(Réponse : B)*

    \item L'optimisation de la vitesse de chargement est critique pour le vocal car :
    \begin{itemize}
        \item[A)] Les assistants vocaux nécessitent des réponses rapides
        \item[B)] Cela n'a pas d'impact
        \item[C)] Cela concerne uniquement le SEO classique
        \item[D)] Cela ralentit les assistants vocaux
    \end{itemize}
    *(Réponse : A)*

    \item Quel pourcentage approximatif des recherches vocales donnent une réponse en featured snippet ?
    \begin{itemize}
        \item[A)] Moins de 10\%
        \item[B)] Environ 30\%
        \item[C)] Environ 40\%
        \item[D)] Plus de 60\%
    \end{itemize}
    *(Réponse : C)*

    \item Les requêtes vocales sont généralement formulées avec :
    \begin{itemize}
        \item[A)] Des mots interrogatifs (qui, quoi, où, comment)
        \item[B)] Des chiffres uniquement
        \item[C)] Des symboles
        \item[D)] Des acronymes
    \end{itemize}
    *(Réponse : A)*

    \item Pour le SEO vocal, il est important de :
    \begin{itemize}
        \item[A)] Ignorer les featured snippets
        \item[B)] Optimiser pour les mots-clés long-tail en format question
        \item[C)] Utiliser uniquement des mots-clés courts
        \item[D)] Éviter le langage naturel
    \end{itemize}
    *(Réponse : B)*

    \item L'essor de la recherche vocale est lié à :
    \begin{itemize}
        \item[A)] La baisse des moteurs de recherche
        \item[B)] L'adoption massive d'assistants vocaux comme Google Assistant et Alexa
        \item[C)] La diminution de l'utilisation d'Internet
        \item[D)] La disparition des écrans tactiles
    \end{itemize}
    *(Réponse : B)*

    \item Les featured snippets sont aussi appelés résultats en position zéro. Cette affirmation est-elle vraie ?
    \begin{itemize}
        \item[A)] Oui
        \item[B)] Non
        \item[C)] Uniquement pour les vidéos
        \item[D)] Uniquement pour les images
    \end{itemize}
    *(Réponse : A)*

    \item Pour optimiser le contenu pour le vocal, il faut privilégier les questions commençant par :
    \begin{itemize}
        \item[A)] "Combien"
        \item[B)] "Qui, quoi, où, quand, comment, pourquoi"
        \item[C)] "Si"
        \item[D)] "Peut-être"
    \end{itemize}
    *(Réponse : B)*

    \item Quelle est la principale différence entre une recherche textuelle et une recherche vocale ?
    \begin{itemize}
        \item[A)] La recherche vocale est plus rapide
        \item[B)] La recherche vocale utilise un langage plus naturel et des phrases complètes
        \item[C)] La recherche vocale utilise des caractères spéciaux
        \item[D)] Il n'y a aucune différence
    \end{itemize}
    *(Réponse : B)*

    \item Quel pourcentage des recherches vocales ont une intention locale ?
    \begin{itemize}
        \item[A)] 22\%
        \item[B)] 58\%
        \item[C)] 75\%
        \item[D)] 90\%
    \end{itemize}
    *(Réponse : B)*

    \item Pour apparaître en réponse vocale, le contenu doit être positionné :
    \begin{itemize}
        \item[A)] En page 2 des résultats
        \item[B)] En position zéro ou top 3 des résultats
        \item[C)] Uniquement en position 1
        \item[D)] Peu importe la position
    \end{itemize}
    *(Réponse : B)*
\end{enumerate}
\section*{Chapitre 30 : A/B Testing et Expérimentation SEO}
\begin{enumerate}
    \item Qu'est-ce que le SEO split testing ?
    \begin{itemize}
        \item[A)] Tester des modifications sur des pages pour mesurer leur impact sur le trafic payant
        \item[B)] Tester des modifications sur des pages pour mesurer leur impact sur le trafic organique
        \item[C)] Tester la vitesse du site
        \item[D)] Tester les backlinks
    \end{itemize}
    *(Réponse : B)*

    \item Quel est le premier type de test SEO mentionné ?
    \begin{itemize}
        \item[A)] Before/After
        \item[B)] A/B test (split test)
        \item[C)] Test apparié
        \item[D)] Test temps réel
    \end{itemize}
    *(Réponse : B)*

    \item Quelle est la durée minimale recommandée pour un test SEO ?
    \begin{itemize}
        \item[A)] 1 semaine
        \item[B)] 2-4 semaines
        \item[C)] 2 mois
        \item[D)] 6 mois
    \end{itemize}
    *(Réponse : B)*

    \item Quel niveau de confiance est considéré comme minimum pour la signification statistique ?
    \begin{itemize}
        \item[A)] 90\%
        \item[B)] 95\%
        \item[C)] 99\%
        \item[D)] 100\%
    \end{itemize}
    *(Réponse : B)*

    \item Quel est l'élément le plus testé en SEO ?
    \begin{itemize}
        \item[A)] Les images
        \item[B)] Les title tags
        \item[C)] Les vidéos
        \item[D)] Les backlinks
    \end{itemize}
    *(Réponse : B)*

    \item Quel type de test compare les performances avant et après un changement ?
    \begin{itemize}
        \item[A)] A/B test
        \item[B)] Before/After
        \item[C)] Test temps réel
        \item[D)] Split test
    \end{itemize}
    *(Réponse : B)*

    \item Que signifie p-value < 0.05 ?
    \begin{itemize}
        \item[A)] Le résultat est invalide
        \item[B)] Le résultat est statistiquement significatif à 95\%
        \item[C)] Le test doit être refait
        \item[D)] L'échantillon est trop petit
    \end{itemize}
    *(Réponse : B)*

    \item Quel outil professionnel de A/B testing SEO est mentionné ?
    \begin{itemize}
        \item[A)] Google Analytics
        \item[B)] SearchPilot
        \item[C)] Microsoft Excel
        \item[D)] Photoshop
    \end{itemize}
    *(Réponse : B)*

    \item Quel est un piège courant du SEO testing ?
    \begin{itemize}
        \item[A)] Tester trop longtemps
        \item[B)] Tester plusieurs changements en même temps
        \item[C)] Avoir trop de trafic
        \item[D)] Utiliser des outils professionnels
    \end{itemize}
    *(Réponse : B)*

    \item Combien d'impressions par variante sont recommandées minimum ?
    \begin{itemize}
        \item[A)] 100
        \item[B)] 1000
        \item[C)] 10000
        \item[D)] 100000
    \end{itemize}
    *(Réponse : B)*

    \item Que faut-il faire des deux versions d'un test pour éviter le contenu dupliqué ?
    \begin{itemize}
        \item[A)] Les indexer toutes les deux
        \item[B)] Ne pas indexer la variante et utiliser une canonique sur la version de contrôle
        \item[C)] Les supprimer toutes les deux
        \item[D)] Les fusionner
    \end{itemize}
    *(Réponse : B)*

    \item Google peut mettre combien de temps à recrawler une page ?
    \begin{itemize}
        \item[A)] 1-2 jours
        \item[B)] 2 à 4 semaines
        \item[C)] 6 mois
        \item[D)] 1 an
    \end{itemize}
    *(Réponse : B)*

    \item Quelle méthode de split est recommandée ?
    \begin{itemize}
        \item[A)] Split par page aléatoire
        \item[B)] Split par URLs (groupe contrôle vs groupe test)
        \item[C)] Split par visiteur
        \item[D)] Split par navigateur
    \end{itemize}
    *(Réponse : B)*

    \item Quel type de schema est testé dans les données structurées ?
    \begin{itemize}
        \item[A)] Person schema uniquement
        \item[B)] FAQ schema, HowTo schema, Review schema, Article schema
        \item[C)] Event schema uniquement
        \item[D)] VideoObject schema uniquement
    \end{itemize}
    *(Réponse : B)*

    \item Quel format de title tag est recommandé pour le test ?
    \begin{itemize}
        \item[A)] Toujours le même format
        \item[B)] "Mot-clé principal -- Secondaire | Marque" vs variations
        \item[C)] Uniquement le nom de la marque
        \item[D)] Uniquement le mot-clé
    \end{itemize}
    *(Réponse : B)*

    \item Quelle est une variante testée pour les meta descriptions ?
    \begin{itemize}
        \item[A)] La couleur
        \item[B)] La présence d'un CTA ("Découvrez" vs "Achetez" vs "Apprenez")
        \item[C)] La police
        \item[D)] La taille du texte
    \end{itemize}
    *(Réponse : B)*

    \item Quel CTA est mentionné comme option de test pour les meta descriptions ?
    \begin{itemize}
        \item[A)] "Cliquez ici"
        \item[B)] "Découvrez", "Achetez", "Apprenez"
        \item[C)] "OK"
        \item[D)] "Suivant"
    \end{itemize}
    *(Réponse : B)*

    \item Les titres avec chiffres génèrent quel impact sur le CTR ?
    \begin{itemize}
        \item[A)] Aucun impact
        \item[B)] +23\% de CTR à position égale
        \item[C)] -10\% de CTR
        \item[D)] +100\% de CTR
    \end{itemize}
    *(Réponse : B)*

    \item Pourquoi tester en SEO ?
    \begin{itemize}
        \item[A)] Pour perdre du temps
        \item[B)] Pour remplacer l'intuition par des données
        \item[C)] Pour augmenter les coûts
        \item[D)] Pour compliquer le travail
    \end{itemize}
    *(Réponse : B)*

    \item Quel facteur peut masquer les effets d'un test SEO ?
    \begin{itemize}
        \item[A)] La météo
        \item[B)] Le bruit algorithmique (fluctuations normales des SERP)
        \item[C)] La phase de la lune
        \item[D)] Le jour de la semaine
    \end{itemize}
    *(Réponse : B)*

    \item Quel outil de test est intégré à GA4 ?
    \begin{itemize}
        \item[A)] SearchPilot
        \item[B)] Google Analytics 4 Experiences A/B natives
        \item[C)] Screaming Frog
        \item[D)] Ahrefs
    \end{itemize}
    *(Réponse : B)*

    \item Les pages testées doivent être comparables en termes de :
    \begin{itemize}
        \item[A)] Couleur
        \item[B)] PageRank et autorité
        \item[C)] Nombre d'images
        \item[D)] Date de création
    \end{itemize}
    *(Réponse : B)*

    \item Quel type de test utilise des pages de catégories similaires ?
    \begin{itemize}
        \item[A)] A/B test
        \item[B)] Test apparié
        \item[C)] Before/After
        \item[D)] Test multivarié
    \end{itemize}
    *(Réponse : B)*

    \item La saisonnalité est un facteur à prendre en compte pour :
    \begin{itemize}
        \item[A)] Ignorer les résultats du test
        \item[B)] Tester sur des périodes similaires
        \item[C)] Tester uniquement en été
        \item[D)] Tester uniquement en hiver
    \end{itemize}
    *(Réponse : B)*

    \item Quel élément de contenu peut être testé en priorité ?
    \begin{itemize}
        \item[A)] La couleur du fond
        \item[B)] La longueur du contenu (1500 vs 3000 mots)
        \item[C)] La langue du site
        \item[D)] Le nom de domaine
    \end{itemize}
    *(Réponse : B)*

    \item Comment s'appelle le test de proportion utilisé pour la signification statistique ?
    \begin{itemize}
        \item[A)] T-test
        \item[B)] Z-test
        \item[C)] Chi-2
        \item[D)] ANOVA
    \end{itemize}
    *(Réponse : B)*

    \item Une valeur z > 1.96 indique :
    \begin{itemize}
        \item[A)] Aucune signification
        \item[B)] Une signification statistique à 95\%
        \item[C)] Une erreur de calcul
        \item[D)] Un échantillon trop petit
    \end{itemize}
    *(Réponse : B)*

    \item Quel est le risque d'un délai insuffisant dans un test SEO ?
    \begin{itemize}
        \item[A)] Aucun risque
        \item[B)] Google peut ne pas avoir recrawlé la page, faussant les résultats
        \item[C)] Le site peut être pénalisé
        \item[D)] Les données sont perdues
    \end{itemize}
    *(Réponse : B)*

    \item Quel élément est prioritaire à tester dans les données structurées ?
    \begin{itemize}
        \item[A)] Person schema
        \item[B)] FAQ schema
        \item[C)] Event schema
        \item[D)] Local Business schema
    \end{itemize}
    *(Réponse : B)*

    \item Une hypothèse de test doit être formulée avant le test pour :
    \begin{itemize}
        \item[A)] Rien, ce n'est pas nécessaire
        \item[B)] Savoir précisément ce qu'on teste et pourquoi
        \item[C)] Impressionner les collègues
        \item[D)] Remplir un document
    \end{itemize}
    *(Réponse : B)*

    \item Quel type de contenu (format) peut être testé ?
    \begin{itemize}
        \item[A)] Couleur uniquement
        \item[B)] Listes vs paragraphes vs guides
        \item[C)] Police de caractères
        \item[D)] Langue
    \end{itemize}
    *(Réponse : B)*

    \item Le monitoring continu avec ajustements dynamiques est appelé :
    \begin{itemize}
        \item[A)] A/B test
        \item[B)] Test temps réel
        \item[C)] Before/After
        \item[D)] Test apparié
    \end{itemize}
    *(Réponse : B)*

    \item Quel outil de test SEO fonctionne au niveau CDN ?
    \begin{itemize}
        \item[A)] SearchPilot
        \item[B)] Cloudflare A/B Testing
        \item[C)] VWO
        \item[D)] AB Tasty
    \end{itemize}
    *(Réponse : B)*

    \item Quel est le résultat du test de longueur de contenu mentionné dans le cas pratique ?
    \begin{itemize}
        \item[A)] +20\% de trafic
        \item[B)] +67\% de trafic pour le groupe B
        \item[C)] -10\% de trafic
        \item[D)] Aucun changement
    \end{itemize}
    *(Réponse : B)*

    \item L'ajout de FAQ schema a généré quel résultat dans le cas pratique ?
    \begin{itemize}
        \item[A)] -15\% de featured snippets
        \item[B)] +15\% de featured snippets, +8\% de trafic
        \item[C)] +50\% de featured snippets
        \item[D)] Aucun changement
    \end{itemize}
    *(Réponse : B)*

    \item Combien de pages sont recommandées minimum pour un split test par URLs ?
    \begin{itemize}
        \item[A)] 10 pages
        \item[B)] 100 pages (50 contrôle + 50 test)
        \item[C)] 5 pages
        \item[D)] 1000 pages
    \end{itemize}
    *(Réponse : B)*

    \item Quel outil de CRO + testing SEO avec intégration GA4 est mentionné ?
    \begin{itemize}
        \item[A)] SearchPilot
        \item[B)] AB Tasty
        \item[C)] Cloudflare
        \item[D)] Screaming Frog
    \end{itemize}
    *(Réponse : B)*

    \item La signification statistique permet d'éviter :
    \begin{itemize}
        \item[A)] Les tests trop longs
        \item[B)] Les conclusions hâtives basées sur le hasard
        \item[C)] L'utilisation d'outils
        \item[D)] La publication de contenu
    \end{itemize}
    *(Réponse : B)*

    \item Comment s'appelle la formule qui calcule la signification statistique d'un test de CTR ?
    \begin{itemize}
        \item[A)] Formule de corrélation
        \item[B)] Z-test pour proportions
        \item[C)] Formule de régression
        \item[D)] Formule de variance
    \end{itemize}
    *(Réponse : B)*

    \item Quelle est la première étape du processus de test rigoureux ?
    \begin{itemize}
        \item[A)] Mesure
        \item[B)] Hypothèse
        \item[C)] Déploiement
        \item[D)] Analyse
    \end{itemize}
    *(Réponse : B)*

    \item Le Before/After est considéré comme :
    \begin{itemize}
        \item[A)] Le plus fiable
        \item[B)] Moins fiable que le split test
        \item[C)] Aussi fiable que le A/B test
        \item[D)] La méthode recommandée
    \end{itemize}
    *(Réponse : B)*

    \item Avant de déployer un changement, il faut :
    \begin{itemize}
        \item[A)] Le déployer immédiatement
        \item[B)] Atteindre la signification statistique
        \item[C)] Attendre au moins 1 jour
        \item[D)] Le tester sur 1 page
    \end{itemize}
    *(Réponse : B)*

    \item L'uplift (amélioration) se calcule comment ?
    \begin{itemize}
\item[A)] CTR$_{\text{variant}}$ - CTR$_{\text{control}}$
\item[B)] (CTR$_{\text{variant}}$ - CTR$_{\text{control}}$) / CTR$_{\text{control}}$ × 100
\item[C)] CTR$_{\text{control}}$ / CTR$_{\text{variant}}$
\item[D)] CTR$_{\text{variant}}$ + CTR$_{\text{control}}$
    \end{itemize}
    *(Réponse : B)*

    \item Quel élément est important à vérifier en pré-test ?
    \begin{itemize}
        \item[A)] La météo
        \item[B)] L'équivalence statistique des groupes
        \item[C)] Le nombre de visiteurs connectés
        \item[D)] La phase lunaire
    \end{itemize}
    *(Réponse : B)*

    \item Les titres interrogatifs vs affirmatifs peuvent être testés pour :
    \begin{itemize}
        \item[A)] La longueur du contenu
        \item[B)] Les title tags
        \item[C)] Les backlinks
        \item[D)] Les images
    \end{itemize}
    *(Réponse : B)*

    \item Quel est le risque d'un canonical incorrect dans un test ?
    \begin{itemize}
        \item[A)] Aucun risque
        \item[B)] Google peut indexer la mauvaise version
        \item[C)] Le site peut être plus rapide
        \item[D)] Les backlinks augmentent
    \end{itemize}
    *(Réponse : B)*

    \item La validation après déploiement consiste à :
    \begin{itemize}
        \item[A)] Supprimer la version gagnante
        \item[B)] Confirmer les résultats sur la durée
        \item[C)] Ignorer les résultats
        \item[D)] Recommencer le test
    \end{itemize}
    *(Réponse : B)*

    \item Quel langage de programmation est utilisé dans l'exemple de test A/B ?
    \begin{itemize}
        \item[A)] JavaScript
        \item[B)] Python
        \item[C)] PHP
        \item[D)] Java
    \end{itemize}
    *(Réponse : B)*

    \item VWO (Visual Website Optimizer) permet :
    \begin{itemize}
        \item[A)] Uniquement le design
        \item[B)] Le testing avec segmentation par source de trafic
        \item[C)] Uniquement l'hébergement
        \item[D)] Uniquement les emails
    \end{itemize}
    *(Réponse : B)*

    \item Quel est le premier signe qu'un test SEO peut être statistiquement significatif ?
    \begin{itemize}
        \item[A)] Un uplift visible à l'oeil nu
        \item[B)] Une p-value inférieure à 0.05
        \item[C)] Un nombre élevé de visiteurs
        \item[D)] Une durée de test d'au moins un mois
    \end{itemize}
    *(Réponse : B)*
\end{enumerate}
\section*{Chapitre 31 : Google Search Console}
\begin{enumerate}
    \item Quel est l'outil gratuit le plus important pour tout professionnel du SEO ?
    \begin{itemize}
        \item[A)] Google Analytics
        \item[B)] Google Search Console
        \item[C)] Google Ads
        \item[D)] Google My Business
    \end{itemize}
    *(Réponse : B)*

    \item Google Search Console fournit des données provenant de :
    \begin{itemize}
        \item[A)] Google Analytics
        \item[B)] L'index Google directement
        \item[C)] Des tiers
        \item[D)] Des utilisateurs
    \end{itemize}
    *(Réponse : B)*

    \item Quel rapport de GSC montre les impressions, clics, CTR et position moyenne par requête ?
    \begin{itemize}
        \item[A)] Rapport de couverture
        \item[B)] Rapport Performance
        \item[C)] Rapport Sitemaps
        \item[D)] Rapport Liens
    \end{itemize}
    *(Réponse : B)*

    \item Quel rapport GSC permet de voir le statut d'indexation des pages ?
    \begin{itemize}
        \item[A)] Rapport Performance
        \item[B)] Rapport Couverture
        \item[C)] Rapport Core Web Vitals
        \item[D)] Rapport Pages mobiles
    \end{itemize}
    *(Réponse : B)*

    \item Quel rapport GSC concerne le LCP, INP et CLS ?
    \begin{itemize}
        \item[A)] Rapport Performance
        \item[B)] Rapport Core Web Vitals
        \item[C)] Rapport Sitemaps
        \item[D)] Rapport AMP
    \end{itemize}
    *(Réponse : B)*

    \item Quel rapport GSC permet de soumettre et vérifier le statut des sitemaps ?
    \begin{itemize}
        \item[A)] Rapport Performance
        \item[B)] Rapport Sitemaps
        \item[C)] Rapport Couverture
        \item[D)] Rapport Liens
    \end{itemize}
    *(Réponse : B)*

    \item Quel rapport GSC montre les pénalités manuelles ?
    \begin{itemize}
        \item[A)] Rapport Liens
        \item[B)] Manual Actions
        \item[C)] Rapport AMP
        \item[D)] Rapport Pages mobiles
    \end{itemize}
    *(Réponse : B)*

    \item Le rapport Performance de GSC inclut toutes ces métriques sauf :
    \begin{itemize}
        \item[A)] Impressions
        \item[B)] Pages vues
        \item[C)] Clics
        \item[D)] Position moyenne
    \end{itemize}
    *(Réponse : B)*

    \item GSC permet de diagnostiquer les problèmes d'utilisabilité mobile via quel rapport ?
    \begin{itemize}
        \item[A)] Rapport Performance
        \item[B)] Rapport Pages mobiles
        \item[C)] Rapport Sitemaps
        \item[D)] Rapport Liens
    \end{itemize}
    *(Réponse : B)*

    \item Le rapport Liens dans GSC montre :
    \begin{itemize}
        \item[A)] Uniquement les liens externes
        \item[B)] Les liens internes et externes
        \item[C)] Uniquement les liens internes
        \item[D)] Uniquement les backlinks
    \end{itemize}
    *(Réponse : B)*

    \item GSC est un outil :
    \begin{itemize}
        \item[A)] Payant
        \item[B)] Gratuit
        \item[C)] Freemium
        \item[D)] Uniquement pour les sites e-commerce
    \end{itemize}
    *(Réponse : B)*

    \item Quel rapport GSC est essentiel pour identifier les erreurs d'indexation ?
    \begin{itemize}
        \item[A)] Performance
        \item[B)] Couverture
        \item[C)] Core Web Vitals
        \item[D)] AMP
    \end{itemize}
    *(Réponse : B)*

    \item Dans GSC, le rapport Performance permet d'analyser :
    \begin{itemize}
        \item[A)] Les backlinks
        \item[B)] Le trafic organique
        \item[C)] Les pages vues
        \item[D)] Le temps passé sur le site
    \end{itemize}
    *(Réponse : B)*

    \item Quelle information peut-on obtenir dans le rapport Couverture de GSC ?
    \begin{itemize}
        \item[A)] Le nombre de visiteurs uniques
        \item[B)] Les erreurs d'indexation et les pages exclues
        \item[C)] Le chiffre d'affaires
        \item[D)] Le nombre d'abonnés
    \end{itemize}
    *(Réponse : B)*

    \item Le rapport AMP dans GSC est utile pour les sites qui utilisent :
    \begin{itemize}
        \item[A)] Les vidéos
        \item[B)] Les pages AMP (Accelerated Mobile Pages)
        \item[C)] Les images
        \item[D)] Les CSS
    \end{itemize}
    *(Réponse : B)*

    \item GSC permet de soumettre quel type de fichier pour aider Google à découvrir les pages ?
    \begin{itemize}
        \item[A)] robots.txt
        \item[B)] Sitemap XML
        \item[C)] .htaccess
        \item[D)] package.json
    \end{itemize}
    *(Réponse : B)*

    \item Combien de propriétés peut-on configurer dans GSC ?
    \begin{itemize}
        \item[A)] Une seule
        \item[B)] Plusieurs (domaines et préfixes d'URL)
        \item[C)] Maximum 5
        \item[D)] Illimité sans restriction
    \end{itemize}
    *(Réponse : B)*

    \item GSC est l'outil indispensable car il fournit :
    \begin{itemize}
        \item[A)] Des données de trafic en temps réel
        \item[B)] Des données directes de l'index Google
        \item[C)] Des informations sur les concurrents
        \item[D)] Des recommandations de design
    \end{itemize}
    *(Réponse : B)*

    \item Le rapport de couverture identifie les pages qui sont :
    \begin{itemize}
        \item[A)] Les plus visitées
        \item[B)] Erreurs, valides, exclues et avertissements
        \item[C)] Les plus lentes
        \item[D)] Les plus partagées
    \end{itemize}
    *(Réponse : B)*

    \item Quel rapport GSC est particulièrement important après une migration de site ?
    \begin{itemize}
        \item[A)] Performance
        \item[B)] Couverture (pour vérifier l'indexation)
        \item[C)] Core Web Vitals
        \item[D)] Liens
    \end{itemize}
    *(Réponse : B)*

    \item Le rapport Core Web Vitals dans GSC groupe les URLs par :
    \begin{itemize}
        \item[A)] Thème
        \item[B)] État (bon, à améliorer, mauvais)
        \item[C)] Date de publication
        \item[D)] Auteur
    \end{itemize}
    *(Réponse : B)*

    \item GSC peut être connecté à quel autre outil Google pour enrichir les données ?
    \begin{itemize}
        \item[A)] Google Ads
        \item[B)] Google Analytics
        \item[C)] Google Tag Manager
        \item[D)] Google Optimize
    \end{itemize}
    *(Réponse : B)*

    \item Quelle action Google Search Console permet de réaliser ?
    \begin{itemize}
        \item[A)] Modifier le contenu du site
        \item[B)] Demander l'indexation d'une URL
        \item[C)] Changer le design du site
        \item[D)] Ajouter des utilisateurs
    \end{itemize}
    *(Réponse : B)*

    \item Le rapport Performance de GSC peut être filtré par :
    \begin{itemize}
        \item[A)] Âge de l'utilisateur
        \item[B)] Requête, page, pays, appareil
        \item[C)] Navigateur
        \item[D)] Système d'exploitation
    \end{itemize}
    *(Réponse : B)*

    \item Quel type d'erreur peut être identifié dans le rapport de couverture ?
    \begin{itemize}
        \item[A)] Erreur 404
        \item[B)] Erreur 500
        \item[C)] Erreur de serveur (5xx), erreur 404, page bloquée par robots.txt
        \item[D)] Erreur 301
    \end{itemize}
    *(Réponse : C)*

    \item GSC permet de voir les backlinks pointant vers un site. Cette affirmation est-elle vraie ?
    \begin{itemize}
        \item[A)] Non, GSC ne montre pas les backlinks
        \item[B)] Oui, via le rapport Liens
        \item[C)] Uniquement pour les sites payants
        \item[D)] Uniquement pour les sites e-commerce
    \end{itemize}
    *(Réponse : B)*

    \item Le rapport Core Web Vitals de GSC mesure :
    \begin{itemize}
        \item[A)] Le nombre de pages visitées
        \item[B)] LCP, INP, CLS par groupe d'URLs
        \item[C)] Le nombre de visiteurs
        \item[D)] Le taux de conversion
    \end{itemize}
    *(Réponse : B)*

    \item GSC est essentiel pour diagnostiquer les problèmes de :
    \begin{itemize}
        \item[A)] Design
        \item[B)] Référencement et indexation
        \item[C)] Email marketing
        \item[D)] Publicité
    \end{itemize}
    *(Réponse : B)*

    \item Quel rapport GSC permet d'analyser les performances par appareil ?
    \begin{itemize}
        \item[A)] Rapport Couverture
        \item[B)] Rapport Performance (avec filtre appareil)
        \item[C)] Rapport Mobile uniquement
        \item[D)] Rapport AMP
    \end{itemize}
    *(Réponse : B)*

    \item GSC permet d'inspecter une URL spécifique pour vérifier :
    \begin{itemize}
        \item[A)] Son design
        \item[B)] Son statut d'indexation et la version indexée
        \item[C)] Son nombre de mots
        \item[D)] Sa date de création
    \end{itemize}
    *(Réponse : B)*

    \item Le rapport Sitemaps dans GSC permet de :
    \begin{itemize}
        \item[A)] Créer un sitemap
        \item[B)] Soumettre un sitemap et voir les erreurs
        \item[C)] Modifier le robots.txt
        \item[D)] Voir les backlinks
    \end{itemize}
    *(Réponse : B)*

    \item GSC est un outil indispensable pour tout professionnel du SEO car il est :
    \begin{itemize}
        \item[A)] Payant et complet
        \item[B)] Gratuit et fournit des données directes de Google
        \item[C)] Complexe à utiliser
        \item[D)] Réservé aux experts
    \end{itemize}
    *(Réponse : B)*

    \item Dans GSC, l'onglet "Efficacité" correspond au rapport :
    \begin{itemize}
        \item[A)] Couverture
        \item[B)] Performance
        \item[C)] Core Web Vitals
        \item[D)] Sitemaps
    \end{itemize}
    *(Réponse : B)*

    \item GSC peut être utilisé pour quel type de site web ?
    \begin{itemize}
        \item[A)] Uniquement les sites e-commerce
        \item[B)] Tout type de site web
        \item[C)] Uniquement les blogs
        \item[D)] Uniquement les sites d'actualité
    \end{itemize}
    *(Réponse : B)*

    \item Comment s'appelle la fonctionnalité de GSC qui permet de signaler des backlinks toxiques ?
    \begin{itemize}
        \item[A)] Link Report
        \item[B)] Disavow (Désaveu de liens)
        \item[C)] Manual Actions
        \item[D)] Spam Report
    \end{itemize}
    *(Réponse : B)*

    \item Le rapport de couverture GSC permet de voir le nombre de pages :
    \begin{itemize}
        \item[A)] Avec le plus de trafic
        \item[B)] Indexées, exclues et avec erreurs
        \item[C)] Les plus récentes
        \item[D)] Avec le plus de commentaires
    \end{itemize}
    *(Réponse : B)*

    \item GSC est utile pour suivre l'évolution du CTR des pages dans les SERP. Cette affirmation est-elle vraie ?
    \begin{itemize}
        \item[A)] Non
        \item[B)] Oui, via le rapport Performance
        \item[C)] Uniquement pour les annonces payantes
        \item[D)] Uniquement pour les images
    \end{itemize}
    *(Réponse : B)*

    \item Quel rapport GSC est utilisé pour vérifier les erreurs d'AMP ?
    \begin{itemize}
        \item[A)] Rapport Performance
        \item[B)] Rapport AMP (si utilisé)
        \item[C)] Rapport Couverture
        \item[D)] Rapport Core Web Vitals
    \end{itemize}
    *(Réponse : B)*

    \item Que permet de faire l'outil d'inspection d'URL dans GSC ?
    \begin{itemize}
        \item[A)] Créer une nouvelle page
        \item[B)] Voir si une URL spécifique est indexée et son état
        \item[C)] Supprimer une page du site
        \item[D)] Modifier le contenu d'une page
    \end{itemize}
    *(Réponse : B)*

    \item GSC peut être configuré de deux manières : par domaine ou par :
    \begin{itemize}
        \item[A)] Par sous-domaine uniquement
        \item[B)] Préfixe d'URL
        \item[C)] Par dossier
        \item[D)] Par fichier
    \end{itemize}
    *(Réponse : B)*

    \item Le rapport Performance de GSC permet d'exporter les données. Cette affirmation est-elle vraie ?
    \begin{itemize}
        \item[A)] Non
        \item[B)] Oui
        \item[C)] Uniquement en PDF
        \item[D)] Uniquement en JSON
    \end{itemize}
    *(Réponse : B)*

    \item GSC permet d'analyser les performances par type de recherche (web, image, vidéo). Cette affirmation est-elle vraie ?
    \begin{itemize}
        \item[A)] Non
        \item[B)] Oui, via le filtre "Type de recherche" dans le rapport Performance
        \item[C)] Uniquement pour les vidéos
        \item[D)] Uniquement pour les images
    \end{itemize}
    *(Réponse : B)*

    \item Le rapport "Pages mobiles" dans GSC vérifie :
    \begin{itemize}
        \item[A)] La vitesse du site sur mobile
        \item[B)] Les problèmes d'utilisabilité mobile
        \item[C)] Le design mobile
        \item[D)] Les applications mobiles
    \end{itemize}
    *(Réponse : B)*

    \item GSC est connecté directement à l'index de Google. Cela signifie que les données sont :
    \begin{itemize}
        \item[A)] Estimées
        \item[B)] Officielles et proviennent de Google
        \item[C)] Fournies par des tiers
        \item[D)] Anonymisées
    \end{itemize}
    *(Réponse : B)*

    \item Quel rapport GSC permet de voir le nombre total de pages indexées ?
    \begin{itemize}
        \item[A)] Performance
        \item[B)] Couverture (onglet "Toutes les pages indexées")
        \item[C)] Core Web Vitals
        \item[D)] Sitemaps
    \end{itemize}
    *(Réponse : B)*

    \item GSC permet de recevoir des notifications par email. Cette affirmation est-elle vraie ?
    \begin{itemize}
        \item[A)] Non
        \item[B)] Oui
        \item[C)] Uniquement pour les comptes premium
        \item[D)] Uniquement pour les erreurs critiques
    \end{itemize}
    *(Réponse : B)*

    \item Dans GSC, on peut configurer le pays cible pour le SEO international. Cette affirmation est-elle vraie ?
    \begin{itemize}
        \item[A)] Non
        \item[B)] Oui, dans les paramètres du site
        \item[C)] Uniquement pour les domaines génériques
        \item[D)] Uniquement pour les sous-répertoires
    \end{itemize}
    *(Réponse : B)*

    \item Quel rapport GSC permet de voir comment Googlebot voit vos pages ?
    \begin{itemize}
        \item[A)] Rapport Performance
        \item[B)] Outil d'inspection d'URL (rendu Google)
        \item[C)] Rapport Core Web Vitals
        \item[D)] Rapport Couverture
    \end{itemize}
    *(Réponse : B)*

    \item GSC est utile pour diagnostiquer les problèmes de :
    \begin{itemize}
        \item[A)] Design responsive
        \item[B)] Référencement, indexation et performances techniques
        \item[C)] Marketing par email
        \item[D)] Réseaux sociaux
    \end{itemize}
    *(Réponse : B)*

    \item Le rapport Liens dans GSC montre les sites qui pointent le plus vers le vôtre. Cette affirmation est-elle vraie ?
    \begin{itemize}
        \item[A)] Non
        \item[B)] Oui
        \item[C)] Uniquement pour les sites payants
        \item[D)] Uniquement pour les sites en HTTPS
    \end{itemize}
    *(Réponse : B)*
\end{enumerate}
\section*{Chapitre 32 : Google Analytics 4 (GA4)}
\begin{enumerate}
    \item GA4 est la dernière version de Google Analytics, basée sur :
    \begin{itemize}
        \item[A)] Les sessions
        \item[B)] Les événements
        \item[C)] Les pages vues
        \item[D)] Les hits
    \end{itemize}
    *(Réponse : B)*

    \item GA4 remplace Universal Analytics depuis quand ?
    \begin{itemize}
        \item[A)] Janvier 2022
        \item[B)] Juillet 2023
        \item[C)] Décembre 2024
        \item[D)] Mars 2020
    \end{itemize}
    *(Réponse : B)*

    \item Quelle métrique GA4 mesure le pourcentage de sessions engagées (10+ secondes, 2+ pages, conversion) ?
    \begin{itemize}
        \item[A)] Taux de rebond
        \item[B)] Engagement rate
        \item[C)] Pages vues par session
        \item[D)] Durée moyenne de session
    \end{itemize}
    *(Réponse : B)*

    \item Comment GA4 définit-il le taux de rebond adapté ?
    \begin{itemize}
        \item[A)] Sessions de 0 seconde
        \item[B)] Sessions non engagées
        \item[C)] Sessions avec 1 page vue
        \item[D)] Sessions sans conversion
    \end{itemize}
    *(Réponse : B)*

    \item Comment GA4 appelle-t-il le trafic provenant des moteurs de recherche ?
    \begin{itemize}
        \item[A)] Trafic direct
        \item[B)] Sessions organiques
        \item[C)] Trafic social
        \item[D)] Trafic email
    \end{itemize}
    *(Réponse : B)*

    \item Quel outil doit être connecté à GA4 pour obtenir les données par mot-clé ?
    \begin{itemize}
        \item[A)] Google Ads
        \item[B)] Google Search Console
        \item[C)] Google Tag Manager
        \item[D)] Google Optimize
    \end{itemize}
    *(Réponse : B)*

    \item Qu'est-ce que l'engagement rate dans GA4 ?
    \begin{itemize}
        \item[A)] Le nombre de clics
        \item[B)] Le pourcentage de sessions engagées
        \item[C)] Le nombre de pages vues
        \item[D)] Le temps de chargement
    \end{itemize}
    *(Réponse : B)*

    \item GA4 utilise un modèle de données basé sur les événements. Cela signifie que :
    \begin{itemize}
        \item[A)] Tout est mesuré en sessions
        \item[B)] Chaque interaction est un événement
        \item[C)] Seules les pages vues sont mesurées
        \item[D)] Les hits sont la seule unité de mesure
    \end{itemize}
    *(Réponse : B)*

    \item Dans GA4, une session engagée est définie comme :
    \begin{itemize}
        \item[A)] Une session de moins de 10 secondes
        \item[B)] Une session de 10+ secondes, 2+ pages, ou avec conversion
        \item[C)] Une session avec une seule page vue
        \item[D)] Une session sans interaction
    \end{itemize}
    *(Réponse : B)*

    \item Les sessions organiques dans GA4 représentent :
    \begin{itemize}
        \item[A)] Le trafic provenant des emails
        \item[B)] Le trafic provenant des moteurs de recherche
        \item[C)] Le trafic provenant des réseaux sociaux
        \item[D)] Le trafic provenant des publicités
    \end{itemize}
    *(Réponse : B)*

    \item Les conversions dans GA4 sont définies comme :
    \begin{itemize}
        \item[A)] Toute page vue
        \item[B)] Des objectifs atteints (achat, inscription, téléchargement)
        \item[C)] Des clics sur les liens
        \item[D)] Des sessions de moins de 10 secondes
    \end{itemize}
    *(Réponse : B)*

    \item La profondeur de navigation est mesurée par :
    \begin{itemize}
        \item[A)] Le nombre de clics
        \item[B)] Les pages vues par session
        \item[C)] Le temps passé
        \item[D)] Le nombre de sessions
    \end{itemize}
    *(Réponse : B)*

    \item Les rapports SEO dans GA4 nécessitent-ils une configuration spécifique ?
    \begin{itemize}
        \item[A)] Non, ils sont automatiques
        \item[B)] Oui, ils nécessitent une configuration spécifique
        \item[C)] Uniquement pour le e-commerce
        \item[D)] Uniquement pour les sites mobiles
    \end{itemize}
    *(Réponse : B)*

    \item Qu'est-ce qui remplace les métriques de session traditionnelles dans GA4 ?
    \begin{itemize}
        \item[A)] Les hits
        \item[B)] L'engagement
        \item[C)] Les pages vues
        \item[D)] Les clics
    \end{itemize}
    *(Réponse : B)*

    \item Quel type de données GA4 peut importer depuis GSC ?
    \begin{itemize}
        \item[A)] Les données démographiques
        \item[B)] Les données de performance par mot-clé (impressions, clics, CTR)
        \item[C)] Les données de vente
        \item[D)] Les données de température
    \end{itemize}
    *(Réponse : B)*

    \item GA4 permet de mesurer les acquisitions par :
    \begin{itemize}
        \item[A)] Uniquement le trafic direct
        \item[B)] Plusieurs canaux (organique, payant, social, email, etc.)
        \item[C)] Uniquement le trafic organique
        \item[D)] Uniquement les publicités
    \end{itemize}
    *(Réponse : B)*

    \item Dans GA4, l'engagement rate remplace en partie :
    \begin{itemize}
        \item[A)] Le nombre de visiteurs
        \item[B)] Le taux de rebond traditionnel
        \item[C)] Le nombre de pages vues
        \item[D)] Le trafic direct
    \end{itemize}
    *(Réponse : B)*

    \item GA4 peut être connecté à Google Search Console pour :
    \begin{itemize}
        \item[A)] Voir les backlinks
        \item[B)] Enrichir les données SEO avec les impressions et clics
        \item[C)] Modifier le robots.txt
        \item[D)] Soumettre des sitemaps
    \end{itemize}
    *(Réponse : B)*

    \item Quelle métrique GA4 est utile pour mesurer la qualité du SEO ?
    \begin{itemize}
        \item[A)] Le taux de clics publicitaires
        \item[B)] L'engagement rate des sessions organiques
        \item[C)] Le nombre d'emails envoyés
        \item[D)] Le nombre d'utilisateurs connectés
    \end{itemize}
    *(Réponse : B)*

    \item GA4 est basé sur les événements plutôt que sur les sessions. Cela permet :
    \begin{itemize}
        \item[A)] Moins de données
        \item[B)] Une mesure plus flexible et précise des interactions utilisateur
        \item[C)] Uniquement le suivi des pages
        \item[D)] Une configuration plus simple
    \end{itemize}
    *(Réponse : B)*

    \item Les pages vues par session mesurent :
    \begin{itemize}
        \item[A)] La vitesse du site
        \item[B)] La profondeur de navigation
        \item[C)] Le nombre de visiteurs
        \item[D)] Le taux de conversion
    \end{itemize}
    *(Réponse : B)*

    \item Dans GA4, un événement peut être :
    \begin{itemize}
        \item[A)] Uniquement un clic
        \item[B)] N'importe quelle interaction : clic, scroll, achat, etc.
        \item[C)] Uniquement une page vue
        \item[D)] Uniquement un formulaire
    \end{itemize}
    *(Réponse : B)*

    \item GA4 remplace Universal Analytics. Cette affirmation est-elle vraie ?
    \begin{itemize}
        \item[A)] Non, ils co-existent
        \item[B)] Oui, depuis juillet 2023
        \item[C)] Non, Universal Analytics est toujours le standard
        \item[D)] Uniquement pour les sites e-commerce
    \end{itemize}
    *(Réponse : B)*

    \item Quel est l'avantage de GA4 par rapport à Universal Analytics pour le SEO ?
    \begin{itemize}
        \item[A)] Il est plus lent
        \item[B)] Il offre un meilleur suivi cross-platform et basé sur les événements
        \item[C)] Il est moins cher
        \item[D)] Il ne nécessite pas de configuration
    \end{itemize}
    *(Réponse : B)*

    \item Les conversions dans GA4 correspondent à :
    \begin{itemize}
        \item[A)] Toute session
        \item[B)] Des événements marqués comme conversions
        \item[C)] Uniquement les achats
        \item[D)] Uniquement les inscriptions
    \end{itemize}
    *(Réponse : B)*

    \item GA4 permet de suivre le comportement des utilisateurs :
    \begin{itemize}
        \item[A)] Uniquement sur un seul appareil
        \item[B)] Sur plusieurs appareils et plateformes
        \item[C)] Uniquement sur mobile
        \item[D)] Uniquement sur desktop
    \end{itemize}
    *(Réponse : B)*

    \item L'engagement rate de GA4 est calculé comme :
    \begin{itemize}
        \item[A)] Sessions engagées / nombre total d'utilisateurs
        \item[B)] Sessions engagées / nombre total de sessions
        \item[C)] Pages vues / sessions
        \item[D)] Conversions / sessions
    \end{itemize}
    *(Réponse : B)*

    \item GA4 utilise le machine learning pour :
    \begin{itemize}
        \item[A)] Créer du contenu
        \item[B)] Combler les lacunes de données et fournir des insights
        \item[C)] Modifier le site
        \item[D)] Envoyer des emails
    \end{itemize}
    *(Réponse : B)*

    \item Quel est l'équivalent GA4 du "taux de rebond" Universal Analytics ?
    \begin{itemize}
        \item[A)] Il n'y a pas d'équivalent
        \item[B)] Le pourcentage de sessions non engagées
        \item[C)] Le nombre de pages vues
        \item[D)] La durée moyenne de session
    \end{itemize}
    *(Réponse : B)*

    \item GA4 est particulièrement utile pour le SEO car il permet de :
    \begin{itemize}
        \item[A)] Créer des backlinks
        \item[B)] Mesurer l'impact du trafic organique sur les conversions
        \item[C)] Modifier les balises title
        \item[D)] Soumettre des sitemaps
    \end{itemize}
    *(Réponse : B)*

    \item Dans GA4, comment s'appelle la métrique qui remplace les "visiteurs uniques" ?
    \begin{itemize}
        \item[A)] Sessions
        \item[B)] Utilisateurs actifs
        \item[C)] Hits
        \item[D)] Pages vues
    \end{itemize}
    *(Réponse : B)*

    \item GA4 peut être intégré à BigQuery. Cette affirmation est-elle vraie ?
    \begin{itemize}
        \item[A)] Non
        \item[B)] Oui
        \item[C)] Uniquement pour les comptes 360
        \item[D)] Uniquement pour les sites e-commerce
    \end{itemize}
    *(Réponse : B)*

    \item Quel type de rapport GA4 est utile pour analyser le trafic SEO ?
    \begin{itemize}
        \item[A)] Rapports en temps réel uniquement
        \item[B)] Rapport "Acquisition" > "Acquisition de trafic"
        \item[C)] Rapports démographiques uniquement
        \item[D)] Rapports technologiques uniquement
    \end{itemize}
    *(Réponse : B)*

    \item Les sessions organiques dans GA4 peuvent être filtrées par page de destination. Cette affirmation est-elle vraie ?
    \begin{itemize}
        \item[A)] Non
        \item[B)] Oui
        \item[C)] Uniquement pour les pages d'accueil
        \item[D)] Uniquement pour les fiches produits
    \end{itemize}
    *(Réponse : B)*

    \item GA4 permet de créer des audiences pour le remarketing. Cette affirmation est-elle vraie ?
    \begin{itemize}
        \item[A)] Non
        \item[B)] Oui
        \item[C)] Uniquement pour Google Ads
        \item[D)] Uniquement pour les sites e-commerce
    \end{itemize}
    *(Réponse : B)*

    \item L'engagement rate est une métrique plus pertinente que le taux de rebond pour :
    \begin{itemize}
        \item[A)] Mesurer la popularité du site
        \item[B)] Mesurer l'engagement réel des utilisateurs
        \item[C)] Compter les visiteurs
        \item[D)] Calculer le chiffre d'affaires
    \end{itemize}
    *(Réponse : B)*

    \item GA4 permet d'exporter les données vers Google Sheets. Cette affirmation est-elle vraie ?
    \begin{itemize}
        \item[A)] Non
        \item[B)] Oui
        \item[C)] Uniquement en CSV
        \item[D)] Uniquement via BigQuery
    \end{itemize}
    *(Réponse : B)*

    \item Quel est l'identifiant principal utilisé par GA4 pour suivre les utilisateurs ?
    \begin{itemize}
        \item[A)] L'adresse IP
        \item[B)] Le Client ID (stocké dans un cookie first-party)
        \item[C)] L'email
        \item[D)] Le numéro de téléphone
    \end{itemize}
    *(Réponse : B)*

    \item GA4 mesure les événements automatiquement. Lesquels sont collectés sans configuration ?
    \begin{itemize}
        \item[A)] Aucun
        \item[B)] page\_view, scroll, click, first\_visit, session\_start, etc.
        \item[C)] Uniquement page\_view
        \item[D)] Uniquement les achats
    \end{itemize}
    *(Réponse : B)*

    \item GA4 propose des "Explorations" qui permettent :
    \begin{itemize}
        \item[A)] De modifier le site web
        \item[B)] De créer des rapports personnalisés avancés
        \item[C)] D'envoyer des newsletters
        \item[D)] De gérer les utilisateurs
    \end{itemize}
    *(Réponse : B)*

    \item GA4 est gratuit pour tous les sites web. Cette affirmation est-elle vraie ?
    \begin{itemize}
        \item[A)] Non, il est payant
        \item[B)] Oui, il est gratuit
        \item[C)] Gratuit jusqu'à 10 millions de hits par mois
        \item[D)] Uniquement pour les petits sites
    \end{itemize}
    *(Réponse : B)*

    \item Dans GA4, le modèle d'attribution par défaut est :
    \begin{itemize}
        \item[A)] Premier clic
        \item[B)] Data-driven (basé sur les données)
        \item[C)] Dernier clic
        \item[D)] Linéaire
    \end{itemize}
    *(Réponse : B)*

    \item GA4 permet de suivre les téléchargements de fichiers comme événements. Cette affirmation est-elle vraie ?
    \begin{itemize}
        \item[A)] Non
        \item[B)] Oui, via la mesure améliorée
        \item[C)] Uniquement pour les PDF
        \item[D)] Uniquement pour les images
    \end{itemize}
    *(Réponse : B)*

    \item GA4 peut être utilisé sans cookie wall. Cette affirmation est-elle vraie ?
    \begin{itemize}
        \item[A)] Non, les cookies sont obligatoires
        \item[B)] Oui, avec des modes de consentement
        \item[C)] Non, GA4 ne fonctionne qu'avec des cookies tiers
        \item[D)] Uniquement dans l'UE
    \end{itemize}
    *(Réponse : B)*

    \item Quel est l'avantage de GA4 pour les sites avec plusieurs domaines ?
    \begin{itemize}
        \item[A)] Aucun avantage spécifique
        \item[B)] Le suivi cross-domain est intégré
        \item[C)] Il faut un compte séparé par domaine
        \item[D)] GA4 ne supporte pas plusieurs domaines
    \end{itemize}
    *(Réponse : B)*

    \item GA4 mesure le trafic en temps réel. Cette affirmation est-elle vraie ?
    \begin{itemize}
        \item[A)] Non
        \item[B)] Oui, avec un rapport en temps réel
        \item[C)] Uniquement avec un délai de 24h
        \item[D)] Uniquement pour les événements de conversion
    \end{itemize}
    *(Réponse : B)*

    \item GA4 est compatible avec Google Tag Manager. Cette affirmation est-elle vraie ?
    \begin{itemize}
        \item[A)] Non
        \item[B)] Oui
        \item[C)] Uniquement pour le remarketing
        \item[D)] Uniquement pour les sites e-commerce
    \end{itemize}
    *(Réponse : B)*

    \item Les "Explorations" dans GA4 remplacent les anciens :
    \begin{itemize}
        \item[A)] Tableaux de bord
        \item[B)] Segments avancés et rapports personnalisés d'Universal Analytics
        \item[C)] Entonnoirs de conversion
        \item[D)] Alertes personnalisées
    \end{itemize}
    *(Réponse : B)*

    \item GA4 nécessite-t-il un marquage spécifique pour le suivi des clics sur les liens externes ?
    \begin{itemize}
        \item[A)] Oui, manuellement
        \item[B)] Non, la mesure améliorée le fait automatiquement
        \item[C)] Uniquement pour les liens d'affiliation
        \item[D)] Uniquement pour les liens sortants
    \end{itemize}
    *(Réponse : B)*

    \item GA4 permet de suivre les interactions des utilisateurs sans modification du code grâce à :
    \begin{itemize}
        \item[A)] Google Tag Manager uniquement
        \item[B)] La mesure améliorée (enhanced measurement)
        \item[C)] Uniquement les événements personnalisés
        \item[D)] Google Ads
    \end{itemize}
    *(Réponse : B)*
\end{enumerate}
\section*{Chapitre 33 : Les KPIs SEO et le Reporting}
\begin{enumerate}
    \item Quelle métrique mesure le nombre de sessions provenant des moteurs de recherche ?
    \begin{itemize}
        \item[A)] Position moyenne
        \item[B)] Trafic organique
        \item[C)] Domain Authority
        \item[D)] Core Web Vitals
    \end{itemize}
    *(Réponse : B)*

    \item Quelle métrique indique la position d'un site par mot-clé dans les SERP ?
    \begin{itemize}
        \item[A)] Trafic organique
        \item[B)] Classement (position moyenne par mot-clé)
        \item[C)] Taux de conversion
        \item[D)] Backlinks
    \end{itemize}
    *(Réponse : B)*

    \item Quel indicateur mesure le nombre d'objectifs atteints divisé par le nombre de sessions ?
    \begin{itemize}
        \item[A)] Trafic organique
        \item[B)] Taux de conversion
        \item[C)] CTR
        \item[D)] Engagement
    \end{itemize}
    *(Réponse : B)*

    \item Quelle métrique mesure l'autorité d'un domaine ?
    \begin{itemize}
        \item[A)] Trafic organique
        \item[B)] Domain Authority, Trust Flow, Citation Flow
        \item[C)] Taux de conversion
        \item[D)] Pages indexées
    \end{itemize}
    *(Réponse : B)*

    \item Quel indicateur mesure le nombre de pages indexées ?
    \begin{itemize}
        \item[A)] Trafic organique
        \item[B)] Indexation (pages indexées, erreurs de crawl)
        \item[C)] CTR
        \item[D)] Backlinks
    \end{itemize}
    *(Réponse : B)*

    \item Les Core Web Vitals mesurent :
    \begin{itemize}
        \item[A)] Le nombre de visiteurs
        \item[B)] LCP, INP, CLS
        \item[C)] Le nombre de pages
        \item[D)] Le nombre de backlinks
    \end{itemize}
    *(Réponse : B)*

    \item Quelle métrique mesure le nombre de domaines référents ?
    \begin{itemize}
        \item[A)] Trafic organique
        \item[B)] Backlinks (nombre de domaines référents)
        \item[C)] CTR
        \item[D)] Position moyenne
    \end{itemize}
    *(Réponse : B)*

    \item Le CTR est le rapport entre :
    \begin{itemize}
        \item[A)] Le nombre de sessions et le nombre d'utilisateurs
        \item[B)] Le nombre de clics et le nombre d'impressions dans les SERP
        \item[C)] Le nombre de pages vues et le nombre de sessions
        \item[D)] Le nombre de conversions et le nombre de sessions
    \end{itemize}
    *(Réponse : B)*

    \item L'engagement dans le tableau de bord SEO mesure :
    \begin{itemize}
        \item[A)] Le nombre de backlinks
        \item[B)] Le temps passé et les pages par session
        \item[C)] Le nombre d'impressions
        \item[D)] La position moyenne
    \end{itemize}
    *(Réponse : B)*

    \item Les KPIs SEO se divisent en quelles catégories principales ?
    \begin{itemize}
        \item[A)] Trafic, design, développement
        \item[B)] Trafic, technique, autorité
        \item[C)] Trafic, email, réseaux sociaux
        \item[D)] Design, marketing, ventes
    \end{itemize}
    *(Réponse : B)*

    \item Un tableau de bord SEO doit être adapté à :
    \begin{itemize}
        \item[A)] Tous les sites de la même manière
        \item[B)] Les objectifs business spécifiques
        \item[C)] Uniquement au trafic
        \item[D)] Uniquement aux backlinks
    \end{itemize}
    *(Réponse : B)*

    \item Le reporting régulier permet de :
    \begin{itemize}
        \item[A)] Remplir des obligations administratives
        \item[B)] Ajuster la stratégie SEO
        \item[C)] Augmenter le trafic automatiquement
        \item[D)] Réduire les coûts
    \end{itemize}
    *(Réponse : B)*

    \item Le Trust Flow est une métrique propriétaire de :
    \begin{itemize}
        \item[A)] Moz
        \item[B)] Majestic
        \item[C)] Ahrefs
        \item[D)] SEMrush
    \end{itemize}
    *(Réponse : B)*

    \item Le Domain Authority (DA) est une métrique propriétaire de :
    \begin{itemize}
        \item[A)] Majestic
        \item[B)] Moz
        \item[C)] Ahrefs
        \item[D)] SEMrush
    \end{itemize}
    *(Réponse : B)*

    \item Le Citation Flow est une métrique propriétaire de :
    \begin{itemize}
        \item[A)] Moz
        \item[B)] Majestic
        \item[C)] Ahrefs
        \item[D)] Google
    \end{itemize}
    *(Réponse : B)*

    \item Combien de catégories de KPIs sont présentées dans le tableau de bord SEO ?
    \begin{itemize}
        \item[A)] 5
        \item[B)] 9 (trafic, classement, conversion, autorité, indexation, Core Web Vitals, backlinks, CTR, engagement)
        \item[C)] 3
        \item[D)] 12
    \end{itemize}
    *(Réponse : B)*

    \item Quelle métrique de trafic est mentionnée dans les KPIs ?
    \begin{itemize}
        \item[A)] Uniquement les sessions
        \item[B)] Sessions, utilisateurs, pages vues
        \item[C)] Uniquement les utilisateurs
        \item[D)] Uniquement les pages vues
    \end{itemize}
    *(Réponse : B)*

    \item Le reporting SEO doit être effectué à quelle fréquence ?
    \begin{itemize}
        \item[A)] Une fois par an
        \item[B)] Régulièrement (mensuel/hebdomadaire selon les besoins)
        \item[C)] Jamais
        \item[D)] Une seule fois
    \end{itemize}
    *(Réponse : B)*

    \item Les Core Web Vitals sont une catégorie de KPIs pour mesurer :
    \begin{itemize}
        \item[A)] Le contenu
        \item[B)] La performance technique et l'expérience utilisateur
        \item[C)] Les backlinks
        \item[D)] Le trafic
    \end{itemize}
    *(Réponse : B)*

    \item Quelle métrique est utile pour mesurer l'autorité d'un site ?
    \begin{itemize}
        \item[A)] Le CTR
        \item[B)] Le Domain Authority
        \item[C)] Les pages vues
        \item[D)] Le taux de conversion
    \end{itemize}
    *(Réponse : B)*

    \item Le nombre de backlinks est un indicateur de :
    \begin{itemize}
        \item[A)] Trafic
        \item[B)] Popularité et autorité
        \item[C)] Performance technique
        \item[D)] Taux de conversion
    \end{itemize}
    *(Réponse : B)*

    \item Un bon reporting SEO doit inclure :
    \begin{itemize}
        \item[A)] Uniquement les chiffres positifs
        \item[B)] Les progrès, les baisses et les recommandations
        \item[C)] Uniquement les résultats
        \item[D)] Uniquement les prévisions
    \end{itemize}
    *(Réponse : B)*

    \item Quelle est la première étape pour construire un tableau de bord SEO ?
    \begin{itemize}
        \item[A)] Choisir un outil
        \item[B)] Définir les objectifs business
        \item[C)] Collecter les données
        \item[D)] Créer des graphiques
    \end{itemize}
    *(Réponse : B)*

    \item Les KPIs SEO doivent être :
    \begin{itemize}
        \item[A)] Identiques pour tous les sites
        \item[B)] Adaptés aux objectifs de chaque site
        \item[C)] Uniquement basés sur le trafic
        \item[D)] Uniquement basés sur les backlinks
    \end{itemize}
    *(Réponse : B)*

    \item Le taux de clic (CTR) est influencé par :
    \begin{itemize}
        \item[A)] Le nombre de pages
        \item[B)] Le title tag et la meta description
        \item[C)] Le nombre de backlinks
        \item[D)] La vitesse du site
    \end{itemize}
    *(Réponse : B)*

    \item Le reporting SEO permet d'ajuster la stratégie. Cette affirmation est-elle vraie ?
    \begin{itemize}
        \item[A)] Non
        \item[B)] Oui
        \item[C)] Uniquement pour les grandes entreprises
        \item[D)] Uniquement pour le e-commerce
    \end{itemize}
    *(Réponse : B)*

    \item Le nombre de pages indexées est un KPI de quelle catégorie ?
    \begin{itemize}
        \item[A)] Trafic
        \item[B)] Indexation / technique
        \item[C)] Autorité
        \item[D)] Conversion
    \end{itemize}
    *(Réponse : B)*

    \item Les nouveaux backlinks et les backlinks perdus sont des KPIS de :
    \begin{itemize}
        \item[A)] Trafic
        \item[B)] Autorité et netlinking
        \item[C)] Performance
        \item[D)] Engagement
    \end{itemize}
    *(Réponse : B)*

    \item Le temps passé sur le site est un indicateur de :
    \begin{itemize}
        \item[A)] Trafic
        \item[B)] Engagement
        \item[C)] Autorité
        \item[D)] Technique
    \end{itemize}
    *(Réponse : B)*

    \item Les KPIs SEO sont essentiels pour :
    \begin{itemize}
        \item[A)] Ignorer les problèmes
        \item[B)] Mesurer l'efficacité des actions SEO
        \item[C)] Remplacer le contenu
        \item[D)] Supprimer les pages
    \end{itemize}
    *(Réponse : B)*

    \item Quelle est la différence entre le trafic organique et le trafic total ?
    \begin{itemize}
        \item[A)] Il n'y a pas de différence
        \item[B)] Le trafic organique vient des moteurs de recherche, le trafic total inclut toutes les sources
        \item[C)] Le trafic organique est payant
        \item[D)] Le trafic total exclut le SEO
    \end{itemize}
    *(Réponse : B)*

    \item Le suivi des positions par mot-clé est important pour mesurer :
    \begin{itemize}
        \item[A)] La vitesse du site
        \item[B)] L'efficacité du référencement
        \item[C)] Le design
        \item[D)] Le nombre de pages
    \end{itemize}
    *(Réponse : B)*

    \item L'analyse du taux de conversion SEO permet de mesurer :
    \begin{itemize}
        \item[A)] La popularité
        \item[B)] La performance business du trafic organique
        \item[C)] Le nombre de pages indexées
        \item[D)] La vitesse du site
    \end{itemize}
    *(Réponse : B)*

    \item Les erreurs de crawl sont un KPI de quelle catégorie ?
    \begin{itemize}
        \item[A)] Autorité
        \item[B)] Indexation / technique
        \item[C)] Engagement
        \item[D)] Trafic
    \end{itemize}
    *(Réponse : B)*

    \item Quel outil est principalement utilisé pour suivre le CTR dans les SERP ?
    \begin{itemize}
        \item[A)] Google Analytics
        \item[B)] Google Search Console
        \item[C)] Google Ads
        \item[D)] Google Tag Manager
    \end{itemize}
    *(Réponse : B)*

    \item Un tableau de bord SEO doit être mis à jour :
    \begin{itemize}
        \item[A)] Une fois par an
        \item[B)] Régulièrement (mensuel recommandé)
        \item[C)] Jamais
        \item[D)] Uniquement en cas de problème
    \end{itemize}
    *(Réponse : B)*

    \item Le Domain Authority (DA) est un score de :
    \begin{itemize}
        \item[A)] 0 à 10
        \item[B)] 1 à 100
        \item[C)] 0 à 1000
        \item[D)] 1 à 10
    \end{itemize}
    *(Réponse : B)*

    \item Quel KPI mesure l'impact des Core Web Vitals sur l'expérience utilisateur ?
    \begin{itemize}
        \item[A)] Trafic organique
        \item[B)] LCP, INP, CLS
        \item[C)] Backlinks
        \item[D)] Taux de conversion
    \end{itemize}
    *(Réponse : B)*

    \item Les pages par session sont un indicateur de :
    \begin{itemize}
        \item[A)] Popularité
        \item[B)] Engagement et profondeur de navigation
        \item[C)] Performance
        \item[D)] Indexation
    \end{itemize}
    *(Réponse : B)*

    \item Un taux de conversion élevé signifie que le trafic est :
    \begin{itemize}
        \item[A)] De faible qualité
        \item[B)] Pertinent et bien ciblé
        \item[C)] Aléatoire
        \item[D)] Peu important
    \end{itemize}
    *(Réponse : B)*

    \item Le Trust Flow mesure :
    \begin{itemize}
        \item[A)] La quantité de backlinks
        \item[B)] La qualité des backlinks basée sur la proximité avec des sites de confiance
        \item[C)] Le nombre de pages
        \item[D)] La vitesse du site
    \end{itemize}
    *(Réponse : B)*

    \item Le reporting SEO implique de communiquer les résultats à :
    \begin{itemize}
        \item[A)] Personne
        \item[B)] Aux parties prenantes (clients, direction, équipe)
        \item[C)] Uniquement aux concurrents
        \item[D)] Uniquement à Google
    \end{itemize}
    *(Réponse : B)*

    \item Quelle est la différence entre Trust Flow et Citation Flow ?
    \begin{itemize}
        \item[A)] Aucune différence
        \item[B)] Trust Flow mesure la qualité, Citation Flow mesure la quantité de backlinks
        \item[C)] Ce sont les mêmes métriques
        \item[D)] Citation Flow mesure la qualité
    \end{itemize}
    *(Réponse : B)*

    \item L'analyse du taux de rebond adapté GA4 est un KPI de :
    \begin{itemize}
        \item[A)] Trafic
        \item[B)] Engagement
        \item[C)] Autorité
        \item[D)] Technique
    \end{itemize}
    *(Réponse : B)*

    \item Un bon KPI SEO doit être :
    \begin{itemize}
        \item[A)] Complexe à mesurer
        \item[B)] Spécifique, mesurable, atteignable, pertinent et temporel (SMART)
        \item[C)] Aléatoire
        \item[D)] Impossible à atteindre
    \end{itemize}
    *(Réponse : B)*

    \item Les KPIs SEO sont utiles pour justifier le budget SEO. Cette affirmation est-elle vraie ?
    \begin{itemize}
        \item[A)] Non
        \item[B)] Oui
        \item[C)] Uniquement pour les petites entreprises
        \item[D)] Uniquement pour les agences
    \end{itemize}
    *(Réponse : B)*

    \item Le nombre de mots-clés en top 3 est un KPI de :
    \begin{itemize}
        \item[A)] Trafic
        \item[B)] Classement / performance SERP
        \item[C)] Technique
        \item[D)] Engagement
    \end{itemize}
    *(Réponse : B)*

    \item Le taux de rebond adapté GA4 est calculé comme :
    \begin{itemize}
        \item[A)] Le nombre de sessions non engagées
        \item[B)] Le pourcentage de sessions non engagées
        \item[C)] Le nombre de pages vues
        \item[D)] Le nombre de clics
    \end{itemize}
    *(Réponse : B)*

    \item Les KPIs SEO permettent de prioriser les actions à entreprendre. Cette affirmation est-elle vraie ?
    \begin{itemize}
        \item[A)] Non
        \item[B)] Oui
        \item[C)] Uniquement pour les backlinks
        \item[D)] Uniquement pour le contenu
    \end{itemize}
    *(Réponse : B)*

    \item Pour construire un tableau de bord SEO efficace, il faut d'abord :
    \begin{itemize}
        \item[A)] Choisir un outil de reporting
        \item[B)] Définir les objectifs business et les KPIs alignés
        \item[C)] Collecter toutes les données disponibles
        \item[D)] Créer des graphiques
    \end{itemize}
    *(Réponse : B)*
\end{enumerate}
\section*{Chapitre 34 : L'Audit SEO Complet}
\begin{enumerate}
    \item Combien de phases comprend un audit SEO complet selon la méthodologie présentée ?
    \begin{itemize}
        \item[A)] 2 phases
        \item[B)] 3 phases (technique, On-Page, Off-Page)
        \item[C)] 5 phases
        \item[D)] 1 phase
    \end{itemize}
    *(Réponse : B)*

    \item Que examine la Phase 1 de l'audit SEO ?
    \begin{itemize}
        \item[A)] Le contenu
        \item[B)] L'audit technique (crawling, indexation, Core Web Vitals)
        \item[C)] Les backlinks
        \item[D)] La concurrence
    \end{itemize}
    *(Réponse : B)*

    \item Que vérifie-t-on dans l'analyse du crawling ?
    \begin{itemize}
        \item[A)] Le design
        \item[B)] Robots.txt, sitemap, structure
        \item[C)] Les images
        \item[D)] Les vidéos
    \end{itemize}
    *(Réponse : B)*

    \item Quel outil est mentionné pour l'audit technique avec une version gratuite limitée à 500 URLs ?
    \begin{itemize}
        \item[A)] Ahrefs
        \item[B)] Screaming Frog
        \item[C)] SEMrush
        \item[D)] Moz
    \end{itemize}
    *(Réponse : B)*

    \item L'analyse de la vitesse utilise quels outils ?
    \begin{itemize}
        \item[A)] Google Analytics et Google Ads
        \item[B)] PageSpeed Insights et Lighthouse
        \item[C)] Google Search Console et Google My Business
        \item[D)] Google Tag Manager et Google Optimize
    \end{itemize}
    *(Réponse : B)*

    \item La Phase 2 de l'audit SEO est l'audit :
    \begin{itemize}
        \item[A)] Technique
        \item[B)] On-Page
        \item[C)] Off-Page
        \item[D)] Concurrentiel
    \end{itemize}
    *(Réponse : B)*

    \item Que vérifie-t-on dans l'audit On-Page concernant les images ?
    \begin{itemize}
        \item[A)] La couleur
        \item[B)] L'alt text, la taille et le format
        \item[C)] La date de création
        \item[D)] L'auteur
    \end{itemize}
    *(Réponse : B)*

    \item La Phase 3 de l'audit SEO est l'audit :
    \begin{itemize}
        \item[A)] Technique
        \item[B)] On-Page
        \item[C)] Off-Page
        \item[D)] De contenu
    \end{itemize}
    *(Réponse : C)*

    \item Que vérifie-t-on dans l'audit Off-Page ?
    \begin{itemize}
        \item[A)] Les balises title
        \item[B)] Le profil de backlinks et les liens toxiques
        \item[C)] Le robots.txt
        \item[D)] La vitesse du site
    \end{itemize}
    *(Réponse : B)*

    \item Quels sont les trois aspects couverts par un audit SEO complet ?
    \begin{itemize}
        \item[A)] Design, développement, marketing
        \item[B)] Technique, On-Page, Off-Page
        \item[C)] Contenu, images, vidéos
        \item[D)] SEO, SEA, SMO
    \end{itemize}
    *(Réponse : B)*

    \item La vérification du mobile-first est incluse dans quelle phase ?
    \begin{itemize}
        \item[A)] Phase 2 (On-Page)
        \item[B)] Phase 1 (technique)
        \item[C)] Phase 3 (Off-Page)
        \item[D)] Phase 4
    \end{itemize}
    *(Réponse : B)*

    \item L'analyse des données structurées fait partie de :
    \begin{itemize}
        \item[A)] Phase 2 (On-Page)
        \item[B)] Phase 1 (technique)
        \item[C)] Phase 3 (Off-Page)
        \item[D)] Aucune phase
    \end{itemize}
    *(Réponse : B)*

    \item L'analyse du maillage interne est incluse dans quelle phase ?
    \begin{itemize}
        \item[A)] Phase 1 (technique)
        \item[B)] Phase 2 (On-Page)
        \item[C)] Phase 3 (Off-Page)
        \item[D)] Aucune phase
    \end{itemize}
    *(Réponse : B)*

    \item L'analyse de la concurrence est incluse dans quelle phase ?
    \begin{itemize}
        \item[A)] Phase 1 (technique)
        \item[B)] Phase 2 (On-Page)
        \item[C)] Phase 3 (Off-Page)
        \item[D)] Aucune phase
    \end{itemize}
    *(Réponse : C)*

    \item L'évaluation de la qualité du contenu fait partie de :
    \begin{itemize}
        \item[A)] Phase 1 (technique)
        \item[B)] Phase 2 (On-Page)
        \item[C)] Phase 3 (Off-Page)
        \item[D)] Phase 4
    \end{itemize}
    *(Réponse : B)*

    \item L'identification des liens toxiques fait partie de :
    \begin{itemize}
        \item[A)] Phase 1 (technique)
        \item[B)] Phase 2 (On-Page)
        \item[C)] Phase 3 (Off-Page)
        \item[D)] Aucune phase
    \end{itemize}
    *(Réponse : C)*

    \item L'inspection du rendu JavaScript est incluse dans :
    \begin{itemize}
        \item[A)] Phase 2 (On-Page)
        \item[B)] Phase 1 (technique)
        \item[C)] Phase 3 (Off-Page)
        \item[D)] Aucune phase
    \end{itemize}
    *(Réponse : B)*

    \item La vérification des balises title et meta description fait partie de :
    \begin{itemize}
        \item[A)] Phase 1 (technique)
        \item[B)] Phase 2 (On-Page)
        \item[C)] Phase 3 (Off-Page)
        \item[D)] Phase 4
    \end{itemize}
    *(Réponse : B)*

    \item L'analyse des mots-clés par page fait partie de :
    \begin{itemize}
        \item[A)] Phase 1 (technique)
        \item[B)] Phase 2 (On-Page)
        \item[C)] Phase 3 (Off-Page)
        \item[D)] Aucune phase
    \end{itemize}
    *(Réponse : B)*

    \item La priorisation des actions est cruciale après quelle étape ?
    \begin{itemize}
        \item[A)] Avant l'audit
        \item[B)] Après l'audit
        \item[C)] Pendant l'audit
        \item[D)] Aucune priorisation nécessaire
    \end{itemize}
    *(Réponse : B)*

    \item L'audit doit être répété à quelle fréquence ?
    \begin{itemize}
        \item[A)] Une seule fois
        \item[B)] Régulièrement
        \item[C)] Tous les 5 ans
        \item[D)] Jamais
    \end{itemize}
    *(Réponse : B)*

    \item La vérification de robots.txt est une étape de :
    \begin{itemize}
        \item[A)] L'audit On-Page
        \item[B)] L'audit technique
        \item[C)] L'audit Off-Page
        \item[D)] L'audit concurrentiel
    \end{itemize}
    *(Réponse : B)*

    \item L'analyse de la structure des headings (H1, H2, H3) fait partie de :
    \begin{itemize}
        \item[A)] L'audit technique
        \item[B)] L'audit On-Page
        \item[C)] L'audit Off-Page
        \item[D)] L'audit de backlinks
    \end{itemize}
    *(Réponse : B)*

    \item L'analyse du profil de backlinks est une étape de :
    \begin{itemize}
        \item[A)] L'audit technique
        \item[B)] L'audit On-Page
        \item[C)] L'audit Off-Page
        \item[D)] L'audit mobile
    \end{itemize}
    *(Réponse : C)*

    \item L'évaluation de l'autorité du domaine fait partie de :
    \begin{itemize}
        \item[A)] L'audit technique
        \item[B)] L'audit On-Page
        \item[C)] L'audit Off-Page
        \item[D)] L'audit de contenu
    \end{itemize}
    *(Réponse : C)*

    \item La vérification des mentions de marque fait partie de :
    \begin{itemize}
        \item[A)] L'audit technique
        \item[B)] L'audit On-Page
        \item[C)] L'audit Off-Page
        \item[D)] L'audit de vitesse
    \end{itemize}
    *(Réponse : C)*

    \item L'analyse du temps de chargement est une étape de :
    \begin{itemize}
        \item[A)] L'audit technique
        \item[B)] L'audit On-Page
        \item[C)] L'audit Off-Page
        \item[D)] L'audit de backlinks
    \end{itemize}
    *(Réponse : A)*

    \item Le rapport d'audit SEO doit inclure :
    \begin{itemize}
        \item[A)] Uniquement les problèmes
        \item[B)] Constats, recommandations, priorisation
        \item[C)] Uniquement les solutions
        \item[D)] Uniquement le résumé
    \end{itemize}
    *(Réponse : B)*

    \item La vérification du protocole HTTPS fait partie de :
    \begin{itemize}
        \item[A)] L'audit On-Page
        \item[B)] L'audit technique
        \item[C)] L'audit Off-Page
        \item[D)] L'audit de contenu
    \end{itemize}
    *(Réponse : B)*

    \item La vérification de la présence et de la syntaxe du fichier robots.txt est une étape de :
    \begin{itemize}
        \item[A)] L'audit On-Page
        \item[B)] L'audit technique
        \item[C)] L'audit Off-Page
        \item[D)] L'audit de la concurrence
    \end{itemize}
    *(Réponse : B)*

    \item L'analyse du sitemap XML fait partie de l'audit technique. Cette affirmation est-elle vraie ?
    \begin{itemize}
        \item[A)] Non
        \item[B)] Oui
        \item[C)] Uniquement pour les grands sites
        \item[D)] Uniquement pour le e-commerce
    \end{itemize}
    *(Réponse : B)*

    \item L'analyse du contenu dupliqué fait partie de quel audit ?
    \begin{itemize}
        \item[A)] Technique
        \item[B)] On-Page
        \item[C)] Off-Page
        \item[D)] Concurrentiel
    \end{itemize}
    *(Réponse : B)*

    \item Le test des données structurées peut être réalisé avec quel outil Google ?
    \begin{itemize}
        \item[A)] Google Analytics
        \item[B)] Test de données structurées Google
        \item[C)] Google Ads
        \item[D)] Google Tag Manager
    \end{itemize}
    *(Réponse : B)*

    \item L'audit SEO complet permet d'identifier :
    \begin{itemize}
        \item[A)] Uniquement les problèmes techniques
        \item[B)] Les problèmes techniques, de contenu et de popularité
        \item[C)] Uniquement les problèmes de backlinks
        \item[D)] Uniquement les problèmes de design
    \end{itemize}
    *(Réponse : B)*

    \item Quelle est la première étape de l'audit technique ?
    \begin{itemize}
        \item[A)] Analyse de la vitesse
        \item[B)] Analyse du crawling (robots.txt, sitemap, structure)
        \item[C)] Test mobile-first
        \item[D)] Analyse des données structurées
    \end{itemize}
    *(Réponse : B)*

    \item La vérification de l'indexation (couverture, erreurs) fait partie de :
    \begin{itemize}
        \item[A)] L'audit On-Page
        \item[B)] L'audit technique
        \item[C)] L'audit Off-Page
        \item[D)] L'audit concurrentiel
    \end{itemize}
    *(Réponse : B)*

    \item L'audit SEO complet doit examiner tous les aspects d'un site. Cette affirmation est-elle vraie ?
    \begin{itemize}
        \item[A)] Non, seulement les aspects techniques
        \item[B)] Oui, techniques et stratégiques
        \item[C)] Non, seulement le contenu
        \item[D)] Non, seulement les backlinks
    \end{itemize}
    *(Réponse : B)*

    \item Screaming Frog permet d'analyser quels éléments en audit ?
    \begin{itemize}
        \item[A)] Uniquement les URLs
        \item[B)] Balises titre, meta descriptions, headings, URLs, temps de chargement
        \item[C)] Uniquement les images
        \item[D)] Uniquement les backlinks
    \end{itemize}
    *(Réponse : B)*

    \item L'audit SEO doit être réalisé avant toute stratégie SEO. Cette affirmation est-elle vraie ?
    \begin{itemize}
        \item[A)] Non
        \item[B)] Oui, c'est la première étape
        \item[C)] Après la stratégie
        \item[D)] Uniquement en cas de problème
    \end{itemize}
    *(Réponse : B)*

    \item La vérification des erreurs 404 fait partie de l'audit technique. Cette affirmation est-elle vraie ?
    \begin{itemize}
        \item[A)] Non
        \item[B)] Oui
        \item[C)] Uniquement pour les sites e-commerce
        \item[D)] Uniquement pour les blogs
    \end{itemize}
    *(Réponse : B)*

    \item L'audit de la concurrence est une étape de quelle phase ?
    \begin{itemize}
        \item[A)] Phase 1 (technique)
        \item[B)] Phase 3 (Off-Page)
        \item[C)] Phase 2 (On-Page)
        \item[D)] Aucune
    \end{itemize}
    *(Réponse : B)*

    \item L'évaluation de la qualité du contenu E-E-A-T fait partie de :
    \begin{itemize}
        \item[A)] L'audit technique
        \item[B)] L'audit On-Page
        \item[C)] L'audit Off-Page
        \item[D)] L'audit de vitesse
    \end{itemize}
    *(Réponse : B)*

    \item La vérification des redirections 301 fait partie de l'audit technique. Cette affirmation est-elle vraie ?
    \begin{itemize}
        \item[A)] Non
        \item[B)] Oui
        \item[C)] Uniquement après une migration
        \item[D)] Uniquement pour les sites multilingues
    \end{itemize}
    *(Réponse : B)*

    \item La priorisation des actions après l'audit doit tenir compte de :
    \begin{itemize}
        \item[A)] Uniquement du coût
        \item[B)] De l'impact et de l'effort nécessaire
        \item[C)] Uniquement de l'impact
        \item[D)] Uniquement de la difficulté
    \end{itemize}
    *(Réponse : B)*

    \item Quelle phase de l'audit vérifie le rendu JavaScript ?
    \begin{itemize}
        \item[A)] Phase 2 (On-Page)
        \item[B)] Phase 1 (technique)
        \item[C)] Phase 3 (Off-Page)
        \item[D)] Phase 4
    \end{itemize}
    *(Réponse : B)*

    \item L'audit SEO est essentiel pour :
    \begin{itemize}
        \item[A)] Ignorer les problèmes
        \item[B)] Établir un diagnostic complet et prioriser les actions
        \item[C)] Augmenter le trafic immédiatement
        \item[D)] Créer du contenu
    \end{itemize}
    *(Réponse : B)*

    \item Dans l'exercice d'audit, quel outil est recommandé en version gratuite ?
    \begin{itemize}
        \item[A)] Ahrefs (limité)
        \item[B)] Screaming Frog (max 500 URLs)
        \item[C)] SEMrush (limité)
        \item[D)] Moz (limité)
    \end{itemize}
    *(Réponse : B)*

    \item Que doit contenir le rapport d'audit ?
    \begin{itemize}
        \item[A)] Uniquement une liste de problèmes
        \item[B)] Des constats détaillés, des recommandations et une priorisation
        \item[C)] Uniquement les solutions
        \item[D)] Uniquement le budget nécessaire
    \end{itemize}
    *(Réponse : B)*

    \item L'audit SEO complet est une action ponctuelle. Cette affirmation est-elle vraie ?
    \begin{itemize}
        \item[A)] Oui
        \item[B)] Non, il doit être répété régulièrement
        \item[C)] Uniquement après une pénalité
        \item[D)] Uniquement pour un nouveau site
    \end{itemize}
    *(Réponse : B)*

    \item L'audit SEO complet est important pour établir un diagnostic avant :
    \begin{itemize}
        \item[A)] La création du site
        \item[B)] Toute stratégie SEO
        \item[C)] L'achat de backlinks
        \item[D)] Le lancement des publicités
    \end{itemize}
    *(Réponse : B)*
\end{enumerate}
\section*{Chapitre 35 : Le Plan d'Action Stratégique}
\begin{enumerate}
    \item Quel est l'objectif du plan d'action stratégique en SEO ?
    \begin{itemize}
        \item[A)] Lister toutes les erreurs techniques
        \item[B)] Transformer les recommandations en actions concrètes
        \item[C)] Créer du contenu
        \item[D)] Analyser les concurrents
    \end{itemize}
    *(Réponse : B)*

    \item Dans la matrice de priorisation, quel niveau correspond aux problèmes d'indexation, erreurs 404 et blocages robots.txt ?
    \begin{itemize}
        \item[A)] Stratégique (Priorité 3)
        \item[B)] Critique (Priorité 1)
        \item[C)] Important (Priorité 2)
        \item[D)] Optimisation (Priorité 4)
    \end{itemize}
    *(Réponse : B)*

    \item Quelle priorité est assignée aux Core Web Vitals et au contenu dupliqué ?
    \begin{itemize}
        \item[A)] Critique (Priorité 1)
        \item[B)] Important (Priorité 2)
        \item[C)] Stratégique (Priorité 3)
        \item[D)] Optimisation (Priorité 4)
    \end{itemize}
    *(Réponse : B)*

    \item Le link building et les Topical Clusters sont classés en :
    \begin{itemize}
        \item[A)] Critique (Priorité 1)
        \item[B)] Important (Priorité 2)
        \item[C)] Stratégique (Priorité 3)
        \item[D)] Optimisation (Priorité 4)
    \end{itemize}
    *(Réponse : C)*

    \item L'amélioration continue et les tests A/B sont classés en :
    \begin{itemize}
        \item[A)] Critique (Priorité 1)
        \item[B)] Important (Priorité 2)
        \item[C)] Stratégique (Priorité 3)
        \item[D)] Optimisation (Priorité 4)
    \end{itemize}
    *(Réponse : D)*

    \item Combien de niveaux de priorité comporte la matrice de priorisation SEO ?
    \begin{itemize}
        \item[A)] 2
        \item[B)] 4
        \item[C)] 3
        \item[D)] 5
    \end{itemize}
    *(Réponse : B)*

    \item Que doit-on faire pendant les jours 1-7 du calendrier SEO type ?
    \begin{itemize}
        \item[A)] Stratégie de contenu
        \item[B)] Audit technique et correction des erreurs critiques
        \item[C)] Link building
        \item[D)] Test A/B
    \end{itemize}
    *(Réponse : B)*

    \item Que doit-on faire pendant les semaines 2-4 du calendrier SEO type ?
    \begin{itemize}
        \item[A)] Audit technique
        \item[B)] Optimisation On-Page (balises, contenu, images)
        \item[C)] Link building
        \item[D)] Mesure et analyse
    \end{itemize}
    *(Réponse : B)*

    \item Que doit-on faire pendant les mois 2-3 du calendrier SEO type ?
    \begin{itemize}
        \item[A)] Audit technique
        \item[B)] Stratégie de contenu (Topical Clusters, SEO sémantique)
        \item[C)] Correction des erreurs critiques
        \item[D)] Test A/B
    \end{itemize}
    *(Réponse : B)*

    \item Que doit-on faire pendant les mois 3-6 du calendrier SEO type ?
    \begin{itemize}
        \item[A)] Audit technique
        \item[B)] Campagne de link building et Digital PR
        \item[C)] Optimisation On-Page
        \item[D)] Test A/B
    \end{itemize}
    *(Réponse : B)*

    \item Quelle action est continue dans le calendrier SEO type ?
    \begin{itemize}
        \item[A)] Audit technique
        \item[B)] Mesure, analyse, ajustement itératif
        \item[C)] Correction des erreurs critiques
        \item[D)] Link building
    \end{itemize}
    *(Réponse : B)*

    \item Le suivi des actions est aussi important que :
    \begin{itemize}
        \item[A)] Leur définition
        \item[B)] Le budget alloué
        \item[C)] Le nombre de pages
        \item[D)] Le design
    \end{itemize}
    *(Réponse : A)*

    \item Les recommandations SEO doivent être priorisées par :
    \begin{itemize}
        \item[A)] Ordre alphabétique
        \item[B)] Impact et effort
        \item[C)] Date de création
        \item[D)] Popularité
    \end{itemize}
    *(Réponse : B)*

    \item Un calendrier SEO type s'étend sur combien de temps ?
    \begin{itemize}
        \item[A)] 1 mois
        \item[B)] 3 à 6 mois
        \item[C)] 1 an
        \item[D)] 2 semaines
    \end{itemize}
    *(Réponse : B)*

    \item Les problèmes d'indexation sont considérés comme :
    \begin{itemize}
        \item[A)] Peu importants
        \item[B)] Critiques (Priorité 1)
        \item[C)] Optionnels
        \item[D)] À traiter plus tard
    \end{itemize}
    *(Réponse : B)*

    \item Le E-E-A-T est classé dans quelle priorité ?
    \begin{itemize}
        \item[A)] Critique (Priorité 1)
        \item[B)] Important (Priorité 2)
        \item[C)] Stratégique (Priorité 3)
        \item[D)] Optimisation (Priorité 4)
    \end{itemize}
    *(Réponse : C)*

    \item Les balises manquantes sont classées en :
    \begin{itemize}
        \item[A)] Critique (Priorité 1)
        \item[B)] Important (Priorité 2)
        \item[C)] Stratégique (Priorité 3)
        \item[D)] Optimisation (Priorité 4)
    \end{itemize}
    *(Réponse : B)*

    \item Les tests A/B sont classés en :
    \begin{itemize}
        \item[A)] Critique (Priorité 1)
        \item[B)] Important (Priorité 2)
        \item[C)] Stratégique (Priorité 3)
        \item[D)] Optimisation (Priorité 4)
    \end{itemize}
    *(Réponse : D)*

    \item La première étape du plan d'action stratégique consiste à :
    \begin{itemize}
        \item[A)] Contacter des influenceurs
        \item[B)] Corriger les erreurs critiques d'indexation
        \item[C)] Créer du contenu
        \item[D)] Lancer des campagnes publicitaires
    \end{itemize}
    *(Réponse : B)*

    \item L'optimisation On-Page inclut l'optimisation de :
    \begin{itemize}
        \item[A)] Uniquement les backlinks
        \item[B)] Balises, contenu et images
        \item[C)] Uniquement les vidéos
        \item[D)] Uniquement les URLs
    \end{itemize}
    *(Réponse : B)*

    \item La stratégie de contenu inclut les Topical Clusters et :
    \begin{itemize}
        \item[A)] Le link building
        \item[B)] Le SEO sémantique
        \item[C)] L'audit technique
        \item[D)] Les tests A/B
    \end{itemize}
    *(Réponse : B)*

    \item La campagne de link building est prévue à partir de quel mois ?
    \begin{itemize}
        \item[A)] Mois 1
        \item[B)] Mois 3-6
        \item[C)] Mois 2
        \item[D)] Mois 12
    \end{itemize}
    *(Réponse : B)*

    \item Les erreurs 404 sont classées dans quelle priorité ?
    \begin{itemize}
        \item[A)] Priorité 4 (Optimisation)
        \item[B)] Priorité 1 (Critique)
        \item[C)] Priorité 2 (Important)
        \item[D)] Priorité 3 (Stratégique)
    \end{itemize}
    *(Réponse : B)*

    \item Les raffinements continus sont classés en :
    \begin{itemize}
        \item[A)] Priorité 1
        \item[B)] Priorité 4
        \item[C)] Priorité 2
        \item[D)] Priorité 3
    \end{itemize}
    *(Réponse : B)*

    \item Le plan d'action stratégique doit être basé sur :
    \begin{itemize}
        \item[A)] Les intuitions
        \item[B)] Les résultats de l'audit SEO
        \item[C)] Les tendances
        \item[D)] Les préférences personnelles
    \end{itemize}
    *(Réponse : B)*

    \item Un calendrier SEO type de 3 à 6 mois permet de :
    \begin{itemize}
        \item[A)] Ignorer les problèmes
        \item[B)] Planifier et exécuter les actions de manière organisée
        \item[C)] Travailler au hasard
        \item[D)] Attendre les résultats
    \end{itemize}
    *(Réponse : B)*

    \item Le Digital PR est une composante de :
    \begin{itemize}
        \item[A)] L'audit technique
        \item[B)] La campagne de link building
        \item[C)] L'optimisation On-Page
        \item[D)] Le suivi
    \end{itemize}
    *(Réponse : B)*

    \item La mesure et l'analyse sont des actions :
    \begin{itemize}
        \item[A)] Ponctuelles
        \item[B)] Continues
        \item[C)] Optionnelles
        \item[D)] À faire une fois par an
    \end{itemize}
    *(Réponse : B)*

    \item La matrice de priorisation permet de classer les actions par :
    \begin{itemize}
        \item[A)] Difficulté uniquement
        \item[B)] Impact et urgence
        \item[C)] Coût uniquement
        \item[D)] Popularité
    \end{itemize}
    *(Réponse : B)*

    \item Les actions de priorité 1 sont qualifiées de :
    \begin{itemize}
        \item[A)] Stratégiques
        \item[B)] Critiques
        \item[C)] Importantes
        \item[D)] D'optimisation
    \end{itemize}
    *(Réponse : B)*

    \item Les actions de priorité 3 sont qualifiées de :
    \begin{itemize}
        \item[A)] Critiques
        \item[B)] Stratégiques
        \item[C)] Importantes
        \item[D)] D'optimisation
    \end{itemize}
    *(Réponse : B)*

    \item Le plan d'action stratégique transforme les recommandations en :
    \begin{itemize}
        \item[A)] Questions
        \item[B)] Actions concrètes et planifiées
        \item[C)] Rapports
        \item[D]) Idées
    \end{itemize}
    *(Réponse : B)*

    \item L'ajustement itératif fait référence à :
    \begin{itemize}
        \item[A)] Une action ponctuelle
        \item[B)] L'optimisation continue basée sur les résultats
        \item[C)] Un audit unique
        \item[D)] Une campagne unique
    \end{itemize}
    *(Réponse : B)*

    \item Les Topical Clusters sont une composante de la stratégie de contenu à quel mois ?
    \begin{itemize}
        \item[A)] Jours 1-7
        \item[B)] Mois 2-3
        \item[C)] Semaines 2-4
        \item[D)] Mois 3-6
    \end{itemize}
    *(Réponse : B)*

    \item La correction des erreurs critiques doit être effectuée :
    \begin{itemize}
        \item[A)] Après la stratégie de contenu
        \item[B)] En premier (Jours 1-7)
        \item[C)] Après le link building
        \item[D)] À la fin du calendrier
    \end{itemize}
    *(Réponse : B)*

    \item Les problèmes de blocages robots.txt sont de priorité :
    \begin{itemize}
        \item[A)] 3
        \item[B)] 1 (Critique)
        \item[C)] 2
        \item[D)] 4
    \end{itemize}
    *(Réponse : B)*

    \item Le SEO sémantique est abordé dans quelle phase du calendrier ?
    \begin{itemize}
        \item[A)] Jours 1-7
        \item[B)] Mois 2-3 (stratégie de contenu)
        \item[C)] Semaines 2-4
        \item[D)] Mois 3-6
    \end{itemize}
    *(Réponse : B)*

    \item La priorisation des actions est cruciale pour :
    \begin{itemize}
        \item[A)] Éviter le travail
        \item[B)] Maximiser l'impact avec les ressources disponibles
        \item[C)] Ralentir le projet
        \item[D)] Ignorer les problèmes importants
    \end{itemize}
    *(Réponse : B)*

    \item Un bon plan d'action SEO doit être :
    \begin{itemize}
        \item[A)] Vague et flexible
        \item[B)] Spécifique, mesurable et temporel
        \item[C)] Uniquement basé sur le trafic
        \item[D)] Uniquement basé sur les backlinks
    \end{itemize}
    *(Réponse : B)*

    \item Quel pourcentage de temps est alloué à l'audit technique dans le calendrier type ?
    \begin{itemize}
        \item[A)] 50\%
        \item[B)] 1 semaine sur 3-6 mois (environ 5\%)
        \item[C)] 100\%
        \item[D)] 0\%
    \end{itemize}
    *(Réponse : B)*

    \item Le plan d'action doit être revu et ajusté :
    \begin{itemize}
        \item[A)] Jamais
        \item[B)] Régulièrement en fonction des résultats
        \item[C)] Une seule fois
        \item[D)] Tous les 2 ans
    \end{itemize}
    *(Réponse : B)*

    \item Les actions stratégiques (Priorité 3) incluent :
    \begin{itemize}
        \item[A)] Les erreurs 404
        \item[B)] Le link building et les Topical Clusters
        \item[C)] Les balises manquantes
        \item[D)] Les Core Web Vitals
    \end{itemize}
    *(Réponse : B)*

    \item L'optimisation continue (Priorité 4) inclut :
    \begin{itemize}
        \item[A)] Les blocages robots.txt
        \item[B)] Les tests A/B et les raffinements
        \item[C)] Les problèmes d'indexation
        \item[D)] Les erreurs 404
    \end{itemize}
    *(Réponse : B)*

    \item Le plan d'action stratégique est la dernière étape après :
    \begin{itemize}
        \item[A)] La création de contenu
        \item[B)] L'audit SEO complet
        \item[C)] Le link building
        \item[D)] Les tests A/B
    \end{itemize}
    *(Réponse : B)*

    \item Les recommandations SEO doivent être priorisées par impact et effort. Cette méthode est appelée :
    \begin{itemize}
        \item[A)] Matrice SWOT
        \item[B)] Matrice de priorisation SEO
        \item[C)] Matrice BCG
        \item[D)] Diagramme de Gantt
    \end{itemize}
    *(Réponse : B)*

    \item La Digital PR fait partie de la stratégie de :
    \begin{itemize}
        \item[A)] Contenu
        \item[B)] Link building
        \item[C)] Technique
        \item[D)] On-Page
    \end{itemize}
    *(Réponse : B)*

    \item Le calendrier SEO type inclut des actions pour combien de périodes distinctes ?
    \begin{itemize}
        \item[A)] 2
        \item[B)] 5 (J1-7, S2-4, M2-3, M3-6, Continu)
        \item[C)] 3
        \item[D)] 7
    \end{itemize}
    *(Réponse : B)*

    \item L'optimisation On-Page commence après :
    \begin{itemize}
        \item[A)] Le link building
        \item[B)] L'audit technique et la correction des erreurs critiques
        \item[C)] La stratégie de contenu
        \item[D)] Les tests A/B
    \end{itemize}
    *(Réponse : B)*

    \item Le plan d'action stratégique est essentiel pour :
    \begin{itemize}
        \item[A)] Travailler sans objectif
        \item[B)] Structurer et exécuter efficacement la stratégie SEO
        \item[C)] Ignorer les priorités
        \item[D)] Remplacer l'audit SEO
    \end{itemize}
    *(Réponse : B)*

    \item Les actions de priorité 2 sont qualifiées de :
    \begin{itemize}
        \item[A)] Critiques
        \item[B)] Importantes
        \item[C)] Stratégiques
        \item[D)] D'optimisation
    \end{itemize}
    *(Réponse : B)*
\end{enumerate}
\section*{Chapitre 36 : SEO et Intelligence Artificielle / Tendances}
\begin{enumerate}
    \item Que signifie SGE dans le contexte du SEO ?
    \begin{itemize}
        \item[A)] Search Google Engine
        \item[B)] Search Generative Experience
        \item[C)] Simple Google Engine
        \item[D)] Standard Google Experience
    \end{itemize}
    *(Réponse : B)*

    \item Google SGE utilise l'IA générative pour :
    \begin{itemize}
        \item[A)] Créer des publicités
        \item[B)] Fournir des réponses synthétisées directement dans les SERP
        \item[C)] Supprimer des résultats
        \item[D)] Analyser les emails
    \end{itemize}
    *(Réponse : B)*

    \item Quels outils d'IA sont mentionnés pour la génération de contenu SEO ?
    \begin{itemize}
        \item[A)] Midjourney, DALL-E
        \item[B)] ChatGPT, Claude, Gemini
        \item[C)] Photoshop, Illustrator
        \item[D)] Excel, PowerPoint
    \end{itemize}
    *(Réponse : B)*

    \item L'IA peut être utilisée en SEO pour :
    \begin{itemize}
        \item[A)] Uniquement la génération de contenu
        \item[B)] Génération de contenu, analyse de données, automatisation, personnalisation
        \item[C)] Uniquement l'analyse de données
        \item[D)] Uniquement l'automatisation
    \end{itemize}
    *(Réponse : B)*

    \item Quel phénomène désigne les recherches où l'utilisateur trouve la réponse directement dans les SERP ?
    \begin{itemize}
        \item[A)] Click-through search
        \item[B)] Zero-click search
        \item[C)] Deep search
        \item[D)] Voice search
    \end{itemize}
    *(Réponse : B)*

    \item Les featured snippets font partie de quel phénomène croissant ?
    \begin{itemize}
        \item[A)] Le SEO traditionnel
        \item[B)] Les zero-click searches
        \item[C)] Le Black Hat SEO
        \item[D)] Le SEO vidéo
    \end{itemize}
    *(Réponse : B)*

    \item Le Search Everywhere Optimization signifie que l'optimisation ne se limite plus à :
    \begin{itemize}
        \item[A)] YouTube
        \item[B)] Google
        \item[C)] Amazon
        \item[D)] TikTok
    \end{itemize}
    *(Réponse : B)*

    \item Quelles plateformes sont mentionnées dans le Search Everywhere Optimization ?
    \begin{itemize}
        \item[A)] Google, Bing, Yahoo
        \item[B)] Amazon, YouTube, TikTok, Reddit, ChatGPT
        \item[C)] Facebook, Instagram, Twitter
        \item[D)] LinkedIn, Pinterest, Snapchat
    \end{itemize}
    *(Réponse : B)*

    \item Les technologies blockchain impactent le SEO via :
    \begin{itemize}
        \item[A)] Les publicités
        \item[B)] L'indexation décentralisée et l'authenticité du contenu
        \item[C)] Les réseaux sociaux
        \item[D)] Les emails
    \end{itemize}
    *(Réponse : B)*

    \item Combien de prédictions pour le futur du SEO sont mentionnées (2026-2030) ?
    \begin{itemize}
        \item[A)] 3
        \item[B)] 5
        \item[C)] 7
        \item[D)] 10
    \end{itemize}
    *(Réponse : B)*

    \item Quelle est la première prédiction pour le futur du SEO ?
    \begin{itemize}
        \item[A)] Le SEO vocal dominera
        \item[B)] L'IA générative dominera les SERP
        \item[C)] Les backlinks ne compteront plus
        \item[D)] Google disparaîtra
    \end{itemize}
    *(Réponse : B)*

    \item Quel facteur de classement deviendra le principal selon les prédictions ?
    \begin{itemize}
        \item[A)] Le nombre de backlinks
        \item[B)] L'autorité thématique (Topical Authority)
        \item[C)] La vitesse du site
        \item[D)] Le nombre de pages
    \end{itemize}
    *(Réponse : B)*

    \item La recherche multimodale inclut :
    \begin{itemize}
        \item[A)] Uniquement le texte
        \item[B)] Texte, image, vidéo, audio
        \item[C)] Uniquement la vidéo
        \item[D)] Uniquement l'audio
    \end{itemize}
    *(Réponse : B)*

    \item Selon les prédictions, l'UX sera :
    \begin{itemize}
        \item[A)] Séparée du SEO
        \item[B)] Indissociable du SEO
        \item[C)] Moins importante
        \item[D)] Optionnelle
    \end{itemize}
    *(Réponse : B)*

    \item Les données first-party deviendront cruciales à cause de :
    \begin{itemize}
        \item[A)] L'augmentation des backlinks
        \item[B)] La confidentialité et la fin des cookies tiers
        \item[C)] La baisse du trafic
        \item[D)] L'augmentation du contenu
    \end{itemize}
    *(Réponse : B)*

    \item L'IA transforme à la fois les moteurs de recherche et :
    \begin{itemize}
        \item[A)] Les utilisateurs
        \item[B)] Le SEO
        \item[C)] Les designers
        \item[D)] Les développeurs
    \end{itemize}
    *(Réponse : B)*

    \item Google SGE modifie la façon dont les utilisateurs interagissent avec :
    \begin{itemize}
        \item[A)] Les emails
        \item[B)] Les SERP
        \item[C)] Les applications
        \item[D)] Les vidéos
    \end{itemize}
    *(Réponse : B)*

    \item Quelle stratégie est recommandée pour s'adapter aux zero-click searches ?
    \begin{itemize}
        \item[A)] Ignorer les featured snippets
        \item[B)] Optimiser pour les featured snippets et utiliser les données structurées
        \item[C)] Réduire le contenu
        \item[D)] Supprimer les pages
    \end{itemize}
    *(Réponse : B)*

    \item Le Web3 et le SEO décentralisé introduisent :
    \begin{itemize}
        \item[A)] Moins de transparence
        \item[B)] De nouveaux modèles de confiance et d'autorité
        \item[C)] Plus de publicités
        \item[D)] Moins de contenu
    \end{itemize}
    *(Réponse : B)*

    \item L'IA pour le SEO permet l'automatisation de :
    \begin{itemize}
        \item[A)] Uniquement la création de contenu
        \item[B)] Le reporting, l'audit et le suivi des positions
        \item[C)] Uniquement l'audit
        \item[D)] Uniquement le reporting
    \end{itemize}
    *(Réponse : B)*

    \item Les recherches zero-click augmentent constamment. Cette affirmation est-elle vraie ?
    \begin{itemize}
        \item[A)] Non
        \item[B)] Oui
        \item[C)] Uniquement sur mobile
        \item[D)] Uniquement sur desktop
    \end{itemize}
    *(Réponse : B)*

    \item Le SEO devient Search Everywhere Optimization. Cela signifie qu'il faut optimiser pour :
    \begin{itemize}
        \item[A)] Uniquement Google
        \item[B)] Google + autres plateformes (Amazon, YouTube, TikTok, ChatGPT, etc.)
        \item[C)] Uniquement les réseaux sociaux
        \item[D)] Uniquement les moteurs de recherche traditionnels
    \end{itemize}
    *(Réponse : B)*

    \item Quelle est une tendance émergente du SEO ?
    \begin{itemize}
        \item[A)] La diminution de l'importance du contenu
        \item[B)] L'essor des zero-click searches
        \item[C)] La fin du référencement
        \item[D)] La baisse de l'utilisation de Google
    \end{itemize}
    *(Réponse : B)*

    \item Les featured snippets sont aussi appelés :
    \begin{itemize}
        \item[A)] Résultats payants
        \item[B)] Position zéro
        \item[C)] Résultats sponsorisés
        \item[D)] Publicités natives
    \end{itemize}
    *(Réponse : B)*

    \item Pour s'adapter au SGE, il faut :
    \begin{itemize}
        \item[A)] Ignorer l'IA
        \item[B)] Créer du contenu de qualité et bien structuré avec des données structurées
        \item[C)] Réduire la production de contenu
        \item[D)] Utiliser uniquement du contenu généré par IA
    \end{itemize}
    *(Réponse : B)*

    \item L'autorité thématique (Topical Authority) se construit par :
    \begin{itemize}
        \item[A)] L'achat de backlinks
        \item[B)] La publication de contenu exhaustif et interconnecté sur un sujet
        \item[C)] La création de nombreuses pages sans lien entre elles
        \item[D)] Le duplicate content
    \end{itemize}
    *(Réponse : B)*

    \item La recherche multimodale signifie que les utilisateurs chercheront avec :
    \begin{itemize}
        \item[A)] Uniquement la voix
        \item[B)] Plusieurs modes (texte, image, vidéo, audio)
        \item[C)] Uniquement le texte
        \item[D)] Uniquement les images
    \end{itemize}
    *(Réponse : B)*

    \item L'IA générative dans Google SGE synthétise les réponses à partir de :
    \begin{itemize}
        \item[A)] Uniquement des publicités
        \item[B)] Plusieurs sources web
        \item[C)] Uniquement des vidéos
        \item[D)] Uniquement des réseaux sociaux
    \end{itemize}
    *(Réponse : B)*

    \item La personnalisation du contenu grâce à l'IA permet :
    \begin{itemize}
        \item[A)] Un contenu identique pour tous
        \item[B)] Un contenu adapté à l'intention de l'utilisateur
        \item[C)] Moins de pertinence
        \item[D)] Des pages plus longues
    \end{itemize}
    *(Réponse : B)*

    \item Quel est l'impact des zero-click searches sur le trafic traditionnel ?
    \begin{itemize}
        \item[A)] Aucun impact
        \item[B)] Réduction du trafic organique traditionnel
        \item[C)] Augmentation du trafic
        \item[D)] Trafic inchangé
    \end{itemize}
    *(Réponse : B)*

    \item Les utilisateurs cherchent aussi sur Amazon pour :
    \begin{itemize}
        \item[A)] Les actualités
        \item[B)] Les produits
        \item[C)] Les vidéos
        \item[D)] Les réseaux sociaux
    \end{itemize}
    *(Réponse : B)*

    \item TikTok est utilisé comme moteur de recherche pour :
    \begin{itemize}
        \item[A)] Les produits
        \item[B)] Les tendances
        \item[C)] Les actualités
        \item[D)] Les emails
    \end{itemize}
    *(Réponse : B)*

    \item Reddit est utilisé comme moteur de recherche pour :
    \begin{itemize}
        \item[A)] Les produits
        \item[B)] Les avis
        \item[C)] Les vidéos
        \item[D)] Les actualités
    \end{itemize}
    *(Réponse : B)*

    \item ChatGPT et les LLMs sont utilisés pour :
    \begin{itemize}
        \item[A)] Les achats
        \item[B)] Les réponses aux questions
        \item[C)] Les réservations
        \item[D)] Les jeux
    \end{itemize}
    *(Réponse : B)*

    \item Les modèles de confiance et d'autorité sont impactés par :
    \begin{itemize}
        \item[A)] Le SEO traditionnel
        \item[B)] Le Web3 et la blockchain
        \item[C)] Les cookies
        \item[D)] Les publicités
    \end{itemize}
    *(Réponse : B)*

    \item L'adaptation aux nouvelles technologies est la clé de :
    \begin{itemize}
        \item[A)] La stagnation
        \item[B)] La pérennité du SEO
        \item[C)] La régression
        \item[D)] L'échec
    \end{itemize}
    *(Réponse : B)*

    \item Pour mesurer sa visibilité au-delà des clics, il faut :
    \begin{itemize}
        \item[A)] Ignorer les impressions
        \item[B)] Mesurer les impressions et la présence dans les featured snippets
        \item[C)] Compter uniquement les clics
        \item[D)] Ignorer les SERP
    \end{itemize}
    *(Réponse : B)*

    \item Dans le Search Everywhere Optimization, YouTube est utilisé pour chercher :
    \begin{itemize}
        \item[A)] Des produits
        \item[B)] Des vidéos
        \item[C)] Des avis
        \item[D)] Des actualités
    \end{itemize}
    *(Réponse : B)*

    \item L'indexation décentralisée est une caractéristique de :
    \begin{itemize}
        \item[A)] Google
        \item[B)] Web3
        \item[C)] Bing
        \item[D)] Yahoo
    \end{itemize}
    *(Réponse : B)*

    \item L'authenticité du contenu vérifiable est rendue possible par :
    \begin{itemize}
        \item[A)] Les cookies
        \item[B)] La blockchain
        \item[C)] Les publicités
        \item[D)] Les emails
    \end{itemize}
    *(Réponse : B)*

    \item Le SGE de Google utilise l'IA générative. Cette affirmation est-elle vraie ?
    \begin{itemize}
        \item[A)] Non
        \item[B)] Oui
        \item[C)] Uniquement pour les recherches payantes
        \item[D)] Uniquement pour les recherches en anglais
    \end{itemize}
    *(Réponse : B)*

    \item Combien de prédictions sont listées dans le "Futur du SEO : 5 prédictions" ?
    \begin{itemize}
        \item[A)] 3
        \item[B)] 5
        \item[C)] 7
        \item[D)] 10
    \end{itemize}
    *(Réponse : B)*

    \item Les featured snippets réduisent le besoin de cliquer sur les résultats. Cette affirmation est-elle vraie ?
    \begin{itemize}
        \item[A)] Non
        \item[B)] Oui, c'est le principe des zero-click searches
        \item[C)] Uniquement sur mobile
        \item[D)] Uniquement pour les vidéos
    \end{itemize}
    *(Réponse : B)*

    \item L'IA est un outil puissant pour la création et l'optimisation de contenu. Cette affirmation est-elle vraie ?
    \begin{itemize}
        \item[A)] Non
        \item[B)] Oui
        \item[C)] Uniquement pour le e-commerce
        \item[D)] Uniquement pour les réseaux sociaux
    \end{itemize}
    *(Réponse : B)*

    \item Les zero-click searches augmentent le trafic organique traditionnel. Cette affirmation est-elle vraie ?
    \begin{itemize}
        \item[A)] Oui
        \item[B)] Non, elles le réduisent
        \item[C)] Aucun impact
        \item[D)] Uniquement sur mobile
    \end{itemize}
    *(Réponse : B)*

    \item Pour s'adapter aux évolutions 2026-2030, il faut :
    \begin{itemize}
        \item[A)] Ignorer les tendances
        \item[B)] Se préparer à l'IA générative, à l'autorité thématique et à la recherche multimodale
        \item[C)] Arrêter le SEO
        \item[D)] Revenir aux méthodes traditionnelles
    \end{itemize}
    *(Réponse : B)*

    \item Le SEO doit évoluer vers une approche multi-plateforme. Cette affirmation est-elle vraie ?
    \begin{itemize}
        \item[A)] Non
        \item[B)] Oui, Search Everywhere Optimization
        \item[C)] Non, Google reste la seule plateforme importante
        \item[D)] Uniquement pour les grandes entreprises
    \end{itemize}
    *(Réponse : B)*

    \item L'IA peut être utilisée pour l'analyse massive de données SEO. Cette affirmation est-elle vraie ?
    \begin{itemize}
        \item[A)] Non
        \item[B)] Oui
        \item[C)] Uniquement pour les petits sites
        \item[D)] Uniquement pour les données de backlinks
    \end{itemize}
    *(Réponse : B)*

    \item La confidentialité et les données first-party deviennent cruciales car :
    \begin{itemize}
        \item[A)] Les cookies tiers sont plus fiables
        \item[B)] Les cookies tiers sont progressivement supprimés
        \item[C)] Les données first-party sont moins chères
        \item[D)] Les données first-party sont inexactes
    \end{itemize}
    *(Réponse : B)*

    \item L'IA générative et Google SGE transforment les SERP. Cette affirmation est-elle vraie ?
    \begin{itemize}
        \item[A)] Non
        \item[B)] Oui
        \item[C)] Uniquement pour les recherches locales
        \item[D)] Uniquement pour les recherches e-commerce
    \end{itemize}
    *(Réponse : B)*
\end{enumerate}

\end{document}